%%%% DON'T EDIT THIS FILE %%%%

\usepackage{xifthen}% \isempty
\usepackage{etoolbox}% \patchcmd
% \usepackage{stringstrings}% \substring

\usepackage{amsmath, amsthm, amssymb}
\usepackage[margin=1in]{geometry}
\usepackage{xspace}
\usepackage{graphicx}
\usepackage[normalem]{ulem}
\usepackage{adjustbox}
\usepackage{xcolor}
\usepackage{changepage}
\usepackage{framed}
\usepackage{multicol}
\usepackage{enumitem}
% \usepackage{listings} % lstlistings environment

\usepackage{makecmds}           %\provideenvironment

\usepackage{xparse,environ}
\usepackage{letltxmacro}       % \LetLtxMacro
\usepackage{wrapfig}
\usepackage{tikz}

\usetikzlibrary{shapes}
\usetikzlibrary{positioning}
\usetikzlibrary{shapes.misc}
\usetikzlibrary{arrows}

\newtheorem{lemma}{Lemma}
\newtheorem{definition}{Definition}
\newtheorem{theorem}{Theorem}
\newtheorem{corollary}{Corollary}
\newtheorem{claim}{Claim}
\newtheorem{observation}{Observation}
\newtheorem{conjecture}{Conjecture}

% A lemma or theorem restated later with its original number (set with
% \setcounter) goes in a restatement environment, which gives it its own link
% anchor (otherwise it would duplicate the original's).
\newenvironment{restatement}
  {\def\theHlemma{restated.\thelemma}\def\theHtheorem{restated.\thetheorem}}{}

\newenvironment{myframed}
{\smallskip\par\noindent\begin{minipage}{\textwidth}\begin{framed}\vspace*{-5pt}}{\vspace*{-5pt}\end{framed}\end{minipage}\smallskip}


\definecolor{darkblue}{RGB}{0,0,200}

 
\let\oldParagraph\paragraph
\providecommand{\gradescopeColor}{\color{black}}
\renewcommand{\paragraph}[1]{\oldParagraph{\gradescopeColor #1}}
\providecommand{\gradescopeInit}{}

\providecommand{\BLUE}[1]{{\color{darkblue} #1}}
\providecommand{\REPLACEME}{\textcolor{red}{\ensuremath{\star \cdots \star}}\xspace}
\provideenvironment{blue}{\color{darkblue}}{}

\newcommand{\dist}[2]{\ensuremath{\operatorname{\sf dist}(#1, #2)}}
\newcommand{\depth}[2]{\ensuremath{\operatorname{\sf depth}_{#1}(#2)}}
\newcommand{\cost}[1]{\ensuremath{\operatorname{\sf cost}(#1)}}
\newcommand{\weight}[1]{\ensuremath{\operatorname{\sf wt}(#1)}}
% \newcommand{\val}[1]{\ensuremath{\operatorname{\sf val}(#1)}}
\newcommand{\eps}{\varepsilon}
\newcommand{\capac}{\textsf{cap}}
\newcommand{\capOf}[1]{\ensuremath{\operatorname{\sf cap}(#1)}}
\newcommand{\rescap}[1]{\ensuremath{\operatorname{\sf cap}_f(#1)}}
\newcommand{\into}[1]{\ensuremath{\operatorname{\sf in}(#1)}}
\newcommand{\out}[1]{\ensuremath{\operatorname{\sf out}(#1)}}
\newenvironment{LP}[1]{
  \begin{array}[t]{@{}l@{}}
    \displaystyle #1 \\
    \begin{aligned}\\[-10pt]}{\end{aligned}\end{array}}

\newcommand{\opt}[1]{\ensuremath{\operatorname{\sf opt}(#1)}}
\newcommand{\Opt}[1]{\ensuremath{\operatorname{\sf Opt}(#1)}}
\newcommand{\half}{\frac{1}{2}}

\newcommand{\E}{\operatorname{E}}
\newcommand{\giv}{\,|\,}
\newcommand{\Z}{\mathbb{Z}}
\newcommand{\Zp}{\mathbb{Z}_{\ge 0}}
\newcommand{\R}{\mathbb{R}}
% \newcommand{\Rp}{\mathbb{R}_{\ge 0}}
\newcommand{\Rp}{\mathbb{R}_{\tiny +}}
\newcommand{\N}{\mathbb{N}}
\newcommand{\Np}{\mathbb{N}_{\ge 0}}
\newcommand{\calS}{{\cal S}}

\newcommand{\OP}[1]{\text{\small\sc #1}\xspace}
\newcommand{\Insert}{\OP{Insert}}
\newcommand{\FindMin}{\OP{Find-Min}}
\newcommand{\DeleteMin}{\OP{Delete-Min}}
\newcommand{\DecreaseKey}{\OP{Decrease}}
\newcommand{\Merge}{\OP{Merge}}
\newcommand{\Cut}{\OP{Cut}}
% \newcommand{\prob}[1]{\text{\sc #1}}
% a problem name; in running text, lines may break inside it
\DeclareRobustCommand{\prob}[1]{\ifmmode\text{#1}\else#1\fi}
\newcommand{\encode}[1]{\langle #1 \rangle}

\newcommand{\best}[1]{\text{\sf best}(#1)}
\newcommand{\worst}[1]{\text{\sf worst}(#1)}
\newcommand{\floor}[1]{\lfloor #1 \rfloor}
\newcommand{\ceil}[1]{\lceil #1 \rceil}

\setlength{\parskip}{2pt}
\setlength{\parindent}{1em}

% \setlength\fboxsep{0pt}

%%% longFormProof, shortFormProof, block environments
%%% \step

\newlist{steps}{enumerate}{7}
\setlist[steps]{wide,
  topsep=1pt,
  parsep=1pt,
  partopsep=1pt,
  itemsep=0pt,
  align=left,
  labelindent=0pt,
  itemindent=0pt, 
}
\setlist[steps,1]{
  label=\arabic*.,
  ref=\arabic*
}
\setlist[steps,2]{
  label=\arabic{stepsi}.\arabic*.,
  ref=\arabic{stepsi}.\arabic*
}
\setlist[steps,3]{
  label=\arabic{stepsi}.\arabic{stepsii}.\arabic*.,
  ref=\arabic{stepsi}.\arabic{stepsii}.\arabic*
}
\setlist[steps,4]{
  label=\arabic{stepsi}.\arabic{stepsii}.\arabic{stepsiii}.\arabic*.,
  ref=\arabic{stepsi}.\arabic{stepsii}.\arabic{stepsiii}.\arabic*
}
\setlist[steps,5]{
  label=\arabic{stepsi}.\arabic{stepsii}.\arabic{stepsiii}.\arabic{stepsiv}.\arabic*.,
  ref=\arabic{stepsi}.\arabic{stepsii}.\arabic{stepsiii}.\arabic{stepsiv}.\arabic*
}
\setlist[steps,6]{
  label=\arabic{stepsi}.\arabic{stepsii}.\arabic{stepsiii}.\arabic{stepsiv}.\arabic{stepsv}.\arabic*.,
  ref=\arabic{stepsi}.\arabic{stepsii}.\arabic{stepsiii}.\arabic{stepsiv}.\arabic{stepsv}.\arabic*
}
\setlist[steps,7]{
  label=\arabic{stepsi}.\arabic{stepsii}.\arabic{stepsiii}.\arabic{stepsiv}. \arabic{stepsv}. \arabic{stepsvi}.\arabic*.,
  ref=\arabic{stepsi}.\arabic{stepsii}.\arabic{stepsiii}.\arabic{stepsiv}.\arabic{stepsv}. \arabic{stepsvi}.\arabic*
}

\newcommand{\step}[1][]{\item\maybeLabel{#1}}

\makeatletter
% \skipitem steps the list counter itself (for block labels).  Give each one
% its own hyperref link target, since step numbers repeat from proof to proof.
\newcounter{skipitem}
\newcommand{\skipitem}{\stepcounter{skipitem}\refstepcounter{\@enumctr}}
\renewcommand{\theHstepsi}{skipitem.\arabic{skipitem}}
\renewcommand{\theHstepsii}{skipitem.\arabic{skipitem}}
\renewcommand{\theHstepsiii}{skipitem.\arabic{skipitem}}
\renewcommand{\theHstepsiv}{skipitem.\arabic{skipitem}}
\renewcommand{\theHstepsv}{skipitem.\arabic{skipitem}}
\renewcommand{\theHstepsvi}{skipitem.\arabic{skipitem}}
\renewcommand{\theHstepsvii}{skipitem.\arabic{skipitem}}
\makeatother

\newcommand{\maybeLabel}[1]{\ifthenelse{\isempty{#1}}{}{\label{#1}}}

\newenvironment{longFormProof}[1][Proof (long form).]
{\begin{proof}[#1]~\par\begin{steps}}{\end{steps}\vspace*{-3ex}\end{proof}}

\newenvironment{shortFormProof}[1][Proof (short form).]
{\begin{proof}[#1]}{\end{proof}}

\newenvironment{longFormSubProof}[1][Proof.]
{\begin{adjustwidth}{1em}{0em}\begin{proof}[#1]~\par\vspace*{-2pt}\skipitem\begin{steps}}
{\end{steps}\vspace*{-1ex}\end{proof}\end{adjustwidth}\smallskip}

\newenvironment{shortFormSubProof}[1][Proof.]
{\begin{proof}[#1]}
{\end{proof}\vspace*{-3pt}}

\newenvironment{block}[1][]
{\skipitem\maybeLabel{#1}\item[] \begin{steps}}{\end{steps}}

\newcommand{\comment}[1]{\hfill{\footnotesize\emph{#1}}}
\newcommand{\lineacross}{\par\vspace*{-0.7\baselineskip}\noindent\hrulefill\par}

%%% problem

% \newcounter{problem}
% \newcommand{\problem}[1][]{}
% {\ifthenelse{\isempty{#1}}{\refstepcounter{problem}}{}%
%   \paragraph{Problem \ifthenelse{\isempty{#1}}
%     {\bf\arabic{problem}.}
%     {#1}} }

\newcommand{\problemNumbered}[1]{
  \setcounter{problem}{#1}%
  \addtocounter{problem}{-1}%
  \problem
}
\newenvironment{problemTheorem}{%
  \setcounter{theorem}{\value{problem}}%
  \addtocounter{theorem}{-1}%
  \begin{theorem}[Problem \arabic{problem}]
  }{\end{theorem}}

\newcounter{problem}
\newcommand{\problem}[1][]
{\ifthenelse{\isempty{#1}}
  {\refstepcounter{problem}\paragraph{Problem~\arabic{problem}.}}
  {\paragraph{Problem~#1.}}\refreshHeader} 

\newcounter{exercise}
\newcommand{\exercise}[1][]
{\ifthenelse{\isempty{#1}}
  {\refstepcounter{exercise}\paragraph{Exercise~\theexercise.}}
  {\paragraph{Exercise~#1.}}}

\newcommand{\solution}{\paragraph{Solution.}} 

\newenvironment{mylist}[1][]
{\begin{list}{#1}{
      \setlength{\itemsep}{-1pt}
      \setlength{\labelwidth}{0pt}
      \setlength{\leftmargin}{1.25em}
    }}
{\end{list}}

\newcommand{\LNTitle}{\textcolor{red}{UNDEFINED LNTitle}}

% In a standalone lecture note's table of contents, leave room for numbers
% like 7.10 (in bold) and 7.10.2.
\makeatletter
\newcommand{\LNtocWidths}{%
  \patchcmd{\l@section}{1.5em}{2.8em}{}{\PackageWarning{macros}{Could not widen section numbers in the contents}}%
  \renewcommand*{\l@subsection}{\@dottedtocline{2}{2.8em}{3.2em}}%
  \renewcommand*{\l@subsubsection}{\@dottedtocline{3}{6.0em}{4.1em}}}
\makeatother

\newcommand{\doTitle}[3]{
  \pagenumbering{roman}
  \renewcommand{\LNTitle}{Lecture Note #1, #3}
  \renewcommand{\currentLN}{#1}
  \codeListingstrue
  % number sections, theorems etc., exercises, figures, tables and equations
  % within the lecture note (e.g. Section 8.2, Lemma 8.6), as in the book
  \renewcommand{\thesection}{#1.\arabic{section}}
  \renewcommand{\thelemma}{#1.\arabic{lemma}}
  \renewcommand{\thedefinition}{#1.\arabic{definition}}
  \renewcommand{\thetheorem}{#1.\arabic{theorem}}
  \renewcommand{\thecorollary}{#1.\arabic{corollary}}
  \renewcommand{\theclaim}{#1.\arabic{claim}}
  \renewcommand{\theobservation}{#1.\arabic{observation}}
  \renewcommand{\theconjecture}{#1.\arabic{conjecture}}
  \renewcommand{\theexercise}{#1.\arabic{exercise}}
  \renewcommand{\thefigure}{#1.\arabic{figure}}
  \renewcommand{\thetable}{#1.\arabic{table}}
  \renewcommand{\theequation}{#1.\arabic{equation}}
  \LNtocWidths
  \title{Lecture Note #1\\ #2}
  % a note's \LNabstract, if it defines one before \doTitle, goes right below the title
  \date{\ifdefined\LNabstract\parbox{0.85\textwidth}{\normalsize\normalfont\LNabstract}\\[3ex]\fi
    draft of \today, by \href{\authorURL}{N.~Young}\\[2ex] {\small\copyrightNotice}}
  \maketitle
  {\small
    \hypertarget{TOC}{}
    \tableofcontents
  }
  \thispagestyle{empty}
  % the text starts on page 1, after the title and contents (numbered i, ii, ...)
  \clearpage
  \pagenumbering{arabic}
}

\def\clap#1{\hbox to 0pt{\hss#1\hss}}

\newcommand{\refreshHeader}{}
\newcommand{\header}{}

\newcommand{\doHeader}[1]{      % for lecture notes
  \renewcommand{\header}{%
    \clap{\raisebox{0pt}{\protect\hyperlink{TOC}{$\uparrow$\scriptsize\,\textsc{contents}\hspace{1.1in}}}}%
    {\footnotesize \LNTitle{} / Section \thesection, #1}} %, #1, draft of \today}}
  \markboth{\header}{\header}
}

\newcommand{\SID}{}
\newcommand{\setAuthor}[1]{\renewcommand{\author}{#1}}
\newcommand{\setSID}[1]{\renewcommand{\SID}{#1}}

\newcommand{\doHwHeader}[1]{                                         % for homeworks
  \renewcommand{\refreshHeader}{
    \renewcommand{\header}{%
      {\footnotesize \classSection #1 / Problem~\arabic{problem} / draft of \today \hfill \author\ifGradescope~(\SID)\fi~~}}
    \markboth{\header}{\header}
  }
  \renewcommand{\header}{%
    {\footnotesize \classSection #1 / Instructions / draft of \today \hfill \author\ifGradescope~(\SID)\fi~~}}
  \markboth{\header}{\header}
}

% Homework settings, for instructors using the homeworks in a course:
% \classSection is the course name shown in each homework's running head
% (e.g. \renewcommand{\classSection}{CS 3000\xspace}), and \Gradescopetrue turns on
% the instructions for submitting to Gradescope (in its fixed-length mode),
% the hints saying on which page of the PDF each answer should be, and the
% student ID in the running head.
\newcommand{\classSection}{}
\newif\ifGradescope
\Gradescopefalse
\newcommand{\GradescopeComment}[1]{\ifGradescope\comment{#1}\fi}

\pagestyle{myheadings}
\addtolength{\headsep}{0.2in}
\addtolength{\topmargin}{-0.3in}
\addtolength{\textheight}{0.4in}

%%%%%%%%%%%%%%%%%%%%%%%%%%%%%%%%%%%%%%%%%%%%%%%%
%%%%%%%%%%%%%%%%%%%%%%%%%%%%%%%%%%%%%%%%%%%%%%%%

\newcommand{\MF}{\href
  {https://mfleck.cs.illinois.edu/building-blocks/}
  {MF}\xspace}

\newcommand{\KWSlides}{\href
  {https://www.cs.princeton.edu/~wayne/kleinberg-tardos/}
  {KW slides}\xspace}

\newcommand{\LLM}{\href
  {https://courses.csail.mit.edu/6.042/spring18/mcs.pdf}
  {LLM}\xspace}

% \DPV, \KT, \CLRS: textbooks listed (without links) in the table of sources in
% the Preliminaries note; each prints as a link to its entry in that table.
\DeclareRobustCommand{\DPV}{\noteLink{prelim}{source.DPV}{DPV}\xspace}
\DeclareRobustCommand{\KT}{\noteLink{prelim}{source.KT}{KT}\xspace}
\DeclareRobustCommand{\CLRS}{\noteLink{prelim}{source.CLRS}{CLRS}\xspace}
\AtBeginDocument{\pdfstringdefDisableCommands{\def\DPV{DPV }\def\KT{KT }\def\CLRS{CLRS }}}

\newcommand{\JE}{\href
  {https://jeffe.cs.illinois.edu/teaching/algorithms/}
  {JE}\xspace}

\newcommand{\JEGreedy}{\href
  {https://jeffe.cs.illinois.edu/teaching/algorithms/book/04-greedy.pdf}
  {JE Chapter 4}\xspace}

\newcommand{\BBCode}{\href
{https://northeastern.blackboard.com/webapps/cmsmain/webui/courses/CS3000.MERGED.202010/code?action=frameset&subaction=view&uniq=-jj2max&course_id=_2597392_1&mask=\%2Fcourses\%2FCS3000.MERGED.202010}
  {code}\xspace}

% The Python code (src/code).  Each lecture note that has code ends with a
% section listing its files: \codeSection{path, path, ...}.  In the text,
% \pycode{greedy_algorithms/Dijkstra.py} prints the path as a link to its listing
% (in the homeworks, which have no listings, to the file on GitHub), and
% \pycode[text]{...} uses the given text instead.  \pythonCode{path, ...} is the
% standard sentence ending a section whose algorithms have code:
% "Python code: path, ... (Section N.k)."
\newcommand{\codeFolder}{https://github.com/nealeyoung/algs101/tree/main/src/code}
\newcommand{\codeRepo}{https://github.com/nealeyoung/algs101/blob/main/src/code}
\newif\ifcodeListings   % set by \doTitle: this document has the listings
\newcommand{\githubCode}[2][]{\href{\codeRepo/#2}{\ifx\relax#1\relax\nolinkurl{#2}\else#1\fi}}
\newcommand{\pycode}[2][]{%
  \ifcodeListings\hyperref[code: #2]{\ifx\relax#1\relax\nolinkurl{#2}\else#1\fi}%
  \else\githubCode[#1]{#2}\fi}
\makeatletter
\newcommand{\pythonCode}[1]{%
  \def\pythonCode@first{}\def\pythonCode@list{}%
  \forcsvlist{\pythonCode@add}{#1}%
  Python code: \pythonCode@list\ (Section~\ref{code: \pythonCode@first}).}
\newcommand{\pythonCode@add}[1]{%
  \ifx\pythonCode@first\@empty\def\pythonCode@first{#1}\appto\pythonCode@list{\pycode{#1}}%
  \else\appto\pythonCode@list{, \pycode{#1}}\fi}
\makeatother

% Code listings that copy and paste exactly (in PDF viewers that honor the
% ToUnicode maps, e.g. MuPDF-based ones; PDFKit viewers such as Preview and
% Skim drop leading spaces, so each listing's heading links to the file on
% GitHub).  Spaces are typeset as white visible-space glyphs that copy as
% spaces, and hyphens and asterisks as the T1 glyphs, so they copy as themselves.
\usepackage[T1,OT1]{fontenc}   % T1 for the code listings only; OT1 stays the default
\usepackage{listings}
\pdfgentounicode=1
\input glyphtounicode
\pdfglyphtounicode{uni2423}{0020}
\makeatletter
\def\lst@visiblespace{\textcolor{white}{\char32}}
% put an invisible space on each blank line, so that blank lines copy too
% (listings' OnEmptyLine hook runs before the line break, so patch the loop)
\def\lst@DoNewLines{%
    \@whilenum\lst@newlines>\@ne \do {\lst@NewLine \lst@visiblespace}%
    \ifnum\lst@newlines>\z@ \lst@NewLine \fi}
\makeatother
\lstset{language=Python,
  basicstyle=\fontencoding{T1}\fontfamily{lmtt}\selectfont\footnotesize,
  upquote=true, showstringspaces=false, columns=fullflexible, keepspaces=true,
  showspaces=true, literate={-}{{\char45}}1 {*}{{\char42}}1,
  frame=single, framesep=1ex, rulecolor=\color{gray}}
% the listings' paths are relative to the lecture note's folder (in the book, to src/)
\ifdefined\BookMode\newcommand{\codeDir}{code/}\else\newcommand{\codeDir}{../code/}\fi
% \codefile{path}: a listing of src/code/path, headed by a link to it on GitHub
\newcommand{\codefile}[1]{%
  \subsection*{\githubCode{#1}}\label{code: #1}%
  \lstinputlisting{\codeDir#1}}
% \codeSection{path, path, ...}: the lecture note's last section, listing its code
\newcommand{\codeSection}[1]{%
  \newpage
  \section{Python code}\label{\LNtag{\currentLN}: sec: code}\doHeader{Python code}
  This section lists the Python code for this lecture note.
  Click on a file's name to open the file on GitHub.
  Copying code from GitHub is more reliable than copying it from this PDF,
  since some PDF viewers drop the leading spaces (the indentation) of copied lines.
  \forcsvlist{\codefile}{#1}}
% the copyright notice, for title pages and the homework first pages
\newcommand{\licenseURL}{https://github.com/nealeyoung/algs101/blob/main/LICENSE}
\newcommand{\authorURL}{https://www.cs.ucr.edu/\string~neal}
\newcommand{\copyrightNotice}{%
  \copyright\ 2018--2026 \href{\authorURL}{Neal E. Young}.
  Text: CC BY-NC-SA 4.0; code: MIT; except where noted.
  See \url{\licenseURL}.}
% the published site: PDFs of the notes, and the homework assignments and templates
\newcommand{\siteURL}{https://nealeyoung.github.io/algs101/}

\newcommand{\blackboard}{\href
  {https://northeastern.blackboard.com/webapps/blackboard/execute/launcher?type=Course&id=_2597392_1&url=}
  {Blackboard}\xspace}

\newcommand{\JUSlides}{\href
  {http://www.khoury.neu.edu/home/jullman/cs3000f18/schedule.html}
  {Prof.~Ullman's slides}\xspace}

% a URL on its own line, in three layouts
\newcommand{\URL}[1]{\par-- \parbox[t]{0.8\textwidth}{\url{#1}} \medskip\par}
\newcommand{\URLsmall}[1]{\par-- \parbox[t]{0.95\textwidth}{\small\url{#1}} \medskip\par}
\newcommand{\URLplain}[1]{\par\parbox[t]{\textwidth}{\url{#1}}\par}

%%%%%%%%%%%%%%%%%%%%%%%%%%%%%%%%%%%%%%%%%%%%%%%%
%%%%%%%%%%%%%%%%%%%%%%%%%%%%%%%%%%%%%%%%%%%%%%%%

\newcommand{\continued}{\vfill\centerline{\emph{(continued)}}\vfill\newpage}

% load hyperref last to avoid conflicts; long URLs may also break after hyphens
\PassOptionsToPackage{hyphens}{url}
\ifdefined\BookMode\else\usepackage{xr-hyper}\fi   % for \lectureNoteRefs (homeworks), below
\usepackage{hyperref}
\hypersetup{
    colorlinks,
    linkcolor={blue!50!black},
    citecolor={blue!50!black},
    urlcolor={blue!80!black}
}

\providecommand{\answerDirectory}{.}
\newcommand{\inputAnswer}[1]{\refreshHeader\input{\answerDirectory/#1}}
\newcommand{\inputSection}[1]{\input{#1}}

% Refer to a lecture note by its tag (the prefix of its labels), e.g.
% \LN{proofs} prints "Lecture Note 2" (in the book, a link to that chapter).  The table below is the only place the numbers appear.
\makeatletter
\newcommand{\defineLN}[3]{\@namedef{LNnumber@#2}{#1}\@namedef{LNtag@#1}{#2}\@namedef{LNfile@#2}{#3}}
\newcommand{\LNnumber}[1]{%
  \ifcsname LNnumber@#1\endcsname\csname LNnumber@#1\endcsname
  \else\PackageError{macros}{Unknown lecture-note tag `#1'}{}\fi}
\newcommand{\LNtag}[1]{\csname LNtag@#1\endcsname}
\newcommand{\LNfile}[1]{\csname LNfile@#1\endcsname}
\makeatother
\defineLN{1}{prelim}{1_preliminaries}
\defineLN{2}{proofs}{2_proofs}
\defineLN{3}{matching}{3_stable_matching}
\defineLN{4}{dc}{4_divide_and_conquer}
\defineLN{5}{greedy}{5_greedy}
\defineLN{6}{graphs}{6_graph_traversal}
\defineLN{7}{dp}{7_dynamic_programming}
\defineLN{8}{flow}{8_flow_reductions_LP}
\defineLN{9}{np}{9_NP}

% The text is the same in the book and in the standalone lecture notes; only
% links differ.  In the book, links to other lecture notes stay in the book.
% In a standalone note, they open the other note's published PDF.
% \noteLink{tag}{target}{text} links text to \hypertarget{target} in lecture note tag.
% \LN{proofs} prints "Lecture Note 2", linked to that lecture note.
\newcommand{\currentLN}{0}   % in a standalone note, its number (set by \doTitle)
\newcommand{\notePDF}[1]{\siteURL notes/\LNfile{#1}.pdf}
\newcommand{\noteLink}[3]{%
  \ifnum\LNnumber{#1}=\currentLN\relax\hyperlink{#2}{#3}%
  \else\href{\notePDF{#1}\#nameddest=#2}{#3}\fi}
% \bookLink{target}{text} links text to \hypertarget{target} in the book's
% appendices (from a standalone note, in the published book).
\DeclareRobustCommand{\bookLink}[2]{%
  \ifdefined\BookMode\hyperlink{#1}{#2}%
  \else\href{\siteURL lecture_notes_on_algorithms.pdf\#nameddest=#1}{#2}\fi}
\DeclareRobustCommand{\LN}[1]{%
  \ifnum\LNnumber{#1}=\currentLN\relax Lecture Note~\LNnumber{#1}%
  \else\href{\notePDF{#1}}{Lecture Note~\LNnumber{#1}}\fi}
\pdfstringdefDisableCommands{\def\LN#1{Lecture Note \LNnumber{#1}}}

% Homeworks refer to the notes' sections, exercises and lemmas by \ref: each
% homework's main file says \lectureNoteRefs after %%%% DON'T EDIT THIS FILE %%%%

\usepackage{xifthen}% \isempty
\usepackage{etoolbox}% \patchcmd
% \usepackage{stringstrings}% \substring

\usepackage{amsmath, amsthm, amssymb}
\usepackage[margin=1in]{geometry}
\usepackage{xspace}
\usepackage{graphicx}
\usepackage[normalem]{ulem}
\usepackage{adjustbox}
\usepackage{xcolor}
\usepackage{changepage}
\usepackage{framed}
\usepackage{multicol}
\usepackage{enumitem}
% \usepackage{listings} % lstlistings environment

\usepackage{makecmds}           %\provideenvironment

\usepackage{xparse,environ}
\usepackage{letltxmacro}       % \LetLtxMacro
\usepackage{wrapfig}
\usepackage{tikz}

\usetikzlibrary{shapes}
\usetikzlibrary{positioning}
\usetikzlibrary{shapes.misc}
\usetikzlibrary{arrows}

\newtheorem{lemma}{Lemma}
\newtheorem{definition}{Definition}
\newtheorem{theorem}{Theorem}
\newtheorem{corollary}{Corollary}
\newtheorem{claim}{Claim}
\newtheorem{observation}{Observation}
\newtheorem{conjecture}{Conjecture}

% A lemma or theorem restated later with its original number (set with
% \setcounter) goes in a restatement environment, which gives it its own link
% anchor (otherwise it would duplicate the original's).
\newenvironment{restatement}
  {\def\theHlemma{restated.\thelemma}\def\theHtheorem{restated.\thetheorem}}{}

\newenvironment{myframed}
{\smallskip\par\noindent\begin{minipage}{\textwidth}\begin{framed}\vspace*{-5pt}}{\vspace*{-5pt}\end{framed}\end{minipage}\smallskip}


\definecolor{darkblue}{RGB}{0,0,200}

 
\let\oldParagraph\paragraph
\providecommand{\gradescopeColor}{\color{black}}
\renewcommand{\paragraph}[1]{\oldParagraph{\gradescopeColor #1}}
\providecommand{\gradescopeInit}{}

\providecommand{\BLUE}[1]{{\color{darkblue} #1}}
\providecommand{\REPLACEME}{\textcolor{red}{\ensuremath{\star \cdots \star}}\xspace}
\provideenvironment{blue}{\color{darkblue}}{}

\newcommand{\dist}[2]{\ensuremath{\operatorname{\sf dist}(#1, #2)}}
\newcommand{\depth}[2]{\ensuremath{\operatorname{\sf depth}_{#1}(#2)}}
\newcommand{\cost}[1]{\ensuremath{\operatorname{\sf cost}(#1)}}
\newcommand{\weight}[1]{\ensuremath{\operatorname{\sf wt}(#1)}}
% \newcommand{\val}[1]{\ensuremath{\operatorname{\sf val}(#1)}}
\newcommand{\eps}{\varepsilon}
\newcommand{\capac}{\textsf{cap}}
\newcommand{\capOf}[1]{\ensuremath{\operatorname{\sf cap}(#1)}}
\newcommand{\rescap}[1]{\ensuremath{\operatorname{\sf cap}_f(#1)}}
\newcommand{\into}[1]{\ensuremath{\operatorname{\sf in}(#1)}}
\newcommand{\out}[1]{\ensuremath{\operatorname{\sf out}(#1)}}
\newenvironment{LP}[1]{
  \begin{array}[t]{@{}l@{}}
    \displaystyle #1 \\
    \begin{aligned}\\[-10pt]}{\end{aligned}\end{array}}

\newcommand{\opt}[1]{\ensuremath{\operatorname{\sf opt}(#1)}}
\newcommand{\Opt}[1]{\ensuremath{\operatorname{\sf Opt}(#1)}}
\newcommand{\half}{\frac{1}{2}}

\newcommand{\E}{\operatorname{E}}
\newcommand{\giv}{\,|\,}
\newcommand{\Z}{\mathbb{Z}}
\newcommand{\Zp}{\mathbb{Z}_{\ge 0}}
\newcommand{\R}{\mathbb{R}}
% \newcommand{\Rp}{\mathbb{R}_{\ge 0}}
\newcommand{\Rp}{\mathbb{R}_{\tiny +}}
\newcommand{\N}{\mathbb{N}}
\newcommand{\Np}{\mathbb{N}_{\ge 0}}
\newcommand{\calS}{{\cal S}}

\newcommand{\OP}[1]{\text{\small\sc #1}\xspace}
\newcommand{\Insert}{\OP{Insert}}
\newcommand{\FindMin}{\OP{Find-Min}}
\newcommand{\DeleteMin}{\OP{Delete-Min}}
\newcommand{\DecreaseKey}{\OP{Decrease}}
\newcommand{\Merge}{\OP{Merge}}
\newcommand{\Cut}{\OP{Cut}}
% \newcommand{\prob}[1]{\text{\sc #1}}
% a problem name; in running text, lines may break inside it
\DeclareRobustCommand{\prob}[1]{\ifmmode\text{#1}\else#1\fi}
\newcommand{\encode}[1]{\langle #1 \rangle}

\newcommand{\best}[1]{\text{\sf best}(#1)}
\newcommand{\worst}[1]{\text{\sf worst}(#1)}
\newcommand{\floor}[1]{\lfloor #1 \rfloor}
\newcommand{\ceil}[1]{\lceil #1 \rceil}

\setlength{\parskip}{2pt}
\setlength{\parindent}{1em}

% \setlength\fboxsep{0pt}

%%% longFormProof, shortFormProof, block environments
%%% \step

\newlist{steps}{enumerate}{7}
\setlist[steps]{wide,
  topsep=1pt,
  parsep=1pt,
  partopsep=1pt,
  itemsep=0pt,
  align=left,
  labelindent=0pt,
  itemindent=0pt, 
}
\setlist[steps,1]{
  label=\arabic*.,
  ref=\arabic*
}
\setlist[steps,2]{
  label=\arabic{stepsi}.\arabic*.,
  ref=\arabic{stepsi}.\arabic*
}
\setlist[steps,3]{
  label=\arabic{stepsi}.\arabic{stepsii}.\arabic*.,
  ref=\arabic{stepsi}.\arabic{stepsii}.\arabic*
}
\setlist[steps,4]{
  label=\arabic{stepsi}.\arabic{stepsii}.\arabic{stepsiii}.\arabic*.,
  ref=\arabic{stepsi}.\arabic{stepsii}.\arabic{stepsiii}.\arabic*
}
\setlist[steps,5]{
  label=\arabic{stepsi}.\arabic{stepsii}.\arabic{stepsiii}.\arabic{stepsiv}.\arabic*.,
  ref=\arabic{stepsi}.\arabic{stepsii}.\arabic{stepsiii}.\arabic{stepsiv}.\arabic*
}
\setlist[steps,6]{
  label=\arabic{stepsi}.\arabic{stepsii}.\arabic{stepsiii}.\arabic{stepsiv}.\arabic{stepsv}.\arabic*.,
  ref=\arabic{stepsi}.\arabic{stepsii}.\arabic{stepsiii}.\arabic{stepsiv}.\arabic{stepsv}.\arabic*
}
\setlist[steps,7]{
  label=\arabic{stepsi}.\arabic{stepsii}.\arabic{stepsiii}.\arabic{stepsiv}. \arabic{stepsv}. \arabic{stepsvi}.\arabic*.,
  ref=\arabic{stepsi}.\arabic{stepsii}.\arabic{stepsiii}.\arabic{stepsiv}.\arabic{stepsv}. \arabic{stepsvi}.\arabic*
}

\newcommand{\step}[1][]{\item\maybeLabel{#1}}

\makeatletter
% \skipitem steps the list counter itself (for block labels).  Give each one
% its own hyperref link target, since step numbers repeat from proof to proof.
\newcounter{skipitem}
\newcommand{\skipitem}{\stepcounter{skipitem}\refstepcounter{\@enumctr}}
\renewcommand{\theHstepsi}{skipitem.\arabic{skipitem}}
\renewcommand{\theHstepsii}{skipitem.\arabic{skipitem}}
\renewcommand{\theHstepsiii}{skipitem.\arabic{skipitem}}
\renewcommand{\theHstepsiv}{skipitem.\arabic{skipitem}}
\renewcommand{\theHstepsv}{skipitem.\arabic{skipitem}}
\renewcommand{\theHstepsvi}{skipitem.\arabic{skipitem}}
\renewcommand{\theHstepsvii}{skipitem.\arabic{skipitem}}
\makeatother

\newcommand{\maybeLabel}[1]{\ifthenelse{\isempty{#1}}{}{\label{#1}}}

\newenvironment{longFormProof}[1][Proof (long form).]
{\begin{proof}[#1]~\par\begin{steps}}{\end{steps}\vspace*{-3ex}\end{proof}}

\newenvironment{shortFormProof}[1][Proof (short form).]
{\begin{proof}[#1]}{\end{proof}}

\newenvironment{longFormSubProof}[1][Proof.]
{\begin{adjustwidth}{1em}{0em}\begin{proof}[#1]~\par\vspace*{-2pt}\skipitem\begin{steps}}
{\end{steps}\vspace*{-1ex}\end{proof}\end{adjustwidth}\smallskip}

\newenvironment{shortFormSubProof}[1][Proof.]
{\begin{proof}[#1]}
{\end{proof}\vspace*{-3pt}}

\newenvironment{block}[1][]
{\skipitem\maybeLabel{#1}\item[] \begin{steps}}{\end{steps}}

\newcommand{\comment}[1]{\hfill{\footnotesize\emph{#1}}}
\newcommand{\lineacross}{\par\vspace*{-0.7\baselineskip}\noindent\hrulefill\par}

%%% problem

% \newcounter{problem}
% \newcommand{\problem}[1][]{}
% {\ifthenelse{\isempty{#1}}{\refstepcounter{problem}}{}%
%   \paragraph{Problem \ifthenelse{\isempty{#1}}
%     {\bf\arabic{problem}.}
%     {#1}} }

\newcommand{\problemNumbered}[1]{
  \setcounter{problem}{#1}%
  \addtocounter{problem}{-1}%
  \problem
}
\newenvironment{problemTheorem}{%
  \setcounter{theorem}{\value{problem}}%
  \addtocounter{theorem}{-1}%
  \begin{theorem}[Problem \arabic{problem}]
  }{\end{theorem}}

\newcounter{problem}
\newcommand{\problem}[1][]
{\ifthenelse{\isempty{#1}}
  {\refstepcounter{problem}\paragraph{Problem~\arabic{problem}.}}
  {\paragraph{Problem~#1.}}\refreshHeader} 

\newcounter{exercise}
\newcommand{\exercise}[1][]
{\ifthenelse{\isempty{#1}}
  {\refstepcounter{exercise}\paragraph{Exercise~\theexercise.}}
  {\paragraph{Exercise~#1.}}}

\newcommand{\solution}{\paragraph{Solution.}} 

\newenvironment{mylist}[1][]
{\begin{list}{#1}{
      \setlength{\itemsep}{-1pt}
      \setlength{\labelwidth}{0pt}
      \setlength{\leftmargin}{1.25em}
    }}
{\end{list}}

\newcommand{\LNTitle}{\textcolor{red}{UNDEFINED LNTitle}}

% In a standalone lecture note's table of contents, leave room for numbers
% like 7.10 (in bold) and 7.10.2.
\makeatletter
\newcommand{\LNtocWidths}{%
  \patchcmd{\l@section}{1.5em}{2.8em}{}{\PackageWarning{macros}{Could not widen section numbers in the contents}}%
  \renewcommand*{\l@subsection}{\@dottedtocline{2}{2.8em}{3.2em}}%
  \renewcommand*{\l@subsubsection}{\@dottedtocline{3}{6.0em}{4.1em}}}
\makeatother

\newcommand{\doTitle}[3]{
  \pagenumbering{roman}
  \renewcommand{\LNTitle}{Lecture Note #1, #3}
  \renewcommand{\currentLN}{#1}
  \codeListingstrue
  % number sections, theorems etc., exercises, figures, tables and equations
  % within the lecture note (e.g. Section 8.2, Lemma 8.6), as in the book
  \renewcommand{\thesection}{#1.\arabic{section}}
  \renewcommand{\thelemma}{#1.\arabic{lemma}}
  \renewcommand{\thedefinition}{#1.\arabic{definition}}
  \renewcommand{\thetheorem}{#1.\arabic{theorem}}
  \renewcommand{\thecorollary}{#1.\arabic{corollary}}
  \renewcommand{\theclaim}{#1.\arabic{claim}}
  \renewcommand{\theobservation}{#1.\arabic{observation}}
  \renewcommand{\theconjecture}{#1.\arabic{conjecture}}
  \renewcommand{\theexercise}{#1.\arabic{exercise}}
  \renewcommand{\thefigure}{#1.\arabic{figure}}
  \renewcommand{\thetable}{#1.\arabic{table}}
  \renewcommand{\theequation}{#1.\arabic{equation}}
  \LNtocWidths
  \title{Lecture Note #1\\ #2}
  % a note's \LNabstract, if it defines one before \doTitle, goes right below the title
  \date{\ifdefined\LNabstract\parbox{0.85\textwidth}{\normalsize\normalfont\LNabstract}\\[3ex]\fi
    draft of \today, by \href{\authorURL}{N.~Young}\\[2ex] {\small\copyrightNotice}}
  \maketitle
  {\small
    \hypertarget{TOC}{}
    \tableofcontents
  }
  \thispagestyle{empty}
  % the text starts on page 1, after the title and contents (numbered i, ii, ...)
  \clearpage
  \pagenumbering{arabic}
}

\def\clap#1{\hbox to 0pt{\hss#1\hss}}

\newcommand{\refreshHeader}{}
\newcommand{\header}{}

\newcommand{\doHeader}[1]{      % for lecture notes
  \renewcommand{\header}{%
    \clap{\raisebox{0pt}{\protect\hyperlink{TOC}{$\uparrow$\scriptsize\,\textsc{contents}\hspace{1.1in}}}}%
    {\footnotesize \LNTitle{} / Section \thesection, #1}} %, #1, draft of \today}}
  \markboth{\header}{\header}
}

\newcommand{\SID}{}
\newcommand{\setAuthor}[1]{\renewcommand{\author}{#1}}
\newcommand{\setSID}[1]{\renewcommand{\SID}{#1}}

\newcommand{\doHwHeader}[1]{                                         % for homeworks
  \renewcommand{\refreshHeader}{
    \renewcommand{\header}{%
      {\footnotesize \classSection #1 / Problem~\arabic{problem} / draft of \today \hfill \author\ifGradescope~(\SID)\fi~~}}
    \markboth{\header}{\header}
  }
  \renewcommand{\header}{%
    {\footnotesize \classSection #1 / Instructions / draft of \today \hfill \author\ifGradescope~(\SID)\fi~~}}
  \markboth{\header}{\header}
}

% Homework settings, for instructors using the homeworks in a course:
% \classSection is the course name shown in each homework's running head
% (e.g. \renewcommand{\classSection}{CS 3000\xspace}), and \Gradescopetrue turns on
% the instructions for submitting to Gradescope (in its fixed-length mode),
% the hints saying on which page of the PDF each answer should be, and the
% student ID in the running head.
\newcommand{\classSection}{}
\newif\ifGradescope
\Gradescopefalse
\newcommand{\GradescopeComment}[1]{\ifGradescope\comment{#1}\fi}

\pagestyle{myheadings}
\addtolength{\headsep}{0.2in}
\addtolength{\topmargin}{-0.3in}
\addtolength{\textheight}{0.4in}

%%%%%%%%%%%%%%%%%%%%%%%%%%%%%%%%%%%%%%%%%%%%%%%%
%%%%%%%%%%%%%%%%%%%%%%%%%%%%%%%%%%%%%%%%%%%%%%%%

\newcommand{\MF}{\href
  {https://mfleck.cs.illinois.edu/building-blocks/}
  {MF}\xspace}

\newcommand{\KWSlides}{\href
  {https://www.cs.princeton.edu/~wayne/kleinberg-tardos/}
  {KW slides}\xspace}

\newcommand{\LLM}{\href
  {https://courses.csail.mit.edu/6.042/spring18/mcs.pdf}
  {LLM}\xspace}

% \DPV, \KT, \CLRS: textbooks listed (without links) in the table of sources in
% the Preliminaries note; each prints as a link to its entry in that table.
\DeclareRobustCommand{\DPV}{\noteLink{prelim}{source.DPV}{DPV}\xspace}
\DeclareRobustCommand{\KT}{\noteLink{prelim}{source.KT}{KT}\xspace}
\DeclareRobustCommand{\CLRS}{\noteLink{prelim}{source.CLRS}{CLRS}\xspace}
\AtBeginDocument{\pdfstringdefDisableCommands{\def\DPV{DPV }\def\KT{KT }\def\CLRS{CLRS }}}

\newcommand{\JE}{\href
  {https://jeffe.cs.illinois.edu/teaching/algorithms/}
  {JE}\xspace}

\newcommand{\JEGreedy}{\href
  {https://jeffe.cs.illinois.edu/teaching/algorithms/book/04-greedy.pdf}
  {JE Chapter 4}\xspace}

\newcommand{\BBCode}{\href
{https://northeastern.blackboard.com/webapps/cmsmain/webui/courses/CS3000.MERGED.202010/code?action=frameset&subaction=view&uniq=-jj2max&course_id=_2597392_1&mask=\%2Fcourses\%2FCS3000.MERGED.202010}
  {code}\xspace}

% The Python code (src/code).  Each lecture note that has code ends with a
% section listing its files: \codeSection{path, path, ...}.  In the text,
% \pycode{greedy_algorithms/Dijkstra.py} prints the path as a link to its listing
% (in the homeworks, which have no listings, to the file on GitHub), and
% \pycode[text]{...} uses the given text instead.  \pythonCode{path, ...} is the
% standard sentence ending a section whose algorithms have code:
% "Python code: path, ... (Section N.k)."
\newcommand{\codeFolder}{https://github.com/nealeyoung/algs101/tree/main/src/code}
\newcommand{\codeRepo}{https://github.com/nealeyoung/algs101/blob/main/src/code}
\newif\ifcodeListings   % set by \doTitle: this document has the listings
\newcommand{\githubCode}[2][]{\href{\codeRepo/#2}{\ifx\relax#1\relax\nolinkurl{#2}\else#1\fi}}
\newcommand{\pycode}[2][]{%
  \ifcodeListings\hyperref[code: #2]{\ifx\relax#1\relax\nolinkurl{#2}\else#1\fi}%
  \else\githubCode[#1]{#2}\fi}
\makeatletter
\newcommand{\pythonCode}[1]{%
  \def\pythonCode@first{}\def\pythonCode@list{}%
  \forcsvlist{\pythonCode@add}{#1}%
  Python code: \pythonCode@list\ (Section~\ref{code: \pythonCode@first}).}
\newcommand{\pythonCode@add}[1]{%
  \ifx\pythonCode@first\@empty\def\pythonCode@first{#1}\appto\pythonCode@list{\pycode{#1}}%
  \else\appto\pythonCode@list{, \pycode{#1}}\fi}
\makeatother

% Code listings that copy and paste exactly (in PDF viewers that honor the
% ToUnicode maps, e.g. MuPDF-based ones; PDFKit viewers such as Preview and
% Skim drop leading spaces, so each listing's heading links to the file on
% GitHub).  Spaces are typeset as white visible-space glyphs that copy as
% spaces, and hyphens and asterisks as the T1 glyphs, so they copy as themselves.
\usepackage[T1,OT1]{fontenc}   % T1 for the code listings only; OT1 stays the default
\usepackage{listings}
\pdfgentounicode=1
\input glyphtounicode
\pdfglyphtounicode{uni2423}{0020}
\makeatletter
\def\lst@visiblespace{\textcolor{white}{\char32}}
% put an invisible space on each blank line, so that blank lines copy too
% (listings' OnEmptyLine hook runs before the line break, so patch the loop)
\def\lst@DoNewLines{%
    \@whilenum\lst@newlines>\@ne \do {\lst@NewLine \lst@visiblespace}%
    \ifnum\lst@newlines>\z@ \lst@NewLine \fi}
\makeatother
\lstset{language=Python,
  basicstyle=\fontencoding{T1}\fontfamily{lmtt}\selectfont\footnotesize,
  upquote=true, showstringspaces=false, columns=fullflexible, keepspaces=true,
  showspaces=true, literate={-}{{\char45}}1 {*}{{\char42}}1,
  frame=single, framesep=1ex, rulecolor=\color{gray}}
% the listings' paths are relative to the lecture note's folder (in the book, to src/)
\ifdefined\BookMode\newcommand{\codeDir}{code/}\else\newcommand{\codeDir}{../code/}\fi
% \codefile{path}: a listing of src/code/path, headed by a link to it on GitHub
\newcommand{\codefile}[1]{%
  \subsection*{\githubCode{#1}}\label{code: #1}%
  \lstinputlisting{\codeDir#1}}
% \codeSection{path, path, ...}: the lecture note's last section, listing its code
\newcommand{\codeSection}[1]{%
  \newpage
  \section{Python code}\label{\LNtag{\currentLN}: sec: code}\doHeader{Python code}
  This section lists the Python code for this lecture note.
  Click on a file's name to open the file on GitHub.
  Copying code from GitHub is more reliable than copying it from this PDF,
  since some PDF viewers drop the leading spaces (the indentation) of copied lines.
  \forcsvlist{\codefile}{#1}}
% the copyright notice, for title pages and the homework first pages
\newcommand{\licenseURL}{https://github.com/nealeyoung/algs101/blob/main/LICENSE}
\newcommand{\authorURL}{https://www.cs.ucr.edu/\string~neal}
\newcommand{\copyrightNotice}{%
  \copyright\ 2018--2026 \href{\authorURL}{Neal E. Young}.
  Text: CC BY-NC-SA 4.0; code: MIT; except where noted.
  See \url{\licenseURL}.}
% the published site: PDFs of the notes, and the homework assignments and templates
\newcommand{\siteURL}{https://nealeyoung.github.io/algs101/}

\newcommand{\blackboard}{\href
  {https://northeastern.blackboard.com/webapps/blackboard/execute/launcher?type=Course&id=_2597392_1&url=}
  {Blackboard}\xspace}

\newcommand{\JUSlides}{\href
  {http://www.khoury.neu.edu/home/jullman/cs3000f18/schedule.html}
  {Prof.~Ullman's slides}\xspace}

% a URL on its own line, in three layouts
\newcommand{\URL}[1]{\par-- \parbox[t]{0.8\textwidth}{\url{#1}} \medskip\par}
\newcommand{\URLsmall}[1]{\par-- \parbox[t]{0.95\textwidth}{\small\url{#1}} \medskip\par}
\newcommand{\URLplain}[1]{\par\parbox[t]{\textwidth}{\url{#1}}\par}

%%%%%%%%%%%%%%%%%%%%%%%%%%%%%%%%%%%%%%%%%%%%%%%%
%%%%%%%%%%%%%%%%%%%%%%%%%%%%%%%%%%%%%%%%%%%%%%%%

\newcommand{\continued}{\vfill\centerline{\emph{(continued)}}\vfill\newpage}

% load hyperref last to avoid conflicts; long URLs may also break after hyphens
\PassOptionsToPackage{hyphens}{url}
\ifdefined\BookMode\else\usepackage{xr-hyper}\fi   % for \lectureNoteRefs (homeworks), below
\usepackage{hyperref}
\hypersetup{
    colorlinks,
    linkcolor={blue!50!black},
    citecolor={blue!50!black},
    urlcolor={blue!80!black}
}

\providecommand{\answerDirectory}{.}
\newcommand{\inputAnswer}[1]{\refreshHeader\input{\answerDirectory/#1}}
\newcommand{\inputSection}[1]{\input{#1}}

% Refer to a lecture note by its tag (the prefix of its labels), e.g.
% \LN{proofs} prints "Lecture Note 2" (in the book, a link to that chapter).  The table below is the only place the numbers appear.
\makeatletter
\newcommand{\defineLN}[3]{\@namedef{LNnumber@#2}{#1}\@namedef{LNtag@#1}{#2}\@namedef{LNfile@#2}{#3}}
\newcommand{\LNnumber}[1]{%
  \ifcsname LNnumber@#1\endcsname\csname LNnumber@#1\endcsname
  \else\PackageError{macros}{Unknown lecture-note tag `#1'}{}\fi}
\newcommand{\LNtag}[1]{\csname LNtag@#1\endcsname}
\newcommand{\LNfile}[1]{\csname LNfile@#1\endcsname}
\makeatother
\defineLN{1}{prelim}{1_preliminaries}
\defineLN{2}{proofs}{2_proofs}
\defineLN{3}{matching}{3_stable_matching}
\defineLN{4}{dc}{4_divide_and_conquer}
\defineLN{5}{greedy}{5_greedy}
\defineLN{6}{graphs}{6_graph_traversal}
\defineLN{7}{dp}{7_dynamic_programming}
\defineLN{8}{flow}{8_flow_reductions_LP}
\defineLN{9}{np}{9_NP}

% The text is the same in the book and in the standalone lecture notes; only
% links differ.  In the book, links to other lecture notes stay in the book.
% In a standalone note, they open the other note's published PDF.
% \noteLink{tag}{target}{text} links text to \hypertarget{target} in lecture note tag.
% \LN{proofs} prints "Lecture Note 2", linked to that lecture note.
\newcommand{\currentLN}{0}   % in a standalone note, its number (set by \doTitle)
\newcommand{\notePDF}[1]{\siteURL notes/\LNfile{#1}.pdf}
\newcommand{\noteLink}[3]{%
  \ifnum\LNnumber{#1}=\currentLN\relax\hyperlink{#2}{#3}%
  \else\href{\notePDF{#1}\#nameddest=#2}{#3}\fi}
% \bookLink{target}{text} links text to \hypertarget{target} in the book's
% appendices (from a standalone note, in the published book).
\DeclareRobustCommand{\bookLink}[2]{%
  \ifdefined\BookMode\hyperlink{#1}{#2}%
  \else\href{\siteURL lecture_notes_on_algorithms.pdf\#nameddest=#1}{#2}\fi}
\DeclareRobustCommand{\LN}[1]{%
  \ifnum\LNnumber{#1}=\currentLN\relax Lecture Note~\LNnumber{#1}%
  \else\href{\notePDF{#1}}{Lecture Note~\LNnumber{#1}}\fi}
\pdfstringdefDisableCommands{\def\LN#1{Lecture Note \LNnumber{#1}}}

% Homeworks refer to the notes' sections, exercises and lemmas by \ref: each
% homework's main file says \lectureNoteRefs after %%%% DON'T EDIT THIS FILE %%%%

\usepackage{xifthen}% \isempty
\usepackage{etoolbox}% \patchcmd
% \usepackage{stringstrings}% \substring

\usepackage{amsmath, amsthm, amssymb}
\usepackage[margin=1in]{geometry}
\usepackage{xspace}
\usepackage{graphicx}
\usepackage[normalem]{ulem}
\usepackage{adjustbox}
\usepackage{xcolor}
\usepackage{changepage}
\usepackage{framed}
\usepackage{multicol}
\usepackage{enumitem}
% \usepackage{listings} % lstlistings environment

\usepackage{makecmds}           %\provideenvironment

\usepackage{xparse,environ}
\usepackage{letltxmacro}       % \LetLtxMacro
\usepackage{wrapfig}
\usepackage{tikz}

\usetikzlibrary{shapes}
\usetikzlibrary{positioning}
\usetikzlibrary{shapes.misc}
\usetikzlibrary{arrows}

\newtheorem{lemma}{Lemma}
\newtheorem{definition}{Definition}
\newtheorem{theorem}{Theorem}
\newtheorem{corollary}{Corollary}
\newtheorem{claim}{Claim}
\newtheorem{observation}{Observation}
\newtheorem{conjecture}{Conjecture}

% A lemma or theorem restated later with its original number (set with
% \setcounter) goes in a restatement environment, which gives it its own link
% anchor (otherwise it would duplicate the original's).
\newenvironment{restatement}
  {\def\theHlemma{restated.\thelemma}\def\theHtheorem{restated.\thetheorem}}{}

\newenvironment{myframed}
{\smallskip\par\noindent\begin{minipage}{\textwidth}\begin{framed}\vspace*{-5pt}}{\vspace*{-5pt}\end{framed}\end{minipage}\smallskip}


\definecolor{darkblue}{RGB}{0,0,200}

 
\let\oldParagraph\paragraph
\providecommand{\gradescopeColor}{\color{black}}
\renewcommand{\paragraph}[1]{\oldParagraph{\gradescopeColor #1}}
\providecommand{\gradescopeInit}{}

\providecommand{\BLUE}[1]{{\color{darkblue} #1}}
\providecommand{\REPLACEME}{\textcolor{red}{\ensuremath{\star \cdots \star}}\xspace}
\provideenvironment{blue}{\color{darkblue}}{}

\newcommand{\dist}[2]{\ensuremath{\operatorname{\sf dist}(#1, #2)}}
\newcommand{\depth}[2]{\ensuremath{\operatorname{\sf depth}_{#1}(#2)}}
\newcommand{\cost}[1]{\ensuremath{\operatorname{\sf cost}(#1)}}
\newcommand{\weight}[1]{\ensuremath{\operatorname{\sf wt}(#1)}}
% \newcommand{\val}[1]{\ensuremath{\operatorname{\sf val}(#1)}}
\newcommand{\eps}{\varepsilon}
\newcommand{\capac}{\textsf{cap}}
\newcommand{\capOf}[1]{\ensuremath{\operatorname{\sf cap}(#1)}}
\newcommand{\rescap}[1]{\ensuremath{\operatorname{\sf cap}_f(#1)}}
\newcommand{\into}[1]{\ensuremath{\operatorname{\sf in}(#1)}}
\newcommand{\out}[1]{\ensuremath{\operatorname{\sf out}(#1)}}
\newenvironment{LP}[1]{
  \begin{array}[t]{@{}l@{}}
    \displaystyle #1 \\
    \begin{aligned}\\[-10pt]}{\end{aligned}\end{array}}

\newcommand{\opt}[1]{\ensuremath{\operatorname{\sf opt}(#1)}}
\newcommand{\Opt}[1]{\ensuremath{\operatorname{\sf Opt}(#1)}}
\newcommand{\half}{\frac{1}{2}}

\newcommand{\E}{\operatorname{E}}
\newcommand{\giv}{\,|\,}
\newcommand{\Z}{\mathbb{Z}}
\newcommand{\Zp}{\mathbb{Z}_{\ge 0}}
\newcommand{\R}{\mathbb{R}}
% \newcommand{\Rp}{\mathbb{R}_{\ge 0}}
\newcommand{\Rp}{\mathbb{R}_{\tiny +}}
\newcommand{\N}{\mathbb{N}}
\newcommand{\Np}{\mathbb{N}_{\ge 0}}
\newcommand{\calS}{{\cal S}}

\newcommand{\OP}[1]{\text{\small\sc #1}\xspace}
\newcommand{\Insert}{\OP{Insert}}
\newcommand{\FindMin}{\OP{Find-Min}}
\newcommand{\DeleteMin}{\OP{Delete-Min}}
\newcommand{\DecreaseKey}{\OP{Decrease}}
\newcommand{\Merge}{\OP{Merge}}
\newcommand{\Cut}{\OP{Cut}}
% \newcommand{\prob}[1]{\text{\sc #1}}
% a problem name; in running text, lines may break inside it
\DeclareRobustCommand{\prob}[1]{\ifmmode\text{#1}\else#1\fi}
\newcommand{\encode}[1]{\langle #1 \rangle}

\newcommand{\best}[1]{\text{\sf best}(#1)}
\newcommand{\worst}[1]{\text{\sf worst}(#1)}
\newcommand{\floor}[1]{\lfloor #1 \rfloor}
\newcommand{\ceil}[1]{\lceil #1 \rceil}

\setlength{\parskip}{2pt}
\setlength{\parindent}{1em}

% \setlength\fboxsep{0pt}

%%% longFormProof, shortFormProof, block environments
%%% \step

\newlist{steps}{enumerate}{7}
\setlist[steps]{wide,
  topsep=1pt,
  parsep=1pt,
  partopsep=1pt,
  itemsep=0pt,
  align=left,
  labelindent=0pt,
  itemindent=0pt, 
}
\setlist[steps,1]{
  label=\arabic*.,
  ref=\arabic*
}
\setlist[steps,2]{
  label=\arabic{stepsi}.\arabic*.,
  ref=\arabic{stepsi}.\arabic*
}
\setlist[steps,3]{
  label=\arabic{stepsi}.\arabic{stepsii}.\arabic*.,
  ref=\arabic{stepsi}.\arabic{stepsii}.\arabic*
}
\setlist[steps,4]{
  label=\arabic{stepsi}.\arabic{stepsii}.\arabic{stepsiii}.\arabic*.,
  ref=\arabic{stepsi}.\arabic{stepsii}.\arabic{stepsiii}.\arabic*
}
\setlist[steps,5]{
  label=\arabic{stepsi}.\arabic{stepsii}.\arabic{stepsiii}.\arabic{stepsiv}.\arabic*.,
  ref=\arabic{stepsi}.\arabic{stepsii}.\arabic{stepsiii}.\arabic{stepsiv}.\arabic*
}
\setlist[steps,6]{
  label=\arabic{stepsi}.\arabic{stepsii}.\arabic{stepsiii}.\arabic{stepsiv}.\arabic{stepsv}.\arabic*.,
  ref=\arabic{stepsi}.\arabic{stepsii}.\arabic{stepsiii}.\arabic{stepsiv}.\arabic{stepsv}.\arabic*
}
\setlist[steps,7]{
  label=\arabic{stepsi}.\arabic{stepsii}.\arabic{stepsiii}.\arabic{stepsiv}. \arabic{stepsv}. \arabic{stepsvi}.\arabic*.,
  ref=\arabic{stepsi}.\arabic{stepsii}.\arabic{stepsiii}.\arabic{stepsiv}.\arabic{stepsv}. \arabic{stepsvi}.\arabic*
}

\newcommand{\step}[1][]{\item\maybeLabel{#1}}

\makeatletter
% \skipitem steps the list counter itself (for block labels).  Give each one
% its own hyperref link target, since step numbers repeat from proof to proof.
\newcounter{skipitem}
\newcommand{\skipitem}{\stepcounter{skipitem}\refstepcounter{\@enumctr}}
\renewcommand{\theHstepsi}{skipitem.\arabic{skipitem}}
\renewcommand{\theHstepsii}{skipitem.\arabic{skipitem}}
\renewcommand{\theHstepsiii}{skipitem.\arabic{skipitem}}
\renewcommand{\theHstepsiv}{skipitem.\arabic{skipitem}}
\renewcommand{\theHstepsv}{skipitem.\arabic{skipitem}}
\renewcommand{\theHstepsvi}{skipitem.\arabic{skipitem}}
\renewcommand{\theHstepsvii}{skipitem.\arabic{skipitem}}
\makeatother

\newcommand{\maybeLabel}[1]{\ifthenelse{\isempty{#1}}{}{\label{#1}}}

\newenvironment{longFormProof}[1][Proof (long form).]
{\begin{proof}[#1]~\par\begin{steps}}{\end{steps}\vspace*{-3ex}\end{proof}}

\newenvironment{shortFormProof}[1][Proof (short form).]
{\begin{proof}[#1]}{\end{proof}}

\newenvironment{longFormSubProof}[1][Proof.]
{\begin{adjustwidth}{1em}{0em}\begin{proof}[#1]~\par\vspace*{-2pt}\skipitem\begin{steps}}
{\end{steps}\vspace*{-1ex}\end{proof}\end{adjustwidth}\smallskip}

\newenvironment{shortFormSubProof}[1][Proof.]
{\begin{proof}[#1]}
{\end{proof}\vspace*{-3pt}}

\newenvironment{block}[1][]
{\skipitem\maybeLabel{#1}\item[] \begin{steps}}{\end{steps}}

\newcommand{\comment}[1]{\hfill{\footnotesize\emph{#1}}}
\newcommand{\lineacross}{\par\vspace*{-0.7\baselineskip}\noindent\hrulefill\par}

%%% problem

% \newcounter{problem}
% \newcommand{\problem}[1][]{}
% {\ifthenelse{\isempty{#1}}{\refstepcounter{problem}}{}%
%   \paragraph{Problem \ifthenelse{\isempty{#1}}
%     {\bf\arabic{problem}.}
%     {#1}} }

\newcommand{\problemNumbered}[1]{
  \setcounter{problem}{#1}%
  \addtocounter{problem}{-1}%
  \problem
}
\newenvironment{problemTheorem}{%
  \setcounter{theorem}{\value{problem}}%
  \addtocounter{theorem}{-1}%
  \begin{theorem}[Problem \arabic{problem}]
  }{\end{theorem}}

\newcounter{problem}
\newcommand{\problem}[1][]
{\ifthenelse{\isempty{#1}}
  {\refstepcounter{problem}\paragraph{Problem~\arabic{problem}.}}
  {\paragraph{Problem~#1.}}\refreshHeader} 

\newcounter{exercise}
\newcommand{\exercise}[1][]
{\ifthenelse{\isempty{#1}}
  {\refstepcounter{exercise}\paragraph{Exercise~\theexercise.}}
  {\paragraph{Exercise~#1.}}}

\newcommand{\solution}{\paragraph{Solution.}} 

\newenvironment{mylist}[1][]
{\begin{list}{#1}{
      \setlength{\itemsep}{-1pt}
      \setlength{\labelwidth}{0pt}
      \setlength{\leftmargin}{1.25em}
    }}
{\end{list}}

\newcommand{\LNTitle}{\textcolor{red}{UNDEFINED LNTitle}}

% In a standalone lecture note's table of contents, leave room for numbers
% like 7.10 (in bold) and 7.10.2.
\makeatletter
\newcommand{\LNtocWidths}{%
  \patchcmd{\l@section}{1.5em}{2.8em}{}{\PackageWarning{macros}{Could not widen section numbers in the contents}}%
  \renewcommand*{\l@subsection}{\@dottedtocline{2}{2.8em}{3.2em}}%
  \renewcommand*{\l@subsubsection}{\@dottedtocline{3}{6.0em}{4.1em}}}
\makeatother

\newcommand{\doTitle}[3]{
  \pagenumbering{roman}
  \renewcommand{\LNTitle}{Lecture Note #1, #3}
  \renewcommand{\currentLN}{#1}
  \codeListingstrue
  % number sections, theorems etc., exercises, figures, tables and equations
  % within the lecture note (e.g. Section 8.2, Lemma 8.6), as in the book
  \renewcommand{\thesection}{#1.\arabic{section}}
  \renewcommand{\thelemma}{#1.\arabic{lemma}}
  \renewcommand{\thedefinition}{#1.\arabic{definition}}
  \renewcommand{\thetheorem}{#1.\arabic{theorem}}
  \renewcommand{\thecorollary}{#1.\arabic{corollary}}
  \renewcommand{\theclaim}{#1.\arabic{claim}}
  \renewcommand{\theobservation}{#1.\arabic{observation}}
  \renewcommand{\theconjecture}{#1.\arabic{conjecture}}
  \renewcommand{\theexercise}{#1.\arabic{exercise}}
  \renewcommand{\thefigure}{#1.\arabic{figure}}
  \renewcommand{\thetable}{#1.\arabic{table}}
  \renewcommand{\theequation}{#1.\arabic{equation}}
  \LNtocWidths
  \title{Lecture Note #1\\ #2}
  % a note's \LNabstract, if it defines one before \doTitle, goes right below the title
  \date{\ifdefined\LNabstract\parbox{0.85\textwidth}{\normalsize\normalfont\LNabstract}\\[3ex]\fi
    draft of \today, by \href{\authorURL}{N.~Young}\\[2ex] {\small\copyrightNotice}}
  \maketitle
  {\small
    \hypertarget{TOC}{}
    \tableofcontents
  }
  \thispagestyle{empty}
  % the text starts on page 1, after the title and contents (numbered i, ii, ...)
  \clearpage
  \pagenumbering{arabic}
}

\def\clap#1{\hbox to 0pt{\hss#1\hss}}

\newcommand{\refreshHeader}{}
\newcommand{\header}{}

\newcommand{\doHeader}[1]{      % for lecture notes
  \renewcommand{\header}{%
    \clap{\raisebox{0pt}{\protect\hyperlink{TOC}{$\uparrow$\scriptsize\,\textsc{contents}\hspace{1.1in}}}}%
    {\footnotesize \LNTitle{} / Section \thesection, #1}} %, #1, draft of \today}}
  \markboth{\header}{\header}
}

\newcommand{\SID}{}
\newcommand{\setAuthor}[1]{\renewcommand{\author}{#1}}
\newcommand{\setSID}[1]{\renewcommand{\SID}{#1}}

\newcommand{\doHwHeader}[1]{                                         % for homeworks
  \renewcommand{\refreshHeader}{
    \renewcommand{\header}{%
      {\footnotesize \classSection #1 / Problem~\arabic{problem} / draft of \today \hfill \author\ifGradescope~(\SID)\fi~~}}
    \markboth{\header}{\header}
  }
  \renewcommand{\header}{%
    {\footnotesize \classSection #1 / Instructions / draft of \today \hfill \author\ifGradescope~(\SID)\fi~~}}
  \markboth{\header}{\header}
}

% Homework settings, for instructors using the homeworks in a course:
% \classSection is the course name shown in each homework's running head
% (e.g. \renewcommand{\classSection}{CS 3000\xspace}), and \Gradescopetrue turns on
% the instructions for submitting to Gradescope (in its fixed-length mode),
% the hints saying on which page of the PDF each answer should be, and the
% student ID in the running head.
\newcommand{\classSection}{}
\newif\ifGradescope
\Gradescopefalse
\newcommand{\GradescopeComment}[1]{\ifGradescope\comment{#1}\fi}

\pagestyle{myheadings}
\addtolength{\headsep}{0.2in}
\addtolength{\topmargin}{-0.3in}
\addtolength{\textheight}{0.4in}

%%%%%%%%%%%%%%%%%%%%%%%%%%%%%%%%%%%%%%%%%%%%%%%%
%%%%%%%%%%%%%%%%%%%%%%%%%%%%%%%%%%%%%%%%%%%%%%%%

\newcommand{\MF}{\href
  {https://mfleck.cs.illinois.edu/building-blocks/}
  {MF}\xspace}

\newcommand{\KWSlides}{\href
  {https://www.cs.princeton.edu/~wayne/kleinberg-tardos/}
  {KW slides}\xspace}

\newcommand{\LLM}{\href
  {https://courses.csail.mit.edu/6.042/spring18/mcs.pdf}
  {LLM}\xspace}

% \DPV, \KT, \CLRS: textbooks listed (without links) in the table of sources in
% the Preliminaries note; each prints as a link to its entry in that table.
\DeclareRobustCommand{\DPV}{\noteLink{prelim}{source.DPV}{DPV}\xspace}
\DeclareRobustCommand{\KT}{\noteLink{prelim}{source.KT}{KT}\xspace}
\DeclareRobustCommand{\CLRS}{\noteLink{prelim}{source.CLRS}{CLRS}\xspace}
\AtBeginDocument{\pdfstringdefDisableCommands{\def\DPV{DPV }\def\KT{KT }\def\CLRS{CLRS }}}

\newcommand{\JE}{\href
  {https://jeffe.cs.illinois.edu/teaching/algorithms/}
  {JE}\xspace}

\newcommand{\JEGreedy}{\href
  {https://jeffe.cs.illinois.edu/teaching/algorithms/book/04-greedy.pdf}
  {JE Chapter 4}\xspace}

\newcommand{\BBCode}{\href
{https://northeastern.blackboard.com/webapps/cmsmain/webui/courses/CS3000.MERGED.202010/code?action=frameset&subaction=view&uniq=-jj2max&course_id=_2597392_1&mask=\%2Fcourses\%2FCS3000.MERGED.202010}
  {code}\xspace}

% The Python code (src/code).  Each lecture note that has code ends with a
% section listing its files: \codeSection{path, path, ...}.  In the text,
% \pycode{greedy_algorithms/Dijkstra.py} prints the path as a link to its listing
% (in the homeworks, which have no listings, to the file on GitHub), and
% \pycode[text]{...} uses the given text instead.  \pythonCode{path, ...} is the
% standard sentence ending a section whose algorithms have code:
% "Python code: path, ... (Section N.k)."
\newcommand{\codeFolder}{https://github.com/nealeyoung/algs101/tree/main/src/code}
\newcommand{\codeRepo}{https://github.com/nealeyoung/algs101/blob/main/src/code}
\newif\ifcodeListings   % set by \doTitle: this document has the listings
\newcommand{\githubCode}[2][]{\href{\codeRepo/#2}{\ifx\relax#1\relax\nolinkurl{#2}\else#1\fi}}
\newcommand{\pycode}[2][]{%
  \ifcodeListings\hyperref[code: #2]{\ifx\relax#1\relax\nolinkurl{#2}\else#1\fi}%
  \else\githubCode[#1]{#2}\fi}
\makeatletter
\newcommand{\pythonCode}[1]{%
  \def\pythonCode@first{}\def\pythonCode@list{}%
  \forcsvlist{\pythonCode@add}{#1}%
  Python code: \pythonCode@list\ (Section~\ref{code: \pythonCode@first}).}
\newcommand{\pythonCode@add}[1]{%
  \ifx\pythonCode@first\@empty\def\pythonCode@first{#1}\appto\pythonCode@list{\pycode{#1}}%
  \else\appto\pythonCode@list{, \pycode{#1}}\fi}
\makeatother

% Code listings that copy and paste exactly (in PDF viewers that honor the
% ToUnicode maps, e.g. MuPDF-based ones; PDFKit viewers such as Preview and
% Skim drop leading spaces, so each listing's heading links to the file on
% GitHub).  Spaces are typeset as white visible-space glyphs that copy as
% spaces, and hyphens and asterisks as the T1 glyphs, so they copy as themselves.
\usepackage[T1,OT1]{fontenc}   % T1 for the code listings only; OT1 stays the default
\usepackage{listings}
\pdfgentounicode=1
\input glyphtounicode
\pdfglyphtounicode{uni2423}{0020}
\makeatletter
\def\lst@visiblespace{\textcolor{white}{\char32}}
% put an invisible space on each blank line, so that blank lines copy too
% (listings' OnEmptyLine hook runs before the line break, so patch the loop)
\def\lst@DoNewLines{%
    \@whilenum\lst@newlines>\@ne \do {\lst@NewLine \lst@visiblespace}%
    \ifnum\lst@newlines>\z@ \lst@NewLine \fi}
\makeatother
\lstset{language=Python,
  basicstyle=\fontencoding{T1}\fontfamily{lmtt}\selectfont\footnotesize,
  upquote=true, showstringspaces=false, columns=fullflexible, keepspaces=true,
  showspaces=true, literate={-}{{\char45}}1 {*}{{\char42}}1,
  frame=single, framesep=1ex, rulecolor=\color{gray}}
% the listings' paths are relative to the lecture note's folder (in the book, to src/)
\ifdefined\BookMode\newcommand{\codeDir}{code/}\else\newcommand{\codeDir}{../code/}\fi
% \codefile{path}: a listing of src/code/path, headed by a link to it on GitHub
\newcommand{\codefile}[1]{%
  \subsection*{\githubCode{#1}}\label{code: #1}%
  \lstinputlisting{\codeDir#1}}
% \codeSection{path, path, ...}: the lecture note's last section, listing its code
\newcommand{\codeSection}[1]{%
  \newpage
  \section{Python code}\label{\LNtag{\currentLN}: sec: code}\doHeader{Python code}
  This section lists the Python code for this lecture note.
  Click on a file's name to open the file on GitHub.
  Copying code from GitHub is more reliable than copying it from this PDF,
  since some PDF viewers drop the leading spaces (the indentation) of copied lines.
  \forcsvlist{\codefile}{#1}}
% the copyright notice, for title pages and the homework first pages
\newcommand{\licenseURL}{https://github.com/nealeyoung/algs101/blob/main/LICENSE}
\newcommand{\authorURL}{https://www.cs.ucr.edu/\string~neal}
\newcommand{\copyrightNotice}{%
  \copyright\ 2018--2026 \href{\authorURL}{Neal E. Young}.
  Text: CC BY-NC-SA 4.0; code: MIT; except where noted.
  See \url{\licenseURL}.}
% the published site: PDFs of the notes, and the homework assignments and templates
\newcommand{\siteURL}{https://nealeyoung.github.io/algs101/}

\newcommand{\blackboard}{\href
  {https://northeastern.blackboard.com/webapps/blackboard/execute/launcher?type=Course&id=_2597392_1&url=}
  {Blackboard}\xspace}

\newcommand{\JUSlides}{\href
  {http://www.khoury.neu.edu/home/jullman/cs3000f18/schedule.html}
  {Prof.~Ullman's slides}\xspace}

% a URL on its own line, in three layouts
\newcommand{\URL}[1]{\par-- \parbox[t]{0.8\textwidth}{\url{#1}} \medskip\par}
\newcommand{\URLsmall}[1]{\par-- \parbox[t]{0.95\textwidth}{\small\url{#1}} \medskip\par}
\newcommand{\URLplain}[1]{\par\parbox[t]{\textwidth}{\url{#1}}\par}

%%%%%%%%%%%%%%%%%%%%%%%%%%%%%%%%%%%%%%%%%%%%%%%%
%%%%%%%%%%%%%%%%%%%%%%%%%%%%%%%%%%%%%%%%%%%%%%%%

\newcommand{\continued}{\vfill\centerline{\emph{(continued)}}\vfill\newpage}

% load hyperref last to avoid conflicts; long URLs may also break after hyphens
\PassOptionsToPackage{hyphens}{url}
\ifdefined\BookMode\else\usepackage{xr-hyper}\fi   % for \lectureNoteRefs (homeworks), below
\usepackage{hyperref}
\hypersetup{
    colorlinks,
    linkcolor={blue!50!black},
    citecolor={blue!50!black},
    urlcolor={blue!80!black}
}

\providecommand{\answerDirectory}{.}
\newcommand{\inputAnswer}[1]{\refreshHeader\input{\answerDirectory/#1}}
\newcommand{\inputSection}[1]{\input{#1}}

% Refer to a lecture note by its tag (the prefix of its labels), e.g.
% \LN{proofs} prints "Lecture Note 2" (in the book, a link to that chapter).  The table below is the only place the numbers appear.
\makeatletter
\newcommand{\defineLN}[3]{\@namedef{LNnumber@#2}{#1}\@namedef{LNtag@#1}{#2}\@namedef{LNfile@#2}{#3}}
\newcommand{\LNnumber}[1]{%
  \ifcsname LNnumber@#1\endcsname\csname LNnumber@#1\endcsname
  \else\PackageError{macros}{Unknown lecture-note tag `#1'}{}\fi}
\newcommand{\LNtag}[1]{\csname LNtag@#1\endcsname}
\newcommand{\LNfile}[1]{\csname LNfile@#1\endcsname}
\makeatother
\defineLN{1}{prelim}{1_preliminaries}
\defineLN{2}{proofs}{2_proofs}
\defineLN{3}{matching}{3_stable_matching}
\defineLN{4}{dc}{4_divide_and_conquer}
\defineLN{5}{greedy}{5_greedy}
\defineLN{6}{graphs}{6_graph_traversal}
\defineLN{7}{dp}{7_dynamic_programming}
\defineLN{8}{flow}{8_flow_reductions_LP}
\defineLN{9}{np}{9_NP}

% The text is the same in the book and in the standalone lecture notes; only
% links differ.  In the book, links to other lecture notes stay in the book.
% In a standalone note, they open the other note's published PDF.
% \noteLink{tag}{target}{text} links text to \hypertarget{target} in lecture note tag.
% \LN{proofs} prints "Lecture Note 2", linked to that lecture note.
\newcommand{\currentLN}{0}   % in a standalone note, its number (set by \doTitle)
\newcommand{\notePDF}[1]{\siteURL notes/\LNfile{#1}.pdf}
\newcommand{\noteLink}[3]{%
  \ifnum\LNnumber{#1}=\currentLN\relax\hyperlink{#2}{#3}%
  \else\href{\notePDF{#1}\#nameddest=#2}{#3}\fi}
% \bookLink{target}{text} links text to \hypertarget{target} in the book's
% appendices (from a standalone note, in the published book).
\DeclareRobustCommand{\bookLink}[2]{%
  \ifdefined\BookMode\hyperlink{#1}{#2}%
  \else\href{\siteURL lecture_notes_on_algorithms.pdf\#nameddest=#1}{#2}\fi}
\DeclareRobustCommand{\LN}[1]{%
  \ifnum\LNnumber{#1}=\currentLN\relax Lecture Note~\LNnumber{#1}%
  \else\href{\notePDF{#1}}{Lecture Note~\LNnumber{#1}}\fi}
\pdfstringdefDisableCommands{\def\LN#1{Lecture Note \LNnumber{#1}}}

% Homeworks refer to the notes' sections, exercises and lemmas by \ref: each
% homework's main file says \lectureNoteRefs after %%%% DON'T EDIT THIS FILE %%%%

\usepackage{xifthen}% \isempty
\usepackage{etoolbox}% \patchcmd
% \usepackage{stringstrings}% \substring

\usepackage{amsmath, amsthm, amssymb}
\usepackage[margin=1in]{geometry}
\usepackage{xspace}
\usepackage{graphicx}
\usepackage[normalem]{ulem}
\usepackage{adjustbox}
\usepackage{xcolor}
\usepackage{changepage}
\usepackage{framed}
\usepackage{multicol}
\usepackage{enumitem}
% \usepackage{listings} % lstlistings environment

\usepackage{makecmds}           %\provideenvironment

\usepackage{xparse,environ}
\usepackage{letltxmacro}       % \LetLtxMacro
\usepackage{wrapfig}
\usepackage{tikz}

\usetikzlibrary{shapes}
\usetikzlibrary{positioning}
\usetikzlibrary{shapes.misc}
\usetikzlibrary{arrows}

\newtheorem{lemma}{Lemma}
\newtheorem{definition}{Definition}
\newtheorem{theorem}{Theorem}
\newtheorem{corollary}{Corollary}
\newtheorem{claim}{Claim}
\newtheorem{observation}{Observation}
\newtheorem{conjecture}{Conjecture}

% A lemma or theorem restated later with its original number (set with
% \setcounter) goes in a restatement environment, which gives it its own link
% anchor (otherwise it would duplicate the original's).
\newenvironment{restatement}
  {\def\theHlemma{restated.\thelemma}\def\theHtheorem{restated.\thetheorem}}{}

\newenvironment{myframed}
{\smallskip\par\noindent\begin{minipage}{\textwidth}\begin{framed}\vspace*{-5pt}}{\vspace*{-5pt}\end{framed}\end{minipage}\smallskip}


\definecolor{darkblue}{RGB}{0,0,200}

 
\let\oldParagraph\paragraph
\providecommand{\gradescopeColor}{\color{black}}
\renewcommand{\paragraph}[1]{\oldParagraph{\gradescopeColor #1}}
\providecommand{\gradescopeInit}{}

\providecommand{\BLUE}[1]{{\color{darkblue} #1}}
\providecommand{\REPLACEME}{\textcolor{red}{\ensuremath{\star \cdots \star}}\xspace}
\provideenvironment{blue}{\color{darkblue}}{}

\newcommand{\dist}[2]{\ensuremath{\operatorname{\sf dist}(#1, #2)}}
\newcommand{\depth}[2]{\ensuremath{\operatorname{\sf depth}_{#1}(#2)}}
\newcommand{\cost}[1]{\ensuremath{\operatorname{\sf cost}(#1)}}
\newcommand{\weight}[1]{\ensuremath{\operatorname{\sf wt}(#1)}}
% \newcommand{\val}[1]{\ensuremath{\operatorname{\sf val}(#1)}}
\newcommand{\eps}{\varepsilon}
\newcommand{\capac}{\textsf{cap}}
\newcommand{\capOf}[1]{\ensuremath{\operatorname{\sf cap}(#1)}}
\newcommand{\rescap}[1]{\ensuremath{\operatorname{\sf cap}_f(#1)}}
\newcommand{\into}[1]{\ensuremath{\operatorname{\sf in}(#1)}}
\newcommand{\out}[1]{\ensuremath{\operatorname{\sf out}(#1)}}
\newenvironment{LP}[1]{
  \begin{array}[t]{@{}l@{}}
    \displaystyle #1 \\
    \begin{aligned}\\[-10pt]}{\end{aligned}\end{array}}

\newcommand{\opt}[1]{\ensuremath{\operatorname{\sf opt}(#1)}}
\newcommand{\Opt}[1]{\ensuremath{\operatorname{\sf Opt}(#1)}}
\newcommand{\half}{\frac{1}{2}}

\newcommand{\E}{\operatorname{E}}
\newcommand{\giv}{\,|\,}
\newcommand{\Z}{\mathbb{Z}}
\newcommand{\Zp}{\mathbb{Z}_{\ge 0}}
\newcommand{\R}{\mathbb{R}}
% \newcommand{\Rp}{\mathbb{R}_{\ge 0}}
\newcommand{\Rp}{\mathbb{R}_{\tiny +}}
\newcommand{\N}{\mathbb{N}}
\newcommand{\Np}{\mathbb{N}_{\ge 0}}
\newcommand{\calS}{{\cal S}}

\newcommand{\OP}[1]{\text{\small\sc #1}\xspace}
\newcommand{\Insert}{\OP{Insert}}
\newcommand{\FindMin}{\OP{Find-Min}}
\newcommand{\DeleteMin}{\OP{Delete-Min}}
\newcommand{\DecreaseKey}{\OP{Decrease}}
\newcommand{\Merge}{\OP{Merge}}
\newcommand{\Cut}{\OP{Cut}}
% \newcommand{\prob}[1]{\text{\sc #1}}
% a problem name; in running text, lines may break inside it
\DeclareRobustCommand{\prob}[1]{\ifmmode\text{#1}\else#1\fi}
\newcommand{\encode}[1]{\langle #1 \rangle}

\newcommand{\best}[1]{\text{\sf best}(#1)}
\newcommand{\worst}[1]{\text{\sf worst}(#1)}
\newcommand{\floor}[1]{\lfloor #1 \rfloor}
\newcommand{\ceil}[1]{\lceil #1 \rceil}

\setlength{\parskip}{2pt}
\setlength{\parindent}{1em}

% \setlength\fboxsep{0pt}

%%% longFormProof, shortFormProof, block environments
%%% \step

\newlist{steps}{enumerate}{7}
\setlist[steps]{wide,
  topsep=1pt,
  parsep=1pt,
  partopsep=1pt,
  itemsep=0pt,
  align=left,
  labelindent=0pt,
  itemindent=0pt, 
}
\setlist[steps,1]{
  label=\arabic*.,
  ref=\arabic*
}
\setlist[steps,2]{
  label=\arabic{stepsi}.\arabic*.,
  ref=\arabic{stepsi}.\arabic*
}
\setlist[steps,3]{
  label=\arabic{stepsi}.\arabic{stepsii}.\arabic*.,
  ref=\arabic{stepsi}.\arabic{stepsii}.\arabic*
}
\setlist[steps,4]{
  label=\arabic{stepsi}.\arabic{stepsii}.\arabic{stepsiii}.\arabic*.,
  ref=\arabic{stepsi}.\arabic{stepsii}.\arabic{stepsiii}.\arabic*
}
\setlist[steps,5]{
  label=\arabic{stepsi}.\arabic{stepsii}.\arabic{stepsiii}.\arabic{stepsiv}.\arabic*.,
  ref=\arabic{stepsi}.\arabic{stepsii}.\arabic{stepsiii}.\arabic{stepsiv}.\arabic*
}
\setlist[steps,6]{
  label=\arabic{stepsi}.\arabic{stepsii}.\arabic{stepsiii}.\arabic{stepsiv}.\arabic{stepsv}.\arabic*.,
  ref=\arabic{stepsi}.\arabic{stepsii}.\arabic{stepsiii}.\arabic{stepsiv}.\arabic{stepsv}.\arabic*
}
\setlist[steps,7]{
  label=\arabic{stepsi}.\arabic{stepsii}.\arabic{stepsiii}.\arabic{stepsiv}. \arabic{stepsv}. \arabic{stepsvi}.\arabic*.,
  ref=\arabic{stepsi}.\arabic{stepsii}.\arabic{stepsiii}.\arabic{stepsiv}.\arabic{stepsv}. \arabic{stepsvi}.\arabic*
}

\newcommand{\step}[1][]{\item\maybeLabel{#1}}

\makeatletter
% \skipitem steps the list counter itself (for block labels).  Give each one
% its own hyperref link target, since step numbers repeat from proof to proof.
\newcounter{skipitem}
\newcommand{\skipitem}{\stepcounter{skipitem}\refstepcounter{\@enumctr}}
\renewcommand{\theHstepsi}{skipitem.\arabic{skipitem}}
\renewcommand{\theHstepsii}{skipitem.\arabic{skipitem}}
\renewcommand{\theHstepsiii}{skipitem.\arabic{skipitem}}
\renewcommand{\theHstepsiv}{skipitem.\arabic{skipitem}}
\renewcommand{\theHstepsv}{skipitem.\arabic{skipitem}}
\renewcommand{\theHstepsvi}{skipitem.\arabic{skipitem}}
\renewcommand{\theHstepsvii}{skipitem.\arabic{skipitem}}
\makeatother

\newcommand{\maybeLabel}[1]{\ifthenelse{\isempty{#1}}{}{\label{#1}}}

\newenvironment{longFormProof}[1][Proof (long form).]
{\begin{proof}[#1]~\par\begin{steps}}{\end{steps}\vspace*{-3ex}\end{proof}}

\newenvironment{shortFormProof}[1][Proof (short form).]
{\begin{proof}[#1]}{\end{proof}}

\newenvironment{longFormSubProof}[1][Proof.]
{\begin{adjustwidth}{1em}{0em}\begin{proof}[#1]~\par\vspace*{-2pt}\skipitem\begin{steps}}
{\end{steps}\vspace*{-1ex}\end{proof}\end{adjustwidth}\smallskip}

\newenvironment{shortFormSubProof}[1][Proof.]
{\begin{proof}[#1]}
{\end{proof}\vspace*{-3pt}}

\newenvironment{block}[1][]
{\skipitem\maybeLabel{#1}\item[] \begin{steps}}{\end{steps}}

\newcommand{\comment}[1]{\hfill{\footnotesize\emph{#1}}}
\newcommand{\lineacross}{\par\vspace*{-0.7\baselineskip}\noindent\hrulefill\par}

%%% problem

% \newcounter{problem}
% \newcommand{\problem}[1][]{}
% {\ifthenelse{\isempty{#1}}{\refstepcounter{problem}}{}%
%   \paragraph{Problem \ifthenelse{\isempty{#1}}
%     {\bf\arabic{problem}.}
%     {#1}} }

\newcommand{\problemNumbered}[1]{
  \setcounter{problem}{#1}%
  \addtocounter{problem}{-1}%
  \problem
}
\newenvironment{problemTheorem}{%
  \setcounter{theorem}{\value{problem}}%
  \addtocounter{theorem}{-1}%
  \begin{theorem}[Problem \arabic{problem}]
  }{\end{theorem}}

\newcounter{problem}
\newcommand{\problem}[1][]
{\ifthenelse{\isempty{#1}}
  {\refstepcounter{problem}\paragraph{Problem~\arabic{problem}.}}
  {\paragraph{Problem~#1.}}\refreshHeader} 

\newcounter{exercise}
\newcommand{\exercise}[1][]
{\ifthenelse{\isempty{#1}}
  {\refstepcounter{exercise}\paragraph{Exercise~\theexercise.}}
  {\paragraph{Exercise~#1.}}}

\newcommand{\solution}{\paragraph{Solution.}} 

\newenvironment{mylist}[1][]
{\begin{list}{#1}{
      \setlength{\itemsep}{-1pt}
      \setlength{\labelwidth}{0pt}
      \setlength{\leftmargin}{1.25em}
    }}
{\end{list}}

\newcommand{\LNTitle}{\textcolor{red}{UNDEFINED LNTitle}}

% In a standalone lecture note's table of contents, leave room for numbers
% like 7.10 (in bold) and 7.10.2.
\makeatletter
\newcommand{\LNtocWidths}{%
  \patchcmd{\l@section}{1.5em}{2.8em}{}{\PackageWarning{macros}{Could not widen section numbers in the contents}}%
  \renewcommand*{\l@subsection}{\@dottedtocline{2}{2.8em}{3.2em}}%
  \renewcommand*{\l@subsubsection}{\@dottedtocline{3}{6.0em}{4.1em}}}
\makeatother

\newcommand{\doTitle}[3]{
  \pagenumbering{roman}
  \renewcommand{\LNTitle}{Lecture Note #1, #3}
  \renewcommand{\currentLN}{#1}
  \codeListingstrue
  % number sections, theorems etc., exercises, figures, tables and equations
  % within the lecture note (e.g. Section 8.2, Lemma 8.6), as in the book
  \renewcommand{\thesection}{#1.\arabic{section}}
  \renewcommand{\thelemma}{#1.\arabic{lemma}}
  \renewcommand{\thedefinition}{#1.\arabic{definition}}
  \renewcommand{\thetheorem}{#1.\arabic{theorem}}
  \renewcommand{\thecorollary}{#1.\arabic{corollary}}
  \renewcommand{\theclaim}{#1.\arabic{claim}}
  \renewcommand{\theobservation}{#1.\arabic{observation}}
  \renewcommand{\theconjecture}{#1.\arabic{conjecture}}
  \renewcommand{\theexercise}{#1.\arabic{exercise}}
  \renewcommand{\thefigure}{#1.\arabic{figure}}
  \renewcommand{\thetable}{#1.\arabic{table}}
  \renewcommand{\theequation}{#1.\arabic{equation}}
  \LNtocWidths
  \title{Lecture Note #1\\ #2}
  % a note's \LNabstract, if it defines one before \doTitle, goes right below the title
  \date{\ifdefined\LNabstract\parbox{0.85\textwidth}{\normalsize\normalfont\LNabstract}\\[3ex]\fi
    draft of \today, by \href{\authorURL}{N.~Young}\\[2ex] {\small\copyrightNotice}}
  \maketitle
  {\small
    \hypertarget{TOC}{}
    \tableofcontents
  }
  \thispagestyle{empty}
  % the text starts on page 1, after the title and contents (numbered i, ii, ...)
  \clearpage
  \pagenumbering{arabic}
}

\def\clap#1{\hbox to 0pt{\hss#1\hss}}

\newcommand{\refreshHeader}{}
\newcommand{\header}{}

\newcommand{\doHeader}[1]{      % for lecture notes
  \renewcommand{\header}{%
    \clap{\raisebox{0pt}{\protect\hyperlink{TOC}{$\uparrow$\scriptsize\,\textsc{contents}\hspace{1.1in}}}}%
    {\footnotesize \LNTitle{} / Section \thesection, #1}} %, #1, draft of \today}}
  \markboth{\header}{\header}
}

\newcommand{\SID}{}
\newcommand{\setAuthor}[1]{\renewcommand{\author}{#1}}
\newcommand{\setSID}[1]{\renewcommand{\SID}{#1}}

\newcommand{\doHwHeader}[1]{                                         % for homeworks
  \renewcommand{\refreshHeader}{
    \renewcommand{\header}{%
      {\footnotesize \classSection #1 / Problem~\arabic{problem} / draft of \today \hfill \author\ifGradescope~(\SID)\fi~~}}
    \markboth{\header}{\header}
  }
  \renewcommand{\header}{%
    {\footnotesize \classSection #1 / Instructions / draft of \today \hfill \author\ifGradescope~(\SID)\fi~~}}
  \markboth{\header}{\header}
}

% Homework settings, for instructors using the homeworks in a course:
% \classSection is the course name shown in each homework's running head
% (e.g. \renewcommand{\classSection}{CS 3000\xspace}), and \Gradescopetrue turns on
% the instructions for submitting to Gradescope (in its fixed-length mode),
% the hints saying on which page of the PDF each answer should be, and the
% student ID in the running head.
\newcommand{\classSection}{}
\newif\ifGradescope
\Gradescopefalse
\newcommand{\GradescopeComment}[1]{\ifGradescope\comment{#1}\fi}

\pagestyle{myheadings}
\addtolength{\headsep}{0.2in}
\addtolength{\topmargin}{-0.3in}
\addtolength{\textheight}{0.4in}

%%%%%%%%%%%%%%%%%%%%%%%%%%%%%%%%%%%%%%%%%%%%%%%%
%%%%%%%%%%%%%%%%%%%%%%%%%%%%%%%%%%%%%%%%%%%%%%%%

\newcommand{\MF}{\href
  {https://mfleck.cs.illinois.edu/building-blocks/}
  {MF}\xspace}

\newcommand{\KWSlides}{\href
  {https://www.cs.princeton.edu/~wayne/kleinberg-tardos/}
  {KW slides}\xspace}

\newcommand{\LLM}{\href
  {https://courses.csail.mit.edu/6.042/spring18/mcs.pdf}
  {LLM}\xspace}

% \DPV, \KT, \CLRS: textbooks listed (without links) in the table of sources in
% the Preliminaries note; each prints as a link to its entry in that table.
\DeclareRobustCommand{\DPV}{\noteLink{prelim}{source.DPV}{DPV}\xspace}
\DeclareRobustCommand{\KT}{\noteLink{prelim}{source.KT}{KT}\xspace}
\DeclareRobustCommand{\CLRS}{\noteLink{prelim}{source.CLRS}{CLRS}\xspace}
\AtBeginDocument{\pdfstringdefDisableCommands{\def\DPV{DPV }\def\KT{KT }\def\CLRS{CLRS }}}

\newcommand{\JE}{\href
  {https://jeffe.cs.illinois.edu/teaching/algorithms/}
  {JE}\xspace}

\newcommand{\JEGreedy}{\href
  {https://jeffe.cs.illinois.edu/teaching/algorithms/book/04-greedy.pdf}
  {JE Chapter 4}\xspace}

\newcommand{\BBCode}{\href
{https://northeastern.blackboard.com/webapps/cmsmain/webui/courses/CS3000.MERGED.202010/code?action=frameset&subaction=view&uniq=-jj2max&course_id=_2597392_1&mask=\%2Fcourses\%2FCS3000.MERGED.202010}
  {code}\xspace}

% The Python code (src/code).  Each lecture note that has code ends with a
% section listing its files: \codeSection{path, path, ...}.  In the text,
% \pycode{greedy_algorithms/Dijkstra.py} prints the path as a link to its listing
% (in the homeworks, which have no listings, to the file on GitHub), and
% \pycode[text]{...} uses the given text instead.  \pythonCode{path, ...} is the
% standard sentence ending a section whose algorithms have code:
% "Python code: path, ... (Section N.k)."
\newcommand{\codeFolder}{https://github.com/nealeyoung/algs101/tree/main/src/code}
\newcommand{\codeRepo}{https://github.com/nealeyoung/algs101/blob/main/src/code}
\newif\ifcodeListings   % set by \doTitle: this document has the listings
\newcommand{\githubCode}[2][]{\href{\codeRepo/#2}{\ifx\relax#1\relax\nolinkurl{#2}\else#1\fi}}
\newcommand{\pycode}[2][]{%
  \ifcodeListings\hyperref[code: #2]{\ifx\relax#1\relax\nolinkurl{#2}\else#1\fi}%
  \else\githubCode[#1]{#2}\fi}
\makeatletter
\newcommand{\pythonCode}[1]{%
  \def\pythonCode@first{}\def\pythonCode@list{}%
  \forcsvlist{\pythonCode@add}{#1}%
  Python code: \pythonCode@list\ (Section~\ref{code: \pythonCode@first}).}
\newcommand{\pythonCode@add}[1]{%
  \ifx\pythonCode@first\@empty\def\pythonCode@first{#1}\appto\pythonCode@list{\pycode{#1}}%
  \else\appto\pythonCode@list{, \pycode{#1}}\fi}
\makeatother

% Code listings that copy and paste exactly (in PDF viewers that honor the
% ToUnicode maps, e.g. MuPDF-based ones; PDFKit viewers such as Preview and
% Skim drop leading spaces, so each listing's heading links to the file on
% GitHub).  Spaces are typeset as white visible-space glyphs that copy as
% spaces, and hyphens and asterisks as the T1 glyphs, so they copy as themselves.
\usepackage[T1,OT1]{fontenc}   % T1 for the code listings only; OT1 stays the default
\usepackage{listings}
\pdfgentounicode=1
\input glyphtounicode
\pdfglyphtounicode{uni2423}{0020}
\makeatletter
\def\lst@visiblespace{\textcolor{white}{\char32}}
% put an invisible space on each blank line, so that blank lines copy too
% (listings' OnEmptyLine hook runs before the line break, so patch the loop)
\def\lst@DoNewLines{%
    \@whilenum\lst@newlines>\@ne \do {\lst@NewLine \lst@visiblespace}%
    \ifnum\lst@newlines>\z@ \lst@NewLine \fi}
\makeatother
\lstset{language=Python,
  basicstyle=\fontencoding{T1}\fontfamily{lmtt}\selectfont\footnotesize,
  upquote=true, showstringspaces=false, columns=fullflexible, keepspaces=true,
  showspaces=true, literate={-}{{\char45}}1 {*}{{\char42}}1,
  frame=single, framesep=1ex, rulecolor=\color{gray}}
% the listings' paths are relative to the lecture note's folder (in the book, to src/)
\ifdefined\BookMode\newcommand{\codeDir}{code/}\else\newcommand{\codeDir}{../code/}\fi
% \codefile{path}: a listing of src/code/path, headed by a link to it on GitHub
\newcommand{\codefile}[1]{%
  \subsection*{\githubCode{#1}}\label{code: #1}%
  \lstinputlisting{\codeDir#1}}
% \codeSection{path, path, ...}: the lecture note's last section, listing its code
\newcommand{\codeSection}[1]{%
  \newpage
  \section{Python code}\label{\LNtag{\currentLN}: sec: code}\doHeader{Python code}
  This section lists the Python code for this lecture note.
  Click on a file's name to open the file on GitHub.
  Copying code from GitHub is more reliable than copying it from this PDF,
  since some PDF viewers drop the leading spaces (the indentation) of copied lines.
  \forcsvlist{\codefile}{#1}}
% the copyright notice, for title pages and the homework first pages
\newcommand{\licenseURL}{https://github.com/nealeyoung/algs101/blob/main/LICENSE}
\newcommand{\authorURL}{https://www.cs.ucr.edu/\string~neal}
\newcommand{\copyrightNotice}{%
  \copyright\ 2018--2026 \href{\authorURL}{Neal E. Young}.
  Text: CC BY-NC-SA 4.0; code: MIT; except where noted.
  See \url{\licenseURL}.}
% the published site: PDFs of the notes, and the homework assignments and templates
\newcommand{\siteURL}{https://nealeyoung.github.io/algs101/}

\newcommand{\blackboard}{\href
  {https://northeastern.blackboard.com/webapps/blackboard/execute/launcher?type=Course&id=_2597392_1&url=}
  {Blackboard}\xspace}

\newcommand{\JUSlides}{\href
  {http://www.khoury.neu.edu/home/jullman/cs3000f18/schedule.html}
  {Prof.~Ullman's slides}\xspace}

% a URL on its own line, in three layouts
\newcommand{\URL}[1]{\par-- \parbox[t]{0.8\textwidth}{\url{#1}} \medskip\par}
\newcommand{\URLsmall}[1]{\par-- \parbox[t]{0.95\textwidth}{\small\url{#1}} \medskip\par}
\newcommand{\URLplain}[1]{\par\parbox[t]{\textwidth}{\url{#1}}\par}

%%%%%%%%%%%%%%%%%%%%%%%%%%%%%%%%%%%%%%%%%%%%%%%%
%%%%%%%%%%%%%%%%%%%%%%%%%%%%%%%%%%%%%%%%%%%%%%%%

\newcommand{\continued}{\vfill\centerline{\emph{(continued)}}\vfill\newpage}

% load hyperref last to avoid conflicts; long URLs may also break after hyphens
\PassOptionsToPackage{hyphens}{url}
\ifdefined\BookMode\else\usepackage{xr-hyper}\fi   % for \lectureNoteRefs (homeworks), below
\usepackage{hyperref}
\hypersetup{
    colorlinks,
    linkcolor={blue!50!black},
    citecolor={blue!50!black},
    urlcolor={blue!80!black}
}

\providecommand{\answerDirectory}{.}
\newcommand{\inputAnswer}[1]{\refreshHeader\input{\answerDirectory/#1}}
\newcommand{\inputSection}[1]{\input{#1}}

% Refer to a lecture note by its tag (the prefix of its labels), e.g.
% \LN{proofs} prints "Lecture Note 2" (in the book, a link to that chapter).  The table below is the only place the numbers appear.
\makeatletter
\newcommand{\defineLN}[3]{\@namedef{LNnumber@#2}{#1}\@namedef{LNtag@#1}{#2}\@namedef{LNfile@#2}{#3}}
\newcommand{\LNnumber}[1]{%
  \ifcsname LNnumber@#1\endcsname\csname LNnumber@#1\endcsname
  \else\PackageError{macros}{Unknown lecture-note tag `#1'}{}\fi}
\newcommand{\LNtag}[1]{\csname LNtag@#1\endcsname}
\newcommand{\LNfile}[1]{\csname LNfile@#1\endcsname}
\makeatother
\defineLN{1}{prelim}{1_preliminaries}
\defineLN{2}{proofs}{2_proofs}
\defineLN{3}{matching}{3_stable_matching}
\defineLN{4}{dc}{4_divide_and_conquer}
\defineLN{5}{greedy}{5_greedy}
\defineLN{6}{graphs}{6_graph_traversal}
\defineLN{7}{dp}{7_dynamic_programming}
\defineLN{8}{flow}{8_flow_reductions_LP}
\defineLN{9}{np}{9_NP}

% The text is the same in the book and in the standalone lecture notes; only
% links differ.  In the book, links to other lecture notes stay in the book.
% In a standalone note, they open the other note's published PDF.
% \noteLink{tag}{target}{text} links text to \hypertarget{target} in lecture note tag.
% \LN{proofs} prints "Lecture Note 2", linked to that lecture note.
\newcommand{\currentLN}{0}   % in a standalone note, its number (set by \doTitle)
\newcommand{\notePDF}[1]{\siteURL notes/\LNfile{#1}.pdf}
\newcommand{\noteLink}[3]{%
  \ifnum\LNnumber{#1}=\currentLN\relax\hyperlink{#2}{#3}%
  \else\href{\notePDF{#1}\#nameddest=#2}{#3}\fi}
% \bookLink{target}{text} links text to \hypertarget{target} in the book's
% appendices (from a standalone note, in the published book).
\DeclareRobustCommand{\bookLink}[2]{%
  \ifdefined\BookMode\hyperlink{#1}{#2}%
  \else\href{\siteURL lecture_notes_on_algorithms.pdf\#nameddest=#1}{#2}\fi}
\DeclareRobustCommand{\LN}[1]{%
  \ifnum\LNnumber{#1}=\currentLN\relax Lecture Note~\LNnumber{#1}%
  \else\href{\notePDF{#1}}{Lecture Note~\LNnumber{#1}}\fi}
\pdfstringdefDisableCommands{\def\LN#1{Lecture Note \LNnumber{#1}}}

% Homeworks refer to the notes' sections, exercises and lemmas by \ref: each
% homework's main file says \lectureNoteRefs after \input{macros}.  A homework
% folder has a link refs/ to src/refs/, which holds each note's label records
% (made by tools/refs.sh from the notes' builds).  Read with xr-hyper, they make
% \ref{greedy: exercise: ...} give the exercise's number in the published note,
% linked to it there; labels only the book has (its appendices) link to the book.
\ifdefined\BookMode\else
  \makeatletter
  \newcommand{\lectureNoteRefs}{%
    \@for\LN@tag:=prelim,proofs,matching,dc,greedy,graphs,dp,flow,np\do{%
      \externaldocument{refs/\LNfile{\LN@tag}}[\notePDF{\LN@tag}]}%
    \externaldocument{refs/book}[\siteURL lecture_notes_on_algorithms.pdf]}
  % hyperref links a \ref to another document's label as URL#anchor; browsers'
  % PDF viewers need URL#nameddest=anchor (the form \noteLink and \bookLink use)
  \def\Hy@setref@link#1#2#3#4#5#6\@nil#7{%
    \begingroup
      \ifx\\#5\\\toks0={\hyper@@link{#5}{#4}}\else\toks0={\hyper@@link{#5}{nameddest=#4}}\fi
      \toks1=\expandafter{#7{#1}{#2}{#3}{#4}{#5}}%
      \edef\x{\endgroup\the\toks0 {\the\toks1 }}%
    \x}
  \makeatother
\fi

% Book mode: main.tex defines \BookMode before \input{macros} to combine the
% lecture notes into one report-class document, one chapter per note.
% Without \BookMode (each lecture note on its own, and the homeworks) none
% of this applies.
\ifdefined\BookMode
  \usepackage{etoc}

  % links to other lecture notes stay in the book
  \renewcommand{\noteLink}[3]{\hyperlink{#2}{#3}}

  % each chapter is a lecture note, and is called one, as in the standalone notes
  \renewcommand{\chaptername}{Lecture Note}
  \setcounter{secnumdepth}{3}   % number subsubsections, as in the standalone notes

  % global table of contents, for main.tex
  \newcommand{\bookContents}{%
    \hypertarget{TOC}{}%
    \etocsetnexttocdepth{section}%
    \tableofcontents}

  % \doTitle{N}{title}{short title} starts chapter N, with a link to the global
  % table of contents and the chapter's own table of contents
  \renewcommand{\doTitle}[3]{%
    \setcounter{chapter}{\numexpr#1-1\relax}%
    \chapter{\texorpdfstring{#2}{#3}}%
    \label{\LNtag{#1}: chapter}%
    \csgdef{LNshorttitle@#1}{#3}%
    \csgdef{LNtitle@#1}{#2}%
    \renewcommand{\LNTitle}{Lecture Note #1, #3}%
    \codeListingstrue
    \hyperlink{TOC}{$\uparrow$\scriptsize\,\textsc{contents}}%
    {\small
      \hypertarget{TOC-#1}{}%
      \etocsettocstyle{\section*{\contentsname}}{}%
      \etocsetnexttocdepth{subsubsection}%
      \localtableofcontents
    }}

  % the code: each lecture note's \codeSection just records its files, and
  % \codeAppendix (in the book's appendix) lists them all, a section per note
  \makeatletter   % (file names have no spaces, so drop the spaces around them)
  \renewcommand{\codeSection}[1]{\csxdef{LNcode@\arabic{chapter}}{\zap@space#1 \@empty}}
  \makeatother
  % \codeAppendix: first a contents list (each note's topic, then its files,
  % linked to the listings), then a section per note listing its files
  \makeatletter
  \def\LNbasename#1{\LN@base#1/\LN@end}   % a/b/c.py -> c.py
  \def\LN@base#1/#2\LN@end{\ifx\relax#2\relax#1\else\LN@base#2\LN@end\fi}
  \makeatother
  \newcommand{\codeContentsFile}[1]{%
    \edef\LNfilename{\LNbasename{#1}}%
    \item \pycode[\texttt{\expandafter\detokenize\expandafter{\LNfilename}}]{#1}}
  \newcommand{\codeAppendix}{%
    \begin{itemize}
    \foreach \n in {1,...,9} {%
      \ifcsdef{LNcode@\n}{%
        \item \hyperref[code: note \n]{\csuse{LNtitle@\n}} (\LN{\LNtag{\n}})
        \begin{itemize}
        \edef\LNcodelist{\csuse{LNcode@\n}}%
        \expandafter\forcsvlist\expandafter\codeContentsFile\expandafter{\LNcodelist}%
        \end{itemize}}{}}
    \end{itemize}
    \newpage
    \foreach \n in {1,...,9} {%
      \ifcsdef{LNcode@\n}{%
        \section[Lecture Note \n: \csuse{LNtitle@\n}]{%
          \texorpdfstring{\LN{\LNtag{\n}}: \csuse{LNtitle@\n}}{Lecture Note \n: \csuse{LNtitle@\n}}}%
        \label{code: note \n}%
        \edef\LNcodelist{\csuse{LNcode@\n}}%
        \expandafter\forcsvlist\expandafter\codefile\expandafter{\LNcodelist}%
        \newpage}{}}}

  % a lecture note is a chapter: \LN{proofs} prints a link "Lecture Note 2"
  \DeclareRobustCommand{\LN}[1]{\hyperref[#1: chapter]{Lecture Note~\ref*{#1: chapter}}}
  \pdfstringdefDisableCommands{\def\LN#1{Lecture Note \LNnumber{#1}}}

  % running head; "contents" links to the current chapter's table of contents
  \renewcommand{\doHeader}[1]{%
    \renewcommand{\header}{%
      \clap{\raisebox{0pt}{\protect\hyperlink{TOC-\arabic{chapter}}{$\uparrow$\scriptsize\,\textsc{contents}\hspace{1.1in}}}}%
      {\footnotesize \LNTitle{} / Section \thesection, #1}}%
    \markboth{\header}{\header}}

  % number theorems, lemmas, etc. and exercises within chapters (e.g. Lemma 8.1)
  \counterwithin{lemma}{chapter}
  \counterwithin{definition}{chapter}
  \counterwithin{theorem}{chapter}
  \counterwithin{corollary}{chapter}
  \counterwithin{claim}{chapter}
  \counterwithin{observation}{chapter}
  \counterwithin{conjecture}{chapter}
  \counterwithin{exercise}{chapter}
\fi

% %%%%%%%%%%%%%%%%%%%%%%%%%%%%%%%%%%%%%%%%%%%%%%%%%%%%%%%%%%%%%%%%%%%%%%%%
% %%%%%%%%%%%%%%%%%%%%%%%%%%%%%%%%%%%%%%%%%%%%%%%%%%%%%%%%%%%%%%%%%%%%%%%%

\ExplSyntaxOn
\seq_new:N \l_itemize_seq
% \NewDocumentEnvironment{REITEMIZE}{ +b }{
\NewDocumentCommand{\REITEMIZECMD}{ +m }{
  \regex_split:nnN {\c{item}} {#1} \l_itemize_seq
  \seq_pop:NN \l_itemize_seq \l_tmpa_tl% remove the (hopefully) empty first item
  \seq_if_empty:NF \l_itemize_seq {
    \l_tmpa_tl
    \seq_map_inline:Nn \l_itemize_seq {
      \regex_split:nnN {\c{item}} {#1} \l_itemize_seq {
        \regex_extract_once:nnNTF { ^\[([^]]*)\](.*) } { ##1} \l_foo_seq
      { \item[\seq_item:Nn \l_foo_seq {2}]\ITEMIZER{\seq_item:Nn \l_foo_seq {3}} } { \item\ITEMIZER{##1}  }
    }
  }
}
}
\ExplSyntaxOff

\newlength{\ITEMHT}
\setlength{\ITEMHT}{1in}

% \setlength{\fboxsep}{1pt}
\newcommand{\ITEMIZER}[1]{%
  % \framebox{
\adjustbox{valign=t}{\parbox{\textwidth}{\linespread{1.15}\selectfont #1}}~~\adjustbox{valign=t}{\parbox{0pt}{\color{gray}\rule[-\ITEMHT]{0.01pt}{\ITEMHT}}}
% }
  % \begin{tabular}{@{}p{\textwidth}@{\rule[-\ITEMHT]{0pt}{\ITEMHT}}}
  %   #1
  % \end{tabular}
}

\newsavebox\thesmashminipage
\newlength{\thesmashminipagelength}
\newlength{\answerboxlength}
\newlength{\answerboxoverrun}
\newenvironment{smashminipage}[2]
{\setlength\answerboxlength{#2}
\begin{lrbox}{\thesmashminipage}\begin{minipage}[t]{#1}}
{\end{minipage}\end{lrbox}%
\setlength{\thesmashminipagelength}{\dimexpr\ht\thesmashminipage+\dp\thesmashminipage}%
\setlength{\answerboxoverrun}{\dimexpr\thesmashminipagelength-\answerboxlength}%
\ifdim\the\thesmashminipagelength>\the\answerboxlength%
\usebox{\thesmashminipage}%
\smash{\color{red}\hspace*{-0.98\textwidth}\rule[-\answerboxlength]{0.98\textwidth}{1pt}}
\errmessage{Answer exceeds allotted length, please shorten it.}
\else
\usebox{\thesmashminipage}%
\fi
}



\newenvironment{answerBox}[1]{%
  \def\ANSWERBOXHEIGHT{#1}
  \begin{adjustbox}{valign=t}
  \begin{smashminipage}{\textwidth}{\ANSWERBOXHEIGHT}
}{\end{smashminipage}\end{adjustbox}~~~\adjustbox{valign=t}{\parbox{0pt}{\color{gray}\rule[-\ANSWERBOXHEIGHT]{0.01pt}{\ANSWERBOXHEIGHT}\rule[-\ANSWERBOXHEIGHT]{10pt}{0.01pt}}}}

% % \setlist[enumerate]{itemsep=0pt,parsep=0pt,topsep=0pt,leftmargin=100pt}

% \makeatletter 
\LetLtxMacro\SVenumerate\enumerate
\let\endSVenumerate\endenumerate

\newcommand{\injectEnumerate}{
  \let\SValph\alph
  \def\alph{}
  \renewenvironment{enumerate}[1][]
  {\let\alph\SValph \SVenumerate[label=(\alph*),itemsep=2pt,parsep=0pt]\REITEMIZE}{\endREITEMIZE\endSVenumerate}
  % \setlist[enumerate]{before=\Collect@Body\REITEMIZECMD}
}
% \makeatother
\newcommand{\restoreEnumerate}{
  \renewenvironment{enumerate}[1][]
  \let\SValph\alph
  \def\alph{}
  {\let\alph\SValph \SVenumerate[label=(\alph*)]}{\endSVenumerate}
  % \setlist[enumerate]{before=,after=}
}
\newcommand{\injectBlue}{
  \let\SVblue\blue
  \let\endSVblue\endblue
  \def\blue{\SVblue\REITEMIZE}
  \def\endblue{\endREITEMIZE\endSVblue}
}
\newcommand{\restoreBlue}{
  \let\blue\SVblue
  \let\endblue\endSVblue
}
\newcommand{\injectEnumerateBlue}{
  \setlist[enumerate]{before=\injectBlue,after=\restoreBlue,itemsep=2pt,parsep=0pt}
}
\newcommand{\restoreEnumerateBlue}{
  \setlist[enumerate]{before=,after=}
}


% %%%%%%%%%%%%%%%%%%%%%%%%%%%%%%%%%%%%%%%%%%%%%%%%%%%%%%%%%%%%%%%%%%%%%%%%
% %%%%%%%%%%%%%%%%%%%%%%%%%%%%%%%%%%%%%%%%%%%%%%%%%%%%%%%%%%%%%%%%%%%%%%%%

% \NewDocumentEnvironment{env}{ }
% {\Environenv}
% {\endEnvironenv}

% \NewEnviron{REITEMIZE}{\REITEMIZECMD\BODY}

\makeatletter 
\newenvironment{REITEMIZE}{\Collect@Body\REITEMIZECMD}{}
\makeatother 

\newcommand{\NODEone}[1]{\ensuremath{v_{#1}}\xspace}
\newcommand{\NODE}[2]{{\NODEone {#1,#2}}\xspace}
\newcommand{\EDGE}[2]{\ensuremath{(\textsf{#1}, \textsf{#2})}}
\newcommand{\ITEM}{\item}

%%% Local Variables:
%%% mode: latex
%%% End:
.  A homework
% folder has a link refs/ to src/refs/, which holds each note's label records
% (made by tools/refs.sh from the notes' builds).  Read with xr-hyper, they make
% \ref{greedy: exercise: ...} give the exercise's number in the published note,
% linked to it there; labels only the book has (its appendices) link to the book.
\ifdefined\BookMode\else
  \makeatletter
  \newcommand{\lectureNoteRefs}{%
    \@for\LN@tag:=prelim,proofs,matching,dc,greedy,graphs,dp,flow,np\do{%
      \externaldocument{refs/\LNfile{\LN@tag}}[\notePDF{\LN@tag}]}%
    \externaldocument{refs/book}[\siteURL lecture_notes_on_algorithms.pdf]}
  % hyperref links a \ref to another document's label as URL#anchor; browsers'
  % PDF viewers need URL#nameddest=anchor (the form \noteLink and \bookLink use)
  \def\Hy@setref@link#1#2#3#4#5#6\@nil#7{%
    \begingroup
      \ifx\\#5\\\toks0={\hyper@@link{#5}{#4}}\else\toks0={\hyper@@link{#5}{nameddest=#4}}\fi
      \toks1=\expandafter{#7{#1}{#2}{#3}{#4}{#5}}%
      \edef\x{\endgroup\the\toks0 {\the\toks1 }}%
    \x}
  \makeatother
\fi

% Book mode: main.tex defines \BookMode before %%%% DON'T EDIT THIS FILE %%%%

\usepackage{xifthen}% \isempty
\usepackage{etoolbox}% \patchcmd
% \usepackage{stringstrings}% \substring

\usepackage{amsmath, amsthm, amssymb}
\usepackage[margin=1in]{geometry}
\usepackage{xspace}
\usepackage{graphicx}
\usepackage[normalem]{ulem}
\usepackage{adjustbox}
\usepackage{xcolor}
\usepackage{changepage}
\usepackage{framed}
\usepackage{multicol}
\usepackage{enumitem}
% \usepackage{listings} % lstlistings environment

\usepackage{makecmds}           %\provideenvironment

\usepackage{xparse,environ}
\usepackage{letltxmacro}       % \LetLtxMacro
\usepackage{wrapfig}
\usepackage{tikz}

\usetikzlibrary{shapes}
\usetikzlibrary{positioning}
\usetikzlibrary{shapes.misc}
\usetikzlibrary{arrows}

\newtheorem{lemma}{Lemma}
\newtheorem{definition}{Definition}
\newtheorem{theorem}{Theorem}
\newtheorem{corollary}{Corollary}
\newtheorem{claim}{Claim}
\newtheorem{observation}{Observation}
\newtheorem{conjecture}{Conjecture}

% A lemma or theorem restated later with its original number (set with
% \setcounter) goes in a restatement environment, which gives it its own link
% anchor (otherwise it would duplicate the original's).
\newenvironment{restatement}
  {\def\theHlemma{restated.\thelemma}\def\theHtheorem{restated.\thetheorem}}{}

\newenvironment{myframed}
{\smallskip\par\noindent\begin{minipage}{\textwidth}\begin{framed}\vspace*{-5pt}}{\vspace*{-5pt}\end{framed}\end{minipage}\smallskip}


\definecolor{darkblue}{RGB}{0,0,200}

 
\let\oldParagraph\paragraph
\providecommand{\gradescopeColor}{\color{black}}
\renewcommand{\paragraph}[1]{\oldParagraph{\gradescopeColor #1}}
\providecommand{\gradescopeInit}{}

\providecommand{\BLUE}[1]{{\color{darkblue} #1}}
\providecommand{\REPLACEME}{\textcolor{red}{\ensuremath{\star \cdots \star}}\xspace}
\provideenvironment{blue}{\color{darkblue}}{}

\newcommand{\dist}[2]{\ensuremath{\operatorname{\sf dist}(#1, #2)}}
\newcommand{\depth}[2]{\ensuremath{\operatorname{\sf depth}_{#1}(#2)}}
\newcommand{\cost}[1]{\ensuremath{\operatorname{\sf cost}(#1)}}
\newcommand{\weight}[1]{\ensuremath{\operatorname{\sf wt}(#1)}}
% \newcommand{\val}[1]{\ensuremath{\operatorname{\sf val}(#1)}}
\newcommand{\eps}{\varepsilon}
\newcommand{\capac}{\textsf{cap}}
\newcommand{\capOf}[1]{\ensuremath{\operatorname{\sf cap}(#1)}}
\newcommand{\rescap}[1]{\ensuremath{\operatorname{\sf cap}_f(#1)}}
\newcommand{\into}[1]{\ensuremath{\operatorname{\sf in}(#1)}}
\newcommand{\out}[1]{\ensuremath{\operatorname{\sf out}(#1)}}
\newenvironment{LP}[1]{
  \begin{array}[t]{@{}l@{}}
    \displaystyle #1 \\
    \begin{aligned}\\[-10pt]}{\end{aligned}\end{array}}

\newcommand{\opt}[1]{\ensuremath{\operatorname{\sf opt}(#1)}}
\newcommand{\Opt}[1]{\ensuremath{\operatorname{\sf Opt}(#1)}}
\newcommand{\half}{\frac{1}{2}}

\newcommand{\E}{\operatorname{E}}
\newcommand{\giv}{\,|\,}
\newcommand{\Z}{\mathbb{Z}}
\newcommand{\Zp}{\mathbb{Z}_{\ge 0}}
\newcommand{\R}{\mathbb{R}}
% \newcommand{\Rp}{\mathbb{R}_{\ge 0}}
\newcommand{\Rp}{\mathbb{R}_{\tiny +}}
\newcommand{\N}{\mathbb{N}}
\newcommand{\Np}{\mathbb{N}_{\ge 0}}
\newcommand{\calS}{{\cal S}}

\newcommand{\OP}[1]{\text{\small\sc #1}\xspace}
\newcommand{\Insert}{\OP{Insert}}
\newcommand{\FindMin}{\OP{Find-Min}}
\newcommand{\DeleteMin}{\OP{Delete-Min}}
\newcommand{\DecreaseKey}{\OP{Decrease}}
\newcommand{\Merge}{\OP{Merge}}
\newcommand{\Cut}{\OP{Cut}}
% \newcommand{\prob}[1]{\text{\sc #1}}
% a problem name; in running text, lines may break inside it
\DeclareRobustCommand{\prob}[1]{\ifmmode\text{#1}\else#1\fi}
\newcommand{\encode}[1]{\langle #1 \rangle}

\newcommand{\best}[1]{\text{\sf best}(#1)}
\newcommand{\worst}[1]{\text{\sf worst}(#1)}
\newcommand{\floor}[1]{\lfloor #1 \rfloor}
\newcommand{\ceil}[1]{\lceil #1 \rceil}

\setlength{\parskip}{2pt}
\setlength{\parindent}{1em}

% \setlength\fboxsep{0pt}

%%% longFormProof, shortFormProof, block environments
%%% \step

\newlist{steps}{enumerate}{7}
\setlist[steps]{wide,
  topsep=1pt,
  parsep=1pt,
  partopsep=1pt,
  itemsep=0pt,
  align=left,
  labelindent=0pt,
  itemindent=0pt, 
}
\setlist[steps,1]{
  label=\arabic*.,
  ref=\arabic*
}
\setlist[steps,2]{
  label=\arabic{stepsi}.\arabic*.,
  ref=\arabic{stepsi}.\arabic*
}
\setlist[steps,3]{
  label=\arabic{stepsi}.\arabic{stepsii}.\arabic*.,
  ref=\arabic{stepsi}.\arabic{stepsii}.\arabic*
}
\setlist[steps,4]{
  label=\arabic{stepsi}.\arabic{stepsii}.\arabic{stepsiii}.\arabic*.,
  ref=\arabic{stepsi}.\arabic{stepsii}.\arabic{stepsiii}.\arabic*
}
\setlist[steps,5]{
  label=\arabic{stepsi}.\arabic{stepsii}.\arabic{stepsiii}.\arabic{stepsiv}.\arabic*.,
  ref=\arabic{stepsi}.\arabic{stepsii}.\arabic{stepsiii}.\arabic{stepsiv}.\arabic*
}
\setlist[steps,6]{
  label=\arabic{stepsi}.\arabic{stepsii}.\arabic{stepsiii}.\arabic{stepsiv}.\arabic{stepsv}.\arabic*.,
  ref=\arabic{stepsi}.\arabic{stepsii}.\arabic{stepsiii}.\arabic{stepsiv}.\arabic{stepsv}.\arabic*
}
\setlist[steps,7]{
  label=\arabic{stepsi}.\arabic{stepsii}.\arabic{stepsiii}.\arabic{stepsiv}. \arabic{stepsv}. \arabic{stepsvi}.\arabic*.,
  ref=\arabic{stepsi}.\arabic{stepsii}.\arabic{stepsiii}.\arabic{stepsiv}.\arabic{stepsv}. \arabic{stepsvi}.\arabic*
}

\newcommand{\step}[1][]{\item\maybeLabel{#1}}

\makeatletter
% \skipitem steps the list counter itself (for block labels).  Give each one
% its own hyperref link target, since step numbers repeat from proof to proof.
\newcounter{skipitem}
\newcommand{\skipitem}{\stepcounter{skipitem}\refstepcounter{\@enumctr}}
\renewcommand{\theHstepsi}{skipitem.\arabic{skipitem}}
\renewcommand{\theHstepsii}{skipitem.\arabic{skipitem}}
\renewcommand{\theHstepsiii}{skipitem.\arabic{skipitem}}
\renewcommand{\theHstepsiv}{skipitem.\arabic{skipitem}}
\renewcommand{\theHstepsv}{skipitem.\arabic{skipitem}}
\renewcommand{\theHstepsvi}{skipitem.\arabic{skipitem}}
\renewcommand{\theHstepsvii}{skipitem.\arabic{skipitem}}
\makeatother

\newcommand{\maybeLabel}[1]{\ifthenelse{\isempty{#1}}{}{\label{#1}}}

\newenvironment{longFormProof}[1][Proof (long form).]
{\begin{proof}[#1]~\par\begin{steps}}{\end{steps}\vspace*{-3ex}\end{proof}}

\newenvironment{shortFormProof}[1][Proof (short form).]
{\begin{proof}[#1]}{\end{proof}}

\newenvironment{longFormSubProof}[1][Proof.]
{\begin{adjustwidth}{1em}{0em}\begin{proof}[#1]~\par\vspace*{-2pt}\skipitem\begin{steps}}
{\end{steps}\vspace*{-1ex}\end{proof}\end{adjustwidth}\smallskip}

\newenvironment{shortFormSubProof}[1][Proof.]
{\begin{proof}[#1]}
{\end{proof}\vspace*{-3pt}}

\newenvironment{block}[1][]
{\skipitem\maybeLabel{#1}\item[] \begin{steps}}{\end{steps}}

\newcommand{\comment}[1]{\hfill{\footnotesize\emph{#1}}}
\newcommand{\lineacross}{\par\vspace*{-0.7\baselineskip}\noindent\hrulefill\par}

%%% problem

% \newcounter{problem}
% \newcommand{\problem}[1][]{}
% {\ifthenelse{\isempty{#1}}{\refstepcounter{problem}}{}%
%   \paragraph{Problem \ifthenelse{\isempty{#1}}
%     {\bf\arabic{problem}.}
%     {#1}} }

\newcommand{\problemNumbered}[1]{
  \setcounter{problem}{#1}%
  \addtocounter{problem}{-1}%
  \problem
}
\newenvironment{problemTheorem}{%
  \setcounter{theorem}{\value{problem}}%
  \addtocounter{theorem}{-1}%
  \begin{theorem}[Problem \arabic{problem}]
  }{\end{theorem}}

\newcounter{problem}
\newcommand{\problem}[1][]
{\ifthenelse{\isempty{#1}}
  {\refstepcounter{problem}\paragraph{Problem~\arabic{problem}.}}
  {\paragraph{Problem~#1.}}\refreshHeader} 

\newcounter{exercise}
\newcommand{\exercise}[1][]
{\ifthenelse{\isempty{#1}}
  {\refstepcounter{exercise}\paragraph{Exercise~\theexercise.}}
  {\paragraph{Exercise~#1.}}}

\newcommand{\solution}{\paragraph{Solution.}} 

\newenvironment{mylist}[1][]
{\begin{list}{#1}{
      \setlength{\itemsep}{-1pt}
      \setlength{\labelwidth}{0pt}
      \setlength{\leftmargin}{1.25em}
    }}
{\end{list}}

\newcommand{\LNTitle}{\textcolor{red}{UNDEFINED LNTitle}}

% In a standalone lecture note's table of contents, leave room for numbers
% like 7.10 (in bold) and 7.10.2.
\makeatletter
\newcommand{\LNtocWidths}{%
  \patchcmd{\l@section}{1.5em}{2.8em}{}{\PackageWarning{macros}{Could not widen section numbers in the contents}}%
  \renewcommand*{\l@subsection}{\@dottedtocline{2}{2.8em}{3.2em}}%
  \renewcommand*{\l@subsubsection}{\@dottedtocline{3}{6.0em}{4.1em}}}
\makeatother

\newcommand{\doTitle}[3]{
  \pagenumbering{roman}
  \renewcommand{\LNTitle}{Lecture Note #1, #3}
  \renewcommand{\currentLN}{#1}
  \codeListingstrue
  % number sections, theorems etc., exercises, figures, tables and equations
  % within the lecture note (e.g. Section 8.2, Lemma 8.6), as in the book
  \renewcommand{\thesection}{#1.\arabic{section}}
  \renewcommand{\thelemma}{#1.\arabic{lemma}}
  \renewcommand{\thedefinition}{#1.\arabic{definition}}
  \renewcommand{\thetheorem}{#1.\arabic{theorem}}
  \renewcommand{\thecorollary}{#1.\arabic{corollary}}
  \renewcommand{\theclaim}{#1.\arabic{claim}}
  \renewcommand{\theobservation}{#1.\arabic{observation}}
  \renewcommand{\theconjecture}{#1.\arabic{conjecture}}
  \renewcommand{\theexercise}{#1.\arabic{exercise}}
  \renewcommand{\thefigure}{#1.\arabic{figure}}
  \renewcommand{\thetable}{#1.\arabic{table}}
  \renewcommand{\theequation}{#1.\arabic{equation}}
  \LNtocWidths
  \title{Lecture Note #1\\ #2}
  % a note's \LNabstract, if it defines one before \doTitle, goes right below the title
  \date{\ifdefined\LNabstract\parbox{0.85\textwidth}{\normalsize\normalfont\LNabstract}\\[3ex]\fi
    draft of \today, by \href{\authorURL}{N.~Young}\\[2ex] {\small\copyrightNotice}}
  \maketitle
  {\small
    \hypertarget{TOC}{}
    \tableofcontents
  }
  \thispagestyle{empty}
  % the text starts on page 1, after the title and contents (numbered i, ii, ...)
  \clearpage
  \pagenumbering{arabic}
}

\def\clap#1{\hbox to 0pt{\hss#1\hss}}

\newcommand{\refreshHeader}{}
\newcommand{\header}{}

\newcommand{\doHeader}[1]{      % for lecture notes
  \renewcommand{\header}{%
    \clap{\raisebox{0pt}{\protect\hyperlink{TOC}{$\uparrow$\scriptsize\,\textsc{contents}\hspace{1.1in}}}}%
    {\footnotesize \LNTitle{} / Section \thesection, #1}} %, #1, draft of \today}}
  \markboth{\header}{\header}
}

\newcommand{\SID}{}
\newcommand{\setAuthor}[1]{\renewcommand{\author}{#1}}
\newcommand{\setSID}[1]{\renewcommand{\SID}{#1}}

\newcommand{\doHwHeader}[1]{                                         % for homeworks
  \renewcommand{\refreshHeader}{
    \renewcommand{\header}{%
      {\footnotesize \classSection #1 / Problem~\arabic{problem} / draft of \today \hfill \author\ifGradescope~(\SID)\fi~~}}
    \markboth{\header}{\header}
  }
  \renewcommand{\header}{%
    {\footnotesize \classSection #1 / Instructions / draft of \today \hfill \author\ifGradescope~(\SID)\fi~~}}
  \markboth{\header}{\header}
}

% Homework settings, for instructors using the homeworks in a course:
% \classSection is the course name shown in each homework's running head
% (e.g. \renewcommand{\classSection}{CS 3000\xspace}), and \Gradescopetrue turns on
% the instructions for submitting to Gradescope (in its fixed-length mode),
% the hints saying on which page of the PDF each answer should be, and the
% student ID in the running head.
\newcommand{\classSection}{}
\newif\ifGradescope
\Gradescopefalse
\newcommand{\GradescopeComment}[1]{\ifGradescope\comment{#1}\fi}

\pagestyle{myheadings}
\addtolength{\headsep}{0.2in}
\addtolength{\topmargin}{-0.3in}
\addtolength{\textheight}{0.4in}

%%%%%%%%%%%%%%%%%%%%%%%%%%%%%%%%%%%%%%%%%%%%%%%%
%%%%%%%%%%%%%%%%%%%%%%%%%%%%%%%%%%%%%%%%%%%%%%%%

\newcommand{\MF}{\href
  {https://mfleck.cs.illinois.edu/building-blocks/}
  {MF}\xspace}

\newcommand{\KWSlides}{\href
  {https://www.cs.princeton.edu/~wayne/kleinberg-tardos/}
  {KW slides}\xspace}

\newcommand{\LLM}{\href
  {https://courses.csail.mit.edu/6.042/spring18/mcs.pdf}
  {LLM}\xspace}

% \DPV, \KT, \CLRS: textbooks listed (without links) in the table of sources in
% the Preliminaries note; each prints as a link to its entry in that table.
\DeclareRobustCommand{\DPV}{\noteLink{prelim}{source.DPV}{DPV}\xspace}
\DeclareRobustCommand{\KT}{\noteLink{prelim}{source.KT}{KT}\xspace}
\DeclareRobustCommand{\CLRS}{\noteLink{prelim}{source.CLRS}{CLRS}\xspace}
\AtBeginDocument{\pdfstringdefDisableCommands{\def\DPV{DPV }\def\KT{KT }\def\CLRS{CLRS }}}

\newcommand{\JE}{\href
  {https://jeffe.cs.illinois.edu/teaching/algorithms/}
  {JE}\xspace}

\newcommand{\JEGreedy}{\href
  {https://jeffe.cs.illinois.edu/teaching/algorithms/book/04-greedy.pdf}
  {JE Chapter 4}\xspace}

\newcommand{\BBCode}{\href
{https://northeastern.blackboard.com/webapps/cmsmain/webui/courses/CS3000.MERGED.202010/code?action=frameset&subaction=view&uniq=-jj2max&course_id=_2597392_1&mask=\%2Fcourses\%2FCS3000.MERGED.202010}
  {code}\xspace}

% The Python code (src/code).  Each lecture note that has code ends with a
% section listing its files: \codeSection{path, path, ...}.  In the text,
% \pycode{greedy_algorithms/Dijkstra.py} prints the path as a link to its listing
% (in the homeworks, which have no listings, to the file on GitHub), and
% \pycode[text]{...} uses the given text instead.  \pythonCode{path, ...} is the
% standard sentence ending a section whose algorithms have code:
% "Python code: path, ... (Section N.k)."
\newcommand{\codeFolder}{https://github.com/nealeyoung/algs101/tree/main/src/code}
\newcommand{\codeRepo}{https://github.com/nealeyoung/algs101/blob/main/src/code}
\newif\ifcodeListings   % set by \doTitle: this document has the listings
\newcommand{\githubCode}[2][]{\href{\codeRepo/#2}{\ifx\relax#1\relax\nolinkurl{#2}\else#1\fi}}
\newcommand{\pycode}[2][]{%
  \ifcodeListings\hyperref[code: #2]{\ifx\relax#1\relax\nolinkurl{#2}\else#1\fi}%
  \else\githubCode[#1]{#2}\fi}
\makeatletter
\newcommand{\pythonCode}[1]{%
  \def\pythonCode@first{}\def\pythonCode@list{}%
  \forcsvlist{\pythonCode@add}{#1}%
  Python code: \pythonCode@list\ (Section~\ref{code: \pythonCode@first}).}
\newcommand{\pythonCode@add}[1]{%
  \ifx\pythonCode@first\@empty\def\pythonCode@first{#1}\appto\pythonCode@list{\pycode{#1}}%
  \else\appto\pythonCode@list{, \pycode{#1}}\fi}
\makeatother

% Code listings that copy and paste exactly (in PDF viewers that honor the
% ToUnicode maps, e.g. MuPDF-based ones; PDFKit viewers such as Preview and
% Skim drop leading spaces, so each listing's heading links to the file on
% GitHub).  Spaces are typeset as white visible-space glyphs that copy as
% spaces, and hyphens and asterisks as the T1 glyphs, so they copy as themselves.
\usepackage[T1,OT1]{fontenc}   % T1 for the code listings only; OT1 stays the default
\usepackage{listings}
\pdfgentounicode=1
\input glyphtounicode
\pdfglyphtounicode{uni2423}{0020}
\makeatletter
\def\lst@visiblespace{\textcolor{white}{\char32}}
% put an invisible space on each blank line, so that blank lines copy too
% (listings' OnEmptyLine hook runs before the line break, so patch the loop)
\def\lst@DoNewLines{%
    \@whilenum\lst@newlines>\@ne \do {\lst@NewLine \lst@visiblespace}%
    \ifnum\lst@newlines>\z@ \lst@NewLine \fi}
\makeatother
\lstset{language=Python,
  basicstyle=\fontencoding{T1}\fontfamily{lmtt}\selectfont\footnotesize,
  upquote=true, showstringspaces=false, columns=fullflexible, keepspaces=true,
  showspaces=true, literate={-}{{\char45}}1 {*}{{\char42}}1,
  frame=single, framesep=1ex, rulecolor=\color{gray}}
% the listings' paths are relative to the lecture note's folder (in the book, to src/)
\ifdefined\BookMode\newcommand{\codeDir}{code/}\else\newcommand{\codeDir}{../code/}\fi
% \codefile{path}: a listing of src/code/path, headed by a link to it on GitHub
\newcommand{\codefile}[1]{%
  \subsection*{\githubCode{#1}}\label{code: #1}%
  \lstinputlisting{\codeDir#1}}
% \codeSection{path, path, ...}: the lecture note's last section, listing its code
\newcommand{\codeSection}[1]{%
  \newpage
  \section{Python code}\label{\LNtag{\currentLN}: sec: code}\doHeader{Python code}
  This section lists the Python code for this lecture note.
  Click on a file's name to open the file on GitHub.
  Copying code from GitHub is more reliable than copying it from this PDF,
  since some PDF viewers drop the leading spaces (the indentation) of copied lines.
  \forcsvlist{\codefile}{#1}}
% the copyright notice, for title pages and the homework first pages
\newcommand{\licenseURL}{https://github.com/nealeyoung/algs101/blob/main/LICENSE}
\newcommand{\authorURL}{https://www.cs.ucr.edu/\string~neal}
\newcommand{\copyrightNotice}{%
  \copyright\ 2018--2026 \href{\authorURL}{Neal E. Young}.
  Text: CC BY-NC-SA 4.0; code: MIT; except where noted.
  See \url{\licenseURL}.}
% the published site: PDFs of the notes, and the homework assignments and templates
\newcommand{\siteURL}{https://nealeyoung.github.io/algs101/}

\newcommand{\blackboard}{\href
  {https://northeastern.blackboard.com/webapps/blackboard/execute/launcher?type=Course&id=_2597392_1&url=}
  {Blackboard}\xspace}

\newcommand{\JUSlides}{\href
  {http://www.khoury.neu.edu/home/jullman/cs3000f18/schedule.html}
  {Prof.~Ullman's slides}\xspace}

% a URL on its own line, in three layouts
\newcommand{\URL}[1]{\par-- \parbox[t]{0.8\textwidth}{\url{#1}} \medskip\par}
\newcommand{\URLsmall}[1]{\par-- \parbox[t]{0.95\textwidth}{\small\url{#1}} \medskip\par}
\newcommand{\URLplain}[1]{\par\parbox[t]{\textwidth}{\url{#1}}\par}

%%%%%%%%%%%%%%%%%%%%%%%%%%%%%%%%%%%%%%%%%%%%%%%%
%%%%%%%%%%%%%%%%%%%%%%%%%%%%%%%%%%%%%%%%%%%%%%%%

\newcommand{\continued}{\vfill\centerline{\emph{(continued)}}\vfill\newpage}

% load hyperref last to avoid conflicts; long URLs may also break after hyphens
\PassOptionsToPackage{hyphens}{url}
\ifdefined\BookMode\else\usepackage{xr-hyper}\fi   % for \lectureNoteRefs (homeworks), below
\usepackage{hyperref}
\hypersetup{
    colorlinks,
    linkcolor={blue!50!black},
    citecolor={blue!50!black},
    urlcolor={blue!80!black}
}

\providecommand{\answerDirectory}{.}
\newcommand{\inputAnswer}[1]{\refreshHeader\input{\answerDirectory/#1}}
\newcommand{\inputSection}[1]{\input{#1}}

% Refer to a lecture note by its tag (the prefix of its labels), e.g.
% \LN{proofs} prints "Lecture Note 2" (in the book, a link to that chapter).  The table below is the only place the numbers appear.
\makeatletter
\newcommand{\defineLN}[3]{\@namedef{LNnumber@#2}{#1}\@namedef{LNtag@#1}{#2}\@namedef{LNfile@#2}{#3}}
\newcommand{\LNnumber}[1]{%
  \ifcsname LNnumber@#1\endcsname\csname LNnumber@#1\endcsname
  \else\PackageError{macros}{Unknown lecture-note tag `#1'}{}\fi}
\newcommand{\LNtag}[1]{\csname LNtag@#1\endcsname}
\newcommand{\LNfile}[1]{\csname LNfile@#1\endcsname}
\makeatother
\defineLN{1}{prelim}{1_preliminaries}
\defineLN{2}{proofs}{2_proofs}
\defineLN{3}{matching}{3_stable_matching}
\defineLN{4}{dc}{4_divide_and_conquer}
\defineLN{5}{greedy}{5_greedy}
\defineLN{6}{graphs}{6_graph_traversal}
\defineLN{7}{dp}{7_dynamic_programming}
\defineLN{8}{flow}{8_flow_reductions_LP}
\defineLN{9}{np}{9_NP}

% The text is the same in the book and in the standalone lecture notes; only
% links differ.  In the book, links to other lecture notes stay in the book.
% In a standalone note, they open the other note's published PDF.
% \noteLink{tag}{target}{text} links text to \hypertarget{target} in lecture note tag.
% \LN{proofs} prints "Lecture Note 2", linked to that lecture note.
\newcommand{\currentLN}{0}   % in a standalone note, its number (set by \doTitle)
\newcommand{\notePDF}[1]{\siteURL notes/\LNfile{#1}.pdf}
\newcommand{\noteLink}[3]{%
  \ifnum\LNnumber{#1}=\currentLN\relax\hyperlink{#2}{#3}%
  \else\href{\notePDF{#1}\#nameddest=#2}{#3}\fi}
% \bookLink{target}{text} links text to \hypertarget{target} in the book's
% appendices (from a standalone note, in the published book).
\DeclareRobustCommand{\bookLink}[2]{%
  \ifdefined\BookMode\hyperlink{#1}{#2}%
  \else\href{\siteURL lecture_notes_on_algorithms.pdf\#nameddest=#1}{#2}\fi}
\DeclareRobustCommand{\LN}[1]{%
  \ifnum\LNnumber{#1}=\currentLN\relax Lecture Note~\LNnumber{#1}%
  \else\href{\notePDF{#1}}{Lecture Note~\LNnumber{#1}}\fi}
\pdfstringdefDisableCommands{\def\LN#1{Lecture Note \LNnumber{#1}}}

% Homeworks refer to the notes' sections, exercises and lemmas by \ref: each
% homework's main file says \lectureNoteRefs after \input{macros}.  A homework
% folder has a link refs/ to src/refs/, which holds each note's label records
% (made by tools/refs.sh from the notes' builds).  Read with xr-hyper, they make
% \ref{greedy: exercise: ...} give the exercise's number in the published note,
% linked to it there; labels only the book has (its appendices) link to the book.
\ifdefined\BookMode\else
  \makeatletter
  \newcommand{\lectureNoteRefs}{%
    \@for\LN@tag:=prelim,proofs,matching,dc,greedy,graphs,dp,flow,np\do{%
      \externaldocument{refs/\LNfile{\LN@tag}}[\notePDF{\LN@tag}]}%
    \externaldocument{refs/book}[\siteURL lecture_notes_on_algorithms.pdf]}
  % hyperref links a \ref to another document's label as URL#anchor; browsers'
  % PDF viewers need URL#nameddest=anchor (the form \noteLink and \bookLink use)
  \def\Hy@setref@link#1#2#3#4#5#6\@nil#7{%
    \begingroup
      \ifx\\#5\\\toks0={\hyper@@link{#5}{#4}}\else\toks0={\hyper@@link{#5}{nameddest=#4}}\fi
      \toks1=\expandafter{#7{#1}{#2}{#3}{#4}{#5}}%
      \edef\x{\endgroup\the\toks0 {\the\toks1 }}%
    \x}
  \makeatother
\fi

% Book mode: main.tex defines \BookMode before \input{macros} to combine the
% lecture notes into one report-class document, one chapter per note.
% Without \BookMode (each lecture note on its own, and the homeworks) none
% of this applies.
\ifdefined\BookMode
  \usepackage{etoc}

  % links to other lecture notes stay in the book
  \renewcommand{\noteLink}[3]{\hyperlink{#2}{#3}}

  % each chapter is a lecture note, and is called one, as in the standalone notes
  \renewcommand{\chaptername}{Lecture Note}
  \setcounter{secnumdepth}{3}   % number subsubsections, as in the standalone notes

  % global table of contents, for main.tex
  \newcommand{\bookContents}{%
    \hypertarget{TOC}{}%
    \etocsetnexttocdepth{section}%
    \tableofcontents}

  % \doTitle{N}{title}{short title} starts chapter N, with a link to the global
  % table of contents and the chapter's own table of contents
  \renewcommand{\doTitle}[3]{%
    \setcounter{chapter}{\numexpr#1-1\relax}%
    \chapter{\texorpdfstring{#2}{#3}}%
    \label{\LNtag{#1}: chapter}%
    \csgdef{LNshorttitle@#1}{#3}%
    \csgdef{LNtitle@#1}{#2}%
    \renewcommand{\LNTitle}{Lecture Note #1, #3}%
    \codeListingstrue
    \hyperlink{TOC}{$\uparrow$\scriptsize\,\textsc{contents}}%
    {\small
      \hypertarget{TOC-#1}{}%
      \etocsettocstyle{\section*{\contentsname}}{}%
      \etocsetnexttocdepth{subsubsection}%
      \localtableofcontents
    }}

  % the code: each lecture note's \codeSection just records its files, and
  % \codeAppendix (in the book's appendix) lists them all, a section per note
  \makeatletter   % (file names have no spaces, so drop the spaces around them)
  \renewcommand{\codeSection}[1]{\csxdef{LNcode@\arabic{chapter}}{\zap@space#1 \@empty}}
  \makeatother
  % \codeAppendix: first a contents list (each note's topic, then its files,
  % linked to the listings), then a section per note listing its files
  \makeatletter
  \def\LNbasename#1{\LN@base#1/\LN@end}   % a/b/c.py -> c.py
  \def\LN@base#1/#2\LN@end{\ifx\relax#2\relax#1\else\LN@base#2\LN@end\fi}
  \makeatother
  \newcommand{\codeContentsFile}[1]{%
    \edef\LNfilename{\LNbasename{#1}}%
    \item \pycode[\texttt{\expandafter\detokenize\expandafter{\LNfilename}}]{#1}}
  \newcommand{\codeAppendix}{%
    \begin{itemize}
    \foreach \n in {1,...,9} {%
      \ifcsdef{LNcode@\n}{%
        \item \hyperref[code: note \n]{\csuse{LNtitle@\n}} (\LN{\LNtag{\n}})
        \begin{itemize}
        \edef\LNcodelist{\csuse{LNcode@\n}}%
        \expandafter\forcsvlist\expandafter\codeContentsFile\expandafter{\LNcodelist}%
        \end{itemize}}{}}
    \end{itemize}
    \newpage
    \foreach \n in {1,...,9} {%
      \ifcsdef{LNcode@\n}{%
        \section[Lecture Note \n: \csuse{LNtitle@\n}]{%
          \texorpdfstring{\LN{\LNtag{\n}}: \csuse{LNtitle@\n}}{Lecture Note \n: \csuse{LNtitle@\n}}}%
        \label{code: note \n}%
        \edef\LNcodelist{\csuse{LNcode@\n}}%
        \expandafter\forcsvlist\expandafter\codefile\expandafter{\LNcodelist}%
        \newpage}{}}}

  % a lecture note is a chapter: \LN{proofs} prints a link "Lecture Note 2"
  \DeclareRobustCommand{\LN}[1]{\hyperref[#1: chapter]{Lecture Note~\ref*{#1: chapter}}}
  \pdfstringdefDisableCommands{\def\LN#1{Lecture Note \LNnumber{#1}}}

  % running head; "contents" links to the current chapter's table of contents
  \renewcommand{\doHeader}[1]{%
    \renewcommand{\header}{%
      \clap{\raisebox{0pt}{\protect\hyperlink{TOC-\arabic{chapter}}{$\uparrow$\scriptsize\,\textsc{contents}\hspace{1.1in}}}}%
      {\footnotesize \LNTitle{} / Section \thesection, #1}}%
    \markboth{\header}{\header}}

  % number theorems, lemmas, etc. and exercises within chapters (e.g. Lemma 8.1)
  \counterwithin{lemma}{chapter}
  \counterwithin{definition}{chapter}
  \counterwithin{theorem}{chapter}
  \counterwithin{corollary}{chapter}
  \counterwithin{claim}{chapter}
  \counterwithin{observation}{chapter}
  \counterwithin{conjecture}{chapter}
  \counterwithin{exercise}{chapter}
\fi

% %%%%%%%%%%%%%%%%%%%%%%%%%%%%%%%%%%%%%%%%%%%%%%%%%%%%%%%%%%%%%%%%%%%%%%%%
% %%%%%%%%%%%%%%%%%%%%%%%%%%%%%%%%%%%%%%%%%%%%%%%%%%%%%%%%%%%%%%%%%%%%%%%%

\ExplSyntaxOn
\seq_new:N \l_itemize_seq
% \NewDocumentEnvironment{REITEMIZE}{ +b }{
\NewDocumentCommand{\REITEMIZECMD}{ +m }{
  \regex_split:nnN {\c{item}} {#1} \l_itemize_seq
  \seq_pop:NN \l_itemize_seq \l_tmpa_tl% remove the (hopefully) empty first item
  \seq_if_empty:NF \l_itemize_seq {
    \l_tmpa_tl
    \seq_map_inline:Nn \l_itemize_seq {
      \regex_split:nnN {\c{item}} {#1} \l_itemize_seq {
        \regex_extract_once:nnNTF { ^\[([^]]*)\](.*) } { ##1} \l_foo_seq
      { \item[\seq_item:Nn \l_foo_seq {2}]\ITEMIZER{\seq_item:Nn \l_foo_seq {3}} } { \item\ITEMIZER{##1}  }
    }
  }
}
}
\ExplSyntaxOff

\newlength{\ITEMHT}
\setlength{\ITEMHT}{1in}

% \setlength{\fboxsep}{1pt}
\newcommand{\ITEMIZER}[1]{%
  % \framebox{
\adjustbox{valign=t}{\parbox{\textwidth}{\linespread{1.15}\selectfont #1}}~~\adjustbox{valign=t}{\parbox{0pt}{\color{gray}\rule[-\ITEMHT]{0.01pt}{\ITEMHT}}}
% }
  % \begin{tabular}{@{}p{\textwidth}@{\rule[-\ITEMHT]{0pt}{\ITEMHT}}}
  %   #1
  % \end{tabular}
}

\newsavebox\thesmashminipage
\newlength{\thesmashminipagelength}
\newlength{\answerboxlength}
\newlength{\answerboxoverrun}
\newenvironment{smashminipage}[2]
{\setlength\answerboxlength{#2}
\begin{lrbox}{\thesmashminipage}\begin{minipage}[t]{#1}}
{\end{minipage}\end{lrbox}%
\setlength{\thesmashminipagelength}{\dimexpr\ht\thesmashminipage+\dp\thesmashminipage}%
\setlength{\answerboxoverrun}{\dimexpr\thesmashminipagelength-\answerboxlength}%
\ifdim\the\thesmashminipagelength>\the\answerboxlength%
\usebox{\thesmashminipage}%
\smash{\color{red}\hspace*{-0.98\textwidth}\rule[-\answerboxlength]{0.98\textwidth}{1pt}}
\errmessage{Answer exceeds allotted length, please shorten it.}
\else
\usebox{\thesmashminipage}%
\fi
}



\newenvironment{answerBox}[1]{%
  \def\ANSWERBOXHEIGHT{#1}
  \begin{adjustbox}{valign=t}
  \begin{smashminipage}{\textwidth}{\ANSWERBOXHEIGHT}
}{\end{smashminipage}\end{adjustbox}~~~\adjustbox{valign=t}{\parbox{0pt}{\color{gray}\rule[-\ANSWERBOXHEIGHT]{0.01pt}{\ANSWERBOXHEIGHT}\rule[-\ANSWERBOXHEIGHT]{10pt}{0.01pt}}}}

% % \setlist[enumerate]{itemsep=0pt,parsep=0pt,topsep=0pt,leftmargin=100pt}

% \makeatletter 
\LetLtxMacro\SVenumerate\enumerate
\let\endSVenumerate\endenumerate

\newcommand{\injectEnumerate}{
  \let\SValph\alph
  \def\alph{}
  \renewenvironment{enumerate}[1][]
  {\let\alph\SValph \SVenumerate[label=(\alph*),itemsep=2pt,parsep=0pt]\REITEMIZE}{\endREITEMIZE\endSVenumerate}
  % \setlist[enumerate]{before=\Collect@Body\REITEMIZECMD}
}
% \makeatother
\newcommand{\restoreEnumerate}{
  \renewenvironment{enumerate}[1][]
  \let\SValph\alph
  \def\alph{}
  {\let\alph\SValph \SVenumerate[label=(\alph*)]}{\endSVenumerate}
  % \setlist[enumerate]{before=,after=}
}
\newcommand{\injectBlue}{
  \let\SVblue\blue
  \let\endSVblue\endblue
  \def\blue{\SVblue\REITEMIZE}
  \def\endblue{\endREITEMIZE\endSVblue}
}
\newcommand{\restoreBlue}{
  \let\blue\SVblue
  \let\endblue\endSVblue
}
\newcommand{\injectEnumerateBlue}{
  \setlist[enumerate]{before=\injectBlue,after=\restoreBlue,itemsep=2pt,parsep=0pt}
}
\newcommand{\restoreEnumerateBlue}{
  \setlist[enumerate]{before=,after=}
}


% %%%%%%%%%%%%%%%%%%%%%%%%%%%%%%%%%%%%%%%%%%%%%%%%%%%%%%%%%%%%%%%%%%%%%%%%
% %%%%%%%%%%%%%%%%%%%%%%%%%%%%%%%%%%%%%%%%%%%%%%%%%%%%%%%%%%%%%%%%%%%%%%%%

% \NewDocumentEnvironment{env}{ }
% {\Environenv}
% {\endEnvironenv}

% \NewEnviron{REITEMIZE}{\REITEMIZECMD\BODY}

\makeatletter 
\newenvironment{REITEMIZE}{\Collect@Body\REITEMIZECMD}{}
\makeatother 

\newcommand{\NODEone}[1]{\ensuremath{v_{#1}}\xspace}
\newcommand{\NODE}[2]{{\NODEone {#1,#2}}\xspace}
\newcommand{\EDGE}[2]{\ensuremath{(\textsf{#1}, \textsf{#2})}}
\newcommand{\ITEM}{\item}

%%% Local Variables:
%%% mode: latex
%%% End:
 to combine the
% lecture notes into one report-class document, one chapter per note.
% Without \BookMode (each lecture note on its own, and the homeworks) none
% of this applies.
\ifdefined\BookMode
  \usepackage{etoc}

  % links to other lecture notes stay in the book
  \renewcommand{\noteLink}[3]{\hyperlink{#2}{#3}}

  % each chapter is a lecture note, and is called one, as in the standalone notes
  \renewcommand{\chaptername}{Lecture Note}
  \setcounter{secnumdepth}{3}   % number subsubsections, as in the standalone notes

  % global table of contents, for main.tex
  \newcommand{\bookContents}{%
    \hypertarget{TOC}{}%
    \etocsetnexttocdepth{section}%
    \tableofcontents}

  % \doTitle{N}{title}{short title} starts chapter N, with a link to the global
  % table of contents and the chapter's own table of contents
  \renewcommand{\doTitle}[3]{%
    \setcounter{chapter}{\numexpr#1-1\relax}%
    \chapter{\texorpdfstring{#2}{#3}}%
    \label{\LNtag{#1}: chapter}%
    \csgdef{LNshorttitle@#1}{#3}%
    \csgdef{LNtitle@#1}{#2}%
    \renewcommand{\LNTitle}{Lecture Note #1, #3}%
    \codeListingstrue
    \hyperlink{TOC}{$\uparrow$\scriptsize\,\textsc{contents}}%
    {\small
      \hypertarget{TOC-#1}{}%
      \etocsettocstyle{\section*{\contentsname}}{}%
      \etocsetnexttocdepth{subsubsection}%
      \localtableofcontents
    }}

  % the code: each lecture note's \codeSection just records its files, and
  % \codeAppendix (in the book's appendix) lists them all, a section per note
  \makeatletter   % (file names have no spaces, so drop the spaces around them)
  \renewcommand{\codeSection}[1]{\csxdef{LNcode@\arabic{chapter}}{\zap@space#1 \@empty}}
  \makeatother
  % \codeAppendix: first a contents list (each note's topic, then its files,
  % linked to the listings), then a section per note listing its files
  \makeatletter
  \def\LNbasename#1{\LN@base#1/\LN@end}   % a/b/c.py -> c.py
  \def\LN@base#1/#2\LN@end{\ifx\relax#2\relax#1\else\LN@base#2\LN@end\fi}
  \makeatother
  \newcommand{\codeContentsFile}[1]{%
    \edef\LNfilename{\LNbasename{#1}}%
    \item \pycode[\texttt{\expandafter\detokenize\expandafter{\LNfilename}}]{#1}}
  \newcommand{\codeAppendix}{%
    \begin{itemize}
    \foreach \n in {1,...,9} {%
      \ifcsdef{LNcode@\n}{%
        \item \hyperref[code: note \n]{\csuse{LNtitle@\n}} (\LN{\LNtag{\n}})
        \begin{itemize}
        \edef\LNcodelist{\csuse{LNcode@\n}}%
        \expandafter\forcsvlist\expandafter\codeContentsFile\expandafter{\LNcodelist}%
        \end{itemize}}{}}
    \end{itemize}
    \newpage
    \foreach \n in {1,...,9} {%
      \ifcsdef{LNcode@\n}{%
        \section[Lecture Note \n: \csuse{LNtitle@\n}]{%
          \texorpdfstring{\LN{\LNtag{\n}}: \csuse{LNtitle@\n}}{Lecture Note \n: \csuse{LNtitle@\n}}}%
        \label{code: note \n}%
        \edef\LNcodelist{\csuse{LNcode@\n}}%
        \expandafter\forcsvlist\expandafter\codefile\expandafter{\LNcodelist}%
        \newpage}{}}}

  % a lecture note is a chapter: \LN{proofs} prints a link "Lecture Note 2"
  \DeclareRobustCommand{\LN}[1]{\hyperref[#1: chapter]{Lecture Note~\ref*{#1: chapter}}}
  \pdfstringdefDisableCommands{\def\LN#1{Lecture Note \LNnumber{#1}}}

  % running head; "contents" links to the current chapter's table of contents
  \renewcommand{\doHeader}[1]{%
    \renewcommand{\header}{%
      \clap{\raisebox{0pt}{\protect\hyperlink{TOC-\arabic{chapter}}{$\uparrow$\scriptsize\,\textsc{contents}\hspace{1.1in}}}}%
      {\footnotesize \LNTitle{} / Section \thesection, #1}}%
    \markboth{\header}{\header}}

  % number theorems, lemmas, etc. and exercises within chapters (e.g. Lemma 8.1)
  \counterwithin{lemma}{chapter}
  \counterwithin{definition}{chapter}
  \counterwithin{theorem}{chapter}
  \counterwithin{corollary}{chapter}
  \counterwithin{claim}{chapter}
  \counterwithin{observation}{chapter}
  \counterwithin{conjecture}{chapter}
  \counterwithin{exercise}{chapter}
\fi

% %%%%%%%%%%%%%%%%%%%%%%%%%%%%%%%%%%%%%%%%%%%%%%%%%%%%%%%%%%%%%%%%%%%%%%%%
% %%%%%%%%%%%%%%%%%%%%%%%%%%%%%%%%%%%%%%%%%%%%%%%%%%%%%%%%%%%%%%%%%%%%%%%%

\ExplSyntaxOn
\seq_new:N \l_itemize_seq
% \NewDocumentEnvironment{REITEMIZE}{ +b }{
\NewDocumentCommand{\REITEMIZECMD}{ +m }{
  \regex_split:nnN {\c{item}} {#1} \l_itemize_seq
  \seq_pop:NN \l_itemize_seq \l_tmpa_tl% remove the (hopefully) empty first item
  \seq_if_empty:NF \l_itemize_seq {
    \l_tmpa_tl
    \seq_map_inline:Nn \l_itemize_seq {
      \regex_split:nnN {\c{item}} {#1} \l_itemize_seq {
        \regex_extract_once:nnNTF { ^\[([^]]*)\](.*) } { ##1} \l_foo_seq
      { \item[\seq_item:Nn \l_foo_seq {2}]\ITEMIZER{\seq_item:Nn \l_foo_seq {3}} } { \item\ITEMIZER{##1}  }
    }
  }
}
}
\ExplSyntaxOff

\newlength{\ITEMHT}
\setlength{\ITEMHT}{1in}

% \setlength{\fboxsep}{1pt}
\newcommand{\ITEMIZER}[1]{%
  % \framebox{
\adjustbox{valign=t}{\parbox{\textwidth}{\linespread{1.15}\selectfont #1}}~~\adjustbox{valign=t}{\parbox{0pt}{\color{gray}\rule[-\ITEMHT]{0.01pt}{\ITEMHT}}}
% }
  % \begin{tabular}{@{}p{\textwidth}@{\rule[-\ITEMHT]{0pt}{\ITEMHT}}}
  %   #1
  % \end{tabular}
}

\newsavebox\thesmashminipage
\newlength{\thesmashminipagelength}
\newlength{\answerboxlength}
\newlength{\answerboxoverrun}
\newenvironment{smashminipage}[2]
{\setlength\answerboxlength{#2}
\begin{lrbox}{\thesmashminipage}\begin{minipage}[t]{#1}}
{\end{minipage}\end{lrbox}%
\setlength{\thesmashminipagelength}{\dimexpr\ht\thesmashminipage+\dp\thesmashminipage}%
\setlength{\answerboxoverrun}{\dimexpr\thesmashminipagelength-\answerboxlength}%
\ifdim\the\thesmashminipagelength>\the\answerboxlength%
\usebox{\thesmashminipage}%
\smash{\color{red}\hspace*{-0.98\textwidth}\rule[-\answerboxlength]{0.98\textwidth}{1pt}}
\errmessage{Answer exceeds allotted length, please shorten it.}
\else
\usebox{\thesmashminipage}%
\fi
}



\newenvironment{answerBox}[1]{%
  \def\ANSWERBOXHEIGHT{#1}
  \begin{adjustbox}{valign=t}
  \begin{smashminipage}{\textwidth}{\ANSWERBOXHEIGHT}
}{\end{smashminipage}\end{adjustbox}~~~\adjustbox{valign=t}{\parbox{0pt}{\color{gray}\rule[-\ANSWERBOXHEIGHT]{0.01pt}{\ANSWERBOXHEIGHT}\rule[-\ANSWERBOXHEIGHT]{10pt}{0.01pt}}}}

% % \setlist[enumerate]{itemsep=0pt,parsep=0pt,topsep=0pt,leftmargin=100pt}

% \makeatletter 
\LetLtxMacro\SVenumerate\enumerate
\let\endSVenumerate\endenumerate

\newcommand{\injectEnumerate}{
  \let\SValph\alph
  \def\alph{}
  \renewenvironment{enumerate}[1][]
  {\let\alph\SValph \SVenumerate[label=(\alph*),itemsep=2pt,parsep=0pt]\REITEMIZE}{\endREITEMIZE\endSVenumerate}
  % \setlist[enumerate]{before=\Collect@Body\REITEMIZECMD}
}
% \makeatother
\newcommand{\restoreEnumerate}{
  \renewenvironment{enumerate}[1][]
  \let\SValph\alph
  \def\alph{}
  {\let\alph\SValph \SVenumerate[label=(\alph*)]}{\endSVenumerate}
  % \setlist[enumerate]{before=,after=}
}
\newcommand{\injectBlue}{
  \let\SVblue\blue
  \let\endSVblue\endblue
  \def\blue{\SVblue\REITEMIZE}
  \def\endblue{\endREITEMIZE\endSVblue}
}
\newcommand{\restoreBlue}{
  \let\blue\SVblue
  \let\endblue\endSVblue
}
\newcommand{\injectEnumerateBlue}{
  \setlist[enumerate]{before=\injectBlue,after=\restoreBlue,itemsep=2pt,parsep=0pt}
}
\newcommand{\restoreEnumerateBlue}{
  \setlist[enumerate]{before=,after=}
}


% %%%%%%%%%%%%%%%%%%%%%%%%%%%%%%%%%%%%%%%%%%%%%%%%%%%%%%%%%%%%%%%%%%%%%%%%
% %%%%%%%%%%%%%%%%%%%%%%%%%%%%%%%%%%%%%%%%%%%%%%%%%%%%%%%%%%%%%%%%%%%%%%%%

% \NewDocumentEnvironment{env}{ }
% {\Environenv}
% {\endEnvironenv}

% \NewEnviron{REITEMIZE}{\REITEMIZECMD\BODY}

\makeatletter 
\newenvironment{REITEMIZE}{\Collect@Body\REITEMIZECMD}{}
\makeatother 

\newcommand{\NODEone}[1]{\ensuremath{v_{#1}}\xspace}
\newcommand{\NODE}[2]{{\NODEone {#1,#2}}\xspace}
\newcommand{\EDGE}[2]{\ensuremath{(\textsf{#1}, \textsf{#2})}}
\newcommand{\ITEM}{\item}

%%% Local Variables:
%%% mode: latex
%%% End:
.  A homework
% folder has a link refs/ to src/refs/, which holds each note's label records
% (made by tools/refs.sh from the notes' builds).  Read with xr-hyper, they make
% \ref{greedy: exercise: ...} give the exercise's number in the published note,
% linked to it there; labels only the book has (its appendices) link to the book.
\ifdefined\BookMode\else
  \makeatletter
  \newcommand{\lectureNoteRefs}{%
    \@for\LN@tag:=prelim,proofs,matching,dc,greedy,graphs,dp,flow,np\do{%
      \externaldocument{refs/\LNfile{\LN@tag}}[\notePDF{\LN@tag}]}%
    \externaldocument{refs/book}[\siteURL lecture_notes_on_algorithms.pdf]}
  % hyperref links a \ref to another document's label as URL#anchor; browsers'
  % PDF viewers need URL#nameddest=anchor (the form \noteLink and \bookLink use)
  \def\Hy@setref@link#1#2#3#4#5#6\@nil#7{%
    \begingroup
      \ifx\\#5\\\toks0={\hyper@@link{#5}{#4}}\else\toks0={\hyper@@link{#5}{nameddest=#4}}\fi
      \toks1=\expandafter{#7{#1}{#2}{#3}{#4}{#5}}%
      \edef\x{\endgroup\the\toks0 {\the\toks1 }}%
    \x}
  \makeatother
\fi

% Book mode: main.tex defines \BookMode before %%%% DON'T EDIT THIS FILE %%%%

\usepackage{xifthen}% \isempty
\usepackage{etoolbox}% \patchcmd
% \usepackage{stringstrings}% \substring

\usepackage{amsmath, amsthm, amssymb}
\usepackage[margin=1in]{geometry}
\usepackage{xspace}
\usepackage{graphicx}
\usepackage[normalem]{ulem}
\usepackage{adjustbox}
\usepackage{xcolor}
\usepackage{changepage}
\usepackage{framed}
\usepackage{multicol}
\usepackage{enumitem}
% \usepackage{listings} % lstlistings environment

\usepackage{makecmds}           %\provideenvironment

\usepackage{xparse,environ}
\usepackage{letltxmacro}       % \LetLtxMacro
\usepackage{wrapfig}
\usepackage{tikz}

\usetikzlibrary{shapes}
\usetikzlibrary{positioning}
\usetikzlibrary{shapes.misc}
\usetikzlibrary{arrows}

\newtheorem{lemma}{Lemma}
\newtheorem{definition}{Definition}
\newtheorem{theorem}{Theorem}
\newtheorem{corollary}{Corollary}
\newtheorem{claim}{Claim}
\newtheorem{observation}{Observation}
\newtheorem{conjecture}{Conjecture}

% A lemma or theorem restated later with its original number (set with
% \setcounter) goes in a restatement environment, which gives it its own link
% anchor (otherwise it would duplicate the original's).
\newenvironment{restatement}
  {\def\theHlemma{restated.\thelemma}\def\theHtheorem{restated.\thetheorem}}{}

\newenvironment{myframed}
{\smallskip\par\noindent\begin{minipage}{\textwidth}\begin{framed}\vspace*{-5pt}}{\vspace*{-5pt}\end{framed}\end{minipage}\smallskip}


\definecolor{darkblue}{RGB}{0,0,200}

 
\let\oldParagraph\paragraph
\providecommand{\gradescopeColor}{\color{black}}
\renewcommand{\paragraph}[1]{\oldParagraph{\gradescopeColor #1}}
\providecommand{\gradescopeInit}{}

\providecommand{\BLUE}[1]{{\color{darkblue} #1}}
\providecommand{\REPLACEME}{\textcolor{red}{\ensuremath{\star \cdots \star}}\xspace}
\provideenvironment{blue}{\color{darkblue}}{}

\newcommand{\dist}[2]{\ensuremath{\operatorname{\sf dist}(#1, #2)}}
\newcommand{\depth}[2]{\ensuremath{\operatorname{\sf depth}_{#1}(#2)}}
\newcommand{\cost}[1]{\ensuremath{\operatorname{\sf cost}(#1)}}
\newcommand{\weight}[1]{\ensuremath{\operatorname{\sf wt}(#1)}}
% \newcommand{\val}[1]{\ensuremath{\operatorname{\sf val}(#1)}}
\newcommand{\eps}{\varepsilon}
\newcommand{\capac}{\textsf{cap}}
\newcommand{\capOf}[1]{\ensuremath{\operatorname{\sf cap}(#1)}}
\newcommand{\rescap}[1]{\ensuremath{\operatorname{\sf cap}_f(#1)}}
\newcommand{\into}[1]{\ensuremath{\operatorname{\sf in}(#1)}}
\newcommand{\out}[1]{\ensuremath{\operatorname{\sf out}(#1)}}
\newenvironment{LP}[1]{
  \begin{array}[t]{@{}l@{}}
    \displaystyle #1 \\
    \begin{aligned}\\[-10pt]}{\end{aligned}\end{array}}

\newcommand{\opt}[1]{\ensuremath{\operatorname{\sf opt}(#1)}}
\newcommand{\Opt}[1]{\ensuremath{\operatorname{\sf Opt}(#1)}}
\newcommand{\half}{\frac{1}{2}}

\newcommand{\E}{\operatorname{E}}
\newcommand{\giv}{\,|\,}
\newcommand{\Z}{\mathbb{Z}}
\newcommand{\Zp}{\mathbb{Z}_{\ge 0}}
\newcommand{\R}{\mathbb{R}}
% \newcommand{\Rp}{\mathbb{R}_{\ge 0}}
\newcommand{\Rp}{\mathbb{R}_{\tiny +}}
\newcommand{\N}{\mathbb{N}}
\newcommand{\Np}{\mathbb{N}_{\ge 0}}
\newcommand{\calS}{{\cal S}}

\newcommand{\OP}[1]{\text{\small\sc #1}\xspace}
\newcommand{\Insert}{\OP{Insert}}
\newcommand{\FindMin}{\OP{Find-Min}}
\newcommand{\DeleteMin}{\OP{Delete-Min}}
\newcommand{\DecreaseKey}{\OP{Decrease}}
\newcommand{\Merge}{\OP{Merge}}
\newcommand{\Cut}{\OP{Cut}}
% \newcommand{\prob}[1]{\text{\sc #1}}
% a problem name; in running text, lines may break inside it
\DeclareRobustCommand{\prob}[1]{\ifmmode\text{#1}\else#1\fi}
\newcommand{\encode}[1]{\langle #1 \rangle}

\newcommand{\best}[1]{\text{\sf best}(#1)}
\newcommand{\worst}[1]{\text{\sf worst}(#1)}
\newcommand{\floor}[1]{\lfloor #1 \rfloor}
\newcommand{\ceil}[1]{\lceil #1 \rceil}

\setlength{\parskip}{2pt}
\setlength{\parindent}{1em}

% \setlength\fboxsep{0pt}

%%% longFormProof, shortFormProof, block environments
%%% \step

\newlist{steps}{enumerate}{7}
\setlist[steps]{wide,
  topsep=1pt,
  parsep=1pt,
  partopsep=1pt,
  itemsep=0pt,
  align=left,
  labelindent=0pt,
  itemindent=0pt, 
}
\setlist[steps,1]{
  label=\arabic*.,
  ref=\arabic*
}
\setlist[steps,2]{
  label=\arabic{stepsi}.\arabic*.,
  ref=\arabic{stepsi}.\arabic*
}
\setlist[steps,3]{
  label=\arabic{stepsi}.\arabic{stepsii}.\arabic*.,
  ref=\arabic{stepsi}.\arabic{stepsii}.\arabic*
}
\setlist[steps,4]{
  label=\arabic{stepsi}.\arabic{stepsii}.\arabic{stepsiii}.\arabic*.,
  ref=\arabic{stepsi}.\arabic{stepsii}.\arabic{stepsiii}.\arabic*
}
\setlist[steps,5]{
  label=\arabic{stepsi}.\arabic{stepsii}.\arabic{stepsiii}.\arabic{stepsiv}.\arabic*.,
  ref=\arabic{stepsi}.\arabic{stepsii}.\arabic{stepsiii}.\arabic{stepsiv}.\arabic*
}
\setlist[steps,6]{
  label=\arabic{stepsi}.\arabic{stepsii}.\arabic{stepsiii}.\arabic{stepsiv}.\arabic{stepsv}.\arabic*.,
  ref=\arabic{stepsi}.\arabic{stepsii}.\arabic{stepsiii}.\arabic{stepsiv}.\arabic{stepsv}.\arabic*
}
\setlist[steps,7]{
  label=\arabic{stepsi}.\arabic{stepsii}.\arabic{stepsiii}.\arabic{stepsiv}. \arabic{stepsv}. \arabic{stepsvi}.\arabic*.,
  ref=\arabic{stepsi}.\arabic{stepsii}.\arabic{stepsiii}.\arabic{stepsiv}.\arabic{stepsv}. \arabic{stepsvi}.\arabic*
}

\newcommand{\step}[1][]{\item\maybeLabel{#1}}

\makeatletter
% \skipitem steps the list counter itself (for block labels).  Give each one
% its own hyperref link target, since step numbers repeat from proof to proof.
\newcounter{skipitem}
\newcommand{\skipitem}{\stepcounter{skipitem}\refstepcounter{\@enumctr}}
\renewcommand{\theHstepsi}{skipitem.\arabic{skipitem}}
\renewcommand{\theHstepsii}{skipitem.\arabic{skipitem}}
\renewcommand{\theHstepsiii}{skipitem.\arabic{skipitem}}
\renewcommand{\theHstepsiv}{skipitem.\arabic{skipitem}}
\renewcommand{\theHstepsv}{skipitem.\arabic{skipitem}}
\renewcommand{\theHstepsvi}{skipitem.\arabic{skipitem}}
\renewcommand{\theHstepsvii}{skipitem.\arabic{skipitem}}
\makeatother

\newcommand{\maybeLabel}[1]{\ifthenelse{\isempty{#1}}{}{\label{#1}}}

\newenvironment{longFormProof}[1][Proof (long form).]
{\begin{proof}[#1]~\par\begin{steps}}{\end{steps}\vspace*{-3ex}\end{proof}}

\newenvironment{shortFormProof}[1][Proof (short form).]
{\begin{proof}[#1]}{\end{proof}}

\newenvironment{longFormSubProof}[1][Proof.]
{\begin{adjustwidth}{1em}{0em}\begin{proof}[#1]~\par\vspace*{-2pt}\skipitem\begin{steps}}
{\end{steps}\vspace*{-1ex}\end{proof}\end{adjustwidth}\smallskip}

\newenvironment{shortFormSubProof}[1][Proof.]
{\begin{proof}[#1]}
{\end{proof}\vspace*{-3pt}}

\newenvironment{block}[1][]
{\skipitem\maybeLabel{#1}\item[] \begin{steps}}{\end{steps}}

\newcommand{\comment}[1]{\hfill{\footnotesize\emph{#1}}}
\newcommand{\lineacross}{\par\vspace*{-0.7\baselineskip}\noindent\hrulefill\par}

%%% problem

% \newcounter{problem}
% \newcommand{\problem}[1][]{}
% {\ifthenelse{\isempty{#1}}{\refstepcounter{problem}}{}%
%   \paragraph{Problem \ifthenelse{\isempty{#1}}
%     {\bf\arabic{problem}.}
%     {#1}} }

\newcommand{\problemNumbered}[1]{
  \setcounter{problem}{#1}%
  \addtocounter{problem}{-1}%
  \problem
}
\newenvironment{problemTheorem}{%
  \setcounter{theorem}{\value{problem}}%
  \addtocounter{theorem}{-1}%
  \begin{theorem}[Problem \arabic{problem}]
  }{\end{theorem}}

\newcounter{problem}
\newcommand{\problem}[1][]
{\ifthenelse{\isempty{#1}}
  {\refstepcounter{problem}\paragraph{Problem~\arabic{problem}.}}
  {\paragraph{Problem~#1.}}\refreshHeader} 

\newcounter{exercise}
\newcommand{\exercise}[1][]
{\ifthenelse{\isempty{#1}}
  {\refstepcounter{exercise}\paragraph{Exercise~\theexercise.}}
  {\paragraph{Exercise~#1.}}}

\newcommand{\solution}{\paragraph{Solution.}} 

\newenvironment{mylist}[1][]
{\begin{list}{#1}{
      \setlength{\itemsep}{-1pt}
      \setlength{\labelwidth}{0pt}
      \setlength{\leftmargin}{1.25em}
    }}
{\end{list}}

\newcommand{\LNTitle}{\textcolor{red}{UNDEFINED LNTitle}}

% In a standalone lecture note's table of contents, leave room for numbers
% like 7.10 (in bold) and 7.10.2.
\makeatletter
\newcommand{\LNtocWidths}{%
  \patchcmd{\l@section}{1.5em}{2.8em}{}{\PackageWarning{macros}{Could not widen section numbers in the contents}}%
  \renewcommand*{\l@subsection}{\@dottedtocline{2}{2.8em}{3.2em}}%
  \renewcommand*{\l@subsubsection}{\@dottedtocline{3}{6.0em}{4.1em}}}
\makeatother

\newcommand{\doTitle}[3]{
  \pagenumbering{roman}
  \renewcommand{\LNTitle}{Lecture Note #1, #3}
  \renewcommand{\currentLN}{#1}
  \codeListingstrue
  % number sections, theorems etc., exercises, figures, tables and equations
  % within the lecture note (e.g. Section 8.2, Lemma 8.6), as in the book
  \renewcommand{\thesection}{#1.\arabic{section}}
  \renewcommand{\thelemma}{#1.\arabic{lemma}}
  \renewcommand{\thedefinition}{#1.\arabic{definition}}
  \renewcommand{\thetheorem}{#1.\arabic{theorem}}
  \renewcommand{\thecorollary}{#1.\arabic{corollary}}
  \renewcommand{\theclaim}{#1.\arabic{claim}}
  \renewcommand{\theobservation}{#1.\arabic{observation}}
  \renewcommand{\theconjecture}{#1.\arabic{conjecture}}
  \renewcommand{\theexercise}{#1.\arabic{exercise}}
  \renewcommand{\thefigure}{#1.\arabic{figure}}
  \renewcommand{\thetable}{#1.\arabic{table}}
  \renewcommand{\theequation}{#1.\arabic{equation}}
  \LNtocWidths
  \title{Lecture Note #1\\ #2}
  % a note's \LNabstract, if it defines one before \doTitle, goes right below the title
  \date{\ifdefined\LNabstract\parbox{0.85\textwidth}{\normalsize\normalfont\LNabstract}\\[3ex]\fi
    draft of \today, by \href{\authorURL}{N.~Young}\\[2ex] {\small\copyrightNotice}}
  \maketitle
  {\small
    \hypertarget{TOC}{}
    \tableofcontents
  }
  \thispagestyle{empty}
  % the text starts on page 1, after the title and contents (numbered i, ii, ...)
  \clearpage
  \pagenumbering{arabic}
}

\def\clap#1{\hbox to 0pt{\hss#1\hss}}

\newcommand{\refreshHeader}{}
\newcommand{\header}{}

\newcommand{\doHeader}[1]{      % for lecture notes
  \renewcommand{\header}{%
    \clap{\raisebox{0pt}{\protect\hyperlink{TOC}{$\uparrow$\scriptsize\,\textsc{contents}\hspace{1.1in}}}}%
    {\footnotesize \LNTitle{} / Section \thesection, #1}} %, #1, draft of \today}}
  \markboth{\header}{\header}
}

\newcommand{\SID}{}
\newcommand{\setAuthor}[1]{\renewcommand{\author}{#1}}
\newcommand{\setSID}[1]{\renewcommand{\SID}{#1}}

\newcommand{\doHwHeader}[1]{                                         % for homeworks
  \renewcommand{\refreshHeader}{
    \renewcommand{\header}{%
      {\footnotesize \classSection #1 / Problem~\arabic{problem} / draft of \today \hfill \author\ifGradescope~(\SID)\fi~~}}
    \markboth{\header}{\header}
  }
  \renewcommand{\header}{%
    {\footnotesize \classSection #1 / Instructions / draft of \today \hfill \author\ifGradescope~(\SID)\fi~~}}
  \markboth{\header}{\header}
}

% Homework settings, for instructors using the homeworks in a course:
% \classSection is the course name shown in each homework's running head
% (e.g. \renewcommand{\classSection}{CS 3000\xspace}), and \Gradescopetrue turns on
% the instructions for submitting to Gradescope (in its fixed-length mode),
% the hints saying on which page of the PDF each answer should be, and the
% student ID in the running head.
\newcommand{\classSection}{}
\newif\ifGradescope
\Gradescopefalse
\newcommand{\GradescopeComment}[1]{\ifGradescope\comment{#1}\fi}

\pagestyle{myheadings}
\addtolength{\headsep}{0.2in}
\addtolength{\topmargin}{-0.3in}
\addtolength{\textheight}{0.4in}

%%%%%%%%%%%%%%%%%%%%%%%%%%%%%%%%%%%%%%%%%%%%%%%%
%%%%%%%%%%%%%%%%%%%%%%%%%%%%%%%%%%%%%%%%%%%%%%%%

\newcommand{\MF}{\href
  {https://mfleck.cs.illinois.edu/building-blocks/}
  {MF}\xspace}

\newcommand{\KWSlides}{\href
  {https://www.cs.princeton.edu/~wayne/kleinberg-tardos/}
  {KW slides}\xspace}

\newcommand{\LLM}{\href
  {https://courses.csail.mit.edu/6.042/spring18/mcs.pdf}
  {LLM}\xspace}

% \DPV, \KT, \CLRS: textbooks listed (without links) in the table of sources in
% the Preliminaries note; each prints as a link to its entry in that table.
\DeclareRobustCommand{\DPV}{\noteLink{prelim}{source.DPV}{DPV}\xspace}
\DeclareRobustCommand{\KT}{\noteLink{prelim}{source.KT}{KT}\xspace}
\DeclareRobustCommand{\CLRS}{\noteLink{prelim}{source.CLRS}{CLRS}\xspace}
\AtBeginDocument{\pdfstringdefDisableCommands{\def\DPV{DPV }\def\KT{KT }\def\CLRS{CLRS }}}

\newcommand{\JE}{\href
  {https://jeffe.cs.illinois.edu/teaching/algorithms/}
  {JE}\xspace}

\newcommand{\JEGreedy}{\href
  {https://jeffe.cs.illinois.edu/teaching/algorithms/book/04-greedy.pdf}
  {JE Chapter 4}\xspace}

\newcommand{\BBCode}{\href
{https://northeastern.blackboard.com/webapps/cmsmain/webui/courses/CS3000.MERGED.202010/code?action=frameset&subaction=view&uniq=-jj2max&course_id=_2597392_1&mask=\%2Fcourses\%2FCS3000.MERGED.202010}
  {code}\xspace}

% The Python code (src/code).  Each lecture note that has code ends with a
% section listing its files: \codeSection{path, path, ...}.  In the text,
% \pycode{greedy_algorithms/Dijkstra.py} prints the path as a link to its listing
% (in the homeworks, which have no listings, to the file on GitHub), and
% \pycode[text]{...} uses the given text instead.  \pythonCode{path, ...} is the
% standard sentence ending a section whose algorithms have code:
% "Python code: path, ... (Section N.k)."
\newcommand{\codeFolder}{https://github.com/nealeyoung/algs101/tree/main/src/code}
\newcommand{\codeRepo}{https://github.com/nealeyoung/algs101/blob/main/src/code}
\newif\ifcodeListings   % set by \doTitle: this document has the listings
\newcommand{\githubCode}[2][]{\href{\codeRepo/#2}{\ifx\relax#1\relax\nolinkurl{#2}\else#1\fi}}
\newcommand{\pycode}[2][]{%
  \ifcodeListings\hyperref[code: #2]{\ifx\relax#1\relax\nolinkurl{#2}\else#1\fi}%
  \else\githubCode[#1]{#2}\fi}
\makeatletter
\newcommand{\pythonCode}[1]{%
  \def\pythonCode@first{}\def\pythonCode@list{}%
  \forcsvlist{\pythonCode@add}{#1}%
  Python code: \pythonCode@list\ (Section~\ref{code: \pythonCode@first}).}
\newcommand{\pythonCode@add}[1]{%
  \ifx\pythonCode@first\@empty\def\pythonCode@first{#1}\appto\pythonCode@list{\pycode{#1}}%
  \else\appto\pythonCode@list{, \pycode{#1}}\fi}
\makeatother

% Code listings that copy and paste exactly (in PDF viewers that honor the
% ToUnicode maps, e.g. MuPDF-based ones; PDFKit viewers such as Preview and
% Skim drop leading spaces, so each listing's heading links to the file on
% GitHub).  Spaces are typeset as white visible-space glyphs that copy as
% spaces, and hyphens and asterisks as the T1 glyphs, so they copy as themselves.
\usepackage[T1,OT1]{fontenc}   % T1 for the code listings only; OT1 stays the default
\usepackage{listings}
\pdfgentounicode=1
\input glyphtounicode
\pdfglyphtounicode{uni2423}{0020}
\makeatletter
\def\lst@visiblespace{\textcolor{white}{\char32}}
% put an invisible space on each blank line, so that blank lines copy too
% (listings' OnEmptyLine hook runs before the line break, so patch the loop)
\def\lst@DoNewLines{%
    \@whilenum\lst@newlines>\@ne \do {\lst@NewLine \lst@visiblespace}%
    \ifnum\lst@newlines>\z@ \lst@NewLine \fi}
\makeatother
\lstset{language=Python,
  basicstyle=\fontencoding{T1}\fontfamily{lmtt}\selectfont\footnotesize,
  upquote=true, showstringspaces=false, columns=fullflexible, keepspaces=true,
  showspaces=true, literate={-}{{\char45}}1 {*}{{\char42}}1,
  frame=single, framesep=1ex, rulecolor=\color{gray}}
% the listings' paths are relative to the lecture note's folder (in the book, to src/)
\ifdefined\BookMode\newcommand{\codeDir}{code/}\else\newcommand{\codeDir}{../code/}\fi
% \codefile{path}: a listing of src/code/path, headed by a link to it on GitHub
\newcommand{\codefile}[1]{%
  \subsection*{\githubCode{#1}}\label{code: #1}%
  \lstinputlisting{\codeDir#1}}
% \codeSection{path, path, ...}: the lecture note's last section, listing its code
\newcommand{\codeSection}[1]{%
  \newpage
  \section{Python code}\label{\LNtag{\currentLN}: sec: code}\doHeader{Python code}
  This section lists the Python code for this lecture note.
  Click on a file's name to open the file on GitHub.
  Copying code from GitHub is more reliable than copying it from this PDF,
  since some PDF viewers drop the leading spaces (the indentation) of copied lines.
  \forcsvlist{\codefile}{#1}}
% the copyright notice, for title pages and the homework first pages
\newcommand{\licenseURL}{https://github.com/nealeyoung/algs101/blob/main/LICENSE}
\newcommand{\authorURL}{https://www.cs.ucr.edu/\string~neal}
\newcommand{\copyrightNotice}{%
  \copyright\ 2018--2026 \href{\authorURL}{Neal E. Young}.
  Text: CC BY-NC-SA 4.0; code: MIT; except where noted.
  See \url{\licenseURL}.}
% the published site: PDFs of the notes, and the homework assignments and templates
\newcommand{\siteURL}{https://nealeyoung.github.io/algs101/}

\newcommand{\blackboard}{\href
  {https://northeastern.blackboard.com/webapps/blackboard/execute/launcher?type=Course&id=_2597392_1&url=}
  {Blackboard}\xspace}

\newcommand{\JUSlides}{\href
  {http://www.khoury.neu.edu/home/jullman/cs3000f18/schedule.html}
  {Prof.~Ullman's slides}\xspace}

% a URL on its own line, in three layouts
\newcommand{\URL}[1]{\par-- \parbox[t]{0.8\textwidth}{\url{#1}} \medskip\par}
\newcommand{\URLsmall}[1]{\par-- \parbox[t]{0.95\textwidth}{\small\url{#1}} \medskip\par}
\newcommand{\URLplain}[1]{\par\parbox[t]{\textwidth}{\url{#1}}\par}

%%%%%%%%%%%%%%%%%%%%%%%%%%%%%%%%%%%%%%%%%%%%%%%%
%%%%%%%%%%%%%%%%%%%%%%%%%%%%%%%%%%%%%%%%%%%%%%%%

\newcommand{\continued}{\vfill\centerline{\emph{(continued)}}\vfill\newpage}

% load hyperref last to avoid conflicts; long URLs may also break after hyphens
\PassOptionsToPackage{hyphens}{url}
\ifdefined\BookMode\else\usepackage{xr-hyper}\fi   % for \lectureNoteRefs (homeworks), below
\usepackage{hyperref}
\hypersetup{
    colorlinks,
    linkcolor={blue!50!black},
    citecolor={blue!50!black},
    urlcolor={blue!80!black}
}

\providecommand{\answerDirectory}{.}
\newcommand{\inputAnswer}[1]{\refreshHeader\input{\answerDirectory/#1}}
\newcommand{\inputSection}[1]{\input{#1}}

% Refer to a lecture note by its tag (the prefix of its labels), e.g.
% \LN{proofs} prints "Lecture Note 2" (in the book, a link to that chapter).  The table below is the only place the numbers appear.
\makeatletter
\newcommand{\defineLN}[3]{\@namedef{LNnumber@#2}{#1}\@namedef{LNtag@#1}{#2}\@namedef{LNfile@#2}{#3}}
\newcommand{\LNnumber}[1]{%
  \ifcsname LNnumber@#1\endcsname\csname LNnumber@#1\endcsname
  \else\PackageError{macros}{Unknown lecture-note tag `#1'}{}\fi}
\newcommand{\LNtag}[1]{\csname LNtag@#1\endcsname}
\newcommand{\LNfile}[1]{\csname LNfile@#1\endcsname}
\makeatother
\defineLN{1}{prelim}{1_preliminaries}
\defineLN{2}{proofs}{2_proofs}
\defineLN{3}{matching}{3_stable_matching}
\defineLN{4}{dc}{4_divide_and_conquer}
\defineLN{5}{greedy}{5_greedy}
\defineLN{6}{graphs}{6_graph_traversal}
\defineLN{7}{dp}{7_dynamic_programming}
\defineLN{8}{flow}{8_flow_reductions_LP}
\defineLN{9}{np}{9_NP}

% The text is the same in the book and in the standalone lecture notes; only
% links differ.  In the book, links to other lecture notes stay in the book.
% In a standalone note, they open the other note's published PDF.
% \noteLink{tag}{target}{text} links text to \hypertarget{target} in lecture note tag.
% \LN{proofs} prints "Lecture Note 2", linked to that lecture note.
\newcommand{\currentLN}{0}   % in a standalone note, its number (set by \doTitle)
\newcommand{\notePDF}[1]{\siteURL notes/\LNfile{#1}.pdf}
\newcommand{\noteLink}[3]{%
  \ifnum\LNnumber{#1}=\currentLN\relax\hyperlink{#2}{#3}%
  \else\href{\notePDF{#1}\#nameddest=#2}{#3}\fi}
% \bookLink{target}{text} links text to \hypertarget{target} in the book's
% appendices (from a standalone note, in the published book).
\DeclareRobustCommand{\bookLink}[2]{%
  \ifdefined\BookMode\hyperlink{#1}{#2}%
  \else\href{\siteURL lecture_notes_on_algorithms.pdf\#nameddest=#1}{#2}\fi}
\DeclareRobustCommand{\LN}[1]{%
  \ifnum\LNnumber{#1}=\currentLN\relax Lecture Note~\LNnumber{#1}%
  \else\href{\notePDF{#1}}{Lecture Note~\LNnumber{#1}}\fi}
\pdfstringdefDisableCommands{\def\LN#1{Lecture Note \LNnumber{#1}}}

% Homeworks refer to the notes' sections, exercises and lemmas by \ref: each
% homework's main file says \lectureNoteRefs after %%%% DON'T EDIT THIS FILE %%%%

\usepackage{xifthen}% \isempty
\usepackage{etoolbox}% \patchcmd
% \usepackage{stringstrings}% \substring

\usepackage{amsmath, amsthm, amssymb}
\usepackage[margin=1in]{geometry}
\usepackage{xspace}
\usepackage{graphicx}
\usepackage[normalem]{ulem}
\usepackage{adjustbox}
\usepackage{xcolor}
\usepackage{changepage}
\usepackage{framed}
\usepackage{multicol}
\usepackage{enumitem}
% \usepackage{listings} % lstlistings environment

\usepackage{makecmds}           %\provideenvironment

\usepackage{xparse,environ}
\usepackage{letltxmacro}       % \LetLtxMacro
\usepackage{wrapfig}
\usepackage{tikz}

\usetikzlibrary{shapes}
\usetikzlibrary{positioning}
\usetikzlibrary{shapes.misc}
\usetikzlibrary{arrows}

\newtheorem{lemma}{Lemma}
\newtheorem{definition}{Definition}
\newtheorem{theorem}{Theorem}
\newtheorem{corollary}{Corollary}
\newtheorem{claim}{Claim}
\newtheorem{observation}{Observation}
\newtheorem{conjecture}{Conjecture}

% A lemma or theorem restated later with its original number (set with
% \setcounter) goes in a restatement environment, which gives it its own link
% anchor (otherwise it would duplicate the original's).
\newenvironment{restatement}
  {\def\theHlemma{restated.\thelemma}\def\theHtheorem{restated.\thetheorem}}{}

\newenvironment{myframed}
{\smallskip\par\noindent\begin{minipage}{\textwidth}\begin{framed}\vspace*{-5pt}}{\vspace*{-5pt}\end{framed}\end{minipage}\smallskip}


\definecolor{darkblue}{RGB}{0,0,200}

 
\let\oldParagraph\paragraph
\providecommand{\gradescopeColor}{\color{black}}
\renewcommand{\paragraph}[1]{\oldParagraph{\gradescopeColor #1}}
\providecommand{\gradescopeInit}{}

\providecommand{\BLUE}[1]{{\color{darkblue} #1}}
\providecommand{\REPLACEME}{\textcolor{red}{\ensuremath{\star \cdots \star}}\xspace}
\provideenvironment{blue}{\color{darkblue}}{}

\newcommand{\dist}[2]{\ensuremath{\operatorname{\sf dist}(#1, #2)}}
\newcommand{\depth}[2]{\ensuremath{\operatorname{\sf depth}_{#1}(#2)}}
\newcommand{\cost}[1]{\ensuremath{\operatorname{\sf cost}(#1)}}
\newcommand{\weight}[1]{\ensuremath{\operatorname{\sf wt}(#1)}}
% \newcommand{\val}[1]{\ensuremath{\operatorname{\sf val}(#1)}}
\newcommand{\eps}{\varepsilon}
\newcommand{\capac}{\textsf{cap}}
\newcommand{\capOf}[1]{\ensuremath{\operatorname{\sf cap}(#1)}}
\newcommand{\rescap}[1]{\ensuremath{\operatorname{\sf cap}_f(#1)}}
\newcommand{\into}[1]{\ensuremath{\operatorname{\sf in}(#1)}}
\newcommand{\out}[1]{\ensuremath{\operatorname{\sf out}(#1)}}
\newenvironment{LP}[1]{
  \begin{array}[t]{@{}l@{}}
    \displaystyle #1 \\
    \begin{aligned}\\[-10pt]}{\end{aligned}\end{array}}

\newcommand{\opt}[1]{\ensuremath{\operatorname{\sf opt}(#1)}}
\newcommand{\Opt}[1]{\ensuremath{\operatorname{\sf Opt}(#1)}}
\newcommand{\half}{\frac{1}{2}}

\newcommand{\E}{\operatorname{E}}
\newcommand{\giv}{\,|\,}
\newcommand{\Z}{\mathbb{Z}}
\newcommand{\Zp}{\mathbb{Z}_{\ge 0}}
\newcommand{\R}{\mathbb{R}}
% \newcommand{\Rp}{\mathbb{R}_{\ge 0}}
\newcommand{\Rp}{\mathbb{R}_{\tiny +}}
\newcommand{\N}{\mathbb{N}}
\newcommand{\Np}{\mathbb{N}_{\ge 0}}
\newcommand{\calS}{{\cal S}}

\newcommand{\OP}[1]{\text{\small\sc #1}\xspace}
\newcommand{\Insert}{\OP{Insert}}
\newcommand{\FindMin}{\OP{Find-Min}}
\newcommand{\DeleteMin}{\OP{Delete-Min}}
\newcommand{\DecreaseKey}{\OP{Decrease}}
\newcommand{\Merge}{\OP{Merge}}
\newcommand{\Cut}{\OP{Cut}}
% \newcommand{\prob}[1]{\text{\sc #1}}
% a problem name; in running text, lines may break inside it
\DeclareRobustCommand{\prob}[1]{\ifmmode\text{#1}\else#1\fi}
\newcommand{\encode}[1]{\langle #1 \rangle}

\newcommand{\best}[1]{\text{\sf best}(#1)}
\newcommand{\worst}[1]{\text{\sf worst}(#1)}
\newcommand{\floor}[1]{\lfloor #1 \rfloor}
\newcommand{\ceil}[1]{\lceil #1 \rceil}

\setlength{\parskip}{2pt}
\setlength{\parindent}{1em}

% \setlength\fboxsep{0pt}

%%% longFormProof, shortFormProof, block environments
%%% \step

\newlist{steps}{enumerate}{7}
\setlist[steps]{wide,
  topsep=1pt,
  parsep=1pt,
  partopsep=1pt,
  itemsep=0pt,
  align=left,
  labelindent=0pt,
  itemindent=0pt, 
}
\setlist[steps,1]{
  label=\arabic*.,
  ref=\arabic*
}
\setlist[steps,2]{
  label=\arabic{stepsi}.\arabic*.,
  ref=\arabic{stepsi}.\arabic*
}
\setlist[steps,3]{
  label=\arabic{stepsi}.\arabic{stepsii}.\arabic*.,
  ref=\arabic{stepsi}.\arabic{stepsii}.\arabic*
}
\setlist[steps,4]{
  label=\arabic{stepsi}.\arabic{stepsii}.\arabic{stepsiii}.\arabic*.,
  ref=\arabic{stepsi}.\arabic{stepsii}.\arabic{stepsiii}.\arabic*
}
\setlist[steps,5]{
  label=\arabic{stepsi}.\arabic{stepsii}.\arabic{stepsiii}.\arabic{stepsiv}.\arabic*.,
  ref=\arabic{stepsi}.\arabic{stepsii}.\arabic{stepsiii}.\arabic{stepsiv}.\arabic*
}
\setlist[steps,6]{
  label=\arabic{stepsi}.\arabic{stepsii}.\arabic{stepsiii}.\arabic{stepsiv}.\arabic{stepsv}.\arabic*.,
  ref=\arabic{stepsi}.\arabic{stepsii}.\arabic{stepsiii}.\arabic{stepsiv}.\arabic{stepsv}.\arabic*
}
\setlist[steps,7]{
  label=\arabic{stepsi}.\arabic{stepsii}.\arabic{stepsiii}.\arabic{stepsiv}. \arabic{stepsv}. \arabic{stepsvi}.\arabic*.,
  ref=\arabic{stepsi}.\arabic{stepsii}.\arabic{stepsiii}.\arabic{stepsiv}.\arabic{stepsv}. \arabic{stepsvi}.\arabic*
}

\newcommand{\step}[1][]{\item\maybeLabel{#1}}

\makeatletter
% \skipitem steps the list counter itself (for block labels).  Give each one
% its own hyperref link target, since step numbers repeat from proof to proof.
\newcounter{skipitem}
\newcommand{\skipitem}{\stepcounter{skipitem}\refstepcounter{\@enumctr}}
\renewcommand{\theHstepsi}{skipitem.\arabic{skipitem}}
\renewcommand{\theHstepsii}{skipitem.\arabic{skipitem}}
\renewcommand{\theHstepsiii}{skipitem.\arabic{skipitem}}
\renewcommand{\theHstepsiv}{skipitem.\arabic{skipitem}}
\renewcommand{\theHstepsv}{skipitem.\arabic{skipitem}}
\renewcommand{\theHstepsvi}{skipitem.\arabic{skipitem}}
\renewcommand{\theHstepsvii}{skipitem.\arabic{skipitem}}
\makeatother

\newcommand{\maybeLabel}[1]{\ifthenelse{\isempty{#1}}{}{\label{#1}}}

\newenvironment{longFormProof}[1][Proof (long form).]
{\begin{proof}[#1]~\par\begin{steps}}{\end{steps}\vspace*{-3ex}\end{proof}}

\newenvironment{shortFormProof}[1][Proof (short form).]
{\begin{proof}[#1]}{\end{proof}}

\newenvironment{longFormSubProof}[1][Proof.]
{\begin{adjustwidth}{1em}{0em}\begin{proof}[#1]~\par\vspace*{-2pt}\skipitem\begin{steps}}
{\end{steps}\vspace*{-1ex}\end{proof}\end{adjustwidth}\smallskip}

\newenvironment{shortFormSubProof}[1][Proof.]
{\begin{proof}[#1]}
{\end{proof}\vspace*{-3pt}}

\newenvironment{block}[1][]
{\skipitem\maybeLabel{#1}\item[] \begin{steps}}{\end{steps}}

\newcommand{\comment}[1]{\hfill{\footnotesize\emph{#1}}}
\newcommand{\lineacross}{\par\vspace*{-0.7\baselineskip}\noindent\hrulefill\par}

%%% problem

% \newcounter{problem}
% \newcommand{\problem}[1][]{}
% {\ifthenelse{\isempty{#1}}{\refstepcounter{problem}}{}%
%   \paragraph{Problem \ifthenelse{\isempty{#1}}
%     {\bf\arabic{problem}.}
%     {#1}} }

\newcommand{\problemNumbered}[1]{
  \setcounter{problem}{#1}%
  \addtocounter{problem}{-1}%
  \problem
}
\newenvironment{problemTheorem}{%
  \setcounter{theorem}{\value{problem}}%
  \addtocounter{theorem}{-1}%
  \begin{theorem}[Problem \arabic{problem}]
  }{\end{theorem}}

\newcounter{problem}
\newcommand{\problem}[1][]
{\ifthenelse{\isempty{#1}}
  {\refstepcounter{problem}\paragraph{Problem~\arabic{problem}.}}
  {\paragraph{Problem~#1.}}\refreshHeader} 

\newcounter{exercise}
\newcommand{\exercise}[1][]
{\ifthenelse{\isempty{#1}}
  {\refstepcounter{exercise}\paragraph{Exercise~\theexercise.}}
  {\paragraph{Exercise~#1.}}}

\newcommand{\solution}{\paragraph{Solution.}} 

\newenvironment{mylist}[1][]
{\begin{list}{#1}{
      \setlength{\itemsep}{-1pt}
      \setlength{\labelwidth}{0pt}
      \setlength{\leftmargin}{1.25em}
    }}
{\end{list}}

\newcommand{\LNTitle}{\textcolor{red}{UNDEFINED LNTitle}}

% In a standalone lecture note's table of contents, leave room for numbers
% like 7.10 (in bold) and 7.10.2.
\makeatletter
\newcommand{\LNtocWidths}{%
  \patchcmd{\l@section}{1.5em}{2.8em}{}{\PackageWarning{macros}{Could not widen section numbers in the contents}}%
  \renewcommand*{\l@subsection}{\@dottedtocline{2}{2.8em}{3.2em}}%
  \renewcommand*{\l@subsubsection}{\@dottedtocline{3}{6.0em}{4.1em}}}
\makeatother

\newcommand{\doTitle}[3]{
  \pagenumbering{roman}
  \renewcommand{\LNTitle}{Lecture Note #1, #3}
  \renewcommand{\currentLN}{#1}
  \codeListingstrue
  % number sections, theorems etc., exercises, figures, tables and equations
  % within the lecture note (e.g. Section 8.2, Lemma 8.6), as in the book
  \renewcommand{\thesection}{#1.\arabic{section}}
  \renewcommand{\thelemma}{#1.\arabic{lemma}}
  \renewcommand{\thedefinition}{#1.\arabic{definition}}
  \renewcommand{\thetheorem}{#1.\arabic{theorem}}
  \renewcommand{\thecorollary}{#1.\arabic{corollary}}
  \renewcommand{\theclaim}{#1.\arabic{claim}}
  \renewcommand{\theobservation}{#1.\arabic{observation}}
  \renewcommand{\theconjecture}{#1.\arabic{conjecture}}
  \renewcommand{\theexercise}{#1.\arabic{exercise}}
  \renewcommand{\thefigure}{#1.\arabic{figure}}
  \renewcommand{\thetable}{#1.\arabic{table}}
  \renewcommand{\theequation}{#1.\arabic{equation}}
  \LNtocWidths
  \title{Lecture Note #1\\ #2}
  % a note's \LNabstract, if it defines one before \doTitle, goes right below the title
  \date{\ifdefined\LNabstract\parbox{0.85\textwidth}{\normalsize\normalfont\LNabstract}\\[3ex]\fi
    draft of \today, by \href{\authorURL}{N.~Young}\\[2ex] {\small\copyrightNotice}}
  \maketitle
  {\small
    \hypertarget{TOC}{}
    \tableofcontents
  }
  \thispagestyle{empty}
  % the text starts on page 1, after the title and contents (numbered i, ii, ...)
  \clearpage
  \pagenumbering{arabic}
}

\def\clap#1{\hbox to 0pt{\hss#1\hss}}

\newcommand{\refreshHeader}{}
\newcommand{\header}{}

\newcommand{\doHeader}[1]{      % for lecture notes
  \renewcommand{\header}{%
    \clap{\raisebox{0pt}{\protect\hyperlink{TOC}{$\uparrow$\scriptsize\,\textsc{contents}\hspace{1.1in}}}}%
    {\footnotesize \LNTitle{} / Section \thesection, #1}} %, #1, draft of \today}}
  \markboth{\header}{\header}
}

\newcommand{\SID}{}
\newcommand{\setAuthor}[1]{\renewcommand{\author}{#1}}
\newcommand{\setSID}[1]{\renewcommand{\SID}{#1}}

\newcommand{\doHwHeader}[1]{                                         % for homeworks
  \renewcommand{\refreshHeader}{
    \renewcommand{\header}{%
      {\footnotesize \classSection #1 / Problem~\arabic{problem} / draft of \today \hfill \author\ifGradescope~(\SID)\fi~~}}
    \markboth{\header}{\header}
  }
  \renewcommand{\header}{%
    {\footnotesize \classSection #1 / Instructions / draft of \today \hfill \author\ifGradescope~(\SID)\fi~~}}
  \markboth{\header}{\header}
}

% Homework settings, for instructors using the homeworks in a course:
% \classSection is the course name shown in each homework's running head
% (e.g. \renewcommand{\classSection}{CS 3000\xspace}), and \Gradescopetrue turns on
% the instructions for submitting to Gradescope (in its fixed-length mode),
% the hints saying on which page of the PDF each answer should be, and the
% student ID in the running head.
\newcommand{\classSection}{}
\newif\ifGradescope
\Gradescopefalse
\newcommand{\GradescopeComment}[1]{\ifGradescope\comment{#1}\fi}

\pagestyle{myheadings}
\addtolength{\headsep}{0.2in}
\addtolength{\topmargin}{-0.3in}
\addtolength{\textheight}{0.4in}

%%%%%%%%%%%%%%%%%%%%%%%%%%%%%%%%%%%%%%%%%%%%%%%%
%%%%%%%%%%%%%%%%%%%%%%%%%%%%%%%%%%%%%%%%%%%%%%%%

\newcommand{\MF}{\href
  {https://mfleck.cs.illinois.edu/building-blocks/}
  {MF}\xspace}

\newcommand{\KWSlides}{\href
  {https://www.cs.princeton.edu/~wayne/kleinberg-tardos/}
  {KW slides}\xspace}

\newcommand{\LLM}{\href
  {https://courses.csail.mit.edu/6.042/spring18/mcs.pdf}
  {LLM}\xspace}

% \DPV, \KT, \CLRS: textbooks listed (without links) in the table of sources in
% the Preliminaries note; each prints as a link to its entry in that table.
\DeclareRobustCommand{\DPV}{\noteLink{prelim}{source.DPV}{DPV}\xspace}
\DeclareRobustCommand{\KT}{\noteLink{prelim}{source.KT}{KT}\xspace}
\DeclareRobustCommand{\CLRS}{\noteLink{prelim}{source.CLRS}{CLRS}\xspace}
\AtBeginDocument{\pdfstringdefDisableCommands{\def\DPV{DPV }\def\KT{KT }\def\CLRS{CLRS }}}

\newcommand{\JE}{\href
  {https://jeffe.cs.illinois.edu/teaching/algorithms/}
  {JE}\xspace}

\newcommand{\JEGreedy}{\href
  {https://jeffe.cs.illinois.edu/teaching/algorithms/book/04-greedy.pdf}
  {JE Chapter 4}\xspace}

\newcommand{\BBCode}{\href
{https://northeastern.blackboard.com/webapps/cmsmain/webui/courses/CS3000.MERGED.202010/code?action=frameset&subaction=view&uniq=-jj2max&course_id=_2597392_1&mask=\%2Fcourses\%2FCS3000.MERGED.202010}
  {code}\xspace}

% The Python code (src/code).  Each lecture note that has code ends with a
% section listing its files: \codeSection{path, path, ...}.  In the text,
% \pycode{greedy_algorithms/Dijkstra.py} prints the path as a link to its listing
% (in the homeworks, which have no listings, to the file on GitHub), and
% \pycode[text]{...} uses the given text instead.  \pythonCode{path, ...} is the
% standard sentence ending a section whose algorithms have code:
% "Python code: path, ... (Section N.k)."
\newcommand{\codeFolder}{https://github.com/nealeyoung/algs101/tree/main/src/code}
\newcommand{\codeRepo}{https://github.com/nealeyoung/algs101/blob/main/src/code}
\newif\ifcodeListings   % set by \doTitle: this document has the listings
\newcommand{\githubCode}[2][]{\href{\codeRepo/#2}{\ifx\relax#1\relax\nolinkurl{#2}\else#1\fi}}
\newcommand{\pycode}[2][]{%
  \ifcodeListings\hyperref[code: #2]{\ifx\relax#1\relax\nolinkurl{#2}\else#1\fi}%
  \else\githubCode[#1]{#2}\fi}
\makeatletter
\newcommand{\pythonCode}[1]{%
  \def\pythonCode@first{}\def\pythonCode@list{}%
  \forcsvlist{\pythonCode@add}{#1}%
  Python code: \pythonCode@list\ (Section~\ref{code: \pythonCode@first}).}
\newcommand{\pythonCode@add}[1]{%
  \ifx\pythonCode@first\@empty\def\pythonCode@first{#1}\appto\pythonCode@list{\pycode{#1}}%
  \else\appto\pythonCode@list{, \pycode{#1}}\fi}
\makeatother

% Code listings that copy and paste exactly (in PDF viewers that honor the
% ToUnicode maps, e.g. MuPDF-based ones; PDFKit viewers such as Preview and
% Skim drop leading spaces, so each listing's heading links to the file on
% GitHub).  Spaces are typeset as white visible-space glyphs that copy as
% spaces, and hyphens and asterisks as the T1 glyphs, so they copy as themselves.
\usepackage[T1,OT1]{fontenc}   % T1 for the code listings only; OT1 stays the default
\usepackage{listings}
\pdfgentounicode=1
\input glyphtounicode
\pdfglyphtounicode{uni2423}{0020}
\makeatletter
\def\lst@visiblespace{\textcolor{white}{\char32}}
% put an invisible space on each blank line, so that blank lines copy too
% (listings' OnEmptyLine hook runs before the line break, so patch the loop)
\def\lst@DoNewLines{%
    \@whilenum\lst@newlines>\@ne \do {\lst@NewLine \lst@visiblespace}%
    \ifnum\lst@newlines>\z@ \lst@NewLine \fi}
\makeatother
\lstset{language=Python,
  basicstyle=\fontencoding{T1}\fontfamily{lmtt}\selectfont\footnotesize,
  upquote=true, showstringspaces=false, columns=fullflexible, keepspaces=true,
  showspaces=true, literate={-}{{\char45}}1 {*}{{\char42}}1,
  frame=single, framesep=1ex, rulecolor=\color{gray}}
% the listings' paths are relative to the lecture note's folder (in the book, to src/)
\ifdefined\BookMode\newcommand{\codeDir}{code/}\else\newcommand{\codeDir}{../code/}\fi
% \codefile{path}: a listing of src/code/path, headed by a link to it on GitHub
\newcommand{\codefile}[1]{%
  \subsection*{\githubCode{#1}}\label{code: #1}%
  \lstinputlisting{\codeDir#1}}
% \codeSection{path, path, ...}: the lecture note's last section, listing its code
\newcommand{\codeSection}[1]{%
  \newpage
  \section{Python code}\label{\LNtag{\currentLN}: sec: code}\doHeader{Python code}
  This section lists the Python code for this lecture note.
  Click on a file's name to open the file on GitHub.
  Copying code from GitHub is more reliable than copying it from this PDF,
  since some PDF viewers drop the leading spaces (the indentation) of copied lines.
  \forcsvlist{\codefile}{#1}}
% the copyright notice, for title pages and the homework first pages
\newcommand{\licenseURL}{https://github.com/nealeyoung/algs101/blob/main/LICENSE}
\newcommand{\authorURL}{https://www.cs.ucr.edu/\string~neal}
\newcommand{\copyrightNotice}{%
  \copyright\ 2018--2026 \href{\authorURL}{Neal E. Young}.
  Text: CC BY-NC-SA 4.0; code: MIT; except where noted.
  See \url{\licenseURL}.}
% the published site: PDFs of the notes, and the homework assignments and templates
\newcommand{\siteURL}{https://nealeyoung.github.io/algs101/}

\newcommand{\blackboard}{\href
  {https://northeastern.blackboard.com/webapps/blackboard/execute/launcher?type=Course&id=_2597392_1&url=}
  {Blackboard}\xspace}

\newcommand{\JUSlides}{\href
  {http://www.khoury.neu.edu/home/jullman/cs3000f18/schedule.html}
  {Prof.~Ullman's slides}\xspace}

% a URL on its own line, in three layouts
\newcommand{\URL}[1]{\par-- \parbox[t]{0.8\textwidth}{\url{#1}} \medskip\par}
\newcommand{\URLsmall}[1]{\par-- \parbox[t]{0.95\textwidth}{\small\url{#1}} \medskip\par}
\newcommand{\URLplain}[1]{\par\parbox[t]{\textwidth}{\url{#1}}\par}

%%%%%%%%%%%%%%%%%%%%%%%%%%%%%%%%%%%%%%%%%%%%%%%%
%%%%%%%%%%%%%%%%%%%%%%%%%%%%%%%%%%%%%%%%%%%%%%%%

\newcommand{\continued}{\vfill\centerline{\emph{(continued)}}\vfill\newpage}

% load hyperref last to avoid conflicts; long URLs may also break after hyphens
\PassOptionsToPackage{hyphens}{url}
\ifdefined\BookMode\else\usepackage{xr-hyper}\fi   % for \lectureNoteRefs (homeworks), below
\usepackage{hyperref}
\hypersetup{
    colorlinks,
    linkcolor={blue!50!black},
    citecolor={blue!50!black},
    urlcolor={blue!80!black}
}

\providecommand{\answerDirectory}{.}
\newcommand{\inputAnswer}[1]{\refreshHeader\input{\answerDirectory/#1}}
\newcommand{\inputSection}[1]{\input{#1}}

% Refer to a lecture note by its tag (the prefix of its labels), e.g.
% \LN{proofs} prints "Lecture Note 2" (in the book, a link to that chapter).  The table below is the only place the numbers appear.
\makeatletter
\newcommand{\defineLN}[3]{\@namedef{LNnumber@#2}{#1}\@namedef{LNtag@#1}{#2}\@namedef{LNfile@#2}{#3}}
\newcommand{\LNnumber}[1]{%
  \ifcsname LNnumber@#1\endcsname\csname LNnumber@#1\endcsname
  \else\PackageError{macros}{Unknown lecture-note tag `#1'}{}\fi}
\newcommand{\LNtag}[1]{\csname LNtag@#1\endcsname}
\newcommand{\LNfile}[1]{\csname LNfile@#1\endcsname}
\makeatother
\defineLN{1}{prelim}{1_preliminaries}
\defineLN{2}{proofs}{2_proofs}
\defineLN{3}{matching}{3_stable_matching}
\defineLN{4}{dc}{4_divide_and_conquer}
\defineLN{5}{greedy}{5_greedy}
\defineLN{6}{graphs}{6_graph_traversal}
\defineLN{7}{dp}{7_dynamic_programming}
\defineLN{8}{flow}{8_flow_reductions_LP}
\defineLN{9}{np}{9_NP}

% The text is the same in the book and in the standalone lecture notes; only
% links differ.  In the book, links to other lecture notes stay in the book.
% In a standalone note, they open the other note's published PDF.
% \noteLink{tag}{target}{text} links text to \hypertarget{target} in lecture note tag.
% \LN{proofs} prints "Lecture Note 2", linked to that lecture note.
\newcommand{\currentLN}{0}   % in a standalone note, its number (set by \doTitle)
\newcommand{\notePDF}[1]{\siteURL notes/\LNfile{#1}.pdf}
\newcommand{\noteLink}[3]{%
  \ifnum\LNnumber{#1}=\currentLN\relax\hyperlink{#2}{#3}%
  \else\href{\notePDF{#1}\#nameddest=#2}{#3}\fi}
% \bookLink{target}{text} links text to \hypertarget{target} in the book's
% appendices (from a standalone note, in the published book).
\DeclareRobustCommand{\bookLink}[2]{%
  \ifdefined\BookMode\hyperlink{#1}{#2}%
  \else\href{\siteURL lecture_notes_on_algorithms.pdf\#nameddest=#1}{#2}\fi}
\DeclareRobustCommand{\LN}[1]{%
  \ifnum\LNnumber{#1}=\currentLN\relax Lecture Note~\LNnumber{#1}%
  \else\href{\notePDF{#1}}{Lecture Note~\LNnumber{#1}}\fi}
\pdfstringdefDisableCommands{\def\LN#1{Lecture Note \LNnumber{#1}}}

% Homeworks refer to the notes' sections, exercises and lemmas by \ref: each
% homework's main file says \lectureNoteRefs after \input{macros}.  A homework
% folder has a link refs/ to src/refs/, which holds each note's label records
% (made by tools/refs.sh from the notes' builds).  Read with xr-hyper, they make
% \ref{greedy: exercise: ...} give the exercise's number in the published note,
% linked to it there; labels only the book has (its appendices) link to the book.
\ifdefined\BookMode\else
  \makeatletter
  \newcommand{\lectureNoteRefs}{%
    \@for\LN@tag:=prelim,proofs,matching,dc,greedy,graphs,dp,flow,np\do{%
      \externaldocument{refs/\LNfile{\LN@tag}}[\notePDF{\LN@tag}]}%
    \externaldocument{refs/book}[\siteURL lecture_notes_on_algorithms.pdf]}
  % hyperref links a \ref to another document's label as URL#anchor; browsers'
  % PDF viewers need URL#nameddest=anchor (the form \noteLink and \bookLink use)
  \def\Hy@setref@link#1#2#3#4#5#6\@nil#7{%
    \begingroup
      \ifx\\#5\\\toks0={\hyper@@link{#5}{#4}}\else\toks0={\hyper@@link{#5}{nameddest=#4}}\fi
      \toks1=\expandafter{#7{#1}{#2}{#3}{#4}{#5}}%
      \edef\x{\endgroup\the\toks0 {\the\toks1 }}%
    \x}
  \makeatother
\fi

% Book mode: main.tex defines \BookMode before \input{macros} to combine the
% lecture notes into one report-class document, one chapter per note.
% Without \BookMode (each lecture note on its own, and the homeworks) none
% of this applies.
\ifdefined\BookMode
  \usepackage{etoc}

  % links to other lecture notes stay in the book
  \renewcommand{\noteLink}[3]{\hyperlink{#2}{#3}}

  % each chapter is a lecture note, and is called one, as in the standalone notes
  \renewcommand{\chaptername}{Lecture Note}
  \setcounter{secnumdepth}{3}   % number subsubsections, as in the standalone notes

  % global table of contents, for main.tex
  \newcommand{\bookContents}{%
    \hypertarget{TOC}{}%
    \etocsetnexttocdepth{section}%
    \tableofcontents}

  % \doTitle{N}{title}{short title} starts chapter N, with a link to the global
  % table of contents and the chapter's own table of contents
  \renewcommand{\doTitle}[3]{%
    \setcounter{chapter}{\numexpr#1-1\relax}%
    \chapter{\texorpdfstring{#2}{#3}}%
    \label{\LNtag{#1}: chapter}%
    \csgdef{LNshorttitle@#1}{#3}%
    \csgdef{LNtitle@#1}{#2}%
    \renewcommand{\LNTitle}{Lecture Note #1, #3}%
    \codeListingstrue
    \hyperlink{TOC}{$\uparrow$\scriptsize\,\textsc{contents}}%
    {\small
      \hypertarget{TOC-#1}{}%
      \etocsettocstyle{\section*{\contentsname}}{}%
      \etocsetnexttocdepth{subsubsection}%
      \localtableofcontents
    }}

  % the code: each lecture note's \codeSection just records its files, and
  % \codeAppendix (in the book's appendix) lists them all, a section per note
  \makeatletter   % (file names have no spaces, so drop the spaces around them)
  \renewcommand{\codeSection}[1]{\csxdef{LNcode@\arabic{chapter}}{\zap@space#1 \@empty}}
  \makeatother
  % \codeAppendix: first a contents list (each note's topic, then its files,
  % linked to the listings), then a section per note listing its files
  \makeatletter
  \def\LNbasename#1{\LN@base#1/\LN@end}   % a/b/c.py -> c.py
  \def\LN@base#1/#2\LN@end{\ifx\relax#2\relax#1\else\LN@base#2\LN@end\fi}
  \makeatother
  \newcommand{\codeContentsFile}[1]{%
    \edef\LNfilename{\LNbasename{#1}}%
    \item \pycode[\texttt{\expandafter\detokenize\expandafter{\LNfilename}}]{#1}}
  \newcommand{\codeAppendix}{%
    \begin{itemize}
    \foreach \n in {1,...,9} {%
      \ifcsdef{LNcode@\n}{%
        \item \hyperref[code: note \n]{\csuse{LNtitle@\n}} (\LN{\LNtag{\n}})
        \begin{itemize}
        \edef\LNcodelist{\csuse{LNcode@\n}}%
        \expandafter\forcsvlist\expandafter\codeContentsFile\expandafter{\LNcodelist}%
        \end{itemize}}{}}
    \end{itemize}
    \newpage
    \foreach \n in {1,...,9} {%
      \ifcsdef{LNcode@\n}{%
        \section[Lecture Note \n: \csuse{LNtitle@\n}]{%
          \texorpdfstring{\LN{\LNtag{\n}}: \csuse{LNtitle@\n}}{Lecture Note \n: \csuse{LNtitle@\n}}}%
        \label{code: note \n}%
        \edef\LNcodelist{\csuse{LNcode@\n}}%
        \expandafter\forcsvlist\expandafter\codefile\expandafter{\LNcodelist}%
        \newpage}{}}}

  % a lecture note is a chapter: \LN{proofs} prints a link "Lecture Note 2"
  \DeclareRobustCommand{\LN}[1]{\hyperref[#1: chapter]{Lecture Note~\ref*{#1: chapter}}}
  \pdfstringdefDisableCommands{\def\LN#1{Lecture Note \LNnumber{#1}}}

  % running head; "contents" links to the current chapter's table of contents
  \renewcommand{\doHeader}[1]{%
    \renewcommand{\header}{%
      \clap{\raisebox{0pt}{\protect\hyperlink{TOC-\arabic{chapter}}{$\uparrow$\scriptsize\,\textsc{contents}\hspace{1.1in}}}}%
      {\footnotesize \LNTitle{} / Section \thesection, #1}}%
    \markboth{\header}{\header}}

  % number theorems, lemmas, etc. and exercises within chapters (e.g. Lemma 8.1)
  \counterwithin{lemma}{chapter}
  \counterwithin{definition}{chapter}
  \counterwithin{theorem}{chapter}
  \counterwithin{corollary}{chapter}
  \counterwithin{claim}{chapter}
  \counterwithin{observation}{chapter}
  \counterwithin{conjecture}{chapter}
  \counterwithin{exercise}{chapter}
\fi

% %%%%%%%%%%%%%%%%%%%%%%%%%%%%%%%%%%%%%%%%%%%%%%%%%%%%%%%%%%%%%%%%%%%%%%%%
% %%%%%%%%%%%%%%%%%%%%%%%%%%%%%%%%%%%%%%%%%%%%%%%%%%%%%%%%%%%%%%%%%%%%%%%%

\ExplSyntaxOn
\seq_new:N \l_itemize_seq
% \NewDocumentEnvironment{REITEMIZE}{ +b }{
\NewDocumentCommand{\REITEMIZECMD}{ +m }{
  \regex_split:nnN {\c{item}} {#1} \l_itemize_seq
  \seq_pop:NN \l_itemize_seq \l_tmpa_tl% remove the (hopefully) empty first item
  \seq_if_empty:NF \l_itemize_seq {
    \l_tmpa_tl
    \seq_map_inline:Nn \l_itemize_seq {
      \regex_split:nnN {\c{item}} {#1} \l_itemize_seq {
        \regex_extract_once:nnNTF { ^\[([^]]*)\](.*) } { ##1} \l_foo_seq
      { \item[\seq_item:Nn \l_foo_seq {2}]\ITEMIZER{\seq_item:Nn \l_foo_seq {3}} } { \item\ITEMIZER{##1}  }
    }
  }
}
}
\ExplSyntaxOff

\newlength{\ITEMHT}
\setlength{\ITEMHT}{1in}

% \setlength{\fboxsep}{1pt}
\newcommand{\ITEMIZER}[1]{%
  % \framebox{
\adjustbox{valign=t}{\parbox{\textwidth}{\linespread{1.15}\selectfont #1}}~~\adjustbox{valign=t}{\parbox{0pt}{\color{gray}\rule[-\ITEMHT]{0.01pt}{\ITEMHT}}}
% }
  % \begin{tabular}{@{}p{\textwidth}@{\rule[-\ITEMHT]{0pt}{\ITEMHT}}}
  %   #1
  % \end{tabular}
}

\newsavebox\thesmashminipage
\newlength{\thesmashminipagelength}
\newlength{\answerboxlength}
\newlength{\answerboxoverrun}
\newenvironment{smashminipage}[2]
{\setlength\answerboxlength{#2}
\begin{lrbox}{\thesmashminipage}\begin{minipage}[t]{#1}}
{\end{minipage}\end{lrbox}%
\setlength{\thesmashminipagelength}{\dimexpr\ht\thesmashminipage+\dp\thesmashminipage}%
\setlength{\answerboxoverrun}{\dimexpr\thesmashminipagelength-\answerboxlength}%
\ifdim\the\thesmashminipagelength>\the\answerboxlength%
\usebox{\thesmashminipage}%
\smash{\color{red}\hspace*{-0.98\textwidth}\rule[-\answerboxlength]{0.98\textwidth}{1pt}}
\errmessage{Answer exceeds allotted length, please shorten it.}
\else
\usebox{\thesmashminipage}%
\fi
}



\newenvironment{answerBox}[1]{%
  \def\ANSWERBOXHEIGHT{#1}
  \begin{adjustbox}{valign=t}
  \begin{smashminipage}{\textwidth}{\ANSWERBOXHEIGHT}
}{\end{smashminipage}\end{adjustbox}~~~\adjustbox{valign=t}{\parbox{0pt}{\color{gray}\rule[-\ANSWERBOXHEIGHT]{0.01pt}{\ANSWERBOXHEIGHT}\rule[-\ANSWERBOXHEIGHT]{10pt}{0.01pt}}}}

% % \setlist[enumerate]{itemsep=0pt,parsep=0pt,topsep=0pt,leftmargin=100pt}

% \makeatletter 
\LetLtxMacro\SVenumerate\enumerate
\let\endSVenumerate\endenumerate

\newcommand{\injectEnumerate}{
  \let\SValph\alph
  \def\alph{}
  \renewenvironment{enumerate}[1][]
  {\let\alph\SValph \SVenumerate[label=(\alph*),itemsep=2pt,parsep=0pt]\REITEMIZE}{\endREITEMIZE\endSVenumerate}
  % \setlist[enumerate]{before=\Collect@Body\REITEMIZECMD}
}
% \makeatother
\newcommand{\restoreEnumerate}{
  \renewenvironment{enumerate}[1][]
  \let\SValph\alph
  \def\alph{}
  {\let\alph\SValph \SVenumerate[label=(\alph*)]}{\endSVenumerate}
  % \setlist[enumerate]{before=,after=}
}
\newcommand{\injectBlue}{
  \let\SVblue\blue
  \let\endSVblue\endblue
  \def\blue{\SVblue\REITEMIZE}
  \def\endblue{\endREITEMIZE\endSVblue}
}
\newcommand{\restoreBlue}{
  \let\blue\SVblue
  \let\endblue\endSVblue
}
\newcommand{\injectEnumerateBlue}{
  \setlist[enumerate]{before=\injectBlue,after=\restoreBlue,itemsep=2pt,parsep=0pt}
}
\newcommand{\restoreEnumerateBlue}{
  \setlist[enumerate]{before=,after=}
}


% %%%%%%%%%%%%%%%%%%%%%%%%%%%%%%%%%%%%%%%%%%%%%%%%%%%%%%%%%%%%%%%%%%%%%%%%
% %%%%%%%%%%%%%%%%%%%%%%%%%%%%%%%%%%%%%%%%%%%%%%%%%%%%%%%%%%%%%%%%%%%%%%%%

% \NewDocumentEnvironment{env}{ }
% {\Environenv}
% {\endEnvironenv}

% \NewEnviron{REITEMIZE}{\REITEMIZECMD\BODY}

\makeatletter 
\newenvironment{REITEMIZE}{\Collect@Body\REITEMIZECMD}{}
\makeatother 

\newcommand{\NODEone}[1]{\ensuremath{v_{#1}}\xspace}
\newcommand{\NODE}[2]{{\NODEone {#1,#2}}\xspace}
\newcommand{\EDGE}[2]{\ensuremath{(\textsf{#1}, \textsf{#2})}}
\newcommand{\ITEM}{\item}

%%% Local Variables:
%%% mode: latex
%%% End:
.  A homework
% folder has a link refs/ to src/refs/, which holds each note's label records
% (made by tools/refs.sh from the notes' builds).  Read with xr-hyper, they make
% \ref{greedy: exercise: ...} give the exercise's number in the published note,
% linked to it there; labels only the book has (its appendices) link to the book.
\ifdefined\BookMode\else
  \makeatletter
  \newcommand{\lectureNoteRefs}{%
    \@for\LN@tag:=prelim,proofs,matching,dc,greedy,graphs,dp,flow,np\do{%
      \externaldocument{refs/\LNfile{\LN@tag}}[\notePDF{\LN@tag}]}%
    \externaldocument{refs/book}[\siteURL lecture_notes_on_algorithms.pdf]}
  % hyperref links a \ref to another document's label as URL#anchor; browsers'
  % PDF viewers need URL#nameddest=anchor (the form \noteLink and \bookLink use)
  \def\Hy@setref@link#1#2#3#4#5#6\@nil#7{%
    \begingroup
      \ifx\\#5\\\toks0={\hyper@@link{#5}{#4}}\else\toks0={\hyper@@link{#5}{nameddest=#4}}\fi
      \toks1=\expandafter{#7{#1}{#2}{#3}{#4}{#5}}%
      \edef\x{\endgroup\the\toks0 {\the\toks1 }}%
    \x}
  \makeatother
\fi

% Book mode: main.tex defines \BookMode before %%%% DON'T EDIT THIS FILE %%%%

\usepackage{xifthen}% \isempty
\usepackage{etoolbox}% \patchcmd
% \usepackage{stringstrings}% \substring

\usepackage{amsmath, amsthm, amssymb}
\usepackage[margin=1in]{geometry}
\usepackage{xspace}
\usepackage{graphicx}
\usepackage[normalem]{ulem}
\usepackage{adjustbox}
\usepackage{xcolor}
\usepackage{changepage}
\usepackage{framed}
\usepackage{multicol}
\usepackage{enumitem}
% \usepackage{listings} % lstlistings environment

\usepackage{makecmds}           %\provideenvironment

\usepackage{xparse,environ}
\usepackage{letltxmacro}       % \LetLtxMacro
\usepackage{wrapfig}
\usepackage{tikz}

\usetikzlibrary{shapes}
\usetikzlibrary{positioning}
\usetikzlibrary{shapes.misc}
\usetikzlibrary{arrows}

\newtheorem{lemma}{Lemma}
\newtheorem{definition}{Definition}
\newtheorem{theorem}{Theorem}
\newtheorem{corollary}{Corollary}
\newtheorem{claim}{Claim}
\newtheorem{observation}{Observation}
\newtheorem{conjecture}{Conjecture}

% A lemma or theorem restated later with its original number (set with
% \setcounter) goes in a restatement environment, which gives it its own link
% anchor (otherwise it would duplicate the original's).
\newenvironment{restatement}
  {\def\theHlemma{restated.\thelemma}\def\theHtheorem{restated.\thetheorem}}{}

\newenvironment{myframed}
{\smallskip\par\noindent\begin{minipage}{\textwidth}\begin{framed}\vspace*{-5pt}}{\vspace*{-5pt}\end{framed}\end{minipage}\smallskip}


\definecolor{darkblue}{RGB}{0,0,200}

 
\let\oldParagraph\paragraph
\providecommand{\gradescopeColor}{\color{black}}
\renewcommand{\paragraph}[1]{\oldParagraph{\gradescopeColor #1}}
\providecommand{\gradescopeInit}{}

\providecommand{\BLUE}[1]{{\color{darkblue} #1}}
\providecommand{\REPLACEME}{\textcolor{red}{\ensuremath{\star \cdots \star}}\xspace}
\provideenvironment{blue}{\color{darkblue}}{}

\newcommand{\dist}[2]{\ensuremath{\operatorname{\sf dist}(#1, #2)}}
\newcommand{\depth}[2]{\ensuremath{\operatorname{\sf depth}_{#1}(#2)}}
\newcommand{\cost}[1]{\ensuremath{\operatorname{\sf cost}(#1)}}
\newcommand{\weight}[1]{\ensuremath{\operatorname{\sf wt}(#1)}}
% \newcommand{\val}[1]{\ensuremath{\operatorname{\sf val}(#1)}}
\newcommand{\eps}{\varepsilon}
\newcommand{\capac}{\textsf{cap}}
\newcommand{\capOf}[1]{\ensuremath{\operatorname{\sf cap}(#1)}}
\newcommand{\rescap}[1]{\ensuremath{\operatorname{\sf cap}_f(#1)}}
\newcommand{\into}[1]{\ensuremath{\operatorname{\sf in}(#1)}}
\newcommand{\out}[1]{\ensuremath{\operatorname{\sf out}(#1)}}
\newenvironment{LP}[1]{
  \begin{array}[t]{@{}l@{}}
    \displaystyle #1 \\
    \begin{aligned}\\[-10pt]}{\end{aligned}\end{array}}

\newcommand{\opt}[1]{\ensuremath{\operatorname{\sf opt}(#1)}}
\newcommand{\Opt}[1]{\ensuremath{\operatorname{\sf Opt}(#1)}}
\newcommand{\half}{\frac{1}{2}}

\newcommand{\E}{\operatorname{E}}
\newcommand{\giv}{\,|\,}
\newcommand{\Z}{\mathbb{Z}}
\newcommand{\Zp}{\mathbb{Z}_{\ge 0}}
\newcommand{\R}{\mathbb{R}}
% \newcommand{\Rp}{\mathbb{R}_{\ge 0}}
\newcommand{\Rp}{\mathbb{R}_{\tiny +}}
\newcommand{\N}{\mathbb{N}}
\newcommand{\Np}{\mathbb{N}_{\ge 0}}
\newcommand{\calS}{{\cal S}}

\newcommand{\OP}[1]{\text{\small\sc #1}\xspace}
\newcommand{\Insert}{\OP{Insert}}
\newcommand{\FindMin}{\OP{Find-Min}}
\newcommand{\DeleteMin}{\OP{Delete-Min}}
\newcommand{\DecreaseKey}{\OP{Decrease}}
\newcommand{\Merge}{\OP{Merge}}
\newcommand{\Cut}{\OP{Cut}}
% \newcommand{\prob}[1]{\text{\sc #1}}
% a problem name; in running text, lines may break inside it
\DeclareRobustCommand{\prob}[1]{\ifmmode\text{#1}\else#1\fi}
\newcommand{\encode}[1]{\langle #1 \rangle}

\newcommand{\best}[1]{\text{\sf best}(#1)}
\newcommand{\worst}[1]{\text{\sf worst}(#1)}
\newcommand{\floor}[1]{\lfloor #1 \rfloor}
\newcommand{\ceil}[1]{\lceil #1 \rceil}

\setlength{\parskip}{2pt}
\setlength{\parindent}{1em}

% \setlength\fboxsep{0pt}

%%% longFormProof, shortFormProof, block environments
%%% \step

\newlist{steps}{enumerate}{7}
\setlist[steps]{wide,
  topsep=1pt,
  parsep=1pt,
  partopsep=1pt,
  itemsep=0pt,
  align=left,
  labelindent=0pt,
  itemindent=0pt, 
}
\setlist[steps,1]{
  label=\arabic*.,
  ref=\arabic*
}
\setlist[steps,2]{
  label=\arabic{stepsi}.\arabic*.,
  ref=\arabic{stepsi}.\arabic*
}
\setlist[steps,3]{
  label=\arabic{stepsi}.\arabic{stepsii}.\arabic*.,
  ref=\arabic{stepsi}.\arabic{stepsii}.\arabic*
}
\setlist[steps,4]{
  label=\arabic{stepsi}.\arabic{stepsii}.\arabic{stepsiii}.\arabic*.,
  ref=\arabic{stepsi}.\arabic{stepsii}.\arabic{stepsiii}.\arabic*
}
\setlist[steps,5]{
  label=\arabic{stepsi}.\arabic{stepsii}.\arabic{stepsiii}.\arabic{stepsiv}.\arabic*.,
  ref=\arabic{stepsi}.\arabic{stepsii}.\arabic{stepsiii}.\arabic{stepsiv}.\arabic*
}
\setlist[steps,6]{
  label=\arabic{stepsi}.\arabic{stepsii}.\arabic{stepsiii}.\arabic{stepsiv}.\arabic{stepsv}.\arabic*.,
  ref=\arabic{stepsi}.\arabic{stepsii}.\arabic{stepsiii}.\arabic{stepsiv}.\arabic{stepsv}.\arabic*
}
\setlist[steps,7]{
  label=\arabic{stepsi}.\arabic{stepsii}.\arabic{stepsiii}.\arabic{stepsiv}. \arabic{stepsv}. \arabic{stepsvi}.\arabic*.,
  ref=\arabic{stepsi}.\arabic{stepsii}.\arabic{stepsiii}.\arabic{stepsiv}.\arabic{stepsv}. \arabic{stepsvi}.\arabic*
}

\newcommand{\step}[1][]{\item\maybeLabel{#1}}

\makeatletter
% \skipitem steps the list counter itself (for block labels).  Give each one
% its own hyperref link target, since step numbers repeat from proof to proof.
\newcounter{skipitem}
\newcommand{\skipitem}{\stepcounter{skipitem}\refstepcounter{\@enumctr}}
\renewcommand{\theHstepsi}{skipitem.\arabic{skipitem}}
\renewcommand{\theHstepsii}{skipitem.\arabic{skipitem}}
\renewcommand{\theHstepsiii}{skipitem.\arabic{skipitem}}
\renewcommand{\theHstepsiv}{skipitem.\arabic{skipitem}}
\renewcommand{\theHstepsv}{skipitem.\arabic{skipitem}}
\renewcommand{\theHstepsvi}{skipitem.\arabic{skipitem}}
\renewcommand{\theHstepsvii}{skipitem.\arabic{skipitem}}
\makeatother

\newcommand{\maybeLabel}[1]{\ifthenelse{\isempty{#1}}{}{\label{#1}}}

\newenvironment{longFormProof}[1][Proof (long form).]
{\begin{proof}[#1]~\par\begin{steps}}{\end{steps}\vspace*{-3ex}\end{proof}}

\newenvironment{shortFormProof}[1][Proof (short form).]
{\begin{proof}[#1]}{\end{proof}}

\newenvironment{longFormSubProof}[1][Proof.]
{\begin{adjustwidth}{1em}{0em}\begin{proof}[#1]~\par\vspace*{-2pt}\skipitem\begin{steps}}
{\end{steps}\vspace*{-1ex}\end{proof}\end{adjustwidth}\smallskip}

\newenvironment{shortFormSubProof}[1][Proof.]
{\begin{proof}[#1]}
{\end{proof}\vspace*{-3pt}}

\newenvironment{block}[1][]
{\skipitem\maybeLabel{#1}\item[] \begin{steps}}{\end{steps}}

\newcommand{\comment}[1]{\hfill{\footnotesize\emph{#1}}}
\newcommand{\lineacross}{\par\vspace*{-0.7\baselineskip}\noindent\hrulefill\par}

%%% problem

% \newcounter{problem}
% \newcommand{\problem}[1][]{}
% {\ifthenelse{\isempty{#1}}{\refstepcounter{problem}}{}%
%   \paragraph{Problem \ifthenelse{\isempty{#1}}
%     {\bf\arabic{problem}.}
%     {#1}} }

\newcommand{\problemNumbered}[1]{
  \setcounter{problem}{#1}%
  \addtocounter{problem}{-1}%
  \problem
}
\newenvironment{problemTheorem}{%
  \setcounter{theorem}{\value{problem}}%
  \addtocounter{theorem}{-1}%
  \begin{theorem}[Problem \arabic{problem}]
  }{\end{theorem}}

\newcounter{problem}
\newcommand{\problem}[1][]
{\ifthenelse{\isempty{#1}}
  {\refstepcounter{problem}\paragraph{Problem~\arabic{problem}.}}
  {\paragraph{Problem~#1.}}\refreshHeader} 

\newcounter{exercise}
\newcommand{\exercise}[1][]
{\ifthenelse{\isempty{#1}}
  {\refstepcounter{exercise}\paragraph{Exercise~\theexercise.}}
  {\paragraph{Exercise~#1.}}}

\newcommand{\solution}{\paragraph{Solution.}} 

\newenvironment{mylist}[1][]
{\begin{list}{#1}{
      \setlength{\itemsep}{-1pt}
      \setlength{\labelwidth}{0pt}
      \setlength{\leftmargin}{1.25em}
    }}
{\end{list}}

\newcommand{\LNTitle}{\textcolor{red}{UNDEFINED LNTitle}}

% In a standalone lecture note's table of contents, leave room for numbers
% like 7.10 (in bold) and 7.10.2.
\makeatletter
\newcommand{\LNtocWidths}{%
  \patchcmd{\l@section}{1.5em}{2.8em}{}{\PackageWarning{macros}{Could not widen section numbers in the contents}}%
  \renewcommand*{\l@subsection}{\@dottedtocline{2}{2.8em}{3.2em}}%
  \renewcommand*{\l@subsubsection}{\@dottedtocline{3}{6.0em}{4.1em}}}
\makeatother

\newcommand{\doTitle}[3]{
  \pagenumbering{roman}
  \renewcommand{\LNTitle}{Lecture Note #1, #3}
  \renewcommand{\currentLN}{#1}
  \codeListingstrue
  % number sections, theorems etc., exercises, figures, tables and equations
  % within the lecture note (e.g. Section 8.2, Lemma 8.6), as in the book
  \renewcommand{\thesection}{#1.\arabic{section}}
  \renewcommand{\thelemma}{#1.\arabic{lemma}}
  \renewcommand{\thedefinition}{#1.\arabic{definition}}
  \renewcommand{\thetheorem}{#1.\arabic{theorem}}
  \renewcommand{\thecorollary}{#1.\arabic{corollary}}
  \renewcommand{\theclaim}{#1.\arabic{claim}}
  \renewcommand{\theobservation}{#1.\arabic{observation}}
  \renewcommand{\theconjecture}{#1.\arabic{conjecture}}
  \renewcommand{\theexercise}{#1.\arabic{exercise}}
  \renewcommand{\thefigure}{#1.\arabic{figure}}
  \renewcommand{\thetable}{#1.\arabic{table}}
  \renewcommand{\theequation}{#1.\arabic{equation}}
  \LNtocWidths
  \title{Lecture Note #1\\ #2}
  % a note's \LNabstract, if it defines one before \doTitle, goes right below the title
  \date{\ifdefined\LNabstract\parbox{0.85\textwidth}{\normalsize\normalfont\LNabstract}\\[3ex]\fi
    draft of \today, by \href{\authorURL}{N.~Young}\\[2ex] {\small\copyrightNotice}}
  \maketitle
  {\small
    \hypertarget{TOC}{}
    \tableofcontents
  }
  \thispagestyle{empty}
  % the text starts on page 1, after the title and contents (numbered i, ii, ...)
  \clearpage
  \pagenumbering{arabic}
}

\def\clap#1{\hbox to 0pt{\hss#1\hss}}

\newcommand{\refreshHeader}{}
\newcommand{\header}{}

\newcommand{\doHeader}[1]{      % for lecture notes
  \renewcommand{\header}{%
    \clap{\raisebox{0pt}{\protect\hyperlink{TOC}{$\uparrow$\scriptsize\,\textsc{contents}\hspace{1.1in}}}}%
    {\footnotesize \LNTitle{} / Section \thesection, #1}} %, #1, draft of \today}}
  \markboth{\header}{\header}
}

\newcommand{\SID}{}
\newcommand{\setAuthor}[1]{\renewcommand{\author}{#1}}
\newcommand{\setSID}[1]{\renewcommand{\SID}{#1}}

\newcommand{\doHwHeader}[1]{                                         % for homeworks
  \renewcommand{\refreshHeader}{
    \renewcommand{\header}{%
      {\footnotesize \classSection #1 / Problem~\arabic{problem} / draft of \today \hfill \author\ifGradescope~(\SID)\fi~~}}
    \markboth{\header}{\header}
  }
  \renewcommand{\header}{%
    {\footnotesize \classSection #1 / Instructions / draft of \today \hfill \author\ifGradescope~(\SID)\fi~~}}
  \markboth{\header}{\header}
}

% Homework settings, for instructors using the homeworks in a course:
% \classSection is the course name shown in each homework's running head
% (e.g. \renewcommand{\classSection}{CS 3000\xspace}), and \Gradescopetrue turns on
% the instructions for submitting to Gradescope (in its fixed-length mode),
% the hints saying on which page of the PDF each answer should be, and the
% student ID in the running head.
\newcommand{\classSection}{}
\newif\ifGradescope
\Gradescopefalse
\newcommand{\GradescopeComment}[1]{\ifGradescope\comment{#1}\fi}

\pagestyle{myheadings}
\addtolength{\headsep}{0.2in}
\addtolength{\topmargin}{-0.3in}
\addtolength{\textheight}{0.4in}

%%%%%%%%%%%%%%%%%%%%%%%%%%%%%%%%%%%%%%%%%%%%%%%%
%%%%%%%%%%%%%%%%%%%%%%%%%%%%%%%%%%%%%%%%%%%%%%%%

\newcommand{\MF}{\href
  {https://mfleck.cs.illinois.edu/building-blocks/}
  {MF}\xspace}

\newcommand{\KWSlides}{\href
  {https://www.cs.princeton.edu/~wayne/kleinberg-tardos/}
  {KW slides}\xspace}

\newcommand{\LLM}{\href
  {https://courses.csail.mit.edu/6.042/spring18/mcs.pdf}
  {LLM}\xspace}

% \DPV, \KT, \CLRS: textbooks listed (without links) in the table of sources in
% the Preliminaries note; each prints as a link to its entry in that table.
\DeclareRobustCommand{\DPV}{\noteLink{prelim}{source.DPV}{DPV}\xspace}
\DeclareRobustCommand{\KT}{\noteLink{prelim}{source.KT}{KT}\xspace}
\DeclareRobustCommand{\CLRS}{\noteLink{prelim}{source.CLRS}{CLRS}\xspace}
\AtBeginDocument{\pdfstringdefDisableCommands{\def\DPV{DPV }\def\KT{KT }\def\CLRS{CLRS }}}

\newcommand{\JE}{\href
  {https://jeffe.cs.illinois.edu/teaching/algorithms/}
  {JE}\xspace}

\newcommand{\JEGreedy}{\href
  {https://jeffe.cs.illinois.edu/teaching/algorithms/book/04-greedy.pdf}
  {JE Chapter 4}\xspace}

\newcommand{\BBCode}{\href
{https://northeastern.blackboard.com/webapps/cmsmain/webui/courses/CS3000.MERGED.202010/code?action=frameset&subaction=view&uniq=-jj2max&course_id=_2597392_1&mask=\%2Fcourses\%2FCS3000.MERGED.202010}
  {code}\xspace}

% The Python code (src/code).  Each lecture note that has code ends with a
% section listing its files: \codeSection{path, path, ...}.  In the text,
% \pycode{greedy_algorithms/Dijkstra.py} prints the path as a link to its listing
% (in the homeworks, which have no listings, to the file on GitHub), and
% \pycode[text]{...} uses the given text instead.  \pythonCode{path, ...} is the
% standard sentence ending a section whose algorithms have code:
% "Python code: path, ... (Section N.k)."
\newcommand{\codeFolder}{https://github.com/nealeyoung/algs101/tree/main/src/code}
\newcommand{\codeRepo}{https://github.com/nealeyoung/algs101/blob/main/src/code}
\newif\ifcodeListings   % set by \doTitle: this document has the listings
\newcommand{\githubCode}[2][]{\href{\codeRepo/#2}{\ifx\relax#1\relax\nolinkurl{#2}\else#1\fi}}
\newcommand{\pycode}[2][]{%
  \ifcodeListings\hyperref[code: #2]{\ifx\relax#1\relax\nolinkurl{#2}\else#1\fi}%
  \else\githubCode[#1]{#2}\fi}
\makeatletter
\newcommand{\pythonCode}[1]{%
  \def\pythonCode@first{}\def\pythonCode@list{}%
  \forcsvlist{\pythonCode@add}{#1}%
  Python code: \pythonCode@list\ (Section~\ref{code: \pythonCode@first}).}
\newcommand{\pythonCode@add}[1]{%
  \ifx\pythonCode@first\@empty\def\pythonCode@first{#1}\appto\pythonCode@list{\pycode{#1}}%
  \else\appto\pythonCode@list{, \pycode{#1}}\fi}
\makeatother

% Code listings that copy and paste exactly (in PDF viewers that honor the
% ToUnicode maps, e.g. MuPDF-based ones; PDFKit viewers such as Preview and
% Skim drop leading spaces, so each listing's heading links to the file on
% GitHub).  Spaces are typeset as white visible-space glyphs that copy as
% spaces, and hyphens and asterisks as the T1 glyphs, so they copy as themselves.
\usepackage[T1,OT1]{fontenc}   % T1 for the code listings only; OT1 stays the default
\usepackage{listings}
\pdfgentounicode=1
\input glyphtounicode
\pdfglyphtounicode{uni2423}{0020}
\makeatletter
\def\lst@visiblespace{\textcolor{white}{\char32}}
% put an invisible space on each blank line, so that blank lines copy too
% (listings' OnEmptyLine hook runs before the line break, so patch the loop)
\def\lst@DoNewLines{%
    \@whilenum\lst@newlines>\@ne \do {\lst@NewLine \lst@visiblespace}%
    \ifnum\lst@newlines>\z@ \lst@NewLine \fi}
\makeatother
\lstset{language=Python,
  basicstyle=\fontencoding{T1}\fontfamily{lmtt}\selectfont\footnotesize,
  upquote=true, showstringspaces=false, columns=fullflexible, keepspaces=true,
  showspaces=true, literate={-}{{\char45}}1 {*}{{\char42}}1,
  frame=single, framesep=1ex, rulecolor=\color{gray}}
% the listings' paths are relative to the lecture note's folder (in the book, to src/)
\ifdefined\BookMode\newcommand{\codeDir}{code/}\else\newcommand{\codeDir}{../code/}\fi
% \codefile{path}: a listing of src/code/path, headed by a link to it on GitHub
\newcommand{\codefile}[1]{%
  \subsection*{\githubCode{#1}}\label{code: #1}%
  \lstinputlisting{\codeDir#1}}
% \codeSection{path, path, ...}: the lecture note's last section, listing its code
\newcommand{\codeSection}[1]{%
  \newpage
  \section{Python code}\label{\LNtag{\currentLN}: sec: code}\doHeader{Python code}
  This section lists the Python code for this lecture note.
  Click on a file's name to open the file on GitHub.
  Copying code from GitHub is more reliable than copying it from this PDF,
  since some PDF viewers drop the leading spaces (the indentation) of copied lines.
  \forcsvlist{\codefile}{#1}}
% the copyright notice, for title pages and the homework first pages
\newcommand{\licenseURL}{https://github.com/nealeyoung/algs101/blob/main/LICENSE}
\newcommand{\authorURL}{https://www.cs.ucr.edu/\string~neal}
\newcommand{\copyrightNotice}{%
  \copyright\ 2018--2026 \href{\authorURL}{Neal E. Young}.
  Text: CC BY-NC-SA 4.0; code: MIT; except where noted.
  See \url{\licenseURL}.}
% the published site: PDFs of the notes, and the homework assignments and templates
\newcommand{\siteURL}{https://nealeyoung.github.io/algs101/}

\newcommand{\blackboard}{\href
  {https://northeastern.blackboard.com/webapps/blackboard/execute/launcher?type=Course&id=_2597392_1&url=}
  {Blackboard}\xspace}

\newcommand{\JUSlides}{\href
  {http://www.khoury.neu.edu/home/jullman/cs3000f18/schedule.html}
  {Prof.~Ullman's slides}\xspace}

% a URL on its own line, in three layouts
\newcommand{\URL}[1]{\par-- \parbox[t]{0.8\textwidth}{\url{#1}} \medskip\par}
\newcommand{\URLsmall}[1]{\par-- \parbox[t]{0.95\textwidth}{\small\url{#1}} \medskip\par}
\newcommand{\URLplain}[1]{\par\parbox[t]{\textwidth}{\url{#1}}\par}

%%%%%%%%%%%%%%%%%%%%%%%%%%%%%%%%%%%%%%%%%%%%%%%%
%%%%%%%%%%%%%%%%%%%%%%%%%%%%%%%%%%%%%%%%%%%%%%%%

\newcommand{\continued}{\vfill\centerline{\emph{(continued)}}\vfill\newpage}

% load hyperref last to avoid conflicts; long URLs may also break after hyphens
\PassOptionsToPackage{hyphens}{url}
\ifdefined\BookMode\else\usepackage{xr-hyper}\fi   % for \lectureNoteRefs (homeworks), below
\usepackage{hyperref}
\hypersetup{
    colorlinks,
    linkcolor={blue!50!black},
    citecolor={blue!50!black},
    urlcolor={blue!80!black}
}

\providecommand{\answerDirectory}{.}
\newcommand{\inputAnswer}[1]{\refreshHeader\input{\answerDirectory/#1}}
\newcommand{\inputSection}[1]{\input{#1}}

% Refer to a lecture note by its tag (the prefix of its labels), e.g.
% \LN{proofs} prints "Lecture Note 2" (in the book, a link to that chapter).  The table below is the only place the numbers appear.
\makeatletter
\newcommand{\defineLN}[3]{\@namedef{LNnumber@#2}{#1}\@namedef{LNtag@#1}{#2}\@namedef{LNfile@#2}{#3}}
\newcommand{\LNnumber}[1]{%
  \ifcsname LNnumber@#1\endcsname\csname LNnumber@#1\endcsname
  \else\PackageError{macros}{Unknown lecture-note tag `#1'}{}\fi}
\newcommand{\LNtag}[1]{\csname LNtag@#1\endcsname}
\newcommand{\LNfile}[1]{\csname LNfile@#1\endcsname}
\makeatother
\defineLN{1}{prelim}{1_preliminaries}
\defineLN{2}{proofs}{2_proofs}
\defineLN{3}{matching}{3_stable_matching}
\defineLN{4}{dc}{4_divide_and_conquer}
\defineLN{5}{greedy}{5_greedy}
\defineLN{6}{graphs}{6_graph_traversal}
\defineLN{7}{dp}{7_dynamic_programming}
\defineLN{8}{flow}{8_flow_reductions_LP}
\defineLN{9}{np}{9_NP}

% The text is the same in the book and in the standalone lecture notes; only
% links differ.  In the book, links to other lecture notes stay in the book.
% In a standalone note, they open the other note's published PDF.
% \noteLink{tag}{target}{text} links text to \hypertarget{target} in lecture note tag.
% \LN{proofs} prints "Lecture Note 2", linked to that lecture note.
\newcommand{\currentLN}{0}   % in a standalone note, its number (set by \doTitle)
\newcommand{\notePDF}[1]{\siteURL notes/\LNfile{#1}.pdf}
\newcommand{\noteLink}[3]{%
  \ifnum\LNnumber{#1}=\currentLN\relax\hyperlink{#2}{#3}%
  \else\href{\notePDF{#1}\#nameddest=#2}{#3}\fi}
% \bookLink{target}{text} links text to \hypertarget{target} in the book's
% appendices (from a standalone note, in the published book).
\DeclareRobustCommand{\bookLink}[2]{%
  \ifdefined\BookMode\hyperlink{#1}{#2}%
  \else\href{\siteURL lecture_notes_on_algorithms.pdf\#nameddest=#1}{#2}\fi}
\DeclareRobustCommand{\LN}[1]{%
  \ifnum\LNnumber{#1}=\currentLN\relax Lecture Note~\LNnumber{#1}%
  \else\href{\notePDF{#1}}{Lecture Note~\LNnumber{#1}}\fi}
\pdfstringdefDisableCommands{\def\LN#1{Lecture Note \LNnumber{#1}}}

% Homeworks refer to the notes' sections, exercises and lemmas by \ref: each
% homework's main file says \lectureNoteRefs after \input{macros}.  A homework
% folder has a link refs/ to src/refs/, which holds each note's label records
% (made by tools/refs.sh from the notes' builds).  Read with xr-hyper, they make
% \ref{greedy: exercise: ...} give the exercise's number in the published note,
% linked to it there; labels only the book has (its appendices) link to the book.
\ifdefined\BookMode\else
  \makeatletter
  \newcommand{\lectureNoteRefs}{%
    \@for\LN@tag:=prelim,proofs,matching,dc,greedy,graphs,dp,flow,np\do{%
      \externaldocument{refs/\LNfile{\LN@tag}}[\notePDF{\LN@tag}]}%
    \externaldocument{refs/book}[\siteURL lecture_notes_on_algorithms.pdf]}
  % hyperref links a \ref to another document's label as URL#anchor; browsers'
  % PDF viewers need URL#nameddest=anchor (the form \noteLink and \bookLink use)
  \def\Hy@setref@link#1#2#3#4#5#6\@nil#7{%
    \begingroup
      \ifx\\#5\\\toks0={\hyper@@link{#5}{#4}}\else\toks0={\hyper@@link{#5}{nameddest=#4}}\fi
      \toks1=\expandafter{#7{#1}{#2}{#3}{#4}{#5}}%
      \edef\x{\endgroup\the\toks0 {\the\toks1 }}%
    \x}
  \makeatother
\fi

% Book mode: main.tex defines \BookMode before \input{macros} to combine the
% lecture notes into one report-class document, one chapter per note.
% Without \BookMode (each lecture note on its own, and the homeworks) none
% of this applies.
\ifdefined\BookMode
  \usepackage{etoc}

  % links to other lecture notes stay in the book
  \renewcommand{\noteLink}[3]{\hyperlink{#2}{#3}}

  % each chapter is a lecture note, and is called one, as in the standalone notes
  \renewcommand{\chaptername}{Lecture Note}
  \setcounter{secnumdepth}{3}   % number subsubsections, as in the standalone notes

  % global table of contents, for main.tex
  \newcommand{\bookContents}{%
    \hypertarget{TOC}{}%
    \etocsetnexttocdepth{section}%
    \tableofcontents}

  % \doTitle{N}{title}{short title} starts chapter N, with a link to the global
  % table of contents and the chapter's own table of contents
  \renewcommand{\doTitle}[3]{%
    \setcounter{chapter}{\numexpr#1-1\relax}%
    \chapter{\texorpdfstring{#2}{#3}}%
    \label{\LNtag{#1}: chapter}%
    \csgdef{LNshorttitle@#1}{#3}%
    \csgdef{LNtitle@#1}{#2}%
    \renewcommand{\LNTitle}{Lecture Note #1, #3}%
    \codeListingstrue
    \hyperlink{TOC}{$\uparrow$\scriptsize\,\textsc{contents}}%
    {\small
      \hypertarget{TOC-#1}{}%
      \etocsettocstyle{\section*{\contentsname}}{}%
      \etocsetnexttocdepth{subsubsection}%
      \localtableofcontents
    }}

  % the code: each lecture note's \codeSection just records its files, and
  % \codeAppendix (in the book's appendix) lists them all, a section per note
  \makeatletter   % (file names have no spaces, so drop the spaces around them)
  \renewcommand{\codeSection}[1]{\csxdef{LNcode@\arabic{chapter}}{\zap@space#1 \@empty}}
  \makeatother
  % \codeAppendix: first a contents list (each note's topic, then its files,
  % linked to the listings), then a section per note listing its files
  \makeatletter
  \def\LNbasename#1{\LN@base#1/\LN@end}   % a/b/c.py -> c.py
  \def\LN@base#1/#2\LN@end{\ifx\relax#2\relax#1\else\LN@base#2\LN@end\fi}
  \makeatother
  \newcommand{\codeContentsFile}[1]{%
    \edef\LNfilename{\LNbasename{#1}}%
    \item \pycode[\texttt{\expandafter\detokenize\expandafter{\LNfilename}}]{#1}}
  \newcommand{\codeAppendix}{%
    \begin{itemize}
    \foreach \n in {1,...,9} {%
      \ifcsdef{LNcode@\n}{%
        \item \hyperref[code: note \n]{\csuse{LNtitle@\n}} (\LN{\LNtag{\n}})
        \begin{itemize}
        \edef\LNcodelist{\csuse{LNcode@\n}}%
        \expandafter\forcsvlist\expandafter\codeContentsFile\expandafter{\LNcodelist}%
        \end{itemize}}{}}
    \end{itemize}
    \newpage
    \foreach \n in {1,...,9} {%
      \ifcsdef{LNcode@\n}{%
        \section[Lecture Note \n: \csuse{LNtitle@\n}]{%
          \texorpdfstring{\LN{\LNtag{\n}}: \csuse{LNtitle@\n}}{Lecture Note \n: \csuse{LNtitle@\n}}}%
        \label{code: note \n}%
        \edef\LNcodelist{\csuse{LNcode@\n}}%
        \expandafter\forcsvlist\expandafter\codefile\expandafter{\LNcodelist}%
        \newpage}{}}}

  % a lecture note is a chapter: \LN{proofs} prints a link "Lecture Note 2"
  \DeclareRobustCommand{\LN}[1]{\hyperref[#1: chapter]{Lecture Note~\ref*{#1: chapter}}}
  \pdfstringdefDisableCommands{\def\LN#1{Lecture Note \LNnumber{#1}}}

  % running head; "contents" links to the current chapter's table of contents
  \renewcommand{\doHeader}[1]{%
    \renewcommand{\header}{%
      \clap{\raisebox{0pt}{\protect\hyperlink{TOC-\arabic{chapter}}{$\uparrow$\scriptsize\,\textsc{contents}\hspace{1.1in}}}}%
      {\footnotesize \LNTitle{} / Section \thesection, #1}}%
    \markboth{\header}{\header}}

  % number theorems, lemmas, etc. and exercises within chapters (e.g. Lemma 8.1)
  \counterwithin{lemma}{chapter}
  \counterwithin{definition}{chapter}
  \counterwithin{theorem}{chapter}
  \counterwithin{corollary}{chapter}
  \counterwithin{claim}{chapter}
  \counterwithin{observation}{chapter}
  \counterwithin{conjecture}{chapter}
  \counterwithin{exercise}{chapter}
\fi

% %%%%%%%%%%%%%%%%%%%%%%%%%%%%%%%%%%%%%%%%%%%%%%%%%%%%%%%%%%%%%%%%%%%%%%%%
% %%%%%%%%%%%%%%%%%%%%%%%%%%%%%%%%%%%%%%%%%%%%%%%%%%%%%%%%%%%%%%%%%%%%%%%%

\ExplSyntaxOn
\seq_new:N \l_itemize_seq
% \NewDocumentEnvironment{REITEMIZE}{ +b }{
\NewDocumentCommand{\REITEMIZECMD}{ +m }{
  \regex_split:nnN {\c{item}} {#1} \l_itemize_seq
  \seq_pop:NN \l_itemize_seq \l_tmpa_tl% remove the (hopefully) empty first item
  \seq_if_empty:NF \l_itemize_seq {
    \l_tmpa_tl
    \seq_map_inline:Nn \l_itemize_seq {
      \regex_split:nnN {\c{item}} {#1} \l_itemize_seq {
        \regex_extract_once:nnNTF { ^\[([^]]*)\](.*) } { ##1} \l_foo_seq
      { \item[\seq_item:Nn \l_foo_seq {2}]\ITEMIZER{\seq_item:Nn \l_foo_seq {3}} } { \item\ITEMIZER{##1}  }
    }
  }
}
}
\ExplSyntaxOff

\newlength{\ITEMHT}
\setlength{\ITEMHT}{1in}

% \setlength{\fboxsep}{1pt}
\newcommand{\ITEMIZER}[1]{%
  % \framebox{
\adjustbox{valign=t}{\parbox{\textwidth}{\linespread{1.15}\selectfont #1}}~~\adjustbox{valign=t}{\parbox{0pt}{\color{gray}\rule[-\ITEMHT]{0.01pt}{\ITEMHT}}}
% }
  % \begin{tabular}{@{}p{\textwidth}@{\rule[-\ITEMHT]{0pt}{\ITEMHT}}}
  %   #1
  % \end{tabular}
}

\newsavebox\thesmashminipage
\newlength{\thesmashminipagelength}
\newlength{\answerboxlength}
\newlength{\answerboxoverrun}
\newenvironment{smashminipage}[2]
{\setlength\answerboxlength{#2}
\begin{lrbox}{\thesmashminipage}\begin{minipage}[t]{#1}}
{\end{minipage}\end{lrbox}%
\setlength{\thesmashminipagelength}{\dimexpr\ht\thesmashminipage+\dp\thesmashminipage}%
\setlength{\answerboxoverrun}{\dimexpr\thesmashminipagelength-\answerboxlength}%
\ifdim\the\thesmashminipagelength>\the\answerboxlength%
\usebox{\thesmashminipage}%
\smash{\color{red}\hspace*{-0.98\textwidth}\rule[-\answerboxlength]{0.98\textwidth}{1pt}}
\errmessage{Answer exceeds allotted length, please shorten it.}
\else
\usebox{\thesmashminipage}%
\fi
}



\newenvironment{answerBox}[1]{%
  \def\ANSWERBOXHEIGHT{#1}
  \begin{adjustbox}{valign=t}
  \begin{smashminipage}{\textwidth}{\ANSWERBOXHEIGHT}
}{\end{smashminipage}\end{adjustbox}~~~\adjustbox{valign=t}{\parbox{0pt}{\color{gray}\rule[-\ANSWERBOXHEIGHT]{0.01pt}{\ANSWERBOXHEIGHT}\rule[-\ANSWERBOXHEIGHT]{10pt}{0.01pt}}}}

% % \setlist[enumerate]{itemsep=0pt,parsep=0pt,topsep=0pt,leftmargin=100pt}

% \makeatletter 
\LetLtxMacro\SVenumerate\enumerate
\let\endSVenumerate\endenumerate

\newcommand{\injectEnumerate}{
  \let\SValph\alph
  \def\alph{}
  \renewenvironment{enumerate}[1][]
  {\let\alph\SValph \SVenumerate[label=(\alph*),itemsep=2pt,parsep=0pt]\REITEMIZE}{\endREITEMIZE\endSVenumerate}
  % \setlist[enumerate]{before=\Collect@Body\REITEMIZECMD}
}
% \makeatother
\newcommand{\restoreEnumerate}{
  \renewenvironment{enumerate}[1][]
  \let\SValph\alph
  \def\alph{}
  {\let\alph\SValph \SVenumerate[label=(\alph*)]}{\endSVenumerate}
  % \setlist[enumerate]{before=,after=}
}
\newcommand{\injectBlue}{
  \let\SVblue\blue
  \let\endSVblue\endblue
  \def\blue{\SVblue\REITEMIZE}
  \def\endblue{\endREITEMIZE\endSVblue}
}
\newcommand{\restoreBlue}{
  \let\blue\SVblue
  \let\endblue\endSVblue
}
\newcommand{\injectEnumerateBlue}{
  \setlist[enumerate]{before=\injectBlue,after=\restoreBlue,itemsep=2pt,parsep=0pt}
}
\newcommand{\restoreEnumerateBlue}{
  \setlist[enumerate]{before=,after=}
}


% %%%%%%%%%%%%%%%%%%%%%%%%%%%%%%%%%%%%%%%%%%%%%%%%%%%%%%%%%%%%%%%%%%%%%%%%
% %%%%%%%%%%%%%%%%%%%%%%%%%%%%%%%%%%%%%%%%%%%%%%%%%%%%%%%%%%%%%%%%%%%%%%%%

% \NewDocumentEnvironment{env}{ }
% {\Environenv}
% {\endEnvironenv}

% \NewEnviron{REITEMIZE}{\REITEMIZECMD\BODY}

\makeatletter 
\newenvironment{REITEMIZE}{\Collect@Body\REITEMIZECMD}{}
\makeatother 

\newcommand{\NODEone}[1]{\ensuremath{v_{#1}}\xspace}
\newcommand{\NODE}[2]{{\NODEone {#1,#2}}\xspace}
\newcommand{\EDGE}[2]{\ensuremath{(\textsf{#1}, \textsf{#2})}}
\newcommand{\ITEM}{\item}

%%% Local Variables:
%%% mode: latex
%%% End:
 to combine the
% lecture notes into one report-class document, one chapter per note.
% Without \BookMode (each lecture note on its own, and the homeworks) none
% of this applies.
\ifdefined\BookMode
  \usepackage{etoc}

  % links to other lecture notes stay in the book
  \renewcommand{\noteLink}[3]{\hyperlink{#2}{#3}}

  % each chapter is a lecture note, and is called one, as in the standalone notes
  \renewcommand{\chaptername}{Lecture Note}
  \setcounter{secnumdepth}{3}   % number subsubsections, as in the standalone notes

  % global table of contents, for main.tex
  \newcommand{\bookContents}{%
    \hypertarget{TOC}{}%
    \etocsetnexttocdepth{section}%
    \tableofcontents}

  % \doTitle{N}{title}{short title} starts chapter N, with a link to the global
  % table of contents and the chapter's own table of contents
  \renewcommand{\doTitle}[3]{%
    \setcounter{chapter}{\numexpr#1-1\relax}%
    \chapter{\texorpdfstring{#2}{#3}}%
    \label{\LNtag{#1}: chapter}%
    \csgdef{LNshorttitle@#1}{#3}%
    \csgdef{LNtitle@#1}{#2}%
    \renewcommand{\LNTitle}{Lecture Note #1, #3}%
    \codeListingstrue
    \hyperlink{TOC}{$\uparrow$\scriptsize\,\textsc{contents}}%
    {\small
      \hypertarget{TOC-#1}{}%
      \etocsettocstyle{\section*{\contentsname}}{}%
      \etocsetnexttocdepth{subsubsection}%
      \localtableofcontents
    }}

  % the code: each lecture note's \codeSection just records its files, and
  % \codeAppendix (in the book's appendix) lists them all, a section per note
  \makeatletter   % (file names have no spaces, so drop the spaces around them)
  \renewcommand{\codeSection}[1]{\csxdef{LNcode@\arabic{chapter}}{\zap@space#1 \@empty}}
  \makeatother
  % \codeAppendix: first a contents list (each note's topic, then its files,
  % linked to the listings), then a section per note listing its files
  \makeatletter
  \def\LNbasename#1{\LN@base#1/\LN@end}   % a/b/c.py -> c.py
  \def\LN@base#1/#2\LN@end{\ifx\relax#2\relax#1\else\LN@base#2\LN@end\fi}
  \makeatother
  \newcommand{\codeContentsFile}[1]{%
    \edef\LNfilename{\LNbasename{#1}}%
    \item \pycode[\texttt{\expandafter\detokenize\expandafter{\LNfilename}}]{#1}}
  \newcommand{\codeAppendix}{%
    \begin{itemize}
    \foreach \n in {1,...,9} {%
      \ifcsdef{LNcode@\n}{%
        \item \hyperref[code: note \n]{\csuse{LNtitle@\n}} (\LN{\LNtag{\n}})
        \begin{itemize}
        \edef\LNcodelist{\csuse{LNcode@\n}}%
        \expandafter\forcsvlist\expandafter\codeContentsFile\expandafter{\LNcodelist}%
        \end{itemize}}{}}
    \end{itemize}
    \newpage
    \foreach \n in {1,...,9} {%
      \ifcsdef{LNcode@\n}{%
        \section[Lecture Note \n: \csuse{LNtitle@\n}]{%
          \texorpdfstring{\LN{\LNtag{\n}}: \csuse{LNtitle@\n}}{Lecture Note \n: \csuse{LNtitle@\n}}}%
        \label{code: note \n}%
        \edef\LNcodelist{\csuse{LNcode@\n}}%
        \expandafter\forcsvlist\expandafter\codefile\expandafter{\LNcodelist}%
        \newpage}{}}}

  % a lecture note is a chapter: \LN{proofs} prints a link "Lecture Note 2"
  \DeclareRobustCommand{\LN}[1]{\hyperref[#1: chapter]{Lecture Note~\ref*{#1: chapter}}}
  \pdfstringdefDisableCommands{\def\LN#1{Lecture Note \LNnumber{#1}}}

  % running head; "contents" links to the current chapter's table of contents
  \renewcommand{\doHeader}[1]{%
    \renewcommand{\header}{%
      \clap{\raisebox{0pt}{\protect\hyperlink{TOC-\arabic{chapter}}{$\uparrow$\scriptsize\,\textsc{contents}\hspace{1.1in}}}}%
      {\footnotesize \LNTitle{} / Section \thesection, #1}}%
    \markboth{\header}{\header}}

  % number theorems, lemmas, etc. and exercises within chapters (e.g. Lemma 8.1)
  \counterwithin{lemma}{chapter}
  \counterwithin{definition}{chapter}
  \counterwithin{theorem}{chapter}
  \counterwithin{corollary}{chapter}
  \counterwithin{claim}{chapter}
  \counterwithin{observation}{chapter}
  \counterwithin{conjecture}{chapter}
  \counterwithin{exercise}{chapter}
\fi

% %%%%%%%%%%%%%%%%%%%%%%%%%%%%%%%%%%%%%%%%%%%%%%%%%%%%%%%%%%%%%%%%%%%%%%%%
% %%%%%%%%%%%%%%%%%%%%%%%%%%%%%%%%%%%%%%%%%%%%%%%%%%%%%%%%%%%%%%%%%%%%%%%%

\ExplSyntaxOn
\seq_new:N \l_itemize_seq
% \NewDocumentEnvironment{REITEMIZE}{ +b }{
\NewDocumentCommand{\REITEMIZECMD}{ +m }{
  \regex_split:nnN {\c{item}} {#1} \l_itemize_seq
  \seq_pop:NN \l_itemize_seq \l_tmpa_tl% remove the (hopefully) empty first item
  \seq_if_empty:NF \l_itemize_seq {
    \l_tmpa_tl
    \seq_map_inline:Nn \l_itemize_seq {
      \regex_split:nnN {\c{item}} {#1} \l_itemize_seq {
        \regex_extract_once:nnNTF { ^\[([^]]*)\](.*) } { ##1} \l_foo_seq
      { \item[\seq_item:Nn \l_foo_seq {2}]\ITEMIZER{\seq_item:Nn \l_foo_seq {3}} } { \item\ITEMIZER{##1}  }
    }
  }
}
}
\ExplSyntaxOff

\newlength{\ITEMHT}
\setlength{\ITEMHT}{1in}

% \setlength{\fboxsep}{1pt}
\newcommand{\ITEMIZER}[1]{%
  % \framebox{
\adjustbox{valign=t}{\parbox{\textwidth}{\linespread{1.15}\selectfont #1}}~~\adjustbox{valign=t}{\parbox{0pt}{\color{gray}\rule[-\ITEMHT]{0.01pt}{\ITEMHT}}}
% }
  % \begin{tabular}{@{}p{\textwidth}@{\rule[-\ITEMHT]{0pt}{\ITEMHT}}}
  %   #1
  % \end{tabular}
}

\newsavebox\thesmashminipage
\newlength{\thesmashminipagelength}
\newlength{\answerboxlength}
\newlength{\answerboxoverrun}
\newenvironment{smashminipage}[2]
{\setlength\answerboxlength{#2}
\begin{lrbox}{\thesmashminipage}\begin{minipage}[t]{#1}}
{\end{minipage}\end{lrbox}%
\setlength{\thesmashminipagelength}{\dimexpr\ht\thesmashminipage+\dp\thesmashminipage}%
\setlength{\answerboxoverrun}{\dimexpr\thesmashminipagelength-\answerboxlength}%
\ifdim\the\thesmashminipagelength>\the\answerboxlength%
\usebox{\thesmashminipage}%
\smash{\color{red}\hspace*{-0.98\textwidth}\rule[-\answerboxlength]{0.98\textwidth}{1pt}}
\errmessage{Answer exceeds allotted length, please shorten it.}
\else
\usebox{\thesmashminipage}%
\fi
}



\newenvironment{answerBox}[1]{%
  \def\ANSWERBOXHEIGHT{#1}
  \begin{adjustbox}{valign=t}
  \begin{smashminipage}{\textwidth}{\ANSWERBOXHEIGHT}
}{\end{smashminipage}\end{adjustbox}~~~\adjustbox{valign=t}{\parbox{0pt}{\color{gray}\rule[-\ANSWERBOXHEIGHT]{0.01pt}{\ANSWERBOXHEIGHT}\rule[-\ANSWERBOXHEIGHT]{10pt}{0.01pt}}}}

% % \setlist[enumerate]{itemsep=0pt,parsep=0pt,topsep=0pt,leftmargin=100pt}

% \makeatletter 
\LetLtxMacro\SVenumerate\enumerate
\let\endSVenumerate\endenumerate

\newcommand{\injectEnumerate}{
  \let\SValph\alph
  \def\alph{}
  \renewenvironment{enumerate}[1][]
  {\let\alph\SValph \SVenumerate[label=(\alph*),itemsep=2pt,parsep=0pt]\REITEMIZE}{\endREITEMIZE\endSVenumerate}
  % \setlist[enumerate]{before=\Collect@Body\REITEMIZECMD}
}
% \makeatother
\newcommand{\restoreEnumerate}{
  \renewenvironment{enumerate}[1][]
  \let\SValph\alph
  \def\alph{}
  {\let\alph\SValph \SVenumerate[label=(\alph*)]}{\endSVenumerate}
  % \setlist[enumerate]{before=,after=}
}
\newcommand{\injectBlue}{
  \let\SVblue\blue
  \let\endSVblue\endblue
  \def\blue{\SVblue\REITEMIZE}
  \def\endblue{\endREITEMIZE\endSVblue}
}
\newcommand{\restoreBlue}{
  \let\blue\SVblue
  \let\endblue\endSVblue
}
\newcommand{\injectEnumerateBlue}{
  \setlist[enumerate]{before=\injectBlue,after=\restoreBlue,itemsep=2pt,parsep=0pt}
}
\newcommand{\restoreEnumerateBlue}{
  \setlist[enumerate]{before=,after=}
}


% %%%%%%%%%%%%%%%%%%%%%%%%%%%%%%%%%%%%%%%%%%%%%%%%%%%%%%%%%%%%%%%%%%%%%%%%
% %%%%%%%%%%%%%%%%%%%%%%%%%%%%%%%%%%%%%%%%%%%%%%%%%%%%%%%%%%%%%%%%%%%%%%%%

% \NewDocumentEnvironment{env}{ }
% {\Environenv}
% {\endEnvironenv}

% \NewEnviron{REITEMIZE}{\REITEMIZECMD\BODY}

\makeatletter 
\newenvironment{REITEMIZE}{\Collect@Body\REITEMIZECMD}{}
\makeatother 

\newcommand{\NODEone}[1]{\ensuremath{v_{#1}}\xspace}
\newcommand{\NODE}[2]{{\NODEone {#1,#2}}\xspace}
\newcommand{\EDGE}[2]{\ensuremath{(\textsf{#1}, \textsf{#2})}}
\newcommand{\ITEM}{\item}

%%% Local Variables:
%%% mode: latex
%%% End:
 to combine the
% lecture notes into one report-class document, one chapter per note.
% Without \BookMode (each lecture note on its own, and the homeworks) none
% of this applies.
\ifdefined\BookMode
  \usepackage{etoc}

  % links to other lecture notes stay in the book
  \renewcommand{\noteLink}[3]{\hyperlink{#2}{#3}}

  % each chapter is a lecture note, and is called one, as in the standalone notes
  \renewcommand{\chaptername}{Lecture Note}
  \setcounter{secnumdepth}{3}   % number subsubsections, as in the standalone notes

  % global table of contents, for main.tex
  \newcommand{\bookContents}{%
    \hypertarget{TOC}{}%
    \etocsetnexttocdepth{section}%
    \tableofcontents}

  % \doTitle{N}{title}{short title} starts chapter N, with a link to the global
  % table of contents and the chapter's own table of contents
  \renewcommand{\doTitle}[3]{%
    \setcounter{chapter}{\numexpr#1-1\relax}%
    \chapter{\texorpdfstring{#2}{#3}}%
    \label{\LNtag{#1}: chapter}%
    \csgdef{LNshorttitle@#1}{#3}%
    \csgdef{LNtitle@#1}{#2}%
    \renewcommand{\LNTitle}{Lecture Note #1, #3}%
    \codeListingstrue
    \hyperlink{TOC}{$\uparrow$\scriptsize\,\textsc{contents}}%
    {\small
      \hypertarget{TOC-#1}{}%
      \etocsettocstyle{\section*{\contentsname}}{}%
      \etocsetnexttocdepth{subsubsection}%
      \localtableofcontents
    }}

  % the code: each lecture note's \codeSection just records its files, and
  % \codeAppendix (in the book's appendix) lists them all, a section per note
  \makeatletter   % (file names have no spaces, so drop the spaces around them)
  \renewcommand{\codeSection}[1]{\csxdef{LNcode@\arabic{chapter}}{\zap@space#1 \@empty}}
  \makeatother
  % \codeAppendix: first a contents list (each note's topic, then its files,
  % linked to the listings), then a section per note listing its files
  \makeatletter
  \def\LNbasename#1{\LN@base#1/\LN@end}   % a/b/c.py -> c.py
  \def\LN@base#1/#2\LN@end{\ifx\relax#2\relax#1\else\LN@base#2\LN@end\fi}
  \makeatother
  \newcommand{\codeContentsFile}[1]{%
    \edef\LNfilename{\LNbasename{#1}}%
    \item \pycode[\texttt{\expandafter\detokenize\expandafter{\LNfilename}}]{#1}}
  \newcommand{\codeAppendix}{%
    \begin{itemize}
    \foreach \n in {1,...,9} {%
      \ifcsdef{LNcode@\n}{%
        \item \hyperref[code: note \n]{\csuse{LNtitle@\n}} (\LN{\LNtag{\n}})
        \begin{itemize}
        \edef\LNcodelist{\csuse{LNcode@\n}}%
        \expandafter\forcsvlist\expandafter\codeContentsFile\expandafter{\LNcodelist}%
        \end{itemize}}{}}
    \end{itemize}
    \newpage
    \foreach \n in {1,...,9} {%
      \ifcsdef{LNcode@\n}{%
        \section[Lecture Note \n: \csuse{LNtitle@\n}]{%
          \texorpdfstring{\LN{\LNtag{\n}}: \csuse{LNtitle@\n}}{Lecture Note \n: \csuse{LNtitle@\n}}}%
        \label{code: note \n}%
        \edef\LNcodelist{\csuse{LNcode@\n}}%
        \expandafter\forcsvlist\expandafter\codefile\expandafter{\LNcodelist}%
        \newpage}{}}}

  % a lecture note is a chapter: \LN{proofs} prints a link "Lecture Note 2"
  \DeclareRobustCommand{\LN}[1]{\hyperref[#1: chapter]{Lecture Note~\ref*{#1: chapter}}}
  \pdfstringdefDisableCommands{\def\LN#1{Lecture Note \LNnumber{#1}}}

  % running head; "contents" links to the current chapter's table of contents
  \renewcommand{\doHeader}[1]{%
    \renewcommand{\header}{%
      \clap{\raisebox{0pt}{\protect\hyperlink{TOC-\arabic{chapter}}{$\uparrow$\scriptsize\,\textsc{contents}\hspace{1.1in}}}}%
      {\footnotesize \LNTitle{} / Section \thesection, #1}}%
    \markboth{\header}{\header}}

  % number theorems, lemmas, etc. and exercises within chapters (e.g. Lemma 8.1)
  \counterwithin{lemma}{chapter}
  \counterwithin{definition}{chapter}
  \counterwithin{theorem}{chapter}
  \counterwithin{corollary}{chapter}
  \counterwithin{claim}{chapter}
  \counterwithin{observation}{chapter}
  \counterwithin{conjecture}{chapter}
  \counterwithin{exercise}{chapter}
\fi

% %%%%%%%%%%%%%%%%%%%%%%%%%%%%%%%%%%%%%%%%%%%%%%%%%%%%%%%%%%%%%%%%%%%%%%%%
% %%%%%%%%%%%%%%%%%%%%%%%%%%%%%%%%%%%%%%%%%%%%%%%%%%%%%%%%%%%%%%%%%%%%%%%%

\ExplSyntaxOn
\seq_new:N \l_itemize_seq
% \NewDocumentEnvironment{REITEMIZE}{ +b }{
\NewDocumentCommand{\REITEMIZECMD}{ +m }{
  \regex_split:nnN {\c{item}} {#1} \l_itemize_seq
  \seq_pop:NN \l_itemize_seq \l_tmpa_tl% remove the (hopefully) empty first item
  \seq_if_empty:NF \l_itemize_seq {
    \l_tmpa_tl
    \seq_map_inline:Nn \l_itemize_seq {
      \regex_split:nnN {\c{item}} {#1} \l_itemize_seq {
        \regex_extract_once:nnNTF { ^\[([^]]*)\](.*) } { ##1} \l_foo_seq
      { \item[\seq_item:Nn \l_foo_seq {2}]\ITEMIZER{\seq_item:Nn \l_foo_seq {3}} } { \item\ITEMIZER{##1}  }
    }
  }
}
}
\ExplSyntaxOff

\newlength{\ITEMHT}
\setlength{\ITEMHT}{1in}

% \setlength{\fboxsep}{1pt}
\newcommand{\ITEMIZER}[1]{%
  % \framebox{
\adjustbox{valign=t}{\parbox{\textwidth}{\linespread{1.15}\selectfont #1}}~~\adjustbox{valign=t}{\parbox{0pt}{\color{gray}\rule[-\ITEMHT]{0.01pt}{\ITEMHT}}}
% }
  % \begin{tabular}{@{}p{\textwidth}@{\rule[-\ITEMHT]{0pt}{\ITEMHT}}}
  %   #1
  % \end{tabular}
}

\newsavebox\thesmashminipage
\newlength{\thesmashminipagelength}
\newlength{\answerboxlength}
\newlength{\answerboxoverrun}
\newenvironment{smashminipage}[2]
{\setlength\answerboxlength{#2}
\begin{lrbox}{\thesmashminipage}\begin{minipage}[t]{#1}}
{\end{minipage}\end{lrbox}%
\setlength{\thesmashminipagelength}{\dimexpr\ht\thesmashminipage+\dp\thesmashminipage}%
\setlength{\answerboxoverrun}{\dimexpr\thesmashminipagelength-\answerboxlength}%
\ifdim\the\thesmashminipagelength>\the\answerboxlength%
\usebox{\thesmashminipage}%
\smash{\color{red}\hspace*{-0.98\textwidth}\rule[-\answerboxlength]{0.98\textwidth}{1pt}}
\errmessage{Answer exceeds allotted length, please shorten it.}
\else
\usebox{\thesmashminipage}%
\fi
}



\newenvironment{answerBox}[1]{%
  \def\ANSWERBOXHEIGHT{#1}
  \begin{adjustbox}{valign=t}
  \begin{smashminipage}{\textwidth}{\ANSWERBOXHEIGHT}
}{\end{smashminipage}\end{adjustbox}~~~\adjustbox{valign=t}{\parbox{0pt}{\color{gray}\rule[-\ANSWERBOXHEIGHT]{0.01pt}{\ANSWERBOXHEIGHT}\rule[-\ANSWERBOXHEIGHT]{10pt}{0.01pt}}}}

% % \setlist[enumerate]{itemsep=0pt,parsep=0pt,topsep=0pt,leftmargin=100pt}

% \makeatletter 
\LetLtxMacro\SVenumerate\enumerate
\let\endSVenumerate\endenumerate

\newcommand{\injectEnumerate}{
  \let\SValph\alph
  \def\alph{}
  \renewenvironment{enumerate}[1][]
  {\let\alph\SValph \SVenumerate[label=(\alph*),itemsep=2pt,parsep=0pt]\REITEMIZE}{\endREITEMIZE\endSVenumerate}
  % \setlist[enumerate]{before=\Collect@Body\REITEMIZECMD}
}
% \makeatother
\newcommand{\restoreEnumerate}{
  \renewenvironment{enumerate}[1][]
  \let\SValph\alph
  \def\alph{}
  {\let\alph\SValph \SVenumerate[label=(\alph*)]}{\endSVenumerate}
  % \setlist[enumerate]{before=,after=}
}
\newcommand{\injectBlue}{
  \let\SVblue\blue
  \let\endSVblue\endblue
  \def\blue{\SVblue\REITEMIZE}
  \def\endblue{\endREITEMIZE\endSVblue}
}
\newcommand{\restoreBlue}{
  \let\blue\SVblue
  \let\endblue\endSVblue
}
\newcommand{\injectEnumerateBlue}{
  \setlist[enumerate]{before=\injectBlue,after=\restoreBlue,itemsep=2pt,parsep=0pt}
}
\newcommand{\restoreEnumerateBlue}{
  \setlist[enumerate]{before=,after=}
}


% %%%%%%%%%%%%%%%%%%%%%%%%%%%%%%%%%%%%%%%%%%%%%%%%%%%%%%%%%%%%%%%%%%%%%%%%
% %%%%%%%%%%%%%%%%%%%%%%%%%%%%%%%%%%%%%%%%%%%%%%%%%%%%%%%%%%%%%%%%%%%%%%%%

% \NewDocumentEnvironment{env}{ }
% {\Environenv}
% {\endEnvironenv}

% \NewEnviron{REITEMIZE}{\REITEMIZECMD\BODY}

\makeatletter 
\newenvironment{REITEMIZE}{\Collect@Body\REITEMIZECMD}{}
\makeatother 

\newcommand{\NODEone}[1]{\ensuremath{v_{#1}}\xspace}
\newcommand{\NODE}[2]{{\NODEone {#1,#2}}\xspace}
\newcommand{\EDGE}[2]{\ensuremath{(\textsf{#1}, \textsf{#2})}}
\newcommand{\ITEM}{\item}

%%% Local Variables:
%%% mode: latex
%%% End:
.  A homework
% folder has a link refs/ to src/refs/, which holds each note's label records
% (made by tools/refs.sh from the notes' builds).  Read with xr-hyper, they make
% \ref{greedy: exercise: ...} give the exercise's number in the published note,
% linked to it there; labels only the book has (its appendices) link to the book.
\ifdefined\BookMode\else
  \makeatletter
  \newcommand{\lectureNoteRefs}{%
    \@for\LN@tag:=prelim,proofs,matching,dc,greedy,graphs,dp,flow,np\do{%
      \externaldocument{refs/\LNfile{\LN@tag}}[\notePDF{\LN@tag}]}%
    \externaldocument{refs/book}[\siteURL lecture_notes_on_algorithms.pdf]}
  % hyperref links a \ref to another document's label as URL#anchor; browsers'
  % PDF viewers need URL#nameddest=anchor (the form \noteLink and \bookLink use)
  \def\Hy@setref@link#1#2#3#4#5#6\@nil#7{%
    \begingroup
      \ifx\\#5\\\toks0={\hyper@@link{#5}{#4}}\else\toks0={\hyper@@link{#5}{nameddest=#4}}\fi
      \toks1=\expandafter{#7{#1}{#2}{#3}{#4}{#5}}%
      \edef\x{\endgroup\the\toks0 {\the\toks1 }}%
    \x}
  \makeatother
\fi

% Book mode: main.tex defines \BookMode before %%%% DON'T EDIT THIS FILE %%%%

\usepackage{xifthen}% \isempty
\usepackage{etoolbox}% \patchcmd
% \usepackage{stringstrings}% \substring

\usepackage{amsmath, amsthm, amssymb}
\usepackage[margin=1in]{geometry}
\usepackage{xspace}
\usepackage{graphicx}
\usepackage[normalem]{ulem}
\usepackage{adjustbox}
\usepackage{xcolor}
\usepackage{changepage}
\usepackage{framed}
\usepackage{multicol}
\usepackage{enumitem}
% \usepackage{listings} % lstlistings environment

\usepackage{makecmds}           %\provideenvironment

\usepackage{xparse,environ}
\usepackage{letltxmacro}       % \LetLtxMacro
\usepackage{wrapfig}
\usepackage{tikz}

\usetikzlibrary{shapes}
\usetikzlibrary{positioning}
\usetikzlibrary{shapes.misc}
\usetikzlibrary{arrows}

\newtheorem{lemma}{Lemma}
\newtheorem{definition}{Definition}
\newtheorem{theorem}{Theorem}
\newtheorem{corollary}{Corollary}
\newtheorem{claim}{Claim}
\newtheorem{observation}{Observation}
\newtheorem{conjecture}{Conjecture}

% A lemma or theorem restated later with its original number (set with
% \setcounter) goes in a restatement environment, which gives it its own link
% anchor (otherwise it would duplicate the original's).
\newenvironment{restatement}
  {\def\theHlemma{restated.\thelemma}\def\theHtheorem{restated.\thetheorem}}{}

\newenvironment{myframed}
{\smallskip\par\noindent\begin{minipage}{\textwidth}\begin{framed}\vspace*{-5pt}}{\vspace*{-5pt}\end{framed}\end{minipage}\smallskip}


\definecolor{darkblue}{RGB}{0,0,200}

 
\let\oldParagraph\paragraph
\providecommand{\gradescopeColor}{\color{black}}
\renewcommand{\paragraph}[1]{\oldParagraph{\gradescopeColor #1}}
\providecommand{\gradescopeInit}{}

\providecommand{\BLUE}[1]{{\color{darkblue} #1}}
\providecommand{\REPLACEME}{\textcolor{red}{\ensuremath{\star \cdots \star}}\xspace}
\provideenvironment{blue}{\color{darkblue}}{}

\newcommand{\dist}[2]{\ensuremath{\operatorname{\sf dist}(#1, #2)}}
\newcommand{\depth}[2]{\ensuremath{\operatorname{\sf depth}_{#1}(#2)}}
\newcommand{\cost}[1]{\ensuremath{\operatorname{\sf cost}(#1)}}
\newcommand{\weight}[1]{\ensuremath{\operatorname{\sf wt}(#1)}}
% \newcommand{\val}[1]{\ensuremath{\operatorname{\sf val}(#1)}}
\newcommand{\eps}{\varepsilon}
\newcommand{\capac}{\textsf{cap}}
\newcommand{\capOf}[1]{\ensuremath{\operatorname{\sf cap}(#1)}}
\newcommand{\rescap}[1]{\ensuremath{\operatorname{\sf cap}_f(#1)}}
\newcommand{\into}[1]{\ensuremath{\operatorname{\sf in}(#1)}}
\newcommand{\out}[1]{\ensuremath{\operatorname{\sf out}(#1)}}
\newenvironment{LP}[1]{
  \begin{array}[t]{@{}l@{}}
    \displaystyle #1 \\
    \begin{aligned}\\[-10pt]}{\end{aligned}\end{array}}

\newcommand{\opt}[1]{\ensuremath{\operatorname{\sf opt}(#1)}}
\newcommand{\Opt}[1]{\ensuremath{\operatorname{\sf Opt}(#1)}}
\newcommand{\half}{\frac{1}{2}}

\newcommand{\E}{\operatorname{E}}
\newcommand{\giv}{\,|\,}
\newcommand{\Z}{\mathbb{Z}}
\newcommand{\Zp}{\mathbb{Z}_{\ge 0}}
\newcommand{\R}{\mathbb{R}}
% \newcommand{\Rp}{\mathbb{R}_{\ge 0}}
\newcommand{\Rp}{\mathbb{R}_{\tiny +}}
\newcommand{\N}{\mathbb{N}}
\newcommand{\Np}{\mathbb{N}_{\ge 0}}
\newcommand{\calS}{{\cal S}}

\newcommand{\OP}[1]{\text{\small\sc #1}\xspace}
\newcommand{\Insert}{\OP{Insert}}
\newcommand{\FindMin}{\OP{Find-Min}}
\newcommand{\DeleteMin}{\OP{Delete-Min}}
\newcommand{\DecreaseKey}{\OP{Decrease}}
\newcommand{\Merge}{\OP{Merge}}
\newcommand{\Cut}{\OP{Cut}}
% \newcommand{\prob}[1]{\text{\sc #1}}
% a problem name; in running text, lines may break inside it
\DeclareRobustCommand{\prob}[1]{\ifmmode\text{#1}\else#1\fi}
\newcommand{\encode}[1]{\langle #1 \rangle}

\newcommand{\best}[1]{\text{\sf best}(#1)}
\newcommand{\worst}[1]{\text{\sf worst}(#1)}
\newcommand{\floor}[1]{\lfloor #1 \rfloor}
\newcommand{\ceil}[1]{\lceil #1 \rceil}

\setlength{\parskip}{2pt}
\setlength{\parindent}{1em}

% \setlength\fboxsep{0pt}

%%% longFormProof, shortFormProof, block environments
%%% \step

\newlist{steps}{enumerate}{7}
\setlist[steps]{wide,
  topsep=1pt,
  parsep=1pt,
  partopsep=1pt,
  itemsep=0pt,
  align=left,
  labelindent=0pt,
  itemindent=0pt, 
}
\setlist[steps,1]{
  label=\arabic*.,
  ref=\arabic*
}
\setlist[steps,2]{
  label=\arabic{stepsi}.\arabic*.,
  ref=\arabic{stepsi}.\arabic*
}
\setlist[steps,3]{
  label=\arabic{stepsi}.\arabic{stepsii}.\arabic*.,
  ref=\arabic{stepsi}.\arabic{stepsii}.\arabic*
}
\setlist[steps,4]{
  label=\arabic{stepsi}.\arabic{stepsii}.\arabic{stepsiii}.\arabic*.,
  ref=\arabic{stepsi}.\arabic{stepsii}.\arabic{stepsiii}.\arabic*
}
\setlist[steps,5]{
  label=\arabic{stepsi}.\arabic{stepsii}.\arabic{stepsiii}.\arabic{stepsiv}.\arabic*.,
  ref=\arabic{stepsi}.\arabic{stepsii}.\arabic{stepsiii}.\arabic{stepsiv}.\arabic*
}
\setlist[steps,6]{
  label=\arabic{stepsi}.\arabic{stepsii}.\arabic{stepsiii}.\arabic{stepsiv}.\arabic{stepsv}.\arabic*.,
  ref=\arabic{stepsi}.\arabic{stepsii}.\arabic{stepsiii}.\arabic{stepsiv}.\arabic{stepsv}.\arabic*
}
\setlist[steps,7]{
  label=\arabic{stepsi}.\arabic{stepsii}.\arabic{stepsiii}.\arabic{stepsiv}. \arabic{stepsv}. \arabic{stepsvi}.\arabic*.,
  ref=\arabic{stepsi}.\arabic{stepsii}.\arabic{stepsiii}.\arabic{stepsiv}.\arabic{stepsv}. \arabic{stepsvi}.\arabic*
}

\newcommand{\step}[1][]{\item\maybeLabel{#1}}

\makeatletter
% \skipitem steps the list counter itself (for block labels).  Give each one
% its own hyperref link target, since step numbers repeat from proof to proof.
\newcounter{skipitem}
\newcommand{\skipitem}{\stepcounter{skipitem}\refstepcounter{\@enumctr}}
\renewcommand{\theHstepsi}{skipitem.\arabic{skipitem}}
\renewcommand{\theHstepsii}{skipitem.\arabic{skipitem}}
\renewcommand{\theHstepsiii}{skipitem.\arabic{skipitem}}
\renewcommand{\theHstepsiv}{skipitem.\arabic{skipitem}}
\renewcommand{\theHstepsv}{skipitem.\arabic{skipitem}}
\renewcommand{\theHstepsvi}{skipitem.\arabic{skipitem}}
\renewcommand{\theHstepsvii}{skipitem.\arabic{skipitem}}
\makeatother

\newcommand{\maybeLabel}[1]{\ifthenelse{\isempty{#1}}{}{\label{#1}}}

\newenvironment{longFormProof}[1][Proof (long form).]
{\begin{proof}[#1]~\par\begin{steps}}{\end{steps}\vspace*{-3ex}\end{proof}}

\newenvironment{shortFormProof}[1][Proof (short form).]
{\begin{proof}[#1]}{\end{proof}}

\newenvironment{longFormSubProof}[1][Proof.]
{\begin{adjustwidth}{1em}{0em}\begin{proof}[#1]~\par\vspace*{-2pt}\skipitem\begin{steps}}
{\end{steps}\vspace*{-1ex}\end{proof}\end{adjustwidth}\smallskip}

\newenvironment{shortFormSubProof}[1][Proof.]
{\begin{proof}[#1]}
{\end{proof}\vspace*{-3pt}}

\newenvironment{block}[1][]
{\skipitem\maybeLabel{#1}\item[] \begin{steps}}{\end{steps}}

\newcommand{\comment}[1]{\hfill{\footnotesize\emph{#1}}}
\newcommand{\lineacross}{\par\vspace*{-0.7\baselineskip}\noindent\hrulefill\par}

%%% problem

% \newcounter{problem}
% \newcommand{\problem}[1][]{}
% {\ifthenelse{\isempty{#1}}{\refstepcounter{problem}}{}%
%   \paragraph{Problem \ifthenelse{\isempty{#1}}
%     {\bf\arabic{problem}.}
%     {#1}} }

\newcommand{\problemNumbered}[1]{
  \setcounter{problem}{#1}%
  \addtocounter{problem}{-1}%
  \problem
}
\newenvironment{problemTheorem}{%
  \setcounter{theorem}{\value{problem}}%
  \addtocounter{theorem}{-1}%
  \begin{theorem}[Problem \arabic{problem}]
  }{\end{theorem}}

\newcounter{problem}
\newcommand{\problem}[1][]
{\ifthenelse{\isempty{#1}}
  {\refstepcounter{problem}\paragraph{Problem~\arabic{problem}.}}
  {\paragraph{Problem~#1.}}\refreshHeader} 

\newcounter{exercise}
\newcommand{\exercise}[1][]
{\ifthenelse{\isempty{#1}}
  {\refstepcounter{exercise}\paragraph{Exercise~\theexercise.}}
  {\paragraph{Exercise~#1.}}}

\newcommand{\solution}{\paragraph{Solution.}} 

\newenvironment{mylist}[1][]
{\begin{list}{#1}{
      \setlength{\itemsep}{-1pt}
      \setlength{\labelwidth}{0pt}
      \setlength{\leftmargin}{1.25em}
    }}
{\end{list}}

\newcommand{\LNTitle}{\textcolor{red}{UNDEFINED LNTitle}}

% In a standalone lecture note's table of contents, leave room for numbers
% like 7.10 (in bold) and 7.10.2.
\makeatletter
\newcommand{\LNtocWidths}{%
  \patchcmd{\l@section}{1.5em}{2.8em}{}{\PackageWarning{macros}{Could not widen section numbers in the contents}}%
  \renewcommand*{\l@subsection}{\@dottedtocline{2}{2.8em}{3.2em}}%
  \renewcommand*{\l@subsubsection}{\@dottedtocline{3}{6.0em}{4.1em}}}
\makeatother

\newcommand{\doTitle}[3]{
  \pagenumbering{roman}
  \renewcommand{\LNTitle}{Lecture Note #1, #3}
  \renewcommand{\currentLN}{#1}
  \codeListingstrue
  % number sections, theorems etc., exercises, figures, tables and equations
  % within the lecture note (e.g. Section 8.2, Lemma 8.6), as in the book
  \renewcommand{\thesection}{#1.\arabic{section}}
  \renewcommand{\thelemma}{#1.\arabic{lemma}}
  \renewcommand{\thedefinition}{#1.\arabic{definition}}
  \renewcommand{\thetheorem}{#1.\arabic{theorem}}
  \renewcommand{\thecorollary}{#1.\arabic{corollary}}
  \renewcommand{\theclaim}{#1.\arabic{claim}}
  \renewcommand{\theobservation}{#1.\arabic{observation}}
  \renewcommand{\theconjecture}{#1.\arabic{conjecture}}
  \renewcommand{\theexercise}{#1.\arabic{exercise}}
  \renewcommand{\thefigure}{#1.\arabic{figure}}
  \renewcommand{\thetable}{#1.\arabic{table}}
  \renewcommand{\theequation}{#1.\arabic{equation}}
  \LNtocWidths
  \title{Lecture Note #1\\ #2}
  % a note's \LNabstract, if it defines one before \doTitle, goes right below the title
  \date{\ifdefined\LNabstract\parbox{0.85\textwidth}{\normalsize\normalfont\LNabstract}\\[3ex]\fi
    draft of \today, by \href{\authorURL}{N.~Young}\\[2ex] {\small\copyrightNotice}}
  \maketitle
  {\small
    \hypertarget{TOC}{}
    \tableofcontents
  }
  \thispagestyle{empty}
  % the text starts on page 1, after the title and contents (numbered i, ii, ...)
  \clearpage
  \pagenumbering{arabic}
}

\def\clap#1{\hbox to 0pt{\hss#1\hss}}

\newcommand{\refreshHeader}{}
\newcommand{\header}{}

\newcommand{\doHeader}[1]{      % for lecture notes
  \renewcommand{\header}{%
    \clap{\raisebox{0pt}{\protect\hyperlink{TOC}{$\uparrow$\scriptsize\,\textsc{contents}\hspace{1.1in}}}}%
    {\footnotesize \LNTitle{} / Section \thesection, #1}} %, #1, draft of \today}}
  \markboth{\header}{\header}
}

\newcommand{\SID}{}
\newcommand{\setAuthor}[1]{\renewcommand{\author}{#1}}
\newcommand{\setSID}[1]{\renewcommand{\SID}{#1}}

\newcommand{\doHwHeader}[1]{                                         % for homeworks
  \renewcommand{\refreshHeader}{
    \renewcommand{\header}{%
      {\footnotesize \classSection #1 / Problem~\arabic{problem} / draft of \today \hfill \author\ifGradescope~(\SID)\fi~~}}
    \markboth{\header}{\header}
  }
  \renewcommand{\header}{%
    {\footnotesize \classSection #1 / Instructions / draft of \today \hfill \author\ifGradescope~(\SID)\fi~~}}
  \markboth{\header}{\header}
}

% Homework settings, for instructors using the homeworks in a course:
% \classSection is the course name shown in each homework's running head
% (e.g. \renewcommand{\classSection}{CS 3000\xspace}), and \Gradescopetrue turns on
% the instructions for submitting to Gradescope (in its fixed-length mode),
% the hints saying on which page of the PDF each answer should be, and the
% student ID in the running head.
\newcommand{\classSection}{}
\newif\ifGradescope
\Gradescopefalse
\newcommand{\GradescopeComment}[1]{\ifGradescope\comment{#1}\fi}

\pagestyle{myheadings}
\addtolength{\headsep}{0.2in}
\addtolength{\topmargin}{-0.3in}
\addtolength{\textheight}{0.4in}

%%%%%%%%%%%%%%%%%%%%%%%%%%%%%%%%%%%%%%%%%%%%%%%%
%%%%%%%%%%%%%%%%%%%%%%%%%%%%%%%%%%%%%%%%%%%%%%%%

\newcommand{\MF}{\href
  {https://mfleck.cs.illinois.edu/building-blocks/}
  {MF}\xspace}

\newcommand{\KWSlides}{\href
  {https://www.cs.princeton.edu/~wayne/kleinberg-tardos/}
  {KW slides}\xspace}

\newcommand{\LLM}{\href
  {https://courses.csail.mit.edu/6.042/spring18/mcs.pdf}
  {LLM}\xspace}

% \DPV, \KT, \CLRS: textbooks listed (without links) in the table of sources in
% the Preliminaries note; each prints as a link to its entry in that table.
\DeclareRobustCommand{\DPV}{\noteLink{prelim}{source.DPV}{DPV}\xspace}
\DeclareRobustCommand{\KT}{\noteLink{prelim}{source.KT}{KT}\xspace}
\DeclareRobustCommand{\CLRS}{\noteLink{prelim}{source.CLRS}{CLRS}\xspace}
\AtBeginDocument{\pdfstringdefDisableCommands{\def\DPV{DPV }\def\KT{KT }\def\CLRS{CLRS }}}

\newcommand{\JE}{\href
  {https://jeffe.cs.illinois.edu/teaching/algorithms/}
  {JE}\xspace}

\newcommand{\JEGreedy}{\href
  {https://jeffe.cs.illinois.edu/teaching/algorithms/book/04-greedy.pdf}
  {JE Chapter 4}\xspace}

\newcommand{\BBCode}{\href
{https://northeastern.blackboard.com/webapps/cmsmain/webui/courses/CS3000.MERGED.202010/code?action=frameset&subaction=view&uniq=-jj2max&course_id=_2597392_1&mask=\%2Fcourses\%2FCS3000.MERGED.202010}
  {code}\xspace}

% The Python code (src/code).  Each lecture note that has code ends with a
% section listing its files: \codeSection{path, path, ...}.  In the text,
% \pycode{greedy_algorithms/Dijkstra.py} prints the path as a link to its listing
% (in the homeworks, which have no listings, to the file on GitHub), and
% \pycode[text]{...} uses the given text instead.  \pythonCode{path, ...} is the
% standard sentence ending a section whose algorithms have code:
% "Python code: path, ... (Section N.k)."
\newcommand{\codeFolder}{https://github.com/nealeyoung/algs101/tree/main/src/code}
\newcommand{\codeRepo}{https://github.com/nealeyoung/algs101/blob/main/src/code}
\newif\ifcodeListings   % set by \doTitle: this document has the listings
\newcommand{\githubCode}[2][]{\href{\codeRepo/#2}{\ifx\relax#1\relax\nolinkurl{#2}\else#1\fi}}
\newcommand{\pycode}[2][]{%
  \ifcodeListings\hyperref[code: #2]{\ifx\relax#1\relax\nolinkurl{#2}\else#1\fi}%
  \else\githubCode[#1]{#2}\fi}
\makeatletter
\newcommand{\pythonCode}[1]{%
  \def\pythonCode@first{}\def\pythonCode@list{}%
  \forcsvlist{\pythonCode@add}{#1}%
  Python code: \pythonCode@list\ (Section~\ref{code: \pythonCode@first}).}
\newcommand{\pythonCode@add}[1]{%
  \ifx\pythonCode@first\@empty\def\pythonCode@first{#1}\appto\pythonCode@list{\pycode{#1}}%
  \else\appto\pythonCode@list{, \pycode{#1}}\fi}
\makeatother

% Code listings that copy and paste exactly (in PDF viewers that honor the
% ToUnicode maps, e.g. MuPDF-based ones; PDFKit viewers such as Preview and
% Skim drop leading spaces, so each listing's heading links to the file on
% GitHub).  Spaces are typeset as white visible-space glyphs that copy as
% spaces, and hyphens and asterisks as the T1 glyphs, so they copy as themselves.
\usepackage[T1,OT1]{fontenc}   % T1 for the code listings only; OT1 stays the default
\usepackage{listings}
\pdfgentounicode=1
\input glyphtounicode
\pdfglyphtounicode{uni2423}{0020}
\makeatletter
\def\lst@visiblespace{\textcolor{white}{\char32}}
% put an invisible space on each blank line, so that blank lines copy too
% (listings' OnEmptyLine hook runs before the line break, so patch the loop)
\def\lst@DoNewLines{%
    \@whilenum\lst@newlines>\@ne \do {\lst@NewLine \lst@visiblespace}%
    \ifnum\lst@newlines>\z@ \lst@NewLine \fi}
\makeatother
\lstset{language=Python,
  basicstyle=\fontencoding{T1}\fontfamily{lmtt}\selectfont\footnotesize,
  upquote=true, showstringspaces=false, columns=fullflexible, keepspaces=true,
  showspaces=true, literate={-}{{\char45}}1 {*}{{\char42}}1,
  frame=single, framesep=1ex, rulecolor=\color{gray}}
% the listings' paths are relative to the lecture note's folder (in the book, to src/)
\ifdefined\BookMode\newcommand{\codeDir}{code/}\else\newcommand{\codeDir}{../code/}\fi
% \codefile{path}: a listing of src/code/path, headed by a link to it on GitHub
\newcommand{\codefile}[1]{%
  \subsection*{\githubCode{#1}}\label{code: #1}%
  \lstinputlisting{\codeDir#1}}
% \codeSection{path, path, ...}: the lecture note's last section, listing its code
\newcommand{\codeSection}[1]{%
  \newpage
  \section{Python code}\label{\LNtag{\currentLN}: sec: code}\doHeader{Python code}
  This section lists the Python code for this lecture note.
  Click on a file's name to open the file on GitHub.
  Copying code from GitHub is more reliable than copying it from this PDF,
  since some PDF viewers drop the leading spaces (the indentation) of copied lines.
  \forcsvlist{\codefile}{#1}}
% the copyright notice, for title pages and the homework first pages
\newcommand{\licenseURL}{https://github.com/nealeyoung/algs101/blob/main/LICENSE}
\newcommand{\authorURL}{https://www.cs.ucr.edu/\string~neal}
\newcommand{\copyrightNotice}{%
  \copyright\ 2018--2026 \href{\authorURL}{Neal E. Young}.
  Text: CC BY-NC-SA 4.0; code: MIT; except where noted.
  See \url{\licenseURL}.}
% the published site: PDFs of the notes, and the homework assignments and templates
\newcommand{\siteURL}{https://nealeyoung.github.io/algs101/}

\newcommand{\blackboard}{\href
  {https://northeastern.blackboard.com/webapps/blackboard/execute/launcher?type=Course&id=_2597392_1&url=}
  {Blackboard}\xspace}

\newcommand{\JUSlides}{\href
  {http://www.khoury.neu.edu/home/jullman/cs3000f18/schedule.html}
  {Prof.~Ullman's slides}\xspace}

% a URL on its own line, in three layouts
\newcommand{\URL}[1]{\par-- \parbox[t]{0.8\textwidth}{\url{#1}} \medskip\par}
\newcommand{\URLsmall}[1]{\par-- \parbox[t]{0.95\textwidth}{\small\url{#1}} \medskip\par}
\newcommand{\URLplain}[1]{\par\parbox[t]{\textwidth}{\url{#1}}\par}

%%%%%%%%%%%%%%%%%%%%%%%%%%%%%%%%%%%%%%%%%%%%%%%%
%%%%%%%%%%%%%%%%%%%%%%%%%%%%%%%%%%%%%%%%%%%%%%%%

\newcommand{\continued}{\vfill\centerline{\emph{(continued)}}\vfill\newpage}

% load hyperref last to avoid conflicts; long URLs may also break after hyphens
\PassOptionsToPackage{hyphens}{url}
\ifdefined\BookMode\else\usepackage{xr-hyper}\fi   % for \lectureNoteRefs (homeworks), below
\usepackage{hyperref}
\hypersetup{
    colorlinks,
    linkcolor={blue!50!black},
    citecolor={blue!50!black},
    urlcolor={blue!80!black}
}

\providecommand{\answerDirectory}{.}
\newcommand{\inputAnswer}[1]{\refreshHeader\input{\answerDirectory/#1}}
\newcommand{\inputSection}[1]{\input{#1}}

% Refer to a lecture note by its tag (the prefix of its labels), e.g.
% \LN{proofs} prints "Lecture Note 2" (in the book, a link to that chapter).  The table below is the only place the numbers appear.
\makeatletter
\newcommand{\defineLN}[3]{\@namedef{LNnumber@#2}{#1}\@namedef{LNtag@#1}{#2}\@namedef{LNfile@#2}{#3}}
\newcommand{\LNnumber}[1]{%
  \ifcsname LNnumber@#1\endcsname\csname LNnumber@#1\endcsname
  \else\PackageError{macros}{Unknown lecture-note tag `#1'}{}\fi}
\newcommand{\LNtag}[1]{\csname LNtag@#1\endcsname}
\newcommand{\LNfile}[1]{\csname LNfile@#1\endcsname}
\makeatother
\defineLN{1}{prelim}{1_preliminaries}
\defineLN{2}{proofs}{2_proofs}
\defineLN{3}{matching}{3_stable_matching}
\defineLN{4}{dc}{4_divide_and_conquer}
\defineLN{5}{greedy}{5_greedy}
\defineLN{6}{graphs}{6_graph_traversal}
\defineLN{7}{dp}{7_dynamic_programming}
\defineLN{8}{flow}{8_flow_reductions_LP}
\defineLN{9}{np}{9_NP}

% The text is the same in the book and in the standalone lecture notes; only
% links differ.  In the book, links to other lecture notes stay in the book.
% In a standalone note, they open the other note's published PDF.
% \noteLink{tag}{target}{text} links text to \hypertarget{target} in lecture note tag.
% \LN{proofs} prints "Lecture Note 2", linked to that lecture note.
\newcommand{\currentLN}{0}   % in a standalone note, its number (set by \doTitle)
\newcommand{\notePDF}[1]{\siteURL notes/\LNfile{#1}.pdf}
\newcommand{\noteLink}[3]{%
  \ifnum\LNnumber{#1}=\currentLN\relax\hyperlink{#2}{#3}%
  \else\href{\notePDF{#1}\#nameddest=#2}{#3}\fi}
% \bookLink{target}{text} links text to \hypertarget{target} in the book's
% appendices (from a standalone note, in the published book).
\DeclareRobustCommand{\bookLink}[2]{%
  \ifdefined\BookMode\hyperlink{#1}{#2}%
  \else\href{\siteURL lecture_notes_on_algorithms.pdf\#nameddest=#1}{#2}\fi}
\DeclareRobustCommand{\LN}[1]{%
  \ifnum\LNnumber{#1}=\currentLN\relax Lecture Note~\LNnumber{#1}%
  \else\href{\notePDF{#1}}{Lecture Note~\LNnumber{#1}}\fi}
\pdfstringdefDisableCommands{\def\LN#1{Lecture Note \LNnumber{#1}}}

% Homeworks refer to the notes' sections, exercises and lemmas by \ref: each
% homework's main file says \lectureNoteRefs after %%%% DON'T EDIT THIS FILE %%%%

\usepackage{xifthen}% \isempty
\usepackage{etoolbox}% \patchcmd
% \usepackage{stringstrings}% \substring

\usepackage{amsmath, amsthm, amssymb}
\usepackage[margin=1in]{geometry}
\usepackage{xspace}
\usepackage{graphicx}
\usepackage[normalem]{ulem}
\usepackage{adjustbox}
\usepackage{xcolor}
\usepackage{changepage}
\usepackage{framed}
\usepackage{multicol}
\usepackage{enumitem}
% \usepackage{listings} % lstlistings environment

\usepackage{makecmds}           %\provideenvironment

\usepackage{xparse,environ}
\usepackage{letltxmacro}       % \LetLtxMacro
\usepackage{wrapfig}
\usepackage{tikz}

\usetikzlibrary{shapes}
\usetikzlibrary{positioning}
\usetikzlibrary{shapes.misc}
\usetikzlibrary{arrows}

\newtheorem{lemma}{Lemma}
\newtheorem{definition}{Definition}
\newtheorem{theorem}{Theorem}
\newtheorem{corollary}{Corollary}
\newtheorem{claim}{Claim}
\newtheorem{observation}{Observation}
\newtheorem{conjecture}{Conjecture}

% A lemma or theorem restated later with its original number (set with
% \setcounter) goes in a restatement environment, which gives it its own link
% anchor (otherwise it would duplicate the original's).
\newenvironment{restatement}
  {\def\theHlemma{restated.\thelemma}\def\theHtheorem{restated.\thetheorem}}{}

\newenvironment{myframed}
{\smallskip\par\noindent\begin{minipage}{\textwidth}\begin{framed}\vspace*{-5pt}}{\vspace*{-5pt}\end{framed}\end{minipage}\smallskip}


\definecolor{darkblue}{RGB}{0,0,200}

 
\let\oldParagraph\paragraph
\providecommand{\gradescopeColor}{\color{black}}
\renewcommand{\paragraph}[1]{\oldParagraph{\gradescopeColor #1}}
\providecommand{\gradescopeInit}{}

\providecommand{\BLUE}[1]{{\color{darkblue} #1}}
\providecommand{\REPLACEME}{\textcolor{red}{\ensuremath{\star \cdots \star}}\xspace}
\provideenvironment{blue}{\color{darkblue}}{}

\newcommand{\dist}[2]{\ensuremath{\operatorname{\sf dist}(#1, #2)}}
\newcommand{\depth}[2]{\ensuremath{\operatorname{\sf depth}_{#1}(#2)}}
\newcommand{\cost}[1]{\ensuremath{\operatorname{\sf cost}(#1)}}
\newcommand{\weight}[1]{\ensuremath{\operatorname{\sf wt}(#1)}}
% \newcommand{\val}[1]{\ensuremath{\operatorname{\sf val}(#1)}}
\newcommand{\eps}{\varepsilon}
\newcommand{\capac}{\textsf{cap}}
\newcommand{\capOf}[1]{\ensuremath{\operatorname{\sf cap}(#1)}}
\newcommand{\rescap}[1]{\ensuremath{\operatorname{\sf cap}_f(#1)}}
\newcommand{\into}[1]{\ensuremath{\operatorname{\sf in}(#1)}}
\newcommand{\out}[1]{\ensuremath{\operatorname{\sf out}(#1)}}
\newenvironment{LP}[1]{
  \begin{array}[t]{@{}l@{}}
    \displaystyle #1 \\
    \begin{aligned}\\[-10pt]}{\end{aligned}\end{array}}

\newcommand{\opt}[1]{\ensuremath{\operatorname{\sf opt}(#1)}}
\newcommand{\Opt}[1]{\ensuremath{\operatorname{\sf Opt}(#1)}}
\newcommand{\half}{\frac{1}{2}}

\newcommand{\E}{\operatorname{E}}
\newcommand{\giv}{\,|\,}
\newcommand{\Z}{\mathbb{Z}}
\newcommand{\Zp}{\mathbb{Z}_{\ge 0}}
\newcommand{\R}{\mathbb{R}}
% \newcommand{\Rp}{\mathbb{R}_{\ge 0}}
\newcommand{\Rp}{\mathbb{R}_{\tiny +}}
\newcommand{\N}{\mathbb{N}}
\newcommand{\Np}{\mathbb{N}_{\ge 0}}
\newcommand{\calS}{{\cal S}}

\newcommand{\OP}[1]{\text{\small\sc #1}\xspace}
\newcommand{\Insert}{\OP{Insert}}
\newcommand{\FindMin}{\OP{Find-Min}}
\newcommand{\DeleteMin}{\OP{Delete-Min}}
\newcommand{\DecreaseKey}{\OP{Decrease}}
\newcommand{\Merge}{\OP{Merge}}
\newcommand{\Cut}{\OP{Cut}}
% \newcommand{\prob}[1]{\text{\sc #1}}
% a problem name; in running text, lines may break inside it
\DeclareRobustCommand{\prob}[1]{\ifmmode\text{#1}\else#1\fi}
\newcommand{\encode}[1]{\langle #1 \rangle}

\newcommand{\best}[1]{\text{\sf best}(#1)}
\newcommand{\worst}[1]{\text{\sf worst}(#1)}
\newcommand{\floor}[1]{\lfloor #1 \rfloor}
\newcommand{\ceil}[1]{\lceil #1 \rceil}

\setlength{\parskip}{2pt}
\setlength{\parindent}{1em}

% \setlength\fboxsep{0pt}

%%% longFormProof, shortFormProof, block environments
%%% \step

\newlist{steps}{enumerate}{7}
\setlist[steps]{wide,
  topsep=1pt,
  parsep=1pt,
  partopsep=1pt,
  itemsep=0pt,
  align=left,
  labelindent=0pt,
  itemindent=0pt, 
}
\setlist[steps,1]{
  label=\arabic*.,
  ref=\arabic*
}
\setlist[steps,2]{
  label=\arabic{stepsi}.\arabic*.,
  ref=\arabic{stepsi}.\arabic*
}
\setlist[steps,3]{
  label=\arabic{stepsi}.\arabic{stepsii}.\arabic*.,
  ref=\arabic{stepsi}.\arabic{stepsii}.\arabic*
}
\setlist[steps,4]{
  label=\arabic{stepsi}.\arabic{stepsii}.\arabic{stepsiii}.\arabic*.,
  ref=\arabic{stepsi}.\arabic{stepsii}.\arabic{stepsiii}.\arabic*
}
\setlist[steps,5]{
  label=\arabic{stepsi}.\arabic{stepsii}.\arabic{stepsiii}.\arabic{stepsiv}.\arabic*.,
  ref=\arabic{stepsi}.\arabic{stepsii}.\arabic{stepsiii}.\arabic{stepsiv}.\arabic*
}
\setlist[steps,6]{
  label=\arabic{stepsi}.\arabic{stepsii}.\arabic{stepsiii}.\arabic{stepsiv}.\arabic{stepsv}.\arabic*.,
  ref=\arabic{stepsi}.\arabic{stepsii}.\arabic{stepsiii}.\arabic{stepsiv}.\arabic{stepsv}.\arabic*
}
\setlist[steps,7]{
  label=\arabic{stepsi}.\arabic{stepsii}.\arabic{stepsiii}.\arabic{stepsiv}. \arabic{stepsv}. \arabic{stepsvi}.\arabic*.,
  ref=\arabic{stepsi}.\arabic{stepsii}.\arabic{stepsiii}.\arabic{stepsiv}.\arabic{stepsv}. \arabic{stepsvi}.\arabic*
}

\newcommand{\step}[1][]{\item\maybeLabel{#1}}

\makeatletter
% \skipitem steps the list counter itself (for block labels).  Give each one
% its own hyperref link target, since step numbers repeat from proof to proof.
\newcounter{skipitem}
\newcommand{\skipitem}{\stepcounter{skipitem}\refstepcounter{\@enumctr}}
\renewcommand{\theHstepsi}{skipitem.\arabic{skipitem}}
\renewcommand{\theHstepsii}{skipitem.\arabic{skipitem}}
\renewcommand{\theHstepsiii}{skipitem.\arabic{skipitem}}
\renewcommand{\theHstepsiv}{skipitem.\arabic{skipitem}}
\renewcommand{\theHstepsv}{skipitem.\arabic{skipitem}}
\renewcommand{\theHstepsvi}{skipitem.\arabic{skipitem}}
\renewcommand{\theHstepsvii}{skipitem.\arabic{skipitem}}
\makeatother

\newcommand{\maybeLabel}[1]{\ifthenelse{\isempty{#1}}{}{\label{#1}}}

\newenvironment{longFormProof}[1][Proof (long form).]
{\begin{proof}[#1]~\par\begin{steps}}{\end{steps}\vspace*{-3ex}\end{proof}}

\newenvironment{shortFormProof}[1][Proof (short form).]
{\begin{proof}[#1]}{\end{proof}}

\newenvironment{longFormSubProof}[1][Proof.]
{\begin{adjustwidth}{1em}{0em}\begin{proof}[#1]~\par\vspace*{-2pt}\skipitem\begin{steps}}
{\end{steps}\vspace*{-1ex}\end{proof}\end{adjustwidth}\smallskip}

\newenvironment{shortFormSubProof}[1][Proof.]
{\begin{proof}[#1]}
{\end{proof}\vspace*{-3pt}}

\newenvironment{block}[1][]
{\skipitem\maybeLabel{#1}\item[] \begin{steps}}{\end{steps}}

\newcommand{\comment}[1]{\hfill{\footnotesize\emph{#1}}}
\newcommand{\lineacross}{\par\vspace*{-0.7\baselineskip}\noindent\hrulefill\par}

%%% problem

% \newcounter{problem}
% \newcommand{\problem}[1][]{}
% {\ifthenelse{\isempty{#1}}{\refstepcounter{problem}}{}%
%   \paragraph{Problem \ifthenelse{\isempty{#1}}
%     {\bf\arabic{problem}.}
%     {#1}} }

\newcommand{\problemNumbered}[1]{
  \setcounter{problem}{#1}%
  \addtocounter{problem}{-1}%
  \problem
}
\newenvironment{problemTheorem}{%
  \setcounter{theorem}{\value{problem}}%
  \addtocounter{theorem}{-1}%
  \begin{theorem}[Problem \arabic{problem}]
  }{\end{theorem}}

\newcounter{problem}
\newcommand{\problem}[1][]
{\ifthenelse{\isempty{#1}}
  {\refstepcounter{problem}\paragraph{Problem~\arabic{problem}.}}
  {\paragraph{Problem~#1.}}\refreshHeader} 

\newcounter{exercise}
\newcommand{\exercise}[1][]
{\ifthenelse{\isempty{#1}}
  {\refstepcounter{exercise}\paragraph{Exercise~\theexercise.}}
  {\paragraph{Exercise~#1.}}}

\newcommand{\solution}{\paragraph{Solution.}} 

\newenvironment{mylist}[1][]
{\begin{list}{#1}{
      \setlength{\itemsep}{-1pt}
      \setlength{\labelwidth}{0pt}
      \setlength{\leftmargin}{1.25em}
    }}
{\end{list}}

\newcommand{\LNTitle}{\textcolor{red}{UNDEFINED LNTitle}}

% In a standalone lecture note's table of contents, leave room for numbers
% like 7.10 (in bold) and 7.10.2.
\makeatletter
\newcommand{\LNtocWidths}{%
  \patchcmd{\l@section}{1.5em}{2.8em}{}{\PackageWarning{macros}{Could not widen section numbers in the contents}}%
  \renewcommand*{\l@subsection}{\@dottedtocline{2}{2.8em}{3.2em}}%
  \renewcommand*{\l@subsubsection}{\@dottedtocline{3}{6.0em}{4.1em}}}
\makeatother

\newcommand{\doTitle}[3]{
  \pagenumbering{roman}
  \renewcommand{\LNTitle}{Lecture Note #1, #3}
  \renewcommand{\currentLN}{#1}
  \codeListingstrue
  % number sections, theorems etc., exercises, figures, tables and equations
  % within the lecture note (e.g. Section 8.2, Lemma 8.6), as in the book
  \renewcommand{\thesection}{#1.\arabic{section}}
  \renewcommand{\thelemma}{#1.\arabic{lemma}}
  \renewcommand{\thedefinition}{#1.\arabic{definition}}
  \renewcommand{\thetheorem}{#1.\arabic{theorem}}
  \renewcommand{\thecorollary}{#1.\arabic{corollary}}
  \renewcommand{\theclaim}{#1.\arabic{claim}}
  \renewcommand{\theobservation}{#1.\arabic{observation}}
  \renewcommand{\theconjecture}{#1.\arabic{conjecture}}
  \renewcommand{\theexercise}{#1.\arabic{exercise}}
  \renewcommand{\thefigure}{#1.\arabic{figure}}
  \renewcommand{\thetable}{#1.\arabic{table}}
  \renewcommand{\theequation}{#1.\arabic{equation}}
  \LNtocWidths
  \title{Lecture Note #1\\ #2}
  % a note's \LNabstract, if it defines one before \doTitle, goes right below the title
  \date{\ifdefined\LNabstract\parbox{0.85\textwidth}{\normalsize\normalfont\LNabstract}\\[3ex]\fi
    draft of \today, by \href{\authorURL}{N.~Young}\\[2ex] {\small\copyrightNotice}}
  \maketitle
  {\small
    \hypertarget{TOC}{}
    \tableofcontents
  }
  \thispagestyle{empty}
  % the text starts on page 1, after the title and contents (numbered i, ii, ...)
  \clearpage
  \pagenumbering{arabic}
}

\def\clap#1{\hbox to 0pt{\hss#1\hss}}

\newcommand{\refreshHeader}{}
\newcommand{\header}{}

\newcommand{\doHeader}[1]{      % for lecture notes
  \renewcommand{\header}{%
    \clap{\raisebox{0pt}{\protect\hyperlink{TOC}{$\uparrow$\scriptsize\,\textsc{contents}\hspace{1.1in}}}}%
    {\footnotesize \LNTitle{} / Section \thesection, #1}} %, #1, draft of \today}}
  \markboth{\header}{\header}
}

\newcommand{\SID}{}
\newcommand{\setAuthor}[1]{\renewcommand{\author}{#1}}
\newcommand{\setSID}[1]{\renewcommand{\SID}{#1}}

\newcommand{\doHwHeader}[1]{                                         % for homeworks
  \renewcommand{\refreshHeader}{
    \renewcommand{\header}{%
      {\footnotesize \classSection #1 / Problem~\arabic{problem} / draft of \today \hfill \author\ifGradescope~(\SID)\fi~~}}
    \markboth{\header}{\header}
  }
  \renewcommand{\header}{%
    {\footnotesize \classSection #1 / Instructions / draft of \today \hfill \author\ifGradescope~(\SID)\fi~~}}
  \markboth{\header}{\header}
}

% Homework settings, for instructors using the homeworks in a course:
% \classSection is the course name shown in each homework's running head
% (e.g. \renewcommand{\classSection}{CS 3000\xspace}), and \Gradescopetrue turns on
% the instructions for submitting to Gradescope (in its fixed-length mode),
% the hints saying on which page of the PDF each answer should be, and the
% student ID in the running head.
\newcommand{\classSection}{}
\newif\ifGradescope
\Gradescopefalse
\newcommand{\GradescopeComment}[1]{\ifGradescope\comment{#1}\fi}

\pagestyle{myheadings}
\addtolength{\headsep}{0.2in}
\addtolength{\topmargin}{-0.3in}
\addtolength{\textheight}{0.4in}

%%%%%%%%%%%%%%%%%%%%%%%%%%%%%%%%%%%%%%%%%%%%%%%%
%%%%%%%%%%%%%%%%%%%%%%%%%%%%%%%%%%%%%%%%%%%%%%%%

\newcommand{\MF}{\href
  {https://mfleck.cs.illinois.edu/building-blocks/}
  {MF}\xspace}

\newcommand{\KWSlides}{\href
  {https://www.cs.princeton.edu/~wayne/kleinberg-tardos/}
  {KW slides}\xspace}

\newcommand{\LLM}{\href
  {https://courses.csail.mit.edu/6.042/spring18/mcs.pdf}
  {LLM}\xspace}

% \DPV, \KT, \CLRS: textbooks listed (without links) in the table of sources in
% the Preliminaries note; each prints as a link to its entry in that table.
\DeclareRobustCommand{\DPV}{\noteLink{prelim}{source.DPV}{DPV}\xspace}
\DeclareRobustCommand{\KT}{\noteLink{prelim}{source.KT}{KT}\xspace}
\DeclareRobustCommand{\CLRS}{\noteLink{prelim}{source.CLRS}{CLRS}\xspace}
\AtBeginDocument{\pdfstringdefDisableCommands{\def\DPV{DPV }\def\KT{KT }\def\CLRS{CLRS }}}

\newcommand{\JE}{\href
  {https://jeffe.cs.illinois.edu/teaching/algorithms/}
  {JE}\xspace}

\newcommand{\JEGreedy}{\href
  {https://jeffe.cs.illinois.edu/teaching/algorithms/book/04-greedy.pdf}
  {JE Chapter 4}\xspace}

\newcommand{\BBCode}{\href
{https://northeastern.blackboard.com/webapps/cmsmain/webui/courses/CS3000.MERGED.202010/code?action=frameset&subaction=view&uniq=-jj2max&course_id=_2597392_1&mask=\%2Fcourses\%2FCS3000.MERGED.202010}
  {code}\xspace}

% The Python code (src/code).  Each lecture note that has code ends with a
% section listing its files: \codeSection{path, path, ...}.  In the text,
% \pycode{greedy_algorithms/Dijkstra.py} prints the path as a link to its listing
% (in the homeworks, which have no listings, to the file on GitHub), and
% \pycode[text]{...} uses the given text instead.  \pythonCode{path, ...} is the
% standard sentence ending a section whose algorithms have code:
% "Python code: path, ... (Section N.k)."
\newcommand{\codeFolder}{https://github.com/nealeyoung/algs101/tree/main/src/code}
\newcommand{\codeRepo}{https://github.com/nealeyoung/algs101/blob/main/src/code}
\newif\ifcodeListings   % set by \doTitle: this document has the listings
\newcommand{\githubCode}[2][]{\href{\codeRepo/#2}{\ifx\relax#1\relax\nolinkurl{#2}\else#1\fi}}
\newcommand{\pycode}[2][]{%
  \ifcodeListings\hyperref[code: #2]{\ifx\relax#1\relax\nolinkurl{#2}\else#1\fi}%
  \else\githubCode[#1]{#2}\fi}
\makeatletter
\newcommand{\pythonCode}[1]{%
  \def\pythonCode@first{}\def\pythonCode@list{}%
  \forcsvlist{\pythonCode@add}{#1}%
  Python code: \pythonCode@list\ (Section~\ref{code: \pythonCode@first}).}
\newcommand{\pythonCode@add}[1]{%
  \ifx\pythonCode@first\@empty\def\pythonCode@first{#1}\appto\pythonCode@list{\pycode{#1}}%
  \else\appto\pythonCode@list{, \pycode{#1}}\fi}
\makeatother

% Code listings that copy and paste exactly (in PDF viewers that honor the
% ToUnicode maps, e.g. MuPDF-based ones; PDFKit viewers such as Preview and
% Skim drop leading spaces, so each listing's heading links to the file on
% GitHub).  Spaces are typeset as white visible-space glyphs that copy as
% spaces, and hyphens and asterisks as the T1 glyphs, so they copy as themselves.
\usepackage[T1,OT1]{fontenc}   % T1 for the code listings only; OT1 stays the default
\usepackage{listings}
\pdfgentounicode=1
\input glyphtounicode
\pdfglyphtounicode{uni2423}{0020}
\makeatletter
\def\lst@visiblespace{\textcolor{white}{\char32}}
% put an invisible space on each blank line, so that blank lines copy too
% (listings' OnEmptyLine hook runs before the line break, so patch the loop)
\def\lst@DoNewLines{%
    \@whilenum\lst@newlines>\@ne \do {\lst@NewLine \lst@visiblespace}%
    \ifnum\lst@newlines>\z@ \lst@NewLine \fi}
\makeatother
\lstset{language=Python,
  basicstyle=\fontencoding{T1}\fontfamily{lmtt}\selectfont\footnotesize,
  upquote=true, showstringspaces=false, columns=fullflexible, keepspaces=true,
  showspaces=true, literate={-}{{\char45}}1 {*}{{\char42}}1,
  frame=single, framesep=1ex, rulecolor=\color{gray}}
% the listings' paths are relative to the lecture note's folder (in the book, to src/)
\ifdefined\BookMode\newcommand{\codeDir}{code/}\else\newcommand{\codeDir}{../code/}\fi
% \codefile{path}: a listing of src/code/path, headed by a link to it on GitHub
\newcommand{\codefile}[1]{%
  \subsection*{\githubCode{#1}}\label{code: #1}%
  \lstinputlisting{\codeDir#1}}
% \codeSection{path, path, ...}: the lecture note's last section, listing its code
\newcommand{\codeSection}[1]{%
  \newpage
  \section{Python code}\label{\LNtag{\currentLN}: sec: code}\doHeader{Python code}
  This section lists the Python code for this lecture note.
  Click on a file's name to open the file on GitHub.
  Copying code from GitHub is more reliable than copying it from this PDF,
  since some PDF viewers drop the leading spaces (the indentation) of copied lines.
  \forcsvlist{\codefile}{#1}}
% the copyright notice, for title pages and the homework first pages
\newcommand{\licenseURL}{https://github.com/nealeyoung/algs101/blob/main/LICENSE}
\newcommand{\authorURL}{https://www.cs.ucr.edu/\string~neal}
\newcommand{\copyrightNotice}{%
  \copyright\ 2018--2026 \href{\authorURL}{Neal E. Young}.
  Text: CC BY-NC-SA 4.0; code: MIT; except where noted.
  See \url{\licenseURL}.}
% the published site: PDFs of the notes, and the homework assignments and templates
\newcommand{\siteURL}{https://nealeyoung.github.io/algs101/}

\newcommand{\blackboard}{\href
  {https://northeastern.blackboard.com/webapps/blackboard/execute/launcher?type=Course&id=_2597392_1&url=}
  {Blackboard}\xspace}

\newcommand{\JUSlides}{\href
  {http://www.khoury.neu.edu/home/jullman/cs3000f18/schedule.html}
  {Prof.~Ullman's slides}\xspace}

% a URL on its own line, in three layouts
\newcommand{\URL}[1]{\par-- \parbox[t]{0.8\textwidth}{\url{#1}} \medskip\par}
\newcommand{\URLsmall}[1]{\par-- \parbox[t]{0.95\textwidth}{\small\url{#1}} \medskip\par}
\newcommand{\URLplain}[1]{\par\parbox[t]{\textwidth}{\url{#1}}\par}

%%%%%%%%%%%%%%%%%%%%%%%%%%%%%%%%%%%%%%%%%%%%%%%%
%%%%%%%%%%%%%%%%%%%%%%%%%%%%%%%%%%%%%%%%%%%%%%%%

\newcommand{\continued}{\vfill\centerline{\emph{(continued)}}\vfill\newpage}

% load hyperref last to avoid conflicts; long URLs may also break after hyphens
\PassOptionsToPackage{hyphens}{url}
\ifdefined\BookMode\else\usepackage{xr-hyper}\fi   % for \lectureNoteRefs (homeworks), below
\usepackage{hyperref}
\hypersetup{
    colorlinks,
    linkcolor={blue!50!black},
    citecolor={blue!50!black},
    urlcolor={blue!80!black}
}

\providecommand{\answerDirectory}{.}
\newcommand{\inputAnswer}[1]{\refreshHeader\input{\answerDirectory/#1}}
\newcommand{\inputSection}[1]{\input{#1}}

% Refer to a lecture note by its tag (the prefix of its labels), e.g.
% \LN{proofs} prints "Lecture Note 2" (in the book, a link to that chapter).  The table below is the only place the numbers appear.
\makeatletter
\newcommand{\defineLN}[3]{\@namedef{LNnumber@#2}{#1}\@namedef{LNtag@#1}{#2}\@namedef{LNfile@#2}{#3}}
\newcommand{\LNnumber}[1]{%
  \ifcsname LNnumber@#1\endcsname\csname LNnumber@#1\endcsname
  \else\PackageError{macros}{Unknown lecture-note tag `#1'}{}\fi}
\newcommand{\LNtag}[1]{\csname LNtag@#1\endcsname}
\newcommand{\LNfile}[1]{\csname LNfile@#1\endcsname}
\makeatother
\defineLN{1}{prelim}{1_preliminaries}
\defineLN{2}{proofs}{2_proofs}
\defineLN{3}{matching}{3_stable_matching}
\defineLN{4}{dc}{4_divide_and_conquer}
\defineLN{5}{greedy}{5_greedy}
\defineLN{6}{graphs}{6_graph_traversal}
\defineLN{7}{dp}{7_dynamic_programming}
\defineLN{8}{flow}{8_flow_reductions_LP}
\defineLN{9}{np}{9_NP}

% The text is the same in the book and in the standalone lecture notes; only
% links differ.  In the book, links to other lecture notes stay in the book.
% In a standalone note, they open the other note's published PDF.
% \noteLink{tag}{target}{text} links text to \hypertarget{target} in lecture note tag.
% \LN{proofs} prints "Lecture Note 2", linked to that lecture note.
\newcommand{\currentLN}{0}   % in a standalone note, its number (set by \doTitle)
\newcommand{\notePDF}[1]{\siteURL notes/\LNfile{#1}.pdf}
\newcommand{\noteLink}[3]{%
  \ifnum\LNnumber{#1}=\currentLN\relax\hyperlink{#2}{#3}%
  \else\href{\notePDF{#1}\#nameddest=#2}{#3}\fi}
% \bookLink{target}{text} links text to \hypertarget{target} in the book's
% appendices (from a standalone note, in the published book).
\DeclareRobustCommand{\bookLink}[2]{%
  \ifdefined\BookMode\hyperlink{#1}{#2}%
  \else\href{\siteURL lecture_notes_on_algorithms.pdf\#nameddest=#1}{#2}\fi}
\DeclareRobustCommand{\LN}[1]{%
  \ifnum\LNnumber{#1}=\currentLN\relax Lecture Note~\LNnumber{#1}%
  \else\href{\notePDF{#1}}{Lecture Note~\LNnumber{#1}}\fi}
\pdfstringdefDisableCommands{\def\LN#1{Lecture Note \LNnumber{#1}}}

% Homeworks refer to the notes' sections, exercises and lemmas by \ref: each
% homework's main file says \lectureNoteRefs after %%%% DON'T EDIT THIS FILE %%%%

\usepackage{xifthen}% \isempty
\usepackage{etoolbox}% \patchcmd
% \usepackage{stringstrings}% \substring

\usepackage{amsmath, amsthm, amssymb}
\usepackage[margin=1in]{geometry}
\usepackage{xspace}
\usepackage{graphicx}
\usepackage[normalem]{ulem}
\usepackage{adjustbox}
\usepackage{xcolor}
\usepackage{changepage}
\usepackage{framed}
\usepackage{multicol}
\usepackage{enumitem}
% \usepackage{listings} % lstlistings environment

\usepackage{makecmds}           %\provideenvironment

\usepackage{xparse,environ}
\usepackage{letltxmacro}       % \LetLtxMacro
\usepackage{wrapfig}
\usepackage{tikz}

\usetikzlibrary{shapes}
\usetikzlibrary{positioning}
\usetikzlibrary{shapes.misc}
\usetikzlibrary{arrows}

\newtheorem{lemma}{Lemma}
\newtheorem{definition}{Definition}
\newtheorem{theorem}{Theorem}
\newtheorem{corollary}{Corollary}
\newtheorem{claim}{Claim}
\newtheorem{observation}{Observation}
\newtheorem{conjecture}{Conjecture}

% A lemma or theorem restated later with its original number (set with
% \setcounter) goes in a restatement environment, which gives it its own link
% anchor (otherwise it would duplicate the original's).
\newenvironment{restatement}
  {\def\theHlemma{restated.\thelemma}\def\theHtheorem{restated.\thetheorem}}{}

\newenvironment{myframed}
{\smallskip\par\noindent\begin{minipage}{\textwidth}\begin{framed}\vspace*{-5pt}}{\vspace*{-5pt}\end{framed}\end{minipage}\smallskip}


\definecolor{darkblue}{RGB}{0,0,200}

 
\let\oldParagraph\paragraph
\providecommand{\gradescopeColor}{\color{black}}
\renewcommand{\paragraph}[1]{\oldParagraph{\gradescopeColor #1}}
\providecommand{\gradescopeInit}{}

\providecommand{\BLUE}[1]{{\color{darkblue} #1}}
\providecommand{\REPLACEME}{\textcolor{red}{\ensuremath{\star \cdots \star}}\xspace}
\provideenvironment{blue}{\color{darkblue}}{}

\newcommand{\dist}[2]{\ensuremath{\operatorname{\sf dist}(#1, #2)}}
\newcommand{\depth}[2]{\ensuremath{\operatorname{\sf depth}_{#1}(#2)}}
\newcommand{\cost}[1]{\ensuremath{\operatorname{\sf cost}(#1)}}
\newcommand{\weight}[1]{\ensuremath{\operatorname{\sf wt}(#1)}}
% \newcommand{\val}[1]{\ensuremath{\operatorname{\sf val}(#1)}}
\newcommand{\eps}{\varepsilon}
\newcommand{\capac}{\textsf{cap}}
\newcommand{\capOf}[1]{\ensuremath{\operatorname{\sf cap}(#1)}}
\newcommand{\rescap}[1]{\ensuremath{\operatorname{\sf cap}_f(#1)}}
\newcommand{\into}[1]{\ensuremath{\operatorname{\sf in}(#1)}}
\newcommand{\out}[1]{\ensuremath{\operatorname{\sf out}(#1)}}
\newenvironment{LP}[1]{
  \begin{array}[t]{@{}l@{}}
    \displaystyle #1 \\
    \begin{aligned}\\[-10pt]}{\end{aligned}\end{array}}

\newcommand{\opt}[1]{\ensuremath{\operatorname{\sf opt}(#1)}}
\newcommand{\Opt}[1]{\ensuremath{\operatorname{\sf Opt}(#1)}}
\newcommand{\half}{\frac{1}{2}}

\newcommand{\E}{\operatorname{E}}
\newcommand{\giv}{\,|\,}
\newcommand{\Z}{\mathbb{Z}}
\newcommand{\Zp}{\mathbb{Z}_{\ge 0}}
\newcommand{\R}{\mathbb{R}}
% \newcommand{\Rp}{\mathbb{R}_{\ge 0}}
\newcommand{\Rp}{\mathbb{R}_{\tiny +}}
\newcommand{\N}{\mathbb{N}}
\newcommand{\Np}{\mathbb{N}_{\ge 0}}
\newcommand{\calS}{{\cal S}}

\newcommand{\OP}[1]{\text{\small\sc #1}\xspace}
\newcommand{\Insert}{\OP{Insert}}
\newcommand{\FindMin}{\OP{Find-Min}}
\newcommand{\DeleteMin}{\OP{Delete-Min}}
\newcommand{\DecreaseKey}{\OP{Decrease}}
\newcommand{\Merge}{\OP{Merge}}
\newcommand{\Cut}{\OP{Cut}}
% \newcommand{\prob}[1]{\text{\sc #1}}
% a problem name; in running text, lines may break inside it
\DeclareRobustCommand{\prob}[1]{\ifmmode\text{#1}\else#1\fi}
\newcommand{\encode}[1]{\langle #1 \rangle}

\newcommand{\best}[1]{\text{\sf best}(#1)}
\newcommand{\worst}[1]{\text{\sf worst}(#1)}
\newcommand{\floor}[1]{\lfloor #1 \rfloor}
\newcommand{\ceil}[1]{\lceil #1 \rceil}

\setlength{\parskip}{2pt}
\setlength{\parindent}{1em}

% \setlength\fboxsep{0pt}

%%% longFormProof, shortFormProof, block environments
%%% \step

\newlist{steps}{enumerate}{7}
\setlist[steps]{wide,
  topsep=1pt,
  parsep=1pt,
  partopsep=1pt,
  itemsep=0pt,
  align=left,
  labelindent=0pt,
  itemindent=0pt, 
}
\setlist[steps,1]{
  label=\arabic*.,
  ref=\arabic*
}
\setlist[steps,2]{
  label=\arabic{stepsi}.\arabic*.,
  ref=\arabic{stepsi}.\arabic*
}
\setlist[steps,3]{
  label=\arabic{stepsi}.\arabic{stepsii}.\arabic*.,
  ref=\arabic{stepsi}.\arabic{stepsii}.\arabic*
}
\setlist[steps,4]{
  label=\arabic{stepsi}.\arabic{stepsii}.\arabic{stepsiii}.\arabic*.,
  ref=\arabic{stepsi}.\arabic{stepsii}.\arabic{stepsiii}.\arabic*
}
\setlist[steps,5]{
  label=\arabic{stepsi}.\arabic{stepsii}.\arabic{stepsiii}.\arabic{stepsiv}.\arabic*.,
  ref=\arabic{stepsi}.\arabic{stepsii}.\arabic{stepsiii}.\arabic{stepsiv}.\arabic*
}
\setlist[steps,6]{
  label=\arabic{stepsi}.\arabic{stepsii}.\arabic{stepsiii}.\arabic{stepsiv}.\arabic{stepsv}.\arabic*.,
  ref=\arabic{stepsi}.\arabic{stepsii}.\arabic{stepsiii}.\arabic{stepsiv}.\arabic{stepsv}.\arabic*
}
\setlist[steps,7]{
  label=\arabic{stepsi}.\arabic{stepsii}.\arabic{stepsiii}.\arabic{stepsiv}. \arabic{stepsv}. \arabic{stepsvi}.\arabic*.,
  ref=\arabic{stepsi}.\arabic{stepsii}.\arabic{stepsiii}.\arabic{stepsiv}.\arabic{stepsv}. \arabic{stepsvi}.\arabic*
}

\newcommand{\step}[1][]{\item\maybeLabel{#1}}

\makeatletter
% \skipitem steps the list counter itself (for block labels).  Give each one
% its own hyperref link target, since step numbers repeat from proof to proof.
\newcounter{skipitem}
\newcommand{\skipitem}{\stepcounter{skipitem}\refstepcounter{\@enumctr}}
\renewcommand{\theHstepsi}{skipitem.\arabic{skipitem}}
\renewcommand{\theHstepsii}{skipitem.\arabic{skipitem}}
\renewcommand{\theHstepsiii}{skipitem.\arabic{skipitem}}
\renewcommand{\theHstepsiv}{skipitem.\arabic{skipitem}}
\renewcommand{\theHstepsv}{skipitem.\arabic{skipitem}}
\renewcommand{\theHstepsvi}{skipitem.\arabic{skipitem}}
\renewcommand{\theHstepsvii}{skipitem.\arabic{skipitem}}
\makeatother

\newcommand{\maybeLabel}[1]{\ifthenelse{\isempty{#1}}{}{\label{#1}}}

\newenvironment{longFormProof}[1][Proof (long form).]
{\begin{proof}[#1]~\par\begin{steps}}{\end{steps}\vspace*{-3ex}\end{proof}}

\newenvironment{shortFormProof}[1][Proof (short form).]
{\begin{proof}[#1]}{\end{proof}}

\newenvironment{longFormSubProof}[1][Proof.]
{\begin{adjustwidth}{1em}{0em}\begin{proof}[#1]~\par\vspace*{-2pt}\skipitem\begin{steps}}
{\end{steps}\vspace*{-1ex}\end{proof}\end{adjustwidth}\smallskip}

\newenvironment{shortFormSubProof}[1][Proof.]
{\begin{proof}[#1]}
{\end{proof}\vspace*{-3pt}}

\newenvironment{block}[1][]
{\skipitem\maybeLabel{#1}\item[] \begin{steps}}{\end{steps}}

\newcommand{\comment}[1]{\hfill{\footnotesize\emph{#1}}}
\newcommand{\lineacross}{\par\vspace*{-0.7\baselineskip}\noindent\hrulefill\par}

%%% problem

% \newcounter{problem}
% \newcommand{\problem}[1][]{}
% {\ifthenelse{\isempty{#1}}{\refstepcounter{problem}}{}%
%   \paragraph{Problem \ifthenelse{\isempty{#1}}
%     {\bf\arabic{problem}.}
%     {#1}} }

\newcommand{\problemNumbered}[1]{
  \setcounter{problem}{#1}%
  \addtocounter{problem}{-1}%
  \problem
}
\newenvironment{problemTheorem}{%
  \setcounter{theorem}{\value{problem}}%
  \addtocounter{theorem}{-1}%
  \begin{theorem}[Problem \arabic{problem}]
  }{\end{theorem}}

\newcounter{problem}
\newcommand{\problem}[1][]
{\ifthenelse{\isempty{#1}}
  {\refstepcounter{problem}\paragraph{Problem~\arabic{problem}.}}
  {\paragraph{Problem~#1.}}\refreshHeader} 

\newcounter{exercise}
\newcommand{\exercise}[1][]
{\ifthenelse{\isempty{#1}}
  {\refstepcounter{exercise}\paragraph{Exercise~\theexercise.}}
  {\paragraph{Exercise~#1.}}}

\newcommand{\solution}{\paragraph{Solution.}} 

\newenvironment{mylist}[1][]
{\begin{list}{#1}{
      \setlength{\itemsep}{-1pt}
      \setlength{\labelwidth}{0pt}
      \setlength{\leftmargin}{1.25em}
    }}
{\end{list}}

\newcommand{\LNTitle}{\textcolor{red}{UNDEFINED LNTitle}}

% In a standalone lecture note's table of contents, leave room for numbers
% like 7.10 (in bold) and 7.10.2.
\makeatletter
\newcommand{\LNtocWidths}{%
  \patchcmd{\l@section}{1.5em}{2.8em}{}{\PackageWarning{macros}{Could not widen section numbers in the contents}}%
  \renewcommand*{\l@subsection}{\@dottedtocline{2}{2.8em}{3.2em}}%
  \renewcommand*{\l@subsubsection}{\@dottedtocline{3}{6.0em}{4.1em}}}
\makeatother

\newcommand{\doTitle}[3]{
  \pagenumbering{roman}
  \renewcommand{\LNTitle}{Lecture Note #1, #3}
  \renewcommand{\currentLN}{#1}
  \codeListingstrue
  % number sections, theorems etc., exercises, figures, tables and equations
  % within the lecture note (e.g. Section 8.2, Lemma 8.6), as in the book
  \renewcommand{\thesection}{#1.\arabic{section}}
  \renewcommand{\thelemma}{#1.\arabic{lemma}}
  \renewcommand{\thedefinition}{#1.\arabic{definition}}
  \renewcommand{\thetheorem}{#1.\arabic{theorem}}
  \renewcommand{\thecorollary}{#1.\arabic{corollary}}
  \renewcommand{\theclaim}{#1.\arabic{claim}}
  \renewcommand{\theobservation}{#1.\arabic{observation}}
  \renewcommand{\theconjecture}{#1.\arabic{conjecture}}
  \renewcommand{\theexercise}{#1.\arabic{exercise}}
  \renewcommand{\thefigure}{#1.\arabic{figure}}
  \renewcommand{\thetable}{#1.\arabic{table}}
  \renewcommand{\theequation}{#1.\arabic{equation}}
  \LNtocWidths
  \title{Lecture Note #1\\ #2}
  % a note's \LNabstract, if it defines one before \doTitle, goes right below the title
  \date{\ifdefined\LNabstract\parbox{0.85\textwidth}{\normalsize\normalfont\LNabstract}\\[3ex]\fi
    draft of \today, by \href{\authorURL}{N.~Young}\\[2ex] {\small\copyrightNotice}}
  \maketitle
  {\small
    \hypertarget{TOC}{}
    \tableofcontents
  }
  \thispagestyle{empty}
  % the text starts on page 1, after the title and contents (numbered i, ii, ...)
  \clearpage
  \pagenumbering{arabic}
}

\def\clap#1{\hbox to 0pt{\hss#1\hss}}

\newcommand{\refreshHeader}{}
\newcommand{\header}{}

\newcommand{\doHeader}[1]{      % for lecture notes
  \renewcommand{\header}{%
    \clap{\raisebox{0pt}{\protect\hyperlink{TOC}{$\uparrow$\scriptsize\,\textsc{contents}\hspace{1.1in}}}}%
    {\footnotesize \LNTitle{} / Section \thesection, #1}} %, #1, draft of \today}}
  \markboth{\header}{\header}
}

\newcommand{\SID}{}
\newcommand{\setAuthor}[1]{\renewcommand{\author}{#1}}
\newcommand{\setSID}[1]{\renewcommand{\SID}{#1}}

\newcommand{\doHwHeader}[1]{                                         % for homeworks
  \renewcommand{\refreshHeader}{
    \renewcommand{\header}{%
      {\footnotesize \classSection #1 / Problem~\arabic{problem} / draft of \today \hfill \author\ifGradescope~(\SID)\fi~~}}
    \markboth{\header}{\header}
  }
  \renewcommand{\header}{%
    {\footnotesize \classSection #1 / Instructions / draft of \today \hfill \author\ifGradescope~(\SID)\fi~~}}
  \markboth{\header}{\header}
}

% Homework settings, for instructors using the homeworks in a course:
% \classSection is the course name shown in each homework's running head
% (e.g. \renewcommand{\classSection}{CS 3000\xspace}), and \Gradescopetrue turns on
% the instructions for submitting to Gradescope (in its fixed-length mode),
% the hints saying on which page of the PDF each answer should be, and the
% student ID in the running head.
\newcommand{\classSection}{}
\newif\ifGradescope
\Gradescopefalse
\newcommand{\GradescopeComment}[1]{\ifGradescope\comment{#1}\fi}

\pagestyle{myheadings}
\addtolength{\headsep}{0.2in}
\addtolength{\topmargin}{-0.3in}
\addtolength{\textheight}{0.4in}

%%%%%%%%%%%%%%%%%%%%%%%%%%%%%%%%%%%%%%%%%%%%%%%%
%%%%%%%%%%%%%%%%%%%%%%%%%%%%%%%%%%%%%%%%%%%%%%%%

\newcommand{\MF}{\href
  {https://mfleck.cs.illinois.edu/building-blocks/}
  {MF}\xspace}

\newcommand{\KWSlides}{\href
  {https://www.cs.princeton.edu/~wayne/kleinberg-tardos/}
  {KW slides}\xspace}

\newcommand{\LLM}{\href
  {https://courses.csail.mit.edu/6.042/spring18/mcs.pdf}
  {LLM}\xspace}

% \DPV, \KT, \CLRS: textbooks listed (without links) in the table of sources in
% the Preliminaries note; each prints as a link to its entry in that table.
\DeclareRobustCommand{\DPV}{\noteLink{prelim}{source.DPV}{DPV}\xspace}
\DeclareRobustCommand{\KT}{\noteLink{prelim}{source.KT}{KT}\xspace}
\DeclareRobustCommand{\CLRS}{\noteLink{prelim}{source.CLRS}{CLRS}\xspace}
\AtBeginDocument{\pdfstringdefDisableCommands{\def\DPV{DPV }\def\KT{KT }\def\CLRS{CLRS }}}

\newcommand{\JE}{\href
  {https://jeffe.cs.illinois.edu/teaching/algorithms/}
  {JE}\xspace}

\newcommand{\JEGreedy}{\href
  {https://jeffe.cs.illinois.edu/teaching/algorithms/book/04-greedy.pdf}
  {JE Chapter 4}\xspace}

\newcommand{\BBCode}{\href
{https://northeastern.blackboard.com/webapps/cmsmain/webui/courses/CS3000.MERGED.202010/code?action=frameset&subaction=view&uniq=-jj2max&course_id=_2597392_1&mask=\%2Fcourses\%2FCS3000.MERGED.202010}
  {code}\xspace}

% The Python code (src/code).  Each lecture note that has code ends with a
% section listing its files: \codeSection{path, path, ...}.  In the text,
% \pycode{greedy_algorithms/Dijkstra.py} prints the path as a link to its listing
% (in the homeworks, which have no listings, to the file on GitHub), and
% \pycode[text]{...} uses the given text instead.  \pythonCode{path, ...} is the
% standard sentence ending a section whose algorithms have code:
% "Python code: path, ... (Section N.k)."
\newcommand{\codeFolder}{https://github.com/nealeyoung/algs101/tree/main/src/code}
\newcommand{\codeRepo}{https://github.com/nealeyoung/algs101/blob/main/src/code}
\newif\ifcodeListings   % set by \doTitle: this document has the listings
\newcommand{\githubCode}[2][]{\href{\codeRepo/#2}{\ifx\relax#1\relax\nolinkurl{#2}\else#1\fi}}
\newcommand{\pycode}[2][]{%
  \ifcodeListings\hyperref[code: #2]{\ifx\relax#1\relax\nolinkurl{#2}\else#1\fi}%
  \else\githubCode[#1]{#2}\fi}
\makeatletter
\newcommand{\pythonCode}[1]{%
  \def\pythonCode@first{}\def\pythonCode@list{}%
  \forcsvlist{\pythonCode@add}{#1}%
  Python code: \pythonCode@list\ (Section~\ref{code: \pythonCode@first}).}
\newcommand{\pythonCode@add}[1]{%
  \ifx\pythonCode@first\@empty\def\pythonCode@first{#1}\appto\pythonCode@list{\pycode{#1}}%
  \else\appto\pythonCode@list{, \pycode{#1}}\fi}
\makeatother

% Code listings that copy and paste exactly (in PDF viewers that honor the
% ToUnicode maps, e.g. MuPDF-based ones; PDFKit viewers such as Preview and
% Skim drop leading spaces, so each listing's heading links to the file on
% GitHub).  Spaces are typeset as white visible-space glyphs that copy as
% spaces, and hyphens and asterisks as the T1 glyphs, so they copy as themselves.
\usepackage[T1,OT1]{fontenc}   % T1 for the code listings only; OT1 stays the default
\usepackage{listings}
\pdfgentounicode=1
\input glyphtounicode
\pdfglyphtounicode{uni2423}{0020}
\makeatletter
\def\lst@visiblespace{\textcolor{white}{\char32}}
% put an invisible space on each blank line, so that blank lines copy too
% (listings' OnEmptyLine hook runs before the line break, so patch the loop)
\def\lst@DoNewLines{%
    \@whilenum\lst@newlines>\@ne \do {\lst@NewLine \lst@visiblespace}%
    \ifnum\lst@newlines>\z@ \lst@NewLine \fi}
\makeatother
\lstset{language=Python,
  basicstyle=\fontencoding{T1}\fontfamily{lmtt}\selectfont\footnotesize,
  upquote=true, showstringspaces=false, columns=fullflexible, keepspaces=true,
  showspaces=true, literate={-}{{\char45}}1 {*}{{\char42}}1,
  frame=single, framesep=1ex, rulecolor=\color{gray}}
% the listings' paths are relative to the lecture note's folder (in the book, to src/)
\ifdefined\BookMode\newcommand{\codeDir}{code/}\else\newcommand{\codeDir}{../code/}\fi
% \codefile{path}: a listing of src/code/path, headed by a link to it on GitHub
\newcommand{\codefile}[1]{%
  \subsection*{\githubCode{#1}}\label{code: #1}%
  \lstinputlisting{\codeDir#1}}
% \codeSection{path, path, ...}: the lecture note's last section, listing its code
\newcommand{\codeSection}[1]{%
  \newpage
  \section{Python code}\label{\LNtag{\currentLN}: sec: code}\doHeader{Python code}
  This section lists the Python code for this lecture note.
  Click on a file's name to open the file on GitHub.
  Copying code from GitHub is more reliable than copying it from this PDF,
  since some PDF viewers drop the leading spaces (the indentation) of copied lines.
  \forcsvlist{\codefile}{#1}}
% the copyright notice, for title pages and the homework first pages
\newcommand{\licenseURL}{https://github.com/nealeyoung/algs101/blob/main/LICENSE}
\newcommand{\authorURL}{https://www.cs.ucr.edu/\string~neal}
\newcommand{\copyrightNotice}{%
  \copyright\ 2018--2026 \href{\authorURL}{Neal E. Young}.
  Text: CC BY-NC-SA 4.0; code: MIT; except where noted.
  See \url{\licenseURL}.}
% the published site: PDFs of the notes, and the homework assignments and templates
\newcommand{\siteURL}{https://nealeyoung.github.io/algs101/}

\newcommand{\blackboard}{\href
  {https://northeastern.blackboard.com/webapps/blackboard/execute/launcher?type=Course&id=_2597392_1&url=}
  {Blackboard}\xspace}

\newcommand{\JUSlides}{\href
  {http://www.khoury.neu.edu/home/jullman/cs3000f18/schedule.html}
  {Prof.~Ullman's slides}\xspace}

% a URL on its own line, in three layouts
\newcommand{\URL}[1]{\par-- \parbox[t]{0.8\textwidth}{\url{#1}} \medskip\par}
\newcommand{\URLsmall}[1]{\par-- \parbox[t]{0.95\textwidth}{\small\url{#1}} \medskip\par}
\newcommand{\URLplain}[1]{\par\parbox[t]{\textwidth}{\url{#1}}\par}

%%%%%%%%%%%%%%%%%%%%%%%%%%%%%%%%%%%%%%%%%%%%%%%%
%%%%%%%%%%%%%%%%%%%%%%%%%%%%%%%%%%%%%%%%%%%%%%%%

\newcommand{\continued}{\vfill\centerline{\emph{(continued)}}\vfill\newpage}

% load hyperref last to avoid conflicts; long URLs may also break after hyphens
\PassOptionsToPackage{hyphens}{url}
\ifdefined\BookMode\else\usepackage{xr-hyper}\fi   % for \lectureNoteRefs (homeworks), below
\usepackage{hyperref}
\hypersetup{
    colorlinks,
    linkcolor={blue!50!black},
    citecolor={blue!50!black},
    urlcolor={blue!80!black}
}

\providecommand{\answerDirectory}{.}
\newcommand{\inputAnswer}[1]{\refreshHeader\input{\answerDirectory/#1}}
\newcommand{\inputSection}[1]{\input{#1}}

% Refer to a lecture note by its tag (the prefix of its labels), e.g.
% \LN{proofs} prints "Lecture Note 2" (in the book, a link to that chapter).  The table below is the only place the numbers appear.
\makeatletter
\newcommand{\defineLN}[3]{\@namedef{LNnumber@#2}{#1}\@namedef{LNtag@#1}{#2}\@namedef{LNfile@#2}{#3}}
\newcommand{\LNnumber}[1]{%
  \ifcsname LNnumber@#1\endcsname\csname LNnumber@#1\endcsname
  \else\PackageError{macros}{Unknown lecture-note tag `#1'}{}\fi}
\newcommand{\LNtag}[1]{\csname LNtag@#1\endcsname}
\newcommand{\LNfile}[1]{\csname LNfile@#1\endcsname}
\makeatother
\defineLN{1}{prelim}{1_preliminaries}
\defineLN{2}{proofs}{2_proofs}
\defineLN{3}{matching}{3_stable_matching}
\defineLN{4}{dc}{4_divide_and_conquer}
\defineLN{5}{greedy}{5_greedy}
\defineLN{6}{graphs}{6_graph_traversal}
\defineLN{7}{dp}{7_dynamic_programming}
\defineLN{8}{flow}{8_flow_reductions_LP}
\defineLN{9}{np}{9_NP}

% The text is the same in the book and in the standalone lecture notes; only
% links differ.  In the book, links to other lecture notes stay in the book.
% In a standalone note, they open the other note's published PDF.
% \noteLink{tag}{target}{text} links text to \hypertarget{target} in lecture note tag.
% \LN{proofs} prints "Lecture Note 2", linked to that lecture note.
\newcommand{\currentLN}{0}   % in a standalone note, its number (set by \doTitle)
\newcommand{\notePDF}[1]{\siteURL notes/\LNfile{#1}.pdf}
\newcommand{\noteLink}[3]{%
  \ifnum\LNnumber{#1}=\currentLN\relax\hyperlink{#2}{#3}%
  \else\href{\notePDF{#1}\#nameddest=#2}{#3}\fi}
% \bookLink{target}{text} links text to \hypertarget{target} in the book's
% appendices (from a standalone note, in the published book).
\DeclareRobustCommand{\bookLink}[2]{%
  \ifdefined\BookMode\hyperlink{#1}{#2}%
  \else\href{\siteURL lecture_notes_on_algorithms.pdf\#nameddest=#1}{#2}\fi}
\DeclareRobustCommand{\LN}[1]{%
  \ifnum\LNnumber{#1}=\currentLN\relax Lecture Note~\LNnumber{#1}%
  \else\href{\notePDF{#1}}{Lecture Note~\LNnumber{#1}}\fi}
\pdfstringdefDisableCommands{\def\LN#1{Lecture Note \LNnumber{#1}}}

% Homeworks refer to the notes' sections, exercises and lemmas by \ref: each
% homework's main file says \lectureNoteRefs after \input{macros}.  A homework
% folder has a link refs/ to src/refs/, which holds each note's label records
% (made by tools/refs.sh from the notes' builds).  Read with xr-hyper, they make
% \ref{greedy: exercise: ...} give the exercise's number in the published note,
% linked to it there; labels only the book has (its appendices) link to the book.
\ifdefined\BookMode\else
  \makeatletter
  \newcommand{\lectureNoteRefs}{%
    \@for\LN@tag:=prelim,proofs,matching,dc,greedy,graphs,dp,flow,np\do{%
      \externaldocument{refs/\LNfile{\LN@tag}}[\notePDF{\LN@tag}]}%
    \externaldocument{refs/book}[\siteURL lecture_notes_on_algorithms.pdf]}
  % hyperref links a \ref to another document's label as URL#anchor; browsers'
  % PDF viewers need URL#nameddest=anchor (the form \noteLink and \bookLink use)
  \def\Hy@setref@link#1#2#3#4#5#6\@nil#7{%
    \begingroup
      \ifx\\#5\\\toks0={\hyper@@link{#5}{#4}}\else\toks0={\hyper@@link{#5}{nameddest=#4}}\fi
      \toks1=\expandafter{#7{#1}{#2}{#3}{#4}{#5}}%
      \edef\x{\endgroup\the\toks0 {\the\toks1 }}%
    \x}
  \makeatother
\fi

% Book mode: main.tex defines \BookMode before \input{macros} to combine the
% lecture notes into one report-class document, one chapter per note.
% Without \BookMode (each lecture note on its own, and the homeworks) none
% of this applies.
\ifdefined\BookMode
  \usepackage{etoc}

  % links to other lecture notes stay in the book
  \renewcommand{\noteLink}[3]{\hyperlink{#2}{#3}}

  % each chapter is a lecture note, and is called one, as in the standalone notes
  \renewcommand{\chaptername}{Lecture Note}
  \setcounter{secnumdepth}{3}   % number subsubsections, as in the standalone notes

  % global table of contents, for main.tex
  \newcommand{\bookContents}{%
    \hypertarget{TOC}{}%
    \etocsetnexttocdepth{section}%
    \tableofcontents}

  % \doTitle{N}{title}{short title} starts chapter N, with a link to the global
  % table of contents and the chapter's own table of contents
  \renewcommand{\doTitle}[3]{%
    \setcounter{chapter}{\numexpr#1-1\relax}%
    \chapter{\texorpdfstring{#2}{#3}}%
    \label{\LNtag{#1}: chapter}%
    \csgdef{LNshorttitle@#1}{#3}%
    \csgdef{LNtitle@#1}{#2}%
    \renewcommand{\LNTitle}{Lecture Note #1, #3}%
    \codeListingstrue
    \hyperlink{TOC}{$\uparrow$\scriptsize\,\textsc{contents}}%
    {\small
      \hypertarget{TOC-#1}{}%
      \etocsettocstyle{\section*{\contentsname}}{}%
      \etocsetnexttocdepth{subsubsection}%
      \localtableofcontents
    }}

  % the code: each lecture note's \codeSection just records its files, and
  % \codeAppendix (in the book's appendix) lists them all, a section per note
  \makeatletter   % (file names have no spaces, so drop the spaces around them)
  \renewcommand{\codeSection}[1]{\csxdef{LNcode@\arabic{chapter}}{\zap@space#1 \@empty}}
  \makeatother
  % \codeAppendix: first a contents list (each note's topic, then its files,
  % linked to the listings), then a section per note listing its files
  \makeatletter
  \def\LNbasename#1{\LN@base#1/\LN@end}   % a/b/c.py -> c.py
  \def\LN@base#1/#2\LN@end{\ifx\relax#2\relax#1\else\LN@base#2\LN@end\fi}
  \makeatother
  \newcommand{\codeContentsFile}[1]{%
    \edef\LNfilename{\LNbasename{#1}}%
    \item \pycode[\texttt{\expandafter\detokenize\expandafter{\LNfilename}}]{#1}}
  \newcommand{\codeAppendix}{%
    \begin{itemize}
    \foreach \n in {1,...,9} {%
      \ifcsdef{LNcode@\n}{%
        \item \hyperref[code: note \n]{\csuse{LNtitle@\n}} (\LN{\LNtag{\n}})
        \begin{itemize}
        \edef\LNcodelist{\csuse{LNcode@\n}}%
        \expandafter\forcsvlist\expandafter\codeContentsFile\expandafter{\LNcodelist}%
        \end{itemize}}{}}
    \end{itemize}
    \newpage
    \foreach \n in {1,...,9} {%
      \ifcsdef{LNcode@\n}{%
        \section[Lecture Note \n: \csuse{LNtitle@\n}]{%
          \texorpdfstring{\LN{\LNtag{\n}}: \csuse{LNtitle@\n}}{Lecture Note \n: \csuse{LNtitle@\n}}}%
        \label{code: note \n}%
        \edef\LNcodelist{\csuse{LNcode@\n}}%
        \expandafter\forcsvlist\expandafter\codefile\expandafter{\LNcodelist}%
        \newpage}{}}}

  % a lecture note is a chapter: \LN{proofs} prints a link "Lecture Note 2"
  \DeclareRobustCommand{\LN}[1]{\hyperref[#1: chapter]{Lecture Note~\ref*{#1: chapter}}}
  \pdfstringdefDisableCommands{\def\LN#1{Lecture Note \LNnumber{#1}}}

  % running head; "contents" links to the current chapter's table of contents
  \renewcommand{\doHeader}[1]{%
    \renewcommand{\header}{%
      \clap{\raisebox{0pt}{\protect\hyperlink{TOC-\arabic{chapter}}{$\uparrow$\scriptsize\,\textsc{contents}\hspace{1.1in}}}}%
      {\footnotesize \LNTitle{} / Section \thesection, #1}}%
    \markboth{\header}{\header}}

  % number theorems, lemmas, etc. and exercises within chapters (e.g. Lemma 8.1)
  \counterwithin{lemma}{chapter}
  \counterwithin{definition}{chapter}
  \counterwithin{theorem}{chapter}
  \counterwithin{corollary}{chapter}
  \counterwithin{claim}{chapter}
  \counterwithin{observation}{chapter}
  \counterwithin{conjecture}{chapter}
  \counterwithin{exercise}{chapter}
\fi

% %%%%%%%%%%%%%%%%%%%%%%%%%%%%%%%%%%%%%%%%%%%%%%%%%%%%%%%%%%%%%%%%%%%%%%%%
% %%%%%%%%%%%%%%%%%%%%%%%%%%%%%%%%%%%%%%%%%%%%%%%%%%%%%%%%%%%%%%%%%%%%%%%%

\ExplSyntaxOn
\seq_new:N \l_itemize_seq
% \NewDocumentEnvironment{REITEMIZE}{ +b }{
\NewDocumentCommand{\REITEMIZECMD}{ +m }{
  \regex_split:nnN {\c{item}} {#1} \l_itemize_seq
  \seq_pop:NN \l_itemize_seq \l_tmpa_tl% remove the (hopefully) empty first item
  \seq_if_empty:NF \l_itemize_seq {
    \l_tmpa_tl
    \seq_map_inline:Nn \l_itemize_seq {
      \regex_split:nnN {\c{item}} {#1} \l_itemize_seq {
        \regex_extract_once:nnNTF { ^\[([^]]*)\](.*) } { ##1} \l_foo_seq
      { \item[\seq_item:Nn \l_foo_seq {2}]\ITEMIZER{\seq_item:Nn \l_foo_seq {3}} } { \item\ITEMIZER{##1}  }
    }
  }
}
}
\ExplSyntaxOff

\newlength{\ITEMHT}
\setlength{\ITEMHT}{1in}

% \setlength{\fboxsep}{1pt}
\newcommand{\ITEMIZER}[1]{%
  % \framebox{
\adjustbox{valign=t}{\parbox{\textwidth}{\linespread{1.15}\selectfont #1}}~~\adjustbox{valign=t}{\parbox{0pt}{\color{gray}\rule[-\ITEMHT]{0.01pt}{\ITEMHT}}}
% }
  % \begin{tabular}{@{}p{\textwidth}@{\rule[-\ITEMHT]{0pt}{\ITEMHT}}}
  %   #1
  % \end{tabular}
}

\newsavebox\thesmashminipage
\newlength{\thesmashminipagelength}
\newlength{\answerboxlength}
\newlength{\answerboxoverrun}
\newenvironment{smashminipage}[2]
{\setlength\answerboxlength{#2}
\begin{lrbox}{\thesmashminipage}\begin{minipage}[t]{#1}}
{\end{minipage}\end{lrbox}%
\setlength{\thesmashminipagelength}{\dimexpr\ht\thesmashminipage+\dp\thesmashminipage}%
\setlength{\answerboxoverrun}{\dimexpr\thesmashminipagelength-\answerboxlength}%
\ifdim\the\thesmashminipagelength>\the\answerboxlength%
\usebox{\thesmashminipage}%
\smash{\color{red}\hspace*{-0.98\textwidth}\rule[-\answerboxlength]{0.98\textwidth}{1pt}}
\errmessage{Answer exceeds allotted length, please shorten it.}
\else
\usebox{\thesmashminipage}%
\fi
}



\newenvironment{answerBox}[1]{%
  \def\ANSWERBOXHEIGHT{#1}
  \begin{adjustbox}{valign=t}
  \begin{smashminipage}{\textwidth}{\ANSWERBOXHEIGHT}
}{\end{smashminipage}\end{adjustbox}~~~\adjustbox{valign=t}{\parbox{0pt}{\color{gray}\rule[-\ANSWERBOXHEIGHT]{0.01pt}{\ANSWERBOXHEIGHT}\rule[-\ANSWERBOXHEIGHT]{10pt}{0.01pt}}}}

% % \setlist[enumerate]{itemsep=0pt,parsep=0pt,topsep=0pt,leftmargin=100pt}

% \makeatletter 
\LetLtxMacro\SVenumerate\enumerate
\let\endSVenumerate\endenumerate

\newcommand{\injectEnumerate}{
  \let\SValph\alph
  \def\alph{}
  \renewenvironment{enumerate}[1][]
  {\let\alph\SValph \SVenumerate[label=(\alph*),itemsep=2pt,parsep=0pt]\REITEMIZE}{\endREITEMIZE\endSVenumerate}
  % \setlist[enumerate]{before=\Collect@Body\REITEMIZECMD}
}
% \makeatother
\newcommand{\restoreEnumerate}{
  \renewenvironment{enumerate}[1][]
  \let\SValph\alph
  \def\alph{}
  {\let\alph\SValph \SVenumerate[label=(\alph*)]}{\endSVenumerate}
  % \setlist[enumerate]{before=,after=}
}
\newcommand{\injectBlue}{
  \let\SVblue\blue
  \let\endSVblue\endblue
  \def\blue{\SVblue\REITEMIZE}
  \def\endblue{\endREITEMIZE\endSVblue}
}
\newcommand{\restoreBlue}{
  \let\blue\SVblue
  \let\endblue\endSVblue
}
\newcommand{\injectEnumerateBlue}{
  \setlist[enumerate]{before=\injectBlue,after=\restoreBlue,itemsep=2pt,parsep=0pt}
}
\newcommand{\restoreEnumerateBlue}{
  \setlist[enumerate]{before=,after=}
}


% %%%%%%%%%%%%%%%%%%%%%%%%%%%%%%%%%%%%%%%%%%%%%%%%%%%%%%%%%%%%%%%%%%%%%%%%
% %%%%%%%%%%%%%%%%%%%%%%%%%%%%%%%%%%%%%%%%%%%%%%%%%%%%%%%%%%%%%%%%%%%%%%%%

% \NewDocumentEnvironment{env}{ }
% {\Environenv}
% {\endEnvironenv}

% \NewEnviron{REITEMIZE}{\REITEMIZECMD\BODY}

\makeatletter 
\newenvironment{REITEMIZE}{\Collect@Body\REITEMIZECMD}{}
\makeatother 

\newcommand{\NODEone}[1]{\ensuremath{v_{#1}}\xspace}
\newcommand{\NODE}[2]{{\NODEone {#1,#2}}\xspace}
\newcommand{\EDGE}[2]{\ensuremath{(\textsf{#1}, \textsf{#2})}}
\newcommand{\ITEM}{\item}

%%% Local Variables:
%%% mode: latex
%%% End:
.  A homework
% folder has a link refs/ to src/refs/, which holds each note's label records
% (made by tools/refs.sh from the notes' builds).  Read with xr-hyper, they make
% \ref{greedy: exercise: ...} give the exercise's number in the published note,
% linked to it there; labels only the book has (its appendices) link to the book.
\ifdefined\BookMode\else
  \makeatletter
  \newcommand{\lectureNoteRefs}{%
    \@for\LN@tag:=prelim,proofs,matching,dc,greedy,graphs,dp,flow,np\do{%
      \externaldocument{refs/\LNfile{\LN@tag}}[\notePDF{\LN@tag}]}%
    \externaldocument{refs/book}[\siteURL lecture_notes_on_algorithms.pdf]}
  % hyperref links a \ref to another document's label as URL#anchor; browsers'
  % PDF viewers need URL#nameddest=anchor (the form \noteLink and \bookLink use)
  \def\Hy@setref@link#1#2#3#4#5#6\@nil#7{%
    \begingroup
      \ifx\\#5\\\toks0={\hyper@@link{#5}{#4}}\else\toks0={\hyper@@link{#5}{nameddest=#4}}\fi
      \toks1=\expandafter{#7{#1}{#2}{#3}{#4}{#5}}%
      \edef\x{\endgroup\the\toks0 {\the\toks1 }}%
    \x}
  \makeatother
\fi

% Book mode: main.tex defines \BookMode before %%%% DON'T EDIT THIS FILE %%%%

\usepackage{xifthen}% \isempty
\usepackage{etoolbox}% \patchcmd
% \usepackage{stringstrings}% \substring

\usepackage{amsmath, amsthm, amssymb}
\usepackage[margin=1in]{geometry}
\usepackage{xspace}
\usepackage{graphicx}
\usepackage[normalem]{ulem}
\usepackage{adjustbox}
\usepackage{xcolor}
\usepackage{changepage}
\usepackage{framed}
\usepackage{multicol}
\usepackage{enumitem}
% \usepackage{listings} % lstlistings environment

\usepackage{makecmds}           %\provideenvironment

\usepackage{xparse,environ}
\usepackage{letltxmacro}       % \LetLtxMacro
\usepackage{wrapfig}
\usepackage{tikz}

\usetikzlibrary{shapes}
\usetikzlibrary{positioning}
\usetikzlibrary{shapes.misc}
\usetikzlibrary{arrows}

\newtheorem{lemma}{Lemma}
\newtheorem{definition}{Definition}
\newtheorem{theorem}{Theorem}
\newtheorem{corollary}{Corollary}
\newtheorem{claim}{Claim}
\newtheorem{observation}{Observation}
\newtheorem{conjecture}{Conjecture}

% A lemma or theorem restated later with its original number (set with
% \setcounter) goes in a restatement environment, which gives it its own link
% anchor (otherwise it would duplicate the original's).
\newenvironment{restatement}
  {\def\theHlemma{restated.\thelemma}\def\theHtheorem{restated.\thetheorem}}{}

\newenvironment{myframed}
{\smallskip\par\noindent\begin{minipage}{\textwidth}\begin{framed}\vspace*{-5pt}}{\vspace*{-5pt}\end{framed}\end{minipage}\smallskip}


\definecolor{darkblue}{RGB}{0,0,200}

 
\let\oldParagraph\paragraph
\providecommand{\gradescopeColor}{\color{black}}
\renewcommand{\paragraph}[1]{\oldParagraph{\gradescopeColor #1}}
\providecommand{\gradescopeInit}{}

\providecommand{\BLUE}[1]{{\color{darkblue} #1}}
\providecommand{\REPLACEME}{\textcolor{red}{\ensuremath{\star \cdots \star}}\xspace}
\provideenvironment{blue}{\color{darkblue}}{}

\newcommand{\dist}[2]{\ensuremath{\operatorname{\sf dist}(#1, #2)}}
\newcommand{\depth}[2]{\ensuremath{\operatorname{\sf depth}_{#1}(#2)}}
\newcommand{\cost}[1]{\ensuremath{\operatorname{\sf cost}(#1)}}
\newcommand{\weight}[1]{\ensuremath{\operatorname{\sf wt}(#1)}}
% \newcommand{\val}[1]{\ensuremath{\operatorname{\sf val}(#1)}}
\newcommand{\eps}{\varepsilon}
\newcommand{\capac}{\textsf{cap}}
\newcommand{\capOf}[1]{\ensuremath{\operatorname{\sf cap}(#1)}}
\newcommand{\rescap}[1]{\ensuremath{\operatorname{\sf cap}_f(#1)}}
\newcommand{\into}[1]{\ensuremath{\operatorname{\sf in}(#1)}}
\newcommand{\out}[1]{\ensuremath{\operatorname{\sf out}(#1)}}
\newenvironment{LP}[1]{
  \begin{array}[t]{@{}l@{}}
    \displaystyle #1 \\
    \begin{aligned}\\[-10pt]}{\end{aligned}\end{array}}

\newcommand{\opt}[1]{\ensuremath{\operatorname{\sf opt}(#1)}}
\newcommand{\Opt}[1]{\ensuremath{\operatorname{\sf Opt}(#1)}}
\newcommand{\half}{\frac{1}{2}}

\newcommand{\E}{\operatorname{E}}
\newcommand{\giv}{\,|\,}
\newcommand{\Z}{\mathbb{Z}}
\newcommand{\Zp}{\mathbb{Z}_{\ge 0}}
\newcommand{\R}{\mathbb{R}}
% \newcommand{\Rp}{\mathbb{R}_{\ge 0}}
\newcommand{\Rp}{\mathbb{R}_{\tiny +}}
\newcommand{\N}{\mathbb{N}}
\newcommand{\Np}{\mathbb{N}_{\ge 0}}
\newcommand{\calS}{{\cal S}}

\newcommand{\OP}[1]{\text{\small\sc #1}\xspace}
\newcommand{\Insert}{\OP{Insert}}
\newcommand{\FindMin}{\OP{Find-Min}}
\newcommand{\DeleteMin}{\OP{Delete-Min}}
\newcommand{\DecreaseKey}{\OP{Decrease}}
\newcommand{\Merge}{\OP{Merge}}
\newcommand{\Cut}{\OP{Cut}}
% \newcommand{\prob}[1]{\text{\sc #1}}
% a problem name; in running text, lines may break inside it
\DeclareRobustCommand{\prob}[1]{\ifmmode\text{#1}\else#1\fi}
\newcommand{\encode}[1]{\langle #1 \rangle}

\newcommand{\best}[1]{\text{\sf best}(#1)}
\newcommand{\worst}[1]{\text{\sf worst}(#1)}
\newcommand{\floor}[1]{\lfloor #1 \rfloor}
\newcommand{\ceil}[1]{\lceil #1 \rceil}

\setlength{\parskip}{2pt}
\setlength{\parindent}{1em}

% \setlength\fboxsep{0pt}

%%% longFormProof, shortFormProof, block environments
%%% \step

\newlist{steps}{enumerate}{7}
\setlist[steps]{wide,
  topsep=1pt,
  parsep=1pt,
  partopsep=1pt,
  itemsep=0pt,
  align=left,
  labelindent=0pt,
  itemindent=0pt, 
}
\setlist[steps,1]{
  label=\arabic*.,
  ref=\arabic*
}
\setlist[steps,2]{
  label=\arabic{stepsi}.\arabic*.,
  ref=\arabic{stepsi}.\arabic*
}
\setlist[steps,3]{
  label=\arabic{stepsi}.\arabic{stepsii}.\arabic*.,
  ref=\arabic{stepsi}.\arabic{stepsii}.\arabic*
}
\setlist[steps,4]{
  label=\arabic{stepsi}.\arabic{stepsii}.\arabic{stepsiii}.\arabic*.,
  ref=\arabic{stepsi}.\arabic{stepsii}.\arabic{stepsiii}.\arabic*
}
\setlist[steps,5]{
  label=\arabic{stepsi}.\arabic{stepsii}.\arabic{stepsiii}.\arabic{stepsiv}.\arabic*.,
  ref=\arabic{stepsi}.\arabic{stepsii}.\arabic{stepsiii}.\arabic{stepsiv}.\arabic*
}
\setlist[steps,6]{
  label=\arabic{stepsi}.\arabic{stepsii}.\arabic{stepsiii}.\arabic{stepsiv}.\arabic{stepsv}.\arabic*.,
  ref=\arabic{stepsi}.\arabic{stepsii}.\arabic{stepsiii}.\arabic{stepsiv}.\arabic{stepsv}.\arabic*
}
\setlist[steps,7]{
  label=\arabic{stepsi}.\arabic{stepsii}.\arabic{stepsiii}.\arabic{stepsiv}. \arabic{stepsv}. \arabic{stepsvi}.\arabic*.,
  ref=\arabic{stepsi}.\arabic{stepsii}.\arabic{stepsiii}.\arabic{stepsiv}.\arabic{stepsv}. \arabic{stepsvi}.\arabic*
}

\newcommand{\step}[1][]{\item\maybeLabel{#1}}

\makeatletter
% \skipitem steps the list counter itself (for block labels).  Give each one
% its own hyperref link target, since step numbers repeat from proof to proof.
\newcounter{skipitem}
\newcommand{\skipitem}{\stepcounter{skipitem}\refstepcounter{\@enumctr}}
\renewcommand{\theHstepsi}{skipitem.\arabic{skipitem}}
\renewcommand{\theHstepsii}{skipitem.\arabic{skipitem}}
\renewcommand{\theHstepsiii}{skipitem.\arabic{skipitem}}
\renewcommand{\theHstepsiv}{skipitem.\arabic{skipitem}}
\renewcommand{\theHstepsv}{skipitem.\arabic{skipitem}}
\renewcommand{\theHstepsvi}{skipitem.\arabic{skipitem}}
\renewcommand{\theHstepsvii}{skipitem.\arabic{skipitem}}
\makeatother

\newcommand{\maybeLabel}[1]{\ifthenelse{\isempty{#1}}{}{\label{#1}}}

\newenvironment{longFormProof}[1][Proof (long form).]
{\begin{proof}[#1]~\par\begin{steps}}{\end{steps}\vspace*{-3ex}\end{proof}}

\newenvironment{shortFormProof}[1][Proof (short form).]
{\begin{proof}[#1]}{\end{proof}}

\newenvironment{longFormSubProof}[1][Proof.]
{\begin{adjustwidth}{1em}{0em}\begin{proof}[#1]~\par\vspace*{-2pt}\skipitem\begin{steps}}
{\end{steps}\vspace*{-1ex}\end{proof}\end{adjustwidth}\smallskip}

\newenvironment{shortFormSubProof}[1][Proof.]
{\begin{proof}[#1]}
{\end{proof}\vspace*{-3pt}}

\newenvironment{block}[1][]
{\skipitem\maybeLabel{#1}\item[] \begin{steps}}{\end{steps}}

\newcommand{\comment}[1]{\hfill{\footnotesize\emph{#1}}}
\newcommand{\lineacross}{\par\vspace*{-0.7\baselineskip}\noindent\hrulefill\par}

%%% problem

% \newcounter{problem}
% \newcommand{\problem}[1][]{}
% {\ifthenelse{\isempty{#1}}{\refstepcounter{problem}}{}%
%   \paragraph{Problem \ifthenelse{\isempty{#1}}
%     {\bf\arabic{problem}.}
%     {#1}} }

\newcommand{\problemNumbered}[1]{
  \setcounter{problem}{#1}%
  \addtocounter{problem}{-1}%
  \problem
}
\newenvironment{problemTheorem}{%
  \setcounter{theorem}{\value{problem}}%
  \addtocounter{theorem}{-1}%
  \begin{theorem}[Problem \arabic{problem}]
  }{\end{theorem}}

\newcounter{problem}
\newcommand{\problem}[1][]
{\ifthenelse{\isempty{#1}}
  {\refstepcounter{problem}\paragraph{Problem~\arabic{problem}.}}
  {\paragraph{Problem~#1.}}\refreshHeader} 

\newcounter{exercise}
\newcommand{\exercise}[1][]
{\ifthenelse{\isempty{#1}}
  {\refstepcounter{exercise}\paragraph{Exercise~\theexercise.}}
  {\paragraph{Exercise~#1.}}}

\newcommand{\solution}{\paragraph{Solution.}} 

\newenvironment{mylist}[1][]
{\begin{list}{#1}{
      \setlength{\itemsep}{-1pt}
      \setlength{\labelwidth}{0pt}
      \setlength{\leftmargin}{1.25em}
    }}
{\end{list}}

\newcommand{\LNTitle}{\textcolor{red}{UNDEFINED LNTitle}}

% In a standalone lecture note's table of contents, leave room for numbers
% like 7.10 (in bold) and 7.10.2.
\makeatletter
\newcommand{\LNtocWidths}{%
  \patchcmd{\l@section}{1.5em}{2.8em}{}{\PackageWarning{macros}{Could not widen section numbers in the contents}}%
  \renewcommand*{\l@subsection}{\@dottedtocline{2}{2.8em}{3.2em}}%
  \renewcommand*{\l@subsubsection}{\@dottedtocline{3}{6.0em}{4.1em}}}
\makeatother

\newcommand{\doTitle}[3]{
  \pagenumbering{roman}
  \renewcommand{\LNTitle}{Lecture Note #1, #3}
  \renewcommand{\currentLN}{#1}
  \codeListingstrue
  % number sections, theorems etc., exercises, figures, tables and equations
  % within the lecture note (e.g. Section 8.2, Lemma 8.6), as in the book
  \renewcommand{\thesection}{#1.\arabic{section}}
  \renewcommand{\thelemma}{#1.\arabic{lemma}}
  \renewcommand{\thedefinition}{#1.\arabic{definition}}
  \renewcommand{\thetheorem}{#1.\arabic{theorem}}
  \renewcommand{\thecorollary}{#1.\arabic{corollary}}
  \renewcommand{\theclaim}{#1.\arabic{claim}}
  \renewcommand{\theobservation}{#1.\arabic{observation}}
  \renewcommand{\theconjecture}{#1.\arabic{conjecture}}
  \renewcommand{\theexercise}{#1.\arabic{exercise}}
  \renewcommand{\thefigure}{#1.\arabic{figure}}
  \renewcommand{\thetable}{#1.\arabic{table}}
  \renewcommand{\theequation}{#1.\arabic{equation}}
  \LNtocWidths
  \title{Lecture Note #1\\ #2}
  % a note's \LNabstract, if it defines one before \doTitle, goes right below the title
  \date{\ifdefined\LNabstract\parbox{0.85\textwidth}{\normalsize\normalfont\LNabstract}\\[3ex]\fi
    draft of \today, by \href{\authorURL}{N.~Young}\\[2ex] {\small\copyrightNotice}}
  \maketitle
  {\small
    \hypertarget{TOC}{}
    \tableofcontents
  }
  \thispagestyle{empty}
  % the text starts on page 1, after the title and contents (numbered i, ii, ...)
  \clearpage
  \pagenumbering{arabic}
}

\def\clap#1{\hbox to 0pt{\hss#1\hss}}

\newcommand{\refreshHeader}{}
\newcommand{\header}{}

\newcommand{\doHeader}[1]{      % for lecture notes
  \renewcommand{\header}{%
    \clap{\raisebox{0pt}{\protect\hyperlink{TOC}{$\uparrow$\scriptsize\,\textsc{contents}\hspace{1.1in}}}}%
    {\footnotesize \LNTitle{} / Section \thesection, #1}} %, #1, draft of \today}}
  \markboth{\header}{\header}
}

\newcommand{\SID}{}
\newcommand{\setAuthor}[1]{\renewcommand{\author}{#1}}
\newcommand{\setSID}[1]{\renewcommand{\SID}{#1}}

\newcommand{\doHwHeader}[1]{                                         % for homeworks
  \renewcommand{\refreshHeader}{
    \renewcommand{\header}{%
      {\footnotesize \classSection #1 / Problem~\arabic{problem} / draft of \today \hfill \author\ifGradescope~(\SID)\fi~~}}
    \markboth{\header}{\header}
  }
  \renewcommand{\header}{%
    {\footnotesize \classSection #1 / Instructions / draft of \today \hfill \author\ifGradescope~(\SID)\fi~~}}
  \markboth{\header}{\header}
}

% Homework settings, for instructors using the homeworks in a course:
% \classSection is the course name shown in each homework's running head
% (e.g. \renewcommand{\classSection}{CS 3000\xspace}), and \Gradescopetrue turns on
% the instructions for submitting to Gradescope (in its fixed-length mode),
% the hints saying on which page of the PDF each answer should be, and the
% student ID in the running head.
\newcommand{\classSection}{}
\newif\ifGradescope
\Gradescopefalse
\newcommand{\GradescopeComment}[1]{\ifGradescope\comment{#1}\fi}

\pagestyle{myheadings}
\addtolength{\headsep}{0.2in}
\addtolength{\topmargin}{-0.3in}
\addtolength{\textheight}{0.4in}

%%%%%%%%%%%%%%%%%%%%%%%%%%%%%%%%%%%%%%%%%%%%%%%%
%%%%%%%%%%%%%%%%%%%%%%%%%%%%%%%%%%%%%%%%%%%%%%%%

\newcommand{\MF}{\href
  {https://mfleck.cs.illinois.edu/building-blocks/}
  {MF}\xspace}

\newcommand{\KWSlides}{\href
  {https://www.cs.princeton.edu/~wayne/kleinberg-tardos/}
  {KW slides}\xspace}

\newcommand{\LLM}{\href
  {https://courses.csail.mit.edu/6.042/spring18/mcs.pdf}
  {LLM}\xspace}

% \DPV, \KT, \CLRS: textbooks listed (without links) in the table of sources in
% the Preliminaries note; each prints as a link to its entry in that table.
\DeclareRobustCommand{\DPV}{\noteLink{prelim}{source.DPV}{DPV}\xspace}
\DeclareRobustCommand{\KT}{\noteLink{prelim}{source.KT}{KT}\xspace}
\DeclareRobustCommand{\CLRS}{\noteLink{prelim}{source.CLRS}{CLRS}\xspace}
\AtBeginDocument{\pdfstringdefDisableCommands{\def\DPV{DPV }\def\KT{KT }\def\CLRS{CLRS }}}

\newcommand{\JE}{\href
  {https://jeffe.cs.illinois.edu/teaching/algorithms/}
  {JE}\xspace}

\newcommand{\JEGreedy}{\href
  {https://jeffe.cs.illinois.edu/teaching/algorithms/book/04-greedy.pdf}
  {JE Chapter 4}\xspace}

\newcommand{\BBCode}{\href
{https://northeastern.blackboard.com/webapps/cmsmain/webui/courses/CS3000.MERGED.202010/code?action=frameset&subaction=view&uniq=-jj2max&course_id=_2597392_1&mask=\%2Fcourses\%2FCS3000.MERGED.202010}
  {code}\xspace}

% The Python code (src/code).  Each lecture note that has code ends with a
% section listing its files: \codeSection{path, path, ...}.  In the text,
% \pycode{greedy_algorithms/Dijkstra.py} prints the path as a link to its listing
% (in the homeworks, which have no listings, to the file on GitHub), and
% \pycode[text]{...} uses the given text instead.  \pythonCode{path, ...} is the
% standard sentence ending a section whose algorithms have code:
% "Python code: path, ... (Section N.k)."
\newcommand{\codeFolder}{https://github.com/nealeyoung/algs101/tree/main/src/code}
\newcommand{\codeRepo}{https://github.com/nealeyoung/algs101/blob/main/src/code}
\newif\ifcodeListings   % set by \doTitle: this document has the listings
\newcommand{\githubCode}[2][]{\href{\codeRepo/#2}{\ifx\relax#1\relax\nolinkurl{#2}\else#1\fi}}
\newcommand{\pycode}[2][]{%
  \ifcodeListings\hyperref[code: #2]{\ifx\relax#1\relax\nolinkurl{#2}\else#1\fi}%
  \else\githubCode[#1]{#2}\fi}
\makeatletter
\newcommand{\pythonCode}[1]{%
  \def\pythonCode@first{}\def\pythonCode@list{}%
  \forcsvlist{\pythonCode@add}{#1}%
  Python code: \pythonCode@list\ (Section~\ref{code: \pythonCode@first}).}
\newcommand{\pythonCode@add}[1]{%
  \ifx\pythonCode@first\@empty\def\pythonCode@first{#1}\appto\pythonCode@list{\pycode{#1}}%
  \else\appto\pythonCode@list{, \pycode{#1}}\fi}
\makeatother

% Code listings that copy and paste exactly (in PDF viewers that honor the
% ToUnicode maps, e.g. MuPDF-based ones; PDFKit viewers such as Preview and
% Skim drop leading spaces, so each listing's heading links to the file on
% GitHub).  Spaces are typeset as white visible-space glyphs that copy as
% spaces, and hyphens and asterisks as the T1 glyphs, so they copy as themselves.
\usepackage[T1,OT1]{fontenc}   % T1 for the code listings only; OT1 stays the default
\usepackage{listings}
\pdfgentounicode=1
\input glyphtounicode
\pdfglyphtounicode{uni2423}{0020}
\makeatletter
\def\lst@visiblespace{\textcolor{white}{\char32}}
% put an invisible space on each blank line, so that blank lines copy too
% (listings' OnEmptyLine hook runs before the line break, so patch the loop)
\def\lst@DoNewLines{%
    \@whilenum\lst@newlines>\@ne \do {\lst@NewLine \lst@visiblespace}%
    \ifnum\lst@newlines>\z@ \lst@NewLine \fi}
\makeatother
\lstset{language=Python,
  basicstyle=\fontencoding{T1}\fontfamily{lmtt}\selectfont\footnotesize,
  upquote=true, showstringspaces=false, columns=fullflexible, keepspaces=true,
  showspaces=true, literate={-}{{\char45}}1 {*}{{\char42}}1,
  frame=single, framesep=1ex, rulecolor=\color{gray}}
% the listings' paths are relative to the lecture note's folder (in the book, to src/)
\ifdefined\BookMode\newcommand{\codeDir}{code/}\else\newcommand{\codeDir}{../code/}\fi
% \codefile{path}: a listing of src/code/path, headed by a link to it on GitHub
\newcommand{\codefile}[1]{%
  \subsection*{\githubCode{#1}}\label{code: #1}%
  \lstinputlisting{\codeDir#1}}
% \codeSection{path, path, ...}: the lecture note's last section, listing its code
\newcommand{\codeSection}[1]{%
  \newpage
  \section{Python code}\label{\LNtag{\currentLN}: sec: code}\doHeader{Python code}
  This section lists the Python code for this lecture note.
  Click on a file's name to open the file on GitHub.
  Copying code from GitHub is more reliable than copying it from this PDF,
  since some PDF viewers drop the leading spaces (the indentation) of copied lines.
  \forcsvlist{\codefile}{#1}}
% the copyright notice, for title pages and the homework first pages
\newcommand{\licenseURL}{https://github.com/nealeyoung/algs101/blob/main/LICENSE}
\newcommand{\authorURL}{https://www.cs.ucr.edu/\string~neal}
\newcommand{\copyrightNotice}{%
  \copyright\ 2018--2026 \href{\authorURL}{Neal E. Young}.
  Text: CC BY-NC-SA 4.0; code: MIT; except where noted.
  See \url{\licenseURL}.}
% the published site: PDFs of the notes, and the homework assignments and templates
\newcommand{\siteURL}{https://nealeyoung.github.io/algs101/}

\newcommand{\blackboard}{\href
  {https://northeastern.blackboard.com/webapps/blackboard/execute/launcher?type=Course&id=_2597392_1&url=}
  {Blackboard}\xspace}

\newcommand{\JUSlides}{\href
  {http://www.khoury.neu.edu/home/jullman/cs3000f18/schedule.html}
  {Prof.~Ullman's slides}\xspace}

% a URL on its own line, in three layouts
\newcommand{\URL}[1]{\par-- \parbox[t]{0.8\textwidth}{\url{#1}} \medskip\par}
\newcommand{\URLsmall}[1]{\par-- \parbox[t]{0.95\textwidth}{\small\url{#1}} \medskip\par}
\newcommand{\URLplain}[1]{\par\parbox[t]{\textwidth}{\url{#1}}\par}

%%%%%%%%%%%%%%%%%%%%%%%%%%%%%%%%%%%%%%%%%%%%%%%%
%%%%%%%%%%%%%%%%%%%%%%%%%%%%%%%%%%%%%%%%%%%%%%%%

\newcommand{\continued}{\vfill\centerline{\emph{(continued)}}\vfill\newpage}

% load hyperref last to avoid conflicts; long URLs may also break after hyphens
\PassOptionsToPackage{hyphens}{url}
\ifdefined\BookMode\else\usepackage{xr-hyper}\fi   % for \lectureNoteRefs (homeworks), below
\usepackage{hyperref}
\hypersetup{
    colorlinks,
    linkcolor={blue!50!black},
    citecolor={blue!50!black},
    urlcolor={blue!80!black}
}

\providecommand{\answerDirectory}{.}
\newcommand{\inputAnswer}[1]{\refreshHeader\input{\answerDirectory/#1}}
\newcommand{\inputSection}[1]{\input{#1}}

% Refer to a lecture note by its tag (the prefix of its labels), e.g.
% \LN{proofs} prints "Lecture Note 2" (in the book, a link to that chapter).  The table below is the only place the numbers appear.
\makeatletter
\newcommand{\defineLN}[3]{\@namedef{LNnumber@#2}{#1}\@namedef{LNtag@#1}{#2}\@namedef{LNfile@#2}{#3}}
\newcommand{\LNnumber}[1]{%
  \ifcsname LNnumber@#1\endcsname\csname LNnumber@#1\endcsname
  \else\PackageError{macros}{Unknown lecture-note tag `#1'}{}\fi}
\newcommand{\LNtag}[1]{\csname LNtag@#1\endcsname}
\newcommand{\LNfile}[1]{\csname LNfile@#1\endcsname}
\makeatother
\defineLN{1}{prelim}{1_preliminaries}
\defineLN{2}{proofs}{2_proofs}
\defineLN{3}{matching}{3_stable_matching}
\defineLN{4}{dc}{4_divide_and_conquer}
\defineLN{5}{greedy}{5_greedy}
\defineLN{6}{graphs}{6_graph_traversal}
\defineLN{7}{dp}{7_dynamic_programming}
\defineLN{8}{flow}{8_flow_reductions_LP}
\defineLN{9}{np}{9_NP}

% The text is the same in the book and in the standalone lecture notes; only
% links differ.  In the book, links to other lecture notes stay in the book.
% In a standalone note, they open the other note's published PDF.
% \noteLink{tag}{target}{text} links text to \hypertarget{target} in lecture note tag.
% \LN{proofs} prints "Lecture Note 2", linked to that lecture note.
\newcommand{\currentLN}{0}   % in a standalone note, its number (set by \doTitle)
\newcommand{\notePDF}[1]{\siteURL notes/\LNfile{#1}.pdf}
\newcommand{\noteLink}[3]{%
  \ifnum\LNnumber{#1}=\currentLN\relax\hyperlink{#2}{#3}%
  \else\href{\notePDF{#1}\#nameddest=#2}{#3}\fi}
% \bookLink{target}{text} links text to \hypertarget{target} in the book's
% appendices (from a standalone note, in the published book).
\DeclareRobustCommand{\bookLink}[2]{%
  \ifdefined\BookMode\hyperlink{#1}{#2}%
  \else\href{\siteURL lecture_notes_on_algorithms.pdf\#nameddest=#1}{#2}\fi}
\DeclareRobustCommand{\LN}[1]{%
  \ifnum\LNnumber{#1}=\currentLN\relax Lecture Note~\LNnumber{#1}%
  \else\href{\notePDF{#1}}{Lecture Note~\LNnumber{#1}}\fi}
\pdfstringdefDisableCommands{\def\LN#1{Lecture Note \LNnumber{#1}}}

% Homeworks refer to the notes' sections, exercises and lemmas by \ref: each
% homework's main file says \lectureNoteRefs after \input{macros}.  A homework
% folder has a link refs/ to src/refs/, which holds each note's label records
% (made by tools/refs.sh from the notes' builds).  Read with xr-hyper, they make
% \ref{greedy: exercise: ...} give the exercise's number in the published note,
% linked to it there; labels only the book has (its appendices) link to the book.
\ifdefined\BookMode\else
  \makeatletter
  \newcommand{\lectureNoteRefs}{%
    \@for\LN@tag:=prelim,proofs,matching,dc,greedy,graphs,dp,flow,np\do{%
      \externaldocument{refs/\LNfile{\LN@tag}}[\notePDF{\LN@tag}]}%
    \externaldocument{refs/book}[\siteURL lecture_notes_on_algorithms.pdf]}
  % hyperref links a \ref to another document's label as URL#anchor; browsers'
  % PDF viewers need URL#nameddest=anchor (the form \noteLink and \bookLink use)
  \def\Hy@setref@link#1#2#3#4#5#6\@nil#7{%
    \begingroup
      \ifx\\#5\\\toks0={\hyper@@link{#5}{#4}}\else\toks0={\hyper@@link{#5}{nameddest=#4}}\fi
      \toks1=\expandafter{#7{#1}{#2}{#3}{#4}{#5}}%
      \edef\x{\endgroup\the\toks0 {\the\toks1 }}%
    \x}
  \makeatother
\fi

% Book mode: main.tex defines \BookMode before \input{macros} to combine the
% lecture notes into one report-class document, one chapter per note.
% Without \BookMode (each lecture note on its own, and the homeworks) none
% of this applies.
\ifdefined\BookMode
  \usepackage{etoc}

  % links to other lecture notes stay in the book
  \renewcommand{\noteLink}[3]{\hyperlink{#2}{#3}}

  % each chapter is a lecture note, and is called one, as in the standalone notes
  \renewcommand{\chaptername}{Lecture Note}
  \setcounter{secnumdepth}{3}   % number subsubsections, as in the standalone notes

  % global table of contents, for main.tex
  \newcommand{\bookContents}{%
    \hypertarget{TOC}{}%
    \etocsetnexttocdepth{section}%
    \tableofcontents}

  % \doTitle{N}{title}{short title} starts chapter N, with a link to the global
  % table of contents and the chapter's own table of contents
  \renewcommand{\doTitle}[3]{%
    \setcounter{chapter}{\numexpr#1-1\relax}%
    \chapter{\texorpdfstring{#2}{#3}}%
    \label{\LNtag{#1}: chapter}%
    \csgdef{LNshorttitle@#1}{#3}%
    \csgdef{LNtitle@#1}{#2}%
    \renewcommand{\LNTitle}{Lecture Note #1, #3}%
    \codeListingstrue
    \hyperlink{TOC}{$\uparrow$\scriptsize\,\textsc{contents}}%
    {\small
      \hypertarget{TOC-#1}{}%
      \etocsettocstyle{\section*{\contentsname}}{}%
      \etocsetnexttocdepth{subsubsection}%
      \localtableofcontents
    }}

  % the code: each lecture note's \codeSection just records its files, and
  % \codeAppendix (in the book's appendix) lists them all, a section per note
  \makeatletter   % (file names have no spaces, so drop the spaces around them)
  \renewcommand{\codeSection}[1]{\csxdef{LNcode@\arabic{chapter}}{\zap@space#1 \@empty}}
  \makeatother
  % \codeAppendix: first a contents list (each note's topic, then its files,
  % linked to the listings), then a section per note listing its files
  \makeatletter
  \def\LNbasename#1{\LN@base#1/\LN@end}   % a/b/c.py -> c.py
  \def\LN@base#1/#2\LN@end{\ifx\relax#2\relax#1\else\LN@base#2\LN@end\fi}
  \makeatother
  \newcommand{\codeContentsFile}[1]{%
    \edef\LNfilename{\LNbasename{#1}}%
    \item \pycode[\texttt{\expandafter\detokenize\expandafter{\LNfilename}}]{#1}}
  \newcommand{\codeAppendix}{%
    \begin{itemize}
    \foreach \n in {1,...,9} {%
      \ifcsdef{LNcode@\n}{%
        \item \hyperref[code: note \n]{\csuse{LNtitle@\n}} (\LN{\LNtag{\n}})
        \begin{itemize}
        \edef\LNcodelist{\csuse{LNcode@\n}}%
        \expandafter\forcsvlist\expandafter\codeContentsFile\expandafter{\LNcodelist}%
        \end{itemize}}{}}
    \end{itemize}
    \newpage
    \foreach \n in {1,...,9} {%
      \ifcsdef{LNcode@\n}{%
        \section[Lecture Note \n: \csuse{LNtitle@\n}]{%
          \texorpdfstring{\LN{\LNtag{\n}}: \csuse{LNtitle@\n}}{Lecture Note \n: \csuse{LNtitle@\n}}}%
        \label{code: note \n}%
        \edef\LNcodelist{\csuse{LNcode@\n}}%
        \expandafter\forcsvlist\expandafter\codefile\expandafter{\LNcodelist}%
        \newpage}{}}}

  % a lecture note is a chapter: \LN{proofs} prints a link "Lecture Note 2"
  \DeclareRobustCommand{\LN}[1]{\hyperref[#1: chapter]{Lecture Note~\ref*{#1: chapter}}}
  \pdfstringdefDisableCommands{\def\LN#1{Lecture Note \LNnumber{#1}}}

  % running head; "contents" links to the current chapter's table of contents
  \renewcommand{\doHeader}[1]{%
    \renewcommand{\header}{%
      \clap{\raisebox{0pt}{\protect\hyperlink{TOC-\arabic{chapter}}{$\uparrow$\scriptsize\,\textsc{contents}\hspace{1.1in}}}}%
      {\footnotesize \LNTitle{} / Section \thesection, #1}}%
    \markboth{\header}{\header}}

  % number theorems, lemmas, etc. and exercises within chapters (e.g. Lemma 8.1)
  \counterwithin{lemma}{chapter}
  \counterwithin{definition}{chapter}
  \counterwithin{theorem}{chapter}
  \counterwithin{corollary}{chapter}
  \counterwithin{claim}{chapter}
  \counterwithin{observation}{chapter}
  \counterwithin{conjecture}{chapter}
  \counterwithin{exercise}{chapter}
\fi

% %%%%%%%%%%%%%%%%%%%%%%%%%%%%%%%%%%%%%%%%%%%%%%%%%%%%%%%%%%%%%%%%%%%%%%%%
% %%%%%%%%%%%%%%%%%%%%%%%%%%%%%%%%%%%%%%%%%%%%%%%%%%%%%%%%%%%%%%%%%%%%%%%%

\ExplSyntaxOn
\seq_new:N \l_itemize_seq
% \NewDocumentEnvironment{REITEMIZE}{ +b }{
\NewDocumentCommand{\REITEMIZECMD}{ +m }{
  \regex_split:nnN {\c{item}} {#1} \l_itemize_seq
  \seq_pop:NN \l_itemize_seq \l_tmpa_tl% remove the (hopefully) empty first item
  \seq_if_empty:NF \l_itemize_seq {
    \l_tmpa_tl
    \seq_map_inline:Nn \l_itemize_seq {
      \regex_split:nnN {\c{item}} {#1} \l_itemize_seq {
        \regex_extract_once:nnNTF { ^\[([^]]*)\](.*) } { ##1} \l_foo_seq
      { \item[\seq_item:Nn \l_foo_seq {2}]\ITEMIZER{\seq_item:Nn \l_foo_seq {3}} } { \item\ITEMIZER{##1}  }
    }
  }
}
}
\ExplSyntaxOff

\newlength{\ITEMHT}
\setlength{\ITEMHT}{1in}

% \setlength{\fboxsep}{1pt}
\newcommand{\ITEMIZER}[1]{%
  % \framebox{
\adjustbox{valign=t}{\parbox{\textwidth}{\linespread{1.15}\selectfont #1}}~~\adjustbox{valign=t}{\parbox{0pt}{\color{gray}\rule[-\ITEMHT]{0.01pt}{\ITEMHT}}}
% }
  % \begin{tabular}{@{}p{\textwidth}@{\rule[-\ITEMHT]{0pt}{\ITEMHT}}}
  %   #1
  % \end{tabular}
}

\newsavebox\thesmashminipage
\newlength{\thesmashminipagelength}
\newlength{\answerboxlength}
\newlength{\answerboxoverrun}
\newenvironment{smashminipage}[2]
{\setlength\answerboxlength{#2}
\begin{lrbox}{\thesmashminipage}\begin{minipage}[t]{#1}}
{\end{minipage}\end{lrbox}%
\setlength{\thesmashminipagelength}{\dimexpr\ht\thesmashminipage+\dp\thesmashminipage}%
\setlength{\answerboxoverrun}{\dimexpr\thesmashminipagelength-\answerboxlength}%
\ifdim\the\thesmashminipagelength>\the\answerboxlength%
\usebox{\thesmashminipage}%
\smash{\color{red}\hspace*{-0.98\textwidth}\rule[-\answerboxlength]{0.98\textwidth}{1pt}}
\errmessage{Answer exceeds allotted length, please shorten it.}
\else
\usebox{\thesmashminipage}%
\fi
}



\newenvironment{answerBox}[1]{%
  \def\ANSWERBOXHEIGHT{#1}
  \begin{adjustbox}{valign=t}
  \begin{smashminipage}{\textwidth}{\ANSWERBOXHEIGHT}
}{\end{smashminipage}\end{adjustbox}~~~\adjustbox{valign=t}{\parbox{0pt}{\color{gray}\rule[-\ANSWERBOXHEIGHT]{0.01pt}{\ANSWERBOXHEIGHT}\rule[-\ANSWERBOXHEIGHT]{10pt}{0.01pt}}}}

% % \setlist[enumerate]{itemsep=0pt,parsep=0pt,topsep=0pt,leftmargin=100pt}

% \makeatletter 
\LetLtxMacro\SVenumerate\enumerate
\let\endSVenumerate\endenumerate

\newcommand{\injectEnumerate}{
  \let\SValph\alph
  \def\alph{}
  \renewenvironment{enumerate}[1][]
  {\let\alph\SValph \SVenumerate[label=(\alph*),itemsep=2pt,parsep=0pt]\REITEMIZE}{\endREITEMIZE\endSVenumerate}
  % \setlist[enumerate]{before=\Collect@Body\REITEMIZECMD}
}
% \makeatother
\newcommand{\restoreEnumerate}{
  \renewenvironment{enumerate}[1][]
  \let\SValph\alph
  \def\alph{}
  {\let\alph\SValph \SVenumerate[label=(\alph*)]}{\endSVenumerate}
  % \setlist[enumerate]{before=,after=}
}
\newcommand{\injectBlue}{
  \let\SVblue\blue
  \let\endSVblue\endblue
  \def\blue{\SVblue\REITEMIZE}
  \def\endblue{\endREITEMIZE\endSVblue}
}
\newcommand{\restoreBlue}{
  \let\blue\SVblue
  \let\endblue\endSVblue
}
\newcommand{\injectEnumerateBlue}{
  \setlist[enumerate]{before=\injectBlue,after=\restoreBlue,itemsep=2pt,parsep=0pt}
}
\newcommand{\restoreEnumerateBlue}{
  \setlist[enumerate]{before=,after=}
}


% %%%%%%%%%%%%%%%%%%%%%%%%%%%%%%%%%%%%%%%%%%%%%%%%%%%%%%%%%%%%%%%%%%%%%%%%
% %%%%%%%%%%%%%%%%%%%%%%%%%%%%%%%%%%%%%%%%%%%%%%%%%%%%%%%%%%%%%%%%%%%%%%%%

% \NewDocumentEnvironment{env}{ }
% {\Environenv}
% {\endEnvironenv}

% \NewEnviron{REITEMIZE}{\REITEMIZECMD\BODY}

\makeatletter 
\newenvironment{REITEMIZE}{\Collect@Body\REITEMIZECMD}{}
\makeatother 

\newcommand{\NODEone}[1]{\ensuremath{v_{#1}}\xspace}
\newcommand{\NODE}[2]{{\NODEone {#1,#2}}\xspace}
\newcommand{\EDGE}[2]{\ensuremath{(\textsf{#1}, \textsf{#2})}}
\newcommand{\ITEM}{\item}

%%% Local Variables:
%%% mode: latex
%%% End:
 to combine the
% lecture notes into one report-class document, one chapter per note.
% Without \BookMode (each lecture note on its own, and the homeworks) none
% of this applies.
\ifdefined\BookMode
  \usepackage{etoc}

  % links to other lecture notes stay in the book
  \renewcommand{\noteLink}[3]{\hyperlink{#2}{#3}}

  % each chapter is a lecture note, and is called one, as in the standalone notes
  \renewcommand{\chaptername}{Lecture Note}
  \setcounter{secnumdepth}{3}   % number subsubsections, as in the standalone notes

  % global table of contents, for main.tex
  \newcommand{\bookContents}{%
    \hypertarget{TOC}{}%
    \etocsetnexttocdepth{section}%
    \tableofcontents}

  % \doTitle{N}{title}{short title} starts chapter N, with a link to the global
  % table of contents and the chapter's own table of contents
  \renewcommand{\doTitle}[3]{%
    \setcounter{chapter}{\numexpr#1-1\relax}%
    \chapter{\texorpdfstring{#2}{#3}}%
    \label{\LNtag{#1}: chapter}%
    \csgdef{LNshorttitle@#1}{#3}%
    \csgdef{LNtitle@#1}{#2}%
    \renewcommand{\LNTitle}{Lecture Note #1, #3}%
    \codeListingstrue
    \hyperlink{TOC}{$\uparrow$\scriptsize\,\textsc{contents}}%
    {\small
      \hypertarget{TOC-#1}{}%
      \etocsettocstyle{\section*{\contentsname}}{}%
      \etocsetnexttocdepth{subsubsection}%
      \localtableofcontents
    }}

  % the code: each lecture note's \codeSection just records its files, and
  % \codeAppendix (in the book's appendix) lists them all, a section per note
  \makeatletter   % (file names have no spaces, so drop the spaces around them)
  \renewcommand{\codeSection}[1]{\csxdef{LNcode@\arabic{chapter}}{\zap@space#1 \@empty}}
  \makeatother
  % \codeAppendix: first a contents list (each note's topic, then its files,
  % linked to the listings), then a section per note listing its files
  \makeatletter
  \def\LNbasename#1{\LN@base#1/\LN@end}   % a/b/c.py -> c.py
  \def\LN@base#1/#2\LN@end{\ifx\relax#2\relax#1\else\LN@base#2\LN@end\fi}
  \makeatother
  \newcommand{\codeContentsFile}[1]{%
    \edef\LNfilename{\LNbasename{#1}}%
    \item \pycode[\texttt{\expandafter\detokenize\expandafter{\LNfilename}}]{#1}}
  \newcommand{\codeAppendix}{%
    \begin{itemize}
    \foreach \n in {1,...,9} {%
      \ifcsdef{LNcode@\n}{%
        \item \hyperref[code: note \n]{\csuse{LNtitle@\n}} (\LN{\LNtag{\n}})
        \begin{itemize}
        \edef\LNcodelist{\csuse{LNcode@\n}}%
        \expandafter\forcsvlist\expandafter\codeContentsFile\expandafter{\LNcodelist}%
        \end{itemize}}{}}
    \end{itemize}
    \newpage
    \foreach \n in {1,...,9} {%
      \ifcsdef{LNcode@\n}{%
        \section[Lecture Note \n: \csuse{LNtitle@\n}]{%
          \texorpdfstring{\LN{\LNtag{\n}}: \csuse{LNtitle@\n}}{Lecture Note \n: \csuse{LNtitle@\n}}}%
        \label{code: note \n}%
        \edef\LNcodelist{\csuse{LNcode@\n}}%
        \expandafter\forcsvlist\expandafter\codefile\expandafter{\LNcodelist}%
        \newpage}{}}}

  % a lecture note is a chapter: \LN{proofs} prints a link "Lecture Note 2"
  \DeclareRobustCommand{\LN}[1]{\hyperref[#1: chapter]{Lecture Note~\ref*{#1: chapter}}}
  \pdfstringdefDisableCommands{\def\LN#1{Lecture Note \LNnumber{#1}}}

  % running head; "contents" links to the current chapter's table of contents
  \renewcommand{\doHeader}[1]{%
    \renewcommand{\header}{%
      \clap{\raisebox{0pt}{\protect\hyperlink{TOC-\arabic{chapter}}{$\uparrow$\scriptsize\,\textsc{contents}\hspace{1.1in}}}}%
      {\footnotesize \LNTitle{} / Section \thesection, #1}}%
    \markboth{\header}{\header}}

  % number theorems, lemmas, etc. and exercises within chapters (e.g. Lemma 8.1)
  \counterwithin{lemma}{chapter}
  \counterwithin{definition}{chapter}
  \counterwithin{theorem}{chapter}
  \counterwithin{corollary}{chapter}
  \counterwithin{claim}{chapter}
  \counterwithin{observation}{chapter}
  \counterwithin{conjecture}{chapter}
  \counterwithin{exercise}{chapter}
\fi

% %%%%%%%%%%%%%%%%%%%%%%%%%%%%%%%%%%%%%%%%%%%%%%%%%%%%%%%%%%%%%%%%%%%%%%%%
% %%%%%%%%%%%%%%%%%%%%%%%%%%%%%%%%%%%%%%%%%%%%%%%%%%%%%%%%%%%%%%%%%%%%%%%%

\ExplSyntaxOn
\seq_new:N \l_itemize_seq
% \NewDocumentEnvironment{REITEMIZE}{ +b }{
\NewDocumentCommand{\REITEMIZECMD}{ +m }{
  \regex_split:nnN {\c{item}} {#1} \l_itemize_seq
  \seq_pop:NN \l_itemize_seq \l_tmpa_tl% remove the (hopefully) empty first item
  \seq_if_empty:NF \l_itemize_seq {
    \l_tmpa_tl
    \seq_map_inline:Nn \l_itemize_seq {
      \regex_split:nnN {\c{item}} {#1} \l_itemize_seq {
        \regex_extract_once:nnNTF { ^\[([^]]*)\](.*) } { ##1} \l_foo_seq
      { \item[\seq_item:Nn \l_foo_seq {2}]\ITEMIZER{\seq_item:Nn \l_foo_seq {3}} } { \item\ITEMIZER{##1}  }
    }
  }
}
}
\ExplSyntaxOff

\newlength{\ITEMHT}
\setlength{\ITEMHT}{1in}

% \setlength{\fboxsep}{1pt}
\newcommand{\ITEMIZER}[1]{%
  % \framebox{
\adjustbox{valign=t}{\parbox{\textwidth}{\linespread{1.15}\selectfont #1}}~~\adjustbox{valign=t}{\parbox{0pt}{\color{gray}\rule[-\ITEMHT]{0.01pt}{\ITEMHT}}}
% }
  % \begin{tabular}{@{}p{\textwidth}@{\rule[-\ITEMHT]{0pt}{\ITEMHT}}}
  %   #1
  % \end{tabular}
}

\newsavebox\thesmashminipage
\newlength{\thesmashminipagelength}
\newlength{\answerboxlength}
\newlength{\answerboxoverrun}
\newenvironment{smashminipage}[2]
{\setlength\answerboxlength{#2}
\begin{lrbox}{\thesmashminipage}\begin{minipage}[t]{#1}}
{\end{minipage}\end{lrbox}%
\setlength{\thesmashminipagelength}{\dimexpr\ht\thesmashminipage+\dp\thesmashminipage}%
\setlength{\answerboxoverrun}{\dimexpr\thesmashminipagelength-\answerboxlength}%
\ifdim\the\thesmashminipagelength>\the\answerboxlength%
\usebox{\thesmashminipage}%
\smash{\color{red}\hspace*{-0.98\textwidth}\rule[-\answerboxlength]{0.98\textwidth}{1pt}}
\errmessage{Answer exceeds allotted length, please shorten it.}
\else
\usebox{\thesmashminipage}%
\fi
}



\newenvironment{answerBox}[1]{%
  \def\ANSWERBOXHEIGHT{#1}
  \begin{adjustbox}{valign=t}
  \begin{smashminipage}{\textwidth}{\ANSWERBOXHEIGHT}
}{\end{smashminipage}\end{adjustbox}~~~\adjustbox{valign=t}{\parbox{0pt}{\color{gray}\rule[-\ANSWERBOXHEIGHT]{0.01pt}{\ANSWERBOXHEIGHT}\rule[-\ANSWERBOXHEIGHT]{10pt}{0.01pt}}}}

% % \setlist[enumerate]{itemsep=0pt,parsep=0pt,topsep=0pt,leftmargin=100pt}

% \makeatletter 
\LetLtxMacro\SVenumerate\enumerate
\let\endSVenumerate\endenumerate

\newcommand{\injectEnumerate}{
  \let\SValph\alph
  \def\alph{}
  \renewenvironment{enumerate}[1][]
  {\let\alph\SValph \SVenumerate[label=(\alph*),itemsep=2pt,parsep=0pt]\REITEMIZE}{\endREITEMIZE\endSVenumerate}
  % \setlist[enumerate]{before=\Collect@Body\REITEMIZECMD}
}
% \makeatother
\newcommand{\restoreEnumerate}{
  \renewenvironment{enumerate}[1][]
  \let\SValph\alph
  \def\alph{}
  {\let\alph\SValph \SVenumerate[label=(\alph*)]}{\endSVenumerate}
  % \setlist[enumerate]{before=,after=}
}
\newcommand{\injectBlue}{
  \let\SVblue\blue
  \let\endSVblue\endblue
  \def\blue{\SVblue\REITEMIZE}
  \def\endblue{\endREITEMIZE\endSVblue}
}
\newcommand{\restoreBlue}{
  \let\blue\SVblue
  \let\endblue\endSVblue
}
\newcommand{\injectEnumerateBlue}{
  \setlist[enumerate]{before=\injectBlue,after=\restoreBlue,itemsep=2pt,parsep=0pt}
}
\newcommand{\restoreEnumerateBlue}{
  \setlist[enumerate]{before=,after=}
}


% %%%%%%%%%%%%%%%%%%%%%%%%%%%%%%%%%%%%%%%%%%%%%%%%%%%%%%%%%%%%%%%%%%%%%%%%
% %%%%%%%%%%%%%%%%%%%%%%%%%%%%%%%%%%%%%%%%%%%%%%%%%%%%%%%%%%%%%%%%%%%%%%%%

% \NewDocumentEnvironment{env}{ }
% {\Environenv}
% {\endEnvironenv}

% \NewEnviron{REITEMIZE}{\REITEMIZECMD\BODY}

\makeatletter 
\newenvironment{REITEMIZE}{\Collect@Body\REITEMIZECMD}{}
\makeatother 

\newcommand{\NODEone}[1]{\ensuremath{v_{#1}}\xspace}
\newcommand{\NODE}[2]{{\NODEone {#1,#2}}\xspace}
\newcommand{\EDGE}[2]{\ensuremath{(\textsf{#1}, \textsf{#2})}}
\newcommand{\ITEM}{\item}

%%% Local Variables:
%%% mode: latex
%%% End:
.  A homework
% folder has a link refs/ to src/refs/, which holds each note's label records
% (made by tools/refs.sh from the notes' builds).  Read with xr-hyper, they make
% \ref{greedy: exercise: ...} give the exercise's number in the published note,
% linked to it there; labels only the book has (its appendices) link to the book.
\ifdefined\BookMode\else
  \makeatletter
  \newcommand{\lectureNoteRefs}{%
    \@for\LN@tag:=prelim,proofs,matching,dc,greedy,graphs,dp,flow,np\do{%
      \externaldocument{refs/\LNfile{\LN@tag}}[\notePDF{\LN@tag}]}%
    \externaldocument{refs/book}[\siteURL lecture_notes_on_algorithms.pdf]}
  % hyperref links a \ref to another document's label as URL#anchor; browsers'
  % PDF viewers need URL#nameddest=anchor (the form \noteLink and \bookLink use)
  \def\Hy@setref@link#1#2#3#4#5#6\@nil#7{%
    \begingroup
      \ifx\\#5\\\toks0={\hyper@@link{#5}{#4}}\else\toks0={\hyper@@link{#5}{nameddest=#4}}\fi
      \toks1=\expandafter{#7{#1}{#2}{#3}{#4}{#5}}%
      \edef\x{\endgroup\the\toks0 {\the\toks1 }}%
    \x}
  \makeatother
\fi

% Book mode: main.tex defines \BookMode before %%%% DON'T EDIT THIS FILE %%%%

\usepackage{xifthen}% \isempty
\usepackage{etoolbox}% \patchcmd
% \usepackage{stringstrings}% \substring

\usepackage{amsmath, amsthm, amssymb}
\usepackage[margin=1in]{geometry}
\usepackage{xspace}
\usepackage{graphicx}
\usepackage[normalem]{ulem}
\usepackage{adjustbox}
\usepackage{xcolor}
\usepackage{changepage}
\usepackage{framed}
\usepackage{multicol}
\usepackage{enumitem}
% \usepackage{listings} % lstlistings environment

\usepackage{makecmds}           %\provideenvironment

\usepackage{xparse,environ}
\usepackage{letltxmacro}       % \LetLtxMacro
\usepackage{wrapfig}
\usepackage{tikz}

\usetikzlibrary{shapes}
\usetikzlibrary{positioning}
\usetikzlibrary{shapes.misc}
\usetikzlibrary{arrows}

\newtheorem{lemma}{Lemma}
\newtheorem{definition}{Definition}
\newtheorem{theorem}{Theorem}
\newtheorem{corollary}{Corollary}
\newtheorem{claim}{Claim}
\newtheorem{observation}{Observation}
\newtheorem{conjecture}{Conjecture}

% A lemma or theorem restated later with its original number (set with
% \setcounter) goes in a restatement environment, which gives it its own link
% anchor (otherwise it would duplicate the original's).
\newenvironment{restatement}
  {\def\theHlemma{restated.\thelemma}\def\theHtheorem{restated.\thetheorem}}{}

\newenvironment{myframed}
{\smallskip\par\noindent\begin{minipage}{\textwidth}\begin{framed}\vspace*{-5pt}}{\vspace*{-5pt}\end{framed}\end{minipage}\smallskip}


\definecolor{darkblue}{RGB}{0,0,200}

 
\let\oldParagraph\paragraph
\providecommand{\gradescopeColor}{\color{black}}
\renewcommand{\paragraph}[1]{\oldParagraph{\gradescopeColor #1}}
\providecommand{\gradescopeInit}{}

\providecommand{\BLUE}[1]{{\color{darkblue} #1}}
\providecommand{\REPLACEME}{\textcolor{red}{\ensuremath{\star \cdots \star}}\xspace}
\provideenvironment{blue}{\color{darkblue}}{}

\newcommand{\dist}[2]{\ensuremath{\operatorname{\sf dist}(#1, #2)}}
\newcommand{\depth}[2]{\ensuremath{\operatorname{\sf depth}_{#1}(#2)}}
\newcommand{\cost}[1]{\ensuremath{\operatorname{\sf cost}(#1)}}
\newcommand{\weight}[1]{\ensuremath{\operatorname{\sf wt}(#1)}}
% \newcommand{\val}[1]{\ensuremath{\operatorname{\sf val}(#1)}}
\newcommand{\eps}{\varepsilon}
\newcommand{\capac}{\textsf{cap}}
\newcommand{\capOf}[1]{\ensuremath{\operatorname{\sf cap}(#1)}}
\newcommand{\rescap}[1]{\ensuremath{\operatorname{\sf cap}_f(#1)}}
\newcommand{\into}[1]{\ensuremath{\operatorname{\sf in}(#1)}}
\newcommand{\out}[1]{\ensuremath{\operatorname{\sf out}(#1)}}
\newenvironment{LP}[1]{
  \begin{array}[t]{@{}l@{}}
    \displaystyle #1 \\
    \begin{aligned}\\[-10pt]}{\end{aligned}\end{array}}

\newcommand{\opt}[1]{\ensuremath{\operatorname{\sf opt}(#1)}}
\newcommand{\Opt}[1]{\ensuremath{\operatorname{\sf Opt}(#1)}}
\newcommand{\half}{\frac{1}{2}}

\newcommand{\E}{\operatorname{E}}
\newcommand{\giv}{\,|\,}
\newcommand{\Z}{\mathbb{Z}}
\newcommand{\Zp}{\mathbb{Z}_{\ge 0}}
\newcommand{\R}{\mathbb{R}}
% \newcommand{\Rp}{\mathbb{R}_{\ge 0}}
\newcommand{\Rp}{\mathbb{R}_{\tiny +}}
\newcommand{\N}{\mathbb{N}}
\newcommand{\Np}{\mathbb{N}_{\ge 0}}
\newcommand{\calS}{{\cal S}}

\newcommand{\OP}[1]{\text{\small\sc #1}\xspace}
\newcommand{\Insert}{\OP{Insert}}
\newcommand{\FindMin}{\OP{Find-Min}}
\newcommand{\DeleteMin}{\OP{Delete-Min}}
\newcommand{\DecreaseKey}{\OP{Decrease}}
\newcommand{\Merge}{\OP{Merge}}
\newcommand{\Cut}{\OP{Cut}}
% \newcommand{\prob}[1]{\text{\sc #1}}
% a problem name; in running text, lines may break inside it
\DeclareRobustCommand{\prob}[1]{\ifmmode\text{#1}\else#1\fi}
\newcommand{\encode}[1]{\langle #1 \rangle}

\newcommand{\best}[1]{\text{\sf best}(#1)}
\newcommand{\worst}[1]{\text{\sf worst}(#1)}
\newcommand{\floor}[1]{\lfloor #1 \rfloor}
\newcommand{\ceil}[1]{\lceil #1 \rceil}

\setlength{\parskip}{2pt}
\setlength{\parindent}{1em}

% \setlength\fboxsep{0pt}

%%% longFormProof, shortFormProof, block environments
%%% \step

\newlist{steps}{enumerate}{7}
\setlist[steps]{wide,
  topsep=1pt,
  parsep=1pt,
  partopsep=1pt,
  itemsep=0pt,
  align=left,
  labelindent=0pt,
  itemindent=0pt, 
}
\setlist[steps,1]{
  label=\arabic*.,
  ref=\arabic*
}
\setlist[steps,2]{
  label=\arabic{stepsi}.\arabic*.,
  ref=\arabic{stepsi}.\arabic*
}
\setlist[steps,3]{
  label=\arabic{stepsi}.\arabic{stepsii}.\arabic*.,
  ref=\arabic{stepsi}.\arabic{stepsii}.\arabic*
}
\setlist[steps,4]{
  label=\arabic{stepsi}.\arabic{stepsii}.\arabic{stepsiii}.\arabic*.,
  ref=\arabic{stepsi}.\arabic{stepsii}.\arabic{stepsiii}.\arabic*
}
\setlist[steps,5]{
  label=\arabic{stepsi}.\arabic{stepsii}.\arabic{stepsiii}.\arabic{stepsiv}.\arabic*.,
  ref=\arabic{stepsi}.\arabic{stepsii}.\arabic{stepsiii}.\arabic{stepsiv}.\arabic*
}
\setlist[steps,6]{
  label=\arabic{stepsi}.\arabic{stepsii}.\arabic{stepsiii}.\arabic{stepsiv}.\arabic{stepsv}.\arabic*.,
  ref=\arabic{stepsi}.\arabic{stepsii}.\arabic{stepsiii}.\arabic{stepsiv}.\arabic{stepsv}.\arabic*
}
\setlist[steps,7]{
  label=\arabic{stepsi}.\arabic{stepsii}.\arabic{stepsiii}.\arabic{stepsiv}. \arabic{stepsv}. \arabic{stepsvi}.\arabic*.,
  ref=\arabic{stepsi}.\arabic{stepsii}.\arabic{stepsiii}.\arabic{stepsiv}.\arabic{stepsv}. \arabic{stepsvi}.\arabic*
}

\newcommand{\step}[1][]{\item\maybeLabel{#1}}

\makeatletter
% \skipitem steps the list counter itself (for block labels).  Give each one
% its own hyperref link target, since step numbers repeat from proof to proof.
\newcounter{skipitem}
\newcommand{\skipitem}{\stepcounter{skipitem}\refstepcounter{\@enumctr}}
\renewcommand{\theHstepsi}{skipitem.\arabic{skipitem}}
\renewcommand{\theHstepsii}{skipitem.\arabic{skipitem}}
\renewcommand{\theHstepsiii}{skipitem.\arabic{skipitem}}
\renewcommand{\theHstepsiv}{skipitem.\arabic{skipitem}}
\renewcommand{\theHstepsv}{skipitem.\arabic{skipitem}}
\renewcommand{\theHstepsvi}{skipitem.\arabic{skipitem}}
\renewcommand{\theHstepsvii}{skipitem.\arabic{skipitem}}
\makeatother

\newcommand{\maybeLabel}[1]{\ifthenelse{\isempty{#1}}{}{\label{#1}}}

\newenvironment{longFormProof}[1][Proof (long form).]
{\begin{proof}[#1]~\par\begin{steps}}{\end{steps}\vspace*{-3ex}\end{proof}}

\newenvironment{shortFormProof}[1][Proof (short form).]
{\begin{proof}[#1]}{\end{proof}}

\newenvironment{longFormSubProof}[1][Proof.]
{\begin{adjustwidth}{1em}{0em}\begin{proof}[#1]~\par\vspace*{-2pt}\skipitem\begin{steps}}
{\end{steps}\vspace*{-1ex}\end{proof}\end{adjustwidth}\smallskip}

\newenvironment{shortFormSubProof}[1][Proof.]
{\begin{proof}[#1]}
{\end{proof}\vspace*{-3pt}}

\newenvironment{block}[1][]
{\skipitem\maybeLabel{#1}\item[] \begin{steps}}{\end{steps}}

\newcommand{\comment}[1]{\hfill{\footnotesize\emph{#1}}}
\newcommand{\lineacross}{\par\vspace*{-0.7\baselineskip}\noindent\hrulefill\par}

%%% problem

% \newcounter{problem}
% \newcommand{\problem}[1][]{}
% {\ifthenelse{\isempty{#1}}{\refstepcounter{problem}}{}%
%   \paragraph{Problem \ifthenelse{\isempty{#1}}
%     {\bf\arabic{problem}.}
%     {#1}} }

\newcommand{\problemNumbered}[1]{
  \setcounter{problem}{#1}%
  \addtocounter{problem}{-1}%
  \problem
}
\newenvironment{problemTheorem}{%
  \setcounter{theorem}{\value{problem}}%
  \addtocounter{theorem}{-1}%
  \begin{theorem}[Problem \arabic{problem}]
  }{\end{theorem}}

\newcounter{problem}
\newcommand{\problem}[1][]
{\ifthenelse{\isempty{#1}}
  {\refstepcounter{problem}\paragraph{Problem~\arabic{problem}.}}
  {\paragraph{Problem~#1.}}\refreshHeader} 

\newcounter{exercise}
\newcommand{\exercise}[1][]
{\ifthenelse{\isempty{#1}}
  {\refstepcounter{exercise}\paragraph{Exercise~\theexercise.}}
  {\paragraph{Exercise~#1.}}}

\newcommand{\solution}{\paragraph{Solution.}} 

\newenvironment{mylist}[1][]
{\begin{list}{#1}{
      \setlength{\itemsep}{-1pt}
      \setlength{\labelwidth}{0pt}
      \setlength{\leftmargin}{1.25em}
    }}
{\end{list}}

\newcommand{\LNTitle}{\textcolor{red}{UNDEFINED LNTitle}}

% In a standalone lecture note's table of contents, leave room for numbers
% like 7.10 (in bold) and 7.10.2.
\makeatletter
\newcommand{\LNtocWidths}{%
  \patchcmd{\l@section}{1.5em}{2.8em}{}{\PackageWarning{macros}{Could not widen section numbers in the contents}}%
  \renewcommand*{\l@subsection}{\@dottedtocline{2}{2.8em}{3.2em}}%
  \renewcommand*{\l@subsubsection}{\@dottedtocline{3}{6.0em}{4.1em}}}
\makeatother

\newcommand{\doTitle}[3]{
  \pagenumbering{roman}
  \renewcommand{\LNTitle}{Lecture Note #1, #3}
  \renewcommand{\currentLN}{#1}
  \codeListingstrue
  % number sections, theorems etc., exercises, figures, tables and equations
  % within the lecture note (e.g. Section 8.2, Lemma 8.6), as in the book
  \renewcommand{\thesection}{#1.\arabic{section}}
  \renewcommand{\thelemma}{#1.\arabic{lemma}}
  \renewcommand{\thedefinition}{#1.\arabic{definition}}
  \renewcommand{\thetheorem}{#1.\arabic{theorem}}
  \renewcommand{\thecorollary}{#1.\arabic{corollary}}
  \renewcommand{\theclaim}{#1.\arabic{claim}}
  \renewcommand{\theobservation}{#1.\arabic{observation}}
  \renewcommand{\theconjecture}{#1.\arabic{conjecture}}
  \renewcommand{\theexercise}{#1.\arabic{exercise}}
  \renewcommand{\thefigure}{#1.\arabic{figure}}
  \renewcommand{\thetable}{#1.\arabic{table}}
  \renewcommand{\theequation}{#1.\arabic{equation}}
  \LNtocWidths
  \title{Lecture Note #1\\ #2}
  % a note's \LNabstract, if it defines one before \doTitle, goes right below the title
  \date{\ifdefined\LNabstract\parbox{0.85\textwidth}{\normalsize\normalfont\LNabstract}\\[3ex]\fi
    draft of \today, by \href{\authorURL}{N.~Young}\\[2ex] {\small\copyrightNotice}}
  \maketitle
  {\small
    \hypertarget{TOC}{}
    \tableofcontents
  }
  \thispagestyle{empty}
  % the text starts on page 1, after the title and contents (numbered i, ii, ...)
  \clearpage
  \pagenumbering{arabic}
}

\def\clap#1{\hbox to 0pt{\hss#1\hss}}

\newcommand{\refreshHeader}{}
\newcommand{\header}{}

\newcommand{\doHeader}[1]{      % for lecture notes
  \renewcommand{\header}{%
    \clap{\raisebox{0pt}{\protect\hyperlink{TOC}{$\uparrow$\scriptsize\,\textsc{contents}\hspace{1.1in}}}}%
    {\footnotesize \LNTitle{} / Section \thesection, #1}} %, #1, draft of \today}}
  \markboth{\header}{\header}
}

\newcommand{\SID}{}
\newcommand{\setAuthor}[1]{\renewcommand{\author}{#1}}
\newcommand{\setSID}[1]{\renewcommand{\SID}{#1}}

\newcommand{\doHwHeader}[1]{                                         % for homeworks
  \renewcommand{\refreshHeader}{
    \renewcommand{\header}{%
      {\footnotesize \classSection #1 / Problem~\arabic{problem} / draft of \today \hfill \author\ifGradescope~(\SID)\fi~~}}
    \markboth{\header}{\header}
  }
  \renewcommand{\header}{%
    {\footnotesize \classSection #1 / Instructions / draft of \today \hfill \author\ifGradescope~(\SID)\fi~~}}
  \markboth{\header}{\header}
}

% Homework settings, for instructors using the homeworks in a course:
% \classSection is the course name shown in each homework's running head
% (e.g. \renewcommand{\classSection}{CS 3000\xspace}), and \Gradescopetrue turns on
% the instructions for submitting to Gradescope (in its fixed-length mode),
% the hints saying on which page of the PDF each answer should be, and the
% student ID in the running head.
\newcommand{\classSection}{}
\newif\ifGradescope
\Gradescopefalse
\newcommand{\GradescopeComment}[1]{\ifGradescope\comment{#1}\fi}

\pagestyle{myheadings}
\addtolength{\headsep}{0.2in}
\addtolength{\topmargin}{-0.3in}
\addtolength{\textheight}{0.4in}

%%%%%%%%%%%%%%%%%%%%%%%%%%%%%%%%%%%%%%%%%%%%%%%%
%%%%%%%%%%%%%%%%%%%%%%%%%%%%%%%%%%%%%%%%%%%%%%%%

\newcommand{\MF}{\href
  {https://mfleck.cs.illinois.edu/building-blocks/}
  {MF}\xspace}

\newcommand{\KWSlides}{\href
  {https://www.cs.princeton.edu/~wayne/kleinberg-tardos/}
  {KW slides}\xspace}

\newcommand{\LLM}{\href
  {https://courses.csail.mit.edu/6.042/spring18/mcs.pdf}
  {LLM}\xspace}

% \DPV, \KT, \CLRS: textbooks listed (without links) in the table of sources in
% the Preliminaries note; each prints as a link to its entry in that table.
\DeclareRobustCommand{\DPV}{\noteLink{prelim}{source.DPV}{DPV}\xspace}
\DeclareRobustCommand{\KT}{\noteLink{prelim}{source.KT}{KT}\xspace}
\DeclareRobustCommand{\CLRS}{\noteLink{prelim}{source.CLRS}{CLRS}\xspace}
\AtBeginDocument{\pdfstringdefDisableCommands{\def\DPV{DPV }\def\KT{KT }\def\CLRS{CLRS }}}

\newcommand{\JE}{\href
  {https://jeffe.cs.illinois.edu/teaching/algorithms/}
  {JE}\xspace}

\newcommand{\JEGreedy}{\href
  {https://jeffe.cs.illinois.edu/teaching/algorithms/book/04-greedy.pdf}
  {JE Chapter 4}\xspace}

\newcommand{\BBCode}{\href
{https://northeastern.blackboard.com/webapps/cmsmain/webui/courses/CS3000.MERGED.202010/code?action=frameset&subaction=view&uniq=-jj2max&course_id=_2597392_1&mask=\%2Fcourses\%2FCS3000.MERGED.202010}
  {code}\xspace}

% The Python code (src/code).  Each lecture note that has code ends with a
% section listing its files: \codeSection{path, path, ...}.  In the text,
% \pycode{greedy_algorithms/Dijkstra.py} prints the path as a link to its listing
% (in the homeworks, which have no listings, to the file on GitHub), and
% \pycode[text]{...} uses the given text instead.  \pythonCode{path, ...} is the
% standard sentence ending a section whose algorithms have code:
% "Python code: path, ... (Section N.k)."
\newcommand{\codeFolder}{https://github.com/nealeyoung/algs101/tree/main/src/code}
\newcommand{\codeRepo}{https://github.com/nealeyoung/algs101/blob/main/src/code}
\newif\ifcodeListings   % set by \doTitle: this document has the listings
\newcommand{\githubCode}[2][]{\href{\codeRepo/#2}{\ifx\relax#1\relax\nolinkurl{#2}\else#1\fi}}
\newcommand{\pycode}[2][]{%
  \ifcodeListings\hyperref[code: #2]{\ifx\relax#1\relax\nolinkurl{#2}\else#1\fi}%
  \else\githubCode[#1]{#2}\fi}
\makeatletter
\newcommand{\pythonCode}[1]{%
  \def\pythonCode@first{}\def\pythonCode@list{}%
  \forcsvlist{\pythonCode@add}{#1}%
  Python code: \pythonCode@list\ (Section~\ref{code: \pythonCode@first}).}
\newcommand{\pythonCode@add}[1]{%
  \ifx\pythonCode@first\@empty\def\pythonCode@first{#1}\appto\pythonCode@list{\pycode{#1}}%
  \else\appto\pythonCode@list{, \pycode{#1}}\fi}
\makeatother

% Code listings that copy and paste exactly (in PDF viewers that honor the
% ToUnicode maps, e.g. MuPDF-based ones; PDFKit viewers such as Preview and
% Skim drop leading spaces, so each listing's heading links to the file on
% GitHub).  Spaces are typeset as white visible-space glyphs that copy as
% spaces, and hyphens and asterisks as the T1 glyphs, so they copy as themselves.
\usepackage[T1,OT1]{fontenc}   % T1 for the code listings only; OT1 stays the default
\usepackage{listings}
\pdfgentounicode=1
\input glyphtounicode
\pdfglyphtounicode{uni2423}{0020}
\makeatletter
\def\lst@visiblespace{\textcolor{white}{\char32}}
% put an invisible space on each blank line, so that blank lines copy too
% (listings' OnEmptyLine hook runs before the line break, so patch the loop)
\def\lst@DoNewLines{%
    \@whilenum\lst@newlines>\@ne \do {\lst@NewLine \lst@visiblespace}%
    \ifnum\lst@newlines>\z@ \lst@NewLine \fi}
\makeatother
\lstset{language=Python,
  basicstyle=\fontencoding{T1}\fontfamily{lmtt}\selectfont\footnotesize,
  upquote=true, showstringspaces=false, columns=fullflexible, keepspaces=true,
  showspaces=true, literate={-}{{\char45}}1 {*}{{\char42}}1,
  frame=single, framesep=1ex, rulecolor=\color{gray}}
% the listings' paths are relative to the lecture note's folder (in the book, to src/)
\ifdefined\BookMode\newcommand{\codeDir}{code/}\else\newcommand{\codeDir}{../code/}\fi
% \codefile{path}: a listing of src/code/path, headed by a link to it on GitHub
\newcommand{\codefile}[1]{%
  \subsection*{\githubCode{#1}}\label{code: #1}%
  \lstinputlisting{\codeDir#1}}
% \codeSection{path, path, ...}: the lecture note's last section, listing its code
\newcommand{\codeSection}[1]{%
  \newpage
  \section{Python code}\label{\LNtag{\currentLN}: sec: code}\doHeader{Python code}
  This section lists the Python code for this lecture note.
  Click on a file's name to open the file on GitHub.
  Copying code from GitHub is more reliable than copying it from this PDF,
  since some PDF viewers drop the leading spaces (the indentation) of copied lines.
  \forcsvlist{\codefile}{#1}}
% the copyright notice, for title pages and the homework first pages
\newcommand{\licenseURL}{https://github.com/nealeyoung/algs101/blob/main/LICENSE}
\newcommand{\authorURL}{https://www.cs.ucr.edu/\string~neal}
\newcommand{\copyrightNotice}{%
  \copyright\ 2018--2026 \href{\authorURL}{Neal E. Young}.
  Text: CC BY-NC-SA 4.0; code: MIT; except where noted.
  See \url{\licenseURL}.}
% the published site: PDFs of the notes, and the homework assignments and templates
\newcommand{\siteURL}{https://nealeyoung.github.io/algs101/}

\newcommand{\blackboard}{\href
  {https://northeastern.blackboard.com/webapps/blackboard/execute/launcher?type=Course&id=_2597392_1&url=}
  {Blackboard}\xspace}

\newcommand{\JUSlides}{\href
  {http://www.khoury.neu.edu/home/jullman/cs3000f18/schedule.html}
  {Prof.~Ullman's slides}\xspace}

% a URL on its own line, in three layouts
\newcommand{\URL}[1]{\par-- \parbox[t]{0.8\textwidth}{\url{#1}} \medskip\par}
\newcommand{\URLsmall}[1]{\par-- \parbox[t]{0.95\textwidth}{\small\url{#1}} \medskip\par}
\newcommand{\URLplain}[1]{\par\parbox[t]{\textwidth}{\url{#1}}\par}

%%%%%%%%%%%%%%%%%%%%%%%%%%%%%%%%%%%%%%%%%%%%%%%%
%%%%%%%%%%%%%%%%%%%%%%%%%%%%%%%%%%%%%%%%%%%%%%%%

\newcommand{\continued}{\vfill\centerline{\emph{(continued)}}\vfill\newpage}

% load hyperref last to avoid conflicts; long URLs may also break after hyphens
\PassOptionsToPackage{hyphens}{url}
\ifdefined\BookMode\else\usepackage{xr-hyper}\fi   % for \lectureNoteRefs (homeworks), below
\usepackage{hyperref}
\hypersetup{
    colorlinks,
    linkcolor={blue!50!black},
    citecolor={blue!50!black},
    urlcolor={blue!80!black}
}

\providecommand{\answerDirectory}{.}
\newcommand{\inputAnswer}[1]{\refreshHeader\input{\answerDirectory/#1}}
\newcommand{\inputSection}[1]{\input{#1}}

% Refer to a lecture note by its tag (the prefix of its labels), e.g.
% \LN{proofs} prints "Lecture Note 2" (in the book, a link to that chapter).  The table below is the only place the numbers appear.
\makeatletter
\newcommand{\defineLN}[3]{\@namedef{LNnumber@#2}{#1}\@namedef{LNtag@#1}{#2}\@namedef{LNfile@#2}{#3}}
\newcommand{\LNnumber}[1]{%
  \ifcsname LNnumber@#1\endcsname\csname LNnumber@#1\endcsname
  \else\PackageError{macros}{Unknown lecture-note tag `#1'}{}\fi}
\newcommand{\LNtag}[1]{\csname LNtag@#1\endcsname}
\newcommand{\LNfile}[1]{\csname LNfile@#1\endcsname}
\makeatother
\defineLN{1}{prelim}{1_preliminaries}
\defineLN{2}{proofs}{2_proofs}
\defineLN{3}{matching}{3_stable_matching}
\defineLN{4}{dc}{4_divide_and_conquer}
\defineLN{5}{greedy}{5_greedy}
\defineLN{6}{graphs}{6_graph_traversal}
\defineLN{7}{dp}{7_dynamic_programming}
\defineLN{8}{flow}{8_flow_reductions_LP}
\defineLN{9}{np}{9_NP}

% The text is the same in the book and in the standalone lecture notes; only
% links differ.  In the book, links to other lecture notes stay in the book.
% In a standalone note, they open the other note's published PDF.
% \noteLink{tag}{target}{text} links text to \hypertarget{target} in lecture note tag.
% \LN{proofs} prints "Lecture Note 2", linked to that lecture note.
\newcommand{\currentLN}{0}   % in a standalone note, its number (set by \doTitle)
\newcommand{\notePDF}[1]{\siteURL notes/\LNfile{#1}.pdf}
\newcommand{\noteLink}[3]{%
  \ifnum\LNnumber{#1}=\currentLN\relax\hyperlink{#2}{#3}%
  \else\href{\notePDF{#1}\#nameddest=#2}{#3}\fi}
% \bookLink{target}{text} links text to \hypertarget{target} in the book's
% appendices (from a standalone note, in the published book).
\DeclareRobustCommand{\bookLink}[2]{%
  \ifdefined\BookMode\hyperlink{#1}{#2}%
  \else\href{\siteURL lecture_notes_on_algorithms.pdf\#nameddest=#1}{#2}\fi}
\DeclareRobustCommand{\LN}[1]{%
  \ifnum\LNnumber{#1}=\currentLN\relax Lecture Note~\LNnumber{#1}%
  \else\href{\notePDF{#1}}{Lecture Note~\LNnumber{#1}}\fi}
\pdfstringdefDisableCommands{\def\LN#1{Lecture Note \LNnumber{#1}}}

% Homeworks refer to the notes' sections, exercises and lemmas by \ref: each
% homework's main file says \lectureNoteRefs after %%%% DON'T EDIT THIS FILE %%%%

\usepackage{xifthen}% \isempty
\usepackage{etoolbox}% \patchcmd
% \usepackage{stringstrings}% \substring

\usepackage{amsmath, amsthm, amssymb}
\usepackage[margin=1in]{geometry}
\usepackage{xspace}
\usepackage{graphicx}
\usepackage[normalem]{ulem}
\usepackage{adjustbox}
\usepackage{xcolor}
\usepackage{changepage}
\usepackage{framed}
\usepackage{multicol}
\usepackage{enumitem}
% \usepackage{listings} % lstlistings environment

\usepackage{makecmds}           %\provideenvironment

\usepackage{xparse,environ}
\usepackage{letltxmacro}       % \LetLtxMacro
\usepackage{wrapfig}
\usepackage{tikz}

\usetikzlibrary{shapes}
\usetikzlibrary{positioning}
\usetikzlibrary{shapes.misc}
\usetikzlibrary{arrows}

\newtheorem{lemma}{Lemma}
\newtheorem{definition}{Definition}
\newtheorem{theorem}{Theorem}
\newtheorem{corollary}{Corollary}
\newtheorem{claim}{Claim}
\newtheorem{observation}{Observation}
\newtheorem{conjecture}{Conjecture}

% A lemma or theorem restated later with its original number (set with
% \setcounter) goes in a restatement environment, which gives it its own link
% anchor (otherwise it would duplicate the original's).
\newenvironment{restatement}
  {\def\theHlemma{restated.\thelemma}\def\theHtheorem{restated.\thetheorem}}{}

\newenvironment{myframed}
{\smallskip\par\noindent\begin{minipage}{\textwidth}\begin{framed}\vspace*{-5pt}}{\vspace*{-5pt}\end{framed}\end{minipage}\smallskip}


\definecolor{darkblue}{RGB}{0,0,200}

 
\let\oldParagraph\paragraph
\providecommand{\gradescopeColor}{\color{black}}
\renewcommand{\paragraph}[1]{\oldParagraph{\gradescopeColor #1}}
\providecommand{\gradescopeInit}{}

\providecommand{\BLUE}[1]{{\color{darkblue} #1}}
\providecommand{\REPLACEME}{\textcolor{red}{\ensuremath{\star \cdots \star}}\xspace}
\provideenvironment{blue}{\color{darkblue}}{}

\newcommand{\dist}[2]{\ensuremath{\operatorname{\sf dist}(#1, #2)}}
\newcommand{\depth}[2]{\ensuremath{\operatorname{\sf depth}_{#1}(#2)}}
\newcommand{\cost}[1]{\ensuremath{\operatorname{\sf cost}(#1)}}
\newcommand{\weight}[1]{\ensuremath{\operatorname{\sf wt}(#1)}}
% \newcommand{\val}[1]{\ensuremath{\operatorname{\sf val}(#1)}}
\newcommand{\eps}{\varepsilon}
\newcommand{\capac}{\textsf{cap}}
\newcommand{\capOf}[1]{\ensuremath{\operatorname{\sf cap}(#1)}}
\newcommand{\rescap}[1]{\ensuremath{\operatorname{\sf cap}_f(#1)}}
\newcommand{\into}[1]{\ensuremath{\operatorname{\sf in}(#1)}}
\newcommand{\out}[1]{\ensuremath{\operatorname{\sf out}(#1)}}
\newenvironment{LP}[1]{
  \begin{array}[t]{@{}l@{}}
    \displaystyle #1 \\
    \begin{aligned}\\[-10pt]}{\end{aligned}\end{array}}

\newcommand{\opt}[1]{\ensuremath{\operatorname{\sf opt}(#1)}}
\newcommand{\Opt}[1]{\ensuremath{\operatorname{\sf Opt}(#1)}}
\newcommand{\half}{\frac{1}{2}}

\newcommand{\E}{\operatorname{E}}
\newcommand{\giv}{\,|\,}
\newcommand{\Z}{\mathbb{Z}}
\newcommand{\Zp}{\mathbb{Z}_{\ge 0}}
\newcommand{\R}{\mathbb{R}}
% \newcommand{\Rp}{\mathbb{R}_{\ge 0}}
\newcommand{\Rp}{\mathbb{R}_{\tiny +}}
\newcommand{\N}{\mathbb{N}}
\newcommand{\Np}{\mathbb{N}_{\ge 0}}
\newcommand{\calS}{{\cal S}}

\newcommand{\OP}[1]{\text{\small\sc #1}\xspace}
\newcommand{\Insert}{\OP{Insert}}
\newcommand{\FindMin}{\OP{Find-Min}}
\newcommand{\DeleteMin}{\OP{Delete-Min}}
\newcommand{\DecreaseKey}{\OP{Decrease}}
\newcommand{\Merge}{\OP{Merge}}
\newcommand{\Cut}{\OP{Cut}}
% \newcommand{\prob}[1]{\text{\sc #1}}
% a problem name; in running text, lines may break inside it
\DeclareRobustCommand{\prob}[1]{\ifmmode\text{#1}\else#1\fi}
\newcommand{\encode}[1]{\langle #1 \rangle}

\newcommand{\best}[1]{\text{\sf best}(#1)}
\newcommand{\worst}[1]{\text{\sf worst}(#1)}
\newcommand{\floor}[1]{\lfloor #1 \rfloor}
\newcommand{\ceil}[1]{\lceil #1 \rceil}

\setlength{\parskip}{2pt}
\setlength{\parindent}{1em}

% \setlength\fboxsep{0pt}

%%% longFormProof, shortFormProof, block environments
%%% \step

\newlist{steps}{enumerate}{7}
\setlist[steps]{wide,
  topsep=1pt,
  parsep=1pt,
  partopsep=1pt,
  itemsep=0pt,
  align=left,
  labelindent=0pt,
  itemindent=0pt, 
}
\setlist[steps,1]{
  label=\arabic*.,
  ref=\arabic*
}
\setlist[steps,2]{
  label=\arabic{stepsi}.\arabic*.,
  ref=\arabic{stepsi}.\arabic*
}
\setlist[steps,3]{
  label=\arabic{stepsi}.\arabic{stepsii}.\arabic*.,
  ref=\arabic{stepsi}.\arabic{stepsii}.\arabic*
}
\setlist[steps,4]{
  label=\arabic{stepsi}.\arabic{stepsii}.\arabic{stepsiii}.\arabic*.,
  ref=\arabic{stepsi}.\arabic{stepsii}.\arabic{stepsiii}.\arabic*
}
\setlist[steps,5]{
  label=\arabic{stepsi}.\arabic{stepsii}.\arabic{stepsiii}.\arabic{stepsiv}.\arabic*.,
  ref=\arabic{stepsi}.\arabic{stepsii}.\arabic{stepsiii}.\arabic{stepsiv}.\arabic*
}
\setlist[steps,6]{
  label=\arabic{stepsi}.\arabic{stepsii}.\arabic{stepsiii}.\arabic{stepsiv}.\arabic{stepsv}.\arabic*.,
  ref=\arabic{stepsi}.\arabic{stepsii}.\arabic{stepsiii}.\arabic{stepsiv}.\arabic{stepsv}.\arabic*
}
\setlist[steps,7]{
  label=\arabic{stepsi}.\arabic{stepsii}.\arabic{stepsiii}.\arabic{stepsiv}. \arabic{stepsv}. \arabic{stepsvi}.\arabic*.,
  ref=\arabic{stepsi}.\arabic{stepsii}.\arabic{stepsiii}.\arabic{stepsiv}.\arabic{stepsv}. \arabic{stepsvi}.\arabic*
}

\newcommand{\step}[1][]{\item\maybeLabel{#1}}

\makeatletter
% \skipitem steps the list counter itself (for block labels).  Give each one
% its own hyperref link target, since step numbers repeat from proof to proof.
\newcounter{skipitem}
\newcommand{\skipitem}{\stepcounter{skipitem}\refstepcounter{\@enumctr}}
\renewcommand{\theHstepsi}{skipitem.\arabic{skipitem}}
\renewcommand{\theHstepsii}{skipitem.\arabic{skipitem}}
\renewcommand{\theHstepsiii}{skipitem.\arabic{skipitem}}
\renewcommand{\theHstepsiv}{skipitem.\arabic{skipitem}}
\renewcommand{\theHstepsv}{skipitem.\arabic{skipitem}}
\renewcommand{\theHstepsvi}{skipitem.\arabic{skipitem}}
\renewcommand{\theHstepsvii}{skipitem.\arabic{skipitem}}
\makeatother

\newcommand{\maybeLabel}[1]{\ifthenelse{\isempty{#1}}{}{\label{#1}}}

\newenvironment{longFormProof}[1][Proof (long form).]
{\begin{proof}[#1]~\par\begin{steps}}{\end{steps}\vspace*{-3ex}\end{proof}}

\newenvironment{shortFormProof}[1][Proof (short form).]
{\begin{proof}[#1]}{\end{proof}}

\newenvironment{longFormSubProof}[1][Proof.]
{\begin{adjustwidth}{1em}{0em}\begin{proof}[#1]~\par\vspace*{-2pt}\skipitem\begin{steps}}
{\end{steps}\vspace*{-1ex}\end{proof}\end{adjustwidth}\smallskip}

\newenvironment{shortFormSubProof}[1][Proof.]
{\begin{proof}[#1]}
{\end{proof}\vspace*{-3pt}}

\newenvironment{block}[1][]
{\skipitem\maybeLabel{#1}\item[] \begin{steps}}{\end{steps}}

\newcommand{\comment}[1]{\hfill{\footnotesize\emph{#1}}}
\newcommand{\lineacross}{\par\vspace*{-0.7\baselineskip}\noindent\hrulefill\par}

%%% problem

% \newcounter{problem}
% \newcommand{\problem}[1][]{}
% {\ifthenelse{\isempty{#1}}{\refstepcounter{problem}}{}%
%   \paragraph{Problem \ifthenelse{\isempty{#1}}
%     {\bf\arabic{problem}.}
%     {#1}} }

\newcommand{\problemNumbered}[1]{
  \setcounter{problem}{#1}%
  \addtocounter{problem}{-1}%
  \problem
}
\newenvironment{problemTheorem}{%
  \setcounter{theorem}{\value{problem}}%
  \addtocounter{theorem}{-1}%
  \begin{theorem}[Problem \arabic{problem}]
  }{\end{theorem}}

\newcounter{problem}
\newcommand{\problem}[1][]
{\ifthenelse{\isempty{#1}}
  {\refstepcounter{problem}\paragraph{Problem~\arabic{problem}.}}
  {\paragraph{Problem~#1.}}\refreshHeader} 

\newcounter{exercise}
\newcommand{\exercise}[1][]
{\ifthenelse{\isempty{#1}}
  {\refstepcounter{exercise}\paragraph{Exercise~\theexercise.}}
  {\paragraph{Exercise~#1.}}}

\newcommand{\solution}{\paragraph{Solution.}} 

\newenvironment{mylist}[1][]
{\begin{list}{#1}{
      \setlength{\itemsep}{-1pt}
      \setlength{\labelwidth}{0pt}
      \setlength{\leftmargin}{1.25em}
    }}
{\end{list}}

\newcommand{\LNTitle}{\textcolor{red}{UNDEFINED LNTitle}}

% In a standalone lecture note's table of contents, leave room for numbers
% like 7.10 (in bold) and 7.10.2.
\makeatletter
\newcommand{\LNtocWidths}{%
  \patchcmd{\l@section}{1.5em}{2.8em}{}{\PackageWarning{macros}{Could not widen section numbers in the contents}}%
  \renewcommand*{\l@subsection}{\@dottedtocline{2}{2.8em}{3.2em}}%
  \renewcommand*{\l@subsubsection}{\@dottedtocline{3}{6.0em}{4.1em}}}
\makeatother

\newcommand{\doTitle}[3]{
  \pagenumbering{roman}
  \renewcommand{\LNTitle}{Lecture Note #1, #3}
  \renewcommand{\currentLN}{#1}
  \codeListingstrue
  % number sections, theorems etc., exercises, figures, tables and equations
  % within the lecture note (e.g. Section 8.2, Lemma 8.6), as in the book
  \renewcommand{\thesection}{#1.\arabic{section}}
  \renewcommand{\thelemma}{#1.\arabic{lemma}}
  \renewcommand{\thedefinition}{#1.\arabic{definition}}
  \renewcommand{\thetheorem}{#1.\arabic{theorem}}
  \renewcommand{\thecorollary}{#1.\arabic{corollary}}
  \renewcommand{\theclaim}{#1.\arabic{claim}}
  \renewcommand{\theobservation}{#1.\arabic{observation}}
  \renewcommand{\theconjecture}{#1.\arabic{conjecture}}
  \renewcommand{\theexercise}{#1.\arabic{exercise}}
  \renewcommand{\thefigure}{#1.\arabic{figure}}
  \renewcommand{\thetable}{#1.\arabic{table}}
  \renewcommand{\theequation}{#1.\arabic{equation}}
  \LNtocWidths
  \title{Lecture Note #1\\ #2}
  % a note's \LNabstract, if it defines one before \doTitle, goes right below the title
  \date{\ifdefined\LNabstract\parbox{0.85\textwidth}{\normalsize\normalfont\LNabstract}\\[3ex]\fi
    draft of \today, by \href{\authorURL}{N.~Young}\\[2ex] {\small\copyrightNotice}}
  \maketitle
  {\small
    \hypertarget{TOC}{}
    \tableofcontents
  }
  \thispagestyle{empty}
  % the text starts on page 1, after the title and contents (numbered i, ii, ...)
  \clearpage
  \pagenumbering{arabic}
}

\def\clap#1{\hbox to 0pt{\hss#1\hss}}

\newcommand{\refreshHeader}{}
\newcommand{\header}{}

\newcommand{\doHeader}[1]{      % for lecture notes
  \renewcommand{\header}{%
    \clap{\raisebox{0pt}{\protect\hyperlink{TOC}{$\uparrow$\scriptsize\,\textsc{contents}\hspace{1.1in}}}}%
    {\footnotesize \LNTitle{} / Section \thesection, #1}} %, #1, draft of \today}}
  \markboth{\header}{\header}
}

\newcommand{\SID}{}
\newcommand{\setAuthor}[1]{\renewcommand{\author}{#1}}
\newcommand{\setSID}[1]{\renewcommand{\SID}{#1}}

\newcommand{\doHwHeader}[1]{                                         % for homeworks
  \renewcommand{\refreshHeader}{
    \renewcommand{\header}{%
      {\footnotesize \classSection #1 / Problem~\arabic{problem} / draft of \today \hfill \author\ifGradescope~(\SID)\fi~~}}
    \markboth{\header}{\header}
  }
  \renewcommand{\header}{%
    {\footnotesize \classSection #1 / Instructions / draft of \today \hfill \author\ifGradescope~(\SID)\fi~~}}
  \markboth{\header}{\header}
}

% Homework settings, for instructors using the homeworks in a course:
% \classSection is the course name shown in each homework's running head
% (e.g. \renewcommand{\classSection}{CS 3000\xspace}), and \Gradescopetrue turns on
% the instructions for submitting to Gradescope (in its fixed-length mode),
% the hints saying on which page of the PDF each answer should be, and the
% student ID in the running head.
\newcommand{\classSection}{}
\newif\ifGradescope
\Gradescopefalse
\newcommand{\GradescopeComment}[1]{\ifGradescope\comment{#1}\fi}

\pagestyle{myheadings}
\addtolength{\headsep}{0.2in}
\addtolength{\topmargin}{-0.3in}
\addtolength{\textheight}{0.4in}

%%%%%%%%%%%%%%%%%%%%%%%%%%%%%%%%%%%%%%%%%%%%%%%%
%%%%%%%%%%%%%%%%%%%%%%%%%%%%%%%%%%%%%%%%%%%%%%%%

\newcommand{\MF}{\href
  {https://mfleck.cs.illinois.edu/building-blocks/}
  {MF}\xspace}

\newcommand{\KWSlides}{\href
  {https://www.cs.princeton.edu/~wayne/kleinberg-tardos/}
  {KW slides}\xspace}

\newcommand{\LLM}{\href
  {https://courses.csail.mit.edu/6.042/spring18/mcs.pdf}
  {LLM}\xspace}

% \DPV, \KT, \CLRS: textbooks listed (without links) in the table of sources in
% the Preliminaries note; each prints as a link to its entry in that table.
\DeclareRobustCommand{\DPV}{\noteLink{prelim}{source.DPV}{DPV}\xspace}
\DeclareRobustCommand{\KT}{\noteLink{prelim}{source.KT}{KT}\xspace}
\DeclareRobustCommand{\CLRS}{\noteLink{prelim}{source.CLRS}{CLRS}\xspace}
\AtBeginDocument{\pdfstringdefDisableCommands{\def\DPV{DPV }\def\KT{KT }\def\CLRS{CLRS }}}

\newcommand{\JE}{\href
  {https://jeffe.cs.illinois.edu/teaching/algorithms/}
  {JE}\xspace}

\newcommand{\JEGreedy}{\href
  {https://jeffe.cs.illinois.edu/teaching/algorithms/book/04-greedy.pdf}
  {JE Chapter 4}\xspace}

\newcommand{\BBCode}{\href
{https://northeastern.blackboard.com/webapps/cmsmain/webui/courses/CS3000.MERGED.202010/code?action=frameset&subaction=view&uniq=-jj2max&course_id=_2597392_1&mask=\%2Fcourses\%2FCS3000.MERGED.202010}
  {code}\xspace}

% The Python code (src/code).  Each lecture note that has code ends with a
% section listing its files: \codeSection{path, path, ...}.  In the text,
% \pycode{greedy_algorithms/Dijkstra.py} prints the path as a link to its listing
% (in the homeworks, which have no listings, to the file on GitHub), and
% \pycode[text]{...} uses the given text instead.  \pythonCode{path, ...} is the
% standard sentence ending a section whose algorithms have code:
% "Python code: path, ... (Section N.k)."
\newcommand{\codeFolder}{https://github.com/nealeyoung/algs101/tree/main/src/code}
\newcommand{\codeRepo}{https://github.com/nealeyoung/algs101/blob/main/src/code}
\newif\ifcodeListings   % set by \doTitle: this document has the listings
\newcommand{\githubCode}[2][]{\href{\codeRepo/#2}{\ifx\relax#1\relax\nolinkurl{#2}\else#1\fi}}
\newcommand{\pycode}[2][]{%
  \ifcodeListings\hyperref[code: #2]{\ifx\relax#1\relax\nolinkurl{#2}\else#1\fi}%
  \else\githubCode[#1]{#2}\fi}
\makeatletter
\newcommand{\pythonCode}[1]{%
  \def\pythonCode@first{}\def\pythonCode@list{}%
  \forcsvlist{\pythonCode@add}{#1}%
  Python code: \pythonCode@list\ (Section~\ref{code: \pythonCode@first}).}
\newcommand{\pythonCode@add}[1]{%
  \ifx\pythonCode@first\@empty\def\pythonCode@first{#1}\appto\pythonCode@list{\pycode{#1}}%
  \else\appto\pythonCode@list{, \pycode{#1}}\fi}
\makeatother

% Code listings that copy and paste exactly (in PDF viewers that honor the
% ToUnicode maps, e.g. MuPDF-based ones; PDFKit viewers such as Preview and
% Skim drop leading spaces, so each listing's heading links to the file on
% GitHub).  Spaces are typeset as white visible-space glyphs that copy as
% spaces, and hyphens and asterisks as the T1 glyphs, so they copy as themselves.
\usepackage[T1,OT1]{fontenc}   % T1 for the code listings only; OT1 stays the default
\usepackage{listings}
\pdfgentounicode=1
\input glyphtounicode
\pdfglyphtounicode{uni2423}{0020}
\makeatletter
\def\lst@visiblespace{\textcolor{white}{\char32}}
% put an invisible space on each blank line, so that blank lines copy too
% (listings' OnEmptyLine hook runs before the line break, so patch the loop)
\def\lst@DoNewLines{%
    \@whilenum\lst@newlines>\@ne \do {\lst@NewLine \lst@visiblespace}%
    \ifnum\lst@newlines>\z@ \lst@NewLine \fi}
\makeatother
\lstset{language=Python,
  basicstyle=\fontencoding{T1}\fontfamily{lmtt}\selectfont\footnotesize,
  upquote=true, showstringspaces=false, columns=fullflexible, keepspaces=true,
  showspaces=true, literate={-}{{\char45}}1 {*}{{\char42}}1,
  frame=single, framesep=1ex, rulecolor=\color{gray}}
% the listings' paths are relative to the lecture note's folder (in the book, to src/)
\ifdefined\BookMode\newcommand{\codeDir}{code/}\else\newcommand{\codeDir}{../code/}\fi
% \codefile{path}: a listing of src/code/path, headed by a link to it on GitHub
\newcommand{\codefile}[1]{%
  \subsection*{\githubCode{#1}}\label{code: #1}%
  \lstinputlisting{\codeDir#1}}
% \codeSection{path, path, ...}: the lecture note's last section, listing its code
\newcommand{\codeSection}[1]{%
  \newpage
  \section{Python code}\label{\LNtag{\currentLN}: sec: code}\doHeader{Python code}
  This section lists the Python code for this lecture note.
  Click on a file's name to open the file on GitHub.
  Copying code from GitHub is more reliable than copying it from this PDF,
  since some PDF viewers drop the leading spaces (the indentation) of copied lines.
  \forcsvlist{\codefile}{#1}}
% the copyright notice, for title pages and the homework first pages
\newcommand{\licenseURL}{https://github.com/nealeyoung/algs101/blob/main/LICENSE}
\newcommand{\authorURL}{https://www.cs.ucr.edu/\string~neal}
\newcommand{\copyrightNotice}{%
  \copyright\ 2018--2026 \href{\authorURL}{Neal E. Young}.
  Text: CC BY-NC-SA 4.0; code: MIT; except where noted.
  See \url{\licenseURL}.}
% the published site: PDFs of the notes, and the homework assignments and templates
\newcommand{\siteURL}{https://nealeyoung.github.io/algs101/}

\newcommand{\blackboard}{\href
  {https://northeastern.blackboard.com/webapps/blackboard/execute/launcher?type=Course&id=_2597392_1&url=}
  {Blackboard}\xspace}

\newcommand{\JUSlides}{\href
  {http://www.khoury.neu.edu/home/jullman/cs3000f18/schedule.html}
  {Prof.~Ullman's slides}\xspace}

% a URL on its own line, in three layouts
\newcommand{\URL}[1]{\par-- \parbox[t]{0.8\textwidth}{\url{#1}} \medskip\par}
\newcommand{\URLsmall}[1]{\par-- \parbox[t]{0.95\textwidth}{\small\url{#1}} \medskip\par}
\newcommand{\URLplain}[1]{\par\parbox[t]{\textwidth}{\url{#1}}\par}

%%%%%%%%%%%%%%%%%%%%%%%%%%%%%%%%%%%%%%%%%%%%%%%%
%%%%%%%%%%%%%%%%%%%%%%%%%%%%%%%%%%%%%%%%%%%%%%%%

\newcommand{\continued}{\vfill\centerline{\emph{(continued)}}\vfill\newpage}

% load hyperref last to avoid conflicts; long URLs may also break after hyphens
\PassOptionsToPackage{hyphens}{url}
\ifdefined\BookMode\else\usepackage{xr-hyper}\fi   % for \lectureNoteRefs (homeworks), below
\usepackage{hyperref}
\hypersetup{
    colorlinks,
    linkcolor={blue!50!black},
    citecolor={blue!50!black},
    urlcolor={blue!80!black}
}

\providecommand{\answerDirectory}{.}
\newcommand{\inputAnswer}[1]{\refreshHeader\input{\answerDirectory/#1}}
\newcommand{\inputSection}[1]{\input{#1}}

% Refer to a lecture note by its tag (the prefix of its labels), e.g.
% \LN{proofs} prints "Lecture Note 2" (in the book, a link to that chapter).  The table below is the only place the numbers appear.
\makeatletter
\newcommand{\defineLN}[3]{\@namedef{LNnumber@#2}{#1}\@namedef{LNtag@#1}{#2}\@namedef{LNfile@#2}{#3}}
\newcommand{\LNnumber}[1]{%
  \ifcsname LNnumber@#1\endcsname\csname LNnumber@#1\endcsname
  \else\PackageError{macros}{Unknown lecture-note tag `#1'}{}\fi}
\newcommand{\LNtag}[1]{\csname LNtag@#1\endcsname}
\newcommand{\LNfile}[1]{\csname LNfile@#1\endcsname}
\makeatother
\defineLN{1}{prelim}{1_preliminaries}
\defineLN{2}{proofs}{2_proofs}
\defineLN{3}{matching}{3_stable_matching}
\defineLN{4}{dc}{4_divide_and_conquer}
\defineLN{5}{greedy}{5_greedy}
\defineLN{6}{graphs}{6_graph_traversal}
\defineLN{7}{dp}{7_dynamic_programming}
\defineLN{8}{flow}{8_flow_reductions_LP}
\defineLN{9}{np}{9_NP}

% The text is the same in the book and in the standalone lecture notes; only
% links differ.  In the book, links to other lecture notes stay in the book.
% In a standalone note, they open the other note's published PDF.
% \noteLink{tag}{target}{text} links text to \hypertarget{target} in lecture note tag.
% \LN{proofs} prints "Lecture Note 2", linked to that lecture note.
\newcommand{\currentLN}{0}   % in a standalone note, its number (set by \doTitle)
\newcommand{\notePDF}[1]{\siteURL notes/\LNfile{#1}.pdf}
\newcommand{\noteLink}[3]{%
  \ifnum\LNnumber{#1}=\currentLN\relax\hyperlink{#2}{#3}%
  \else\href{\notePDF{#1}\#nameddest=#2}{#3}\fi}
% \bookLink{target}{text} links text to \hypertarget{target} in the book's
% appendices (from a standalone note, in the published book).
\DeclareRobustCommand{\bookLink}[2]{%
  \ifdefined\BookMode\hyperlink{#1}{#2}%
  \else\href{\siteURL lecture_notes_on_algorithms.pdf\#nameddest=#1}{#2}\fi}
\DeclareRobustCommand{\LN}[1]{%
  \ifnum\LNnumber{#1}=\currentLN\relax Lecture Note~\LNnumber{#1}%
  \else\href{\notePDF{#1}}{Lecture Note~\LNnumber{#1}}\fi}
\pdfstringdefDisableCommands{\def\LN#1{Lecture Note \LNnumber{#1}}}

% Homeworks refer to the notes' sections, exercises and lemmas by \ref: each
% homework's main file says \lectureNoteRefs after \input{macros}.  A homework
% folder has a link refs/ to src/refs/, which holds each note's label records
% (made by tools/refs.sh from the notes' builds).  Read with xr-hyper, they make
% \ref{greedy: exercise: ...} give the exercise's number in the published note,
% linked to it there; labels only the book has (its appendices) link to the book.
\ifdefined\BookMode\else
  \makeatletter
  \newcommand{\lectureNoteRefs}{%
    \@for\LN@tag:=prelim,proofs,matching,dc,greedy,graphs,dp,flow,np\do{%
      \externaldocument{refs/\LNfile{\LN@tag}}[\notePDF{\LN@tag}]}%
    \externaldocument{refs/book}[\siteURL lecture_notes_on_algorithms.pdf]}
  % hyperref links a \ref to another document's label as URL#anchor; browsers'
  % PDF viewers need URL#nameddest=anchor (the form \noteLink and \bookLink use)
  \def\Hy@setref@link#1#2#3#4#5#6\@nil#7{%
    \begingroup
      \ifx\\#5\\\toks0={\hyper@@link{#5}{#4}}\else\toks0={\hyper@@link{#5}{nameddest=#4}}\fi
      \toks1=\expandafter{#7{#1}{#2}{#3}{#4}{#5}}%
      \edef\x{\endgroup\the\toks0 {\the\toks1 }}%
    \x}
  \makeatother
\fi

% Book mode: main.tex defines \BookMode before \input{macros} to combine the
% lecture notes into one report-class document, one chapter per note.
% Without \BookMode (each lecture note on its own, and the homeworks) none
% of this applies.
\ifdefined\BookMode
  \usepackage{etoc}

  % links to other lecture notes stay in the book
  \renewcommand{\noteLink}[3]{\hyperlink{#2}{#3}}

  % each chapter is a lecture note, and is called one, as in the standalone notes
  \renewcommand{\chaptername}{Lecture Note}
  \setcounter{secnumdepth}{3}   % number subsubsections, as in the standalone notes

  % global table of contents, for main.tex
  \newcommand{\bookContents}{%
    \hypertarget{TOC}{}%
    \etocsetnexttocdepth{section}%
    \tableofcontents}

  % \doTitle{N}{title}{short title} starts chapter N, with a link to the global
  % table of contents and the chapter's own table of contents
  \renewcommand{\doTitle}[3]{%
    \setcounter{chapter}{\numexpr#1-1\relax}%
    \chapter{\texorpdfstring{#2}{#3}}%
    \label{\LNtag{#1}: chapter}%
    \csgdef{LNshorttitle@#1}{#3}%
    \csgdef{LNtitle@#1}{#2}%
    \renewcommand{\LNTitle}{Lecture Note #1, #3}%
    \codeListingstrue
    \hyperlink{TOC}{$\uparrow$\scriptsize\,\textsc{contents}}%
    {\small
      \hypertarget{TOC-#1}{}%
      \etocsettocstyle{\section*{\contentsname}}{}%
      \etocsetnexttocdepth{subsubsection}%
      \localtableofcontents
    }}

  % the code: each lecture note's \codeSection just records its files, and
  % \codeAppendix (in the book's appendix) lists them all, a section per note
  \makeatletter   % (file names have no spaces, so drop the spaces around them)
  \renewcommand{\codeSection}[1]{\csxdef{LNcode@\arabic{chapter}}{\zap@space#1 \@empty}}
  \makeatother
  % \codeAppendix: first a contents list (each note's topic, then its files,
  % linked to the listings), then a section per note listing its files
  \makeatletter
  \def\LNbasename#1{\LN@base#1/\LN@end}   % a/b/c.py -> c.py
  \def\LN@base#1/#2\LN@end{\ifx\relax#2\relax#1\else\LN@base#2\LN@end\fi}
  \makeatother
  \newcommand{\codeContentsFile}[1]{%
    \edef\LNfilename{\LNbasename{#1}}%
    \item \pycode[\texttt{\expandafter\detokenize\expandafter{\LNfilename}}]{#1}}
  \newcommand{\codeAppendix}{%
    \begin{itemize}
    \foreach \n in {1,...,9} {%
      \ifcsdef{LNcode@\n}{%
        \item \hyperref[code: note \n]{\csuse{LNtitle@\n}} (\LN{\LNtag{\n}})
        \begin{itemize}
        \edef\LNcodelist{\csuse{LNcode@\n}}%
        \expandafter\forcsvlist\expandafter\codeContentsFile\expandafter{\LNcodelist}%
        \end{itemize}}{}}
    \end{itemize}
    \newpage
    \foreach \n in {1,...,9} {%
      \ifcsdef{LNcode@\n}{%
        \section[Lecture Note \n: \csuse{LNtitle@\n}]{%
          \texorpdfstring{\LN{\LNtag{\n}}: \csuse{LNtitle@\n}}{Lecture Note \n: \csuse{LNtitle@\n}}}%
        \label{code: note \n}%
        \edef\LNcodelist{\csuse{LNcode@\n}}%
        \expandafter\forcsvlist\expandafter\codefile\expandafter{\LNcodelist}%
        \newpage}{}}}

  % a lecture note is a chapter: \LN{proofs} prints a link "Lecture Note 2"
  \DeclareRobustCommand{\LN}[1]{\hyperref[#1: chapter]{Lecture Note~\ref*{#1: chapter}}}
  \pdfstringdefDisableCommands{\def\LN#1{Lecture Note \LNnumber{#1}}}

  % running head; "contents" links to the current chapter's table of contents
  \renewcommand{\doHeader}[1]{%
    \renewcommand{\header}{%
      \clap{\raisebox{0pt}{\protect\hyperlink{TOC-\arabic{chapter}}{$\uparrow$\scriptsize\,\textsc{contents}\hspace{1.1in}}}}%
      {\footnotesize \LNTitle{} / Section \thesection, #1}}%
    \markboth{\header}{\header}}

  % number theorems, lemmas, etc. and exercises within chapters (e.g. Lemma 8.1)
  \counterwithin{lemma}{chapter}
  \counterwithin{definition}{chapter}
  \counterwithin{theorem}{chapter}
  \counterwithin{corollary}{chapter}
  \counterwithin{claim}{chapter}
  \counterwithin{observation}{chapter}
  \counterwithin{conjecture}{chapter}
  \counterwithin{exercise}{chapter}
\fi

% %%%%%%%%%%%%%%%%%%%%%%%%%%%%%%%%%%%%%%%%%%%%%%%%%%%%%%%%%%%%%%%%%%%%%%%%
% %%%%%%%%%%%%%%%%%%%%%%%%%%%%%%%%%%%%%%%%%%%%%%%%%%%%%%%%%%%%%%%%%%%%%%%%

\ExplSyntaxOn
\seq_new:N \l_itemize_seq
% \NewDocumentEnvironment{REITEMIZE}{ +b }{
\NewDocumentCommand{\REITEMIZECMD}{ +m }{
  \regex_split:nnN {\c{item}} {#1} \l_itemize_seq
  \seq_pop:NN \l_itemize_seq \l_tmpa_tl% remove the (hopefully) empty first item
  \seq_if_empty:NF \l_itemize_seq {
    \l_tmpa_tl
    \seq_map_inline:Nn \l_itemize_seq {
      \regex_split:nnN {\c{item}} {#1} \l_itemize_seq {
        \regex_extract_once:nnNTF { ^\[([^]]*)\](.*) } { ##1} \l_foo_seq
      { \item[\seq_item:Nn \l_foo_seq {2}]\ITEMIZER{\seq_item:Nn \l_foo_seq {3}} } { \item\ITEMIZER{##1}  }
    }
  }
}
}
\ExplSyntaxOff

\newlength{\ITEMHT}
\setlength{\ITEMHT}{1in}

% \setlength{\fboxsep}{1pt}
\newcommand{\ITEMIZER}[1]{%
  % \framebox{
\adjustbox{valign=t}{\parbox{\textwidth}{\linespread{1.15}\selectfont #1}}~~\adjustbox{valign=t}{\parbox{0pt}{\color{gray}\rule[-\ITEMHT]{0.01pt}{\ITEMHT}}}
% }
  % \begin{tabular}{@{}p{\textwidth}@{\rule[-\ITEMHT]{0pt}{\ITEMHT}}}
  %   #1
  % \end{tabular}
}

\newsavebox\thesmashminipage
\newlength{\thesmashminipagelength}
\newlength{\answerboxlength}
\newlength{\answerboxoverrun}
\newenvironment{smashminipage}[2]
{\setlength\answerboxlength{#2}
\begin{lrbox}{\thesmashminipage}\begin{minipage}[t]{#1}}
{\end{minipage}\end{lrbox}%
\setlength{\thesmashminipagelength}{\dimexpr\ht\thesmashminipage+\dp\thesmashminipage}%
\setlength{\answerboxoverrun}{\dimexpr\thesmashminipagelength-\answerboxlength}%
\ifdim\the\thesmashminipagelength>\the\answerboxlength%
\usebox{\thesmashminipage}%
\smash{\color{red}\hspace*{-0.98\textwidth}\rule[-\answerboxlength]{0.98\textwidth}{1pt}}
\errmessage{Answer exceeds allotted length, please shorten it.}
\else
\usebox{\thesmashminipage}%
\fi
}



\newenvironment{answerBox}[1]{%
  \def\ANSWERBOXHEIGHT{#1}
  \begin{adjustbox}{valign=t}
  \begin{smashminipage}{\textwidth}{\ANSWERBOXHEIGHT}
}{\end{smashminipage}\end{adjustbox}~~~\adjustbox{valign=t}{\parbox{0pt}{\color{gray}\rule[-\ANSWERBOXHEIGHT]{0.01pt}{\ANSWERBOXHEIGHT}\rule[-\ANSWERBOXHEIGHT]{10pt}{0.01pt}}}}

% % \setlist[enumerate]{itemsep=0pt,parsep=0pt,topsep=0pt,leftmargin=100pt}

% \makeatletter 
\LetLtxMacro\SVenumerate\enumerate
\let\endSVenumerate\endenumerate

\newcommand{\injectEnumerate}{
  \let\SValph\alph
  \def\alph{}
  \renewenvironment{enumerate}[1][]
  {\let\alph\SValph \SVenumerate[label=(\alph*),itemsep=2pt,parsep=0pt]\REITEMIZE}{\endREITEMIZE\endSVenumerate}
  % \setlist[enumerate]{before=\Collect@Body\REITEMIZECMD}
}
% \makeatother
\newcommand{\restoreEnumerate}{
  \renewenvironment{enumerate}[1][]
  \let\SValph\alph
  \def\alph{}
  {\let\alph\SValph \SVenumerate[label=(\alph*)]}{\endSVenumerate}
  % \setlist[enumerate]{before=,after=}
}
\newcommand{\injectBlue}{
  \let\SVblue\blue
  \let\endSVblue\endblue
  \def\blue{\SVblue\REITEMIZE}
  \def\endblue{\endREITEMIZE\endSVblue}
}
\newcommand{\restoreBlue}{
  \let\blue\SVblue
  \let\endblue\endSVblue
}
\newcommand{\injectEnumerateBlue}{
  \setlist[enumerate]{before=\injectBlue,after=\restoreBlue,itemsep=2pt,parsep=0pt}
}
\newcommand{\restoreEnumerateBlue}{
  \setlist[enumerate]{before=,after=}
}


% %%%%%%%%%%%%%%%%%%%%%%%%%%%%%%%%%%%%%%%%%%%%%%%%%%%%%%%%%%%%%%%%%%%%%%%%
% %%%%%%%%%%%%%%%%%%%%%%%%%%%%%%%%%%%%%%%%%%%%%%%%%%%%%%%%%%%%%%%%%%%%%%%%

% \NewDocumentEnvironment{env}{ }
% {\Environenv}
% {\endEnvironenv}

% \NewEnviron{REITEMIZE}{\REITEMIZECMD\BODY}

\makeatletter 
\newenvironment{REITEMIZE}{\Collect@Body\REITEMIZECMD}{}
\makeatother 

\newcommand{\NODEone}[1]{\ensuremath{v_{#1}}\xspace}
\newcommand{\NODE}[2]{{\NODEone {#1,#2}}\xspace}
\newcommand{\EDGE}[2]{\ensuremath{(\textsf{#1}, \textsf{#2})}}
\newcommand{\ITEM}{\item}

%%% Local Variables:
%%% mode: latex
%%% End:
.  A homework
% folder has a link refs/ to src/refs/, which holds each note's label records
% (made by tools/refs.sh from the notes' builds).  Read with xr-hyper, they make
% \ref{greedy: exercise: ...} give the exercise's number in the published note,
% linked to it there; labels only the book has (its appendices) link to the book.
\ifdefined\BookMode\else
  \makeatletter
  \newcommand{\lectureNoteRefs}{%
    \@for\LN@tag:=prelim,proofs,matching,dc,greedy,graphs,dp,flow,np\do{%
      \externaldocument{refs/\LNfile{\LN@tag}}[\notePDF{\LN@tag}]}%
    \externaldocument{refs/book}[\siteURL lecture_notes_on_algorithms.pdf]}
  % hyperref links a \ref to another document's label as URL#anchor; browsers'
  % PDF viewers need URL#nameddest=anchor (the form \noteLink and \bookLink use)
  \def\Hy@setref@link#1#2#3#4#5#6\@nil#7{%
    \begingroup
      \ifx\\#5\\\toks0={\hyper@@link{#5}{#4}}\else\toks0={\hyper@@link{#5}{nameddest=#4}}\fi
      \toks1=\expandafter{#7{#1}{#2}{#3}{#4}{#5}}%
      \edef\x{\endgroup\the\toks0 {\the\toks1 }}%
    \x}
  \makeatother
\fi

% Book mode: main.tex defines \BookMode before %%%% DON'T EDIT THIS FILE %%%%

\usepackage{xifthen}% \isempty
\usepackage{etoolbox}% \patchcmd
% \usepackage{stringstrings}% \substring

\usepackage{amsmath, amsthm, amssymb}
\usepackage[margin=1in]{geometry}
\usepackage{xspace}
\usepackage{graphicx}
\usepackage[normalem]{ulem}
\usepackage{adjustbox}
\usepackage{xcolor}
\usepackage{changepage}
\usepackage{framed}
\usepackage{multicol}
\usepackage{enumitem}
% \usepackage{listings} % lstlistings environment

\usepackage{makecmds}           %\provideenvironment

\usepackage{xparse,environ}
\usepackage{letltxmacro}       % \LetLtxMacro
\usepackage{wrapfig}
\usepackage{tikz}

\usetikzlibrary{shapes}
\usetikzlibrary{positioning}
\usetikzlibrary{shapes.misc}
\usetikzlibrary{arrows}

\newtheorem{lemma}{Lemma}
\newtheorem{definition}{Definition}
\newtheorem{theorem}{Theorem}
\newtheorem{corollary}{Corollary}
\newtheorem{claim}{Claim}
\newtheorem{observation}{Observation}
\newtheorem{conjecture}{Conjecture}

% A lemma or theorem restated later with its original number (set with
% \setcounter) goes in a restatement environment, which gives it its own link
% anchor (otherwise it would duplicate the original's).
\newenvironment{restatement}
  {\def\theHlemma{restated.\thelemma}\def\theHtheorem{restated.\thetheorem}}{}

\newenvironment{myframed}
{\smallskip\par\noindent\begin{minipage}{\textwidth}\begin{framed}\vspace*{-5pt}}{\vspace*{-5pt}\end{framed}\end{minipage}\smallskip}


\definecolor{darkblue}{RGB}{0,0,200}

 
\let\oldParagraph\paragraph
\providecommand{\gradescopeColor}{\color{black}}
\renewcommand{\paragraph}[1]{\oldParagraph{\gradescopeColor #1}}
\providecommand{\gradescopeInit}{}

\providecommand{\BLUE}[1]{{\color{darkblue} #1}}
\providecommand{\REPLACEME}{\textcolor{red}{\ensuremath{\star \cdots \star}}\xspace}
\provideenvironment{blue}{\color{darkblue}}{}

\newcommand{\dist}[2]{\ensuremath{\operatorname{\sf dist}(#1, #2)}}
\newcommand{\depth}[2]{\ensuremath{\operatorname{\sf depth}_{#1}(#2)}}
\newcommand{\cost}[1]{\ensuremath{\operatorname{\sf cost}(#1)}}
\newcommand{\weight}[1]{\ensuremath{\operatorname{\sf wt}(#1)}}
% \newcommand{\val}[1]{\ensuremath{\operatorname{\sf val}(#1)}}
\newcommand{\eps}{\varepsilon}
\newcommand{\capac}{\textsf{cap}}
\newcommand{\capOf}[1]{\ensuremath{\operatorname{\sf cap}(#1)}}
\newcommand{\rescap}[1]{\ensuremath{\operatorname{\sf cap}_f(#1)}}
\newcommand{\into}[1]{\ensuremath{\operatorname{\sf in}(#1)}}
\newcommand{\out}[1]{\ensuremath{\operatorname{\sf out}(#1)}}
\newenvironment{LP}[1]{
  \begin{array}[t]{@{}l@{}}
    \displaystyle #1 \\
    \begin{aligned}\\[-10pt]}{\end{aligned}\end{array}}

\newcommand{\opt}[1]{\ensuremath{\operatorname{\sf opt}(#1)}}
\newcommand{\Opt}[1]{\ensuremath{\operatorname{\sf Opt}(#1)}}
\newcommand{\half}{\frac{1}{2}}

\newcommand{\E}{\operatorname{E}}
\newcommand{\giv}{\,|\,}
\newcommand{\Z}{\mathbb{Z}}
\newcommand{\Zp}{\mathbb{Z}_{\ge 0}}
\newcommand{\R}{\mathbb{R}}
% \newcommand{\Rp}{\mathbb{R}_{\ge 0}}
\newcommand{\Rp}{\mathbb{R}_{\tiny +}}
\newcommand{\N}{\mathbb{N}}
\newcommand{\Np}{\mathbb{N}_{\ge 0}}
\newcommand{\calS}{{\cal S}}

\newcommand{\OP}[1]{\text{\small\sc #1}\xspace}
\newcommand{\Insert}{\OP{Insert}}
\newcommand{\FindMin}{\OP{Find-Min}}
\newcommand{\DeleteMin}{\OP{Delete-Min}}
\newcommand{\DecreaseKey}{\OP{Decrease}}
\newcommand{\Merge}{\OP{Merge}}
\newcommand{\Cut}{\OP{Cut}}
% \newcommand{\prob}[1]{\text{\sc #1}}
% a problem name; in running text, lines may break inside it
\DeclareRobustCommand{\prob}[1]{\ifmmode\text{#1}\else#1\fi}
\newcommand{\encode}[1]{\langle #1 \rangle}

\newcommand{\best}[1]{\text{\sf best}(#1)}
\newcommand{\worst}[1]{\text{\sf worst}(#1)}
\newcommand{\floor}[1]{\lfloor #1 \rfloor}
\newcommand{\ceil}[1]{\lceil #1 \rceil}

\setlength{\parskip}{2pt}
\setlength{\parindent}{1em}

% \setlength\fboxsep{0pt}

%%% longFormProof, shortFormProof, block environments
%%% \step

\newlist{steps}{enumerate}{7}
\setlist[steps]{wide,
  topsep=1pt,
  parsep=1pt,
  partopsep=1pt,
  itemsep=0pt,
  align=left,
  labelindent=0pt,
  itemindent=0pt, 
}
\setlist[steps,1]{
  label=\arabic*.,
  ref=\arabic*
}
\setlist[steps,2]{
  label=\arabic{stepsi}.\arabic*.,
  ref=\arabic{stepsi}.\arabic*
}
\setlist[steps,3]{
  label=\arabic{stepsi}.\arabic{stepsii}.\arabic*.,
  ref=\arabic{stepsi}.\arabic{stepsii}.\arabic*
}
\setlist[steps,4]{
  label=\arabic{stepsi}.\arabic{stepsii}.\arabic{stepsiii}.\arabic*.,
  ref=\arabic{stepsi}.\arabic{stepsii}.\arabic{stepsiii}.\arabic*
}
\setlist[steps,5]{
  label=\arabic{stepsi}.\arabic{stepsii}.\arabic{stepsiii}.\arabic{stepsiv}.\arabic*.,
  ref=\arabic{stepsi}.\arabic{stepsii}.\arabic{stepsiii}.\arabic{stepsiv}.\arabic*
}
\setlist[steps,6]{
  label=\arabic{stepsi}.\arabic{stepsii}.\arabic{stepsiii}.\arabic{stepsiv}.\arabic{stepsv}.\arabic*.,
  ref=\arabic{stepsi}.\arabic{stepsii}.\arabic{stepsiii}.\arabic{stepsiv}.\arabic{stepsv}.\arabic*
}
\setlist[steps,7]{
  label=\arabic{stepsi}.\arabic{stepsii}.\arabic{stepsiii}.\arabic{stepsiv}. \arabic{stepsv}. \arabic{stepsvi}.\arabic*.,
  ref=\arabic{stepsi}.\arabic{stepsii}.\arabic{stepsiii}.\arabic{stepsiv}.\arabic{stepsv}. \arabic{stepsvi}.\arabic*
}

\newcommand{\step}[1][]{\item\maybeLabel{#1}}

\makeatletter
% \skipitem steps the list counter itself (for block labels).  Give each one
% its own hyperref link target, since step numbers repeat from proof to proof.
\newcounter{skipitem}
\newcommand{\skipitem}{\stepcounter{skipitem}\refstepcounter{\@enumctr}}
\renewcommand{\theHstepsi}{skipitem.\arabic{skipitem}}
\renewcommand{\theHstepsii}{skipitem.\arabic{skipitem}}
\renewcommand{\theHstepsiii}{skipitem.\arabic{skipitem}}
\renewcommand{\theHstepsiv}{skipitem.\arabic{skipitem}}
\renewcommand{\theHstepsv}{skipitem.\arabic{skipitem}}
\renewcommand{\theHstepsvi}{skipitem.\arabic{skipitem}}
\renewcommand{\theHstepsvii}{skipitem.\arabic{skipitem}}
\makeatother

\newcommand{\maybeLabel}[1]{\ifthenelse{\isempty{#1}}{}{\label{#1}}}

\newenvironment{longFormProof}[1][Proof (long form).]
{\begin{proof}[#1]~\par\begin{steps}}{\end{steps}\vspace*{-3ex}\end{proof}}

\newenvironment{shortFormProof}[1][Proof (short form).]
{\begin{proof}[#1]}{\end{proof}}

\newenvironment{longFormSubProof}[1][Proof.]
{\begin{adjustwidth}{1em}{0em}\begin{proof}[#1]~\par\vspace*{-2pt}\skipitem\begin{steps}}
{\end{steps}\vspace*{-1ex}\end{proof}\end{adjustwidth}\smallskip}

\newenvironment{shortFormSubProof}[1][Proof.]
{\begin{proof}[#1]}
{\end{proof}\vspace*{-3pt}}

\newenvironment{block}[1][]
{\skipitem\maybeLabel{#1}\item[] \begin{steps}}{\end{steps}}

\newcommand{\comment}[1]{\hfill{\footnotesize\emph{#1}}}
\newcommand{\lineacross}{\par\vspace*{-0.7\baselineskip}\noindent\hrulefill\par}

%%% problem

% \newcounter{problem}
% \newcommand{\problem}[1][]{}
% {\ifthenelse{\isempty{#1}}{\refstepcounter{problem}}{}%
%   \paragraph{Problem \ifthenelse{\isempty{#1}}
%     {\bf\arabic{problem}.}
%     {#1}} }

\newcommand{\problemNumbered}[1]{
  \setcounter{problem}{#1}%
  \addtocounter{problem}{-1}%
  \problem
}
\newenvironment{problemTheorem}{%
  \setcounter{theorem}{\value{problem}}%
  \addtocounter{theorem}{-1}%
  \begin{theorem}[Problem \arabic{problem}]
  }{\end{theorem}}

\newcounter{problem}
\newcommand{\problem}[1][]
{\ifthenelse{\isempty{#1}}
  {\refstepcounter{problem}\paragraph{Problem~\arabic{problem}.}}
  {\paragraph{Problem~#1.}}\refreshHeader} 

\newcounter{exercise}
\newcommand{\exercise}[1][]
{\ifthenelse{\isempty{#1}}
  {\refstepcounter{exercise}\paragraph{Exercise~\theexercise.}}
  {\paragraph{Exercise~#1.}}}

\newcommand{\solution}{\paragraph{Solution.}} 

\newenvironment{mylist}[1][]
{\begin{list}{#1}{
      \setlength{\itemsep}{-1pt}
      \setlength{\labelwidth}{0pt}
      \setlength{\leftmargin}{1.25em}
    }}
{\end{list}}

\newcommand{\LNTitle}{\textcolor{red}{UNDEFINED LNTitle}}

% In a standalone lecture note's table of contents, leave room for numbers
% like 7.10 (in bold) and 7.10.2.
\makeatletter
\newcommand{\LNtocWidths}{%
  \patchcmd{\l@section}{1.5em}{2.8em}{}{\PackageWarning{macros}{Could not widen section numbers in the contents}}%
  \renewcommand*{\l@subsection}{\@dottedtocline{2}{2.8em}{3.2em}}%
  \renewcommand*{\l@subsubsection}{\@dottedtocline{3}{6.0em}{4.1em}}}
\makeatother

\newcommand{\doTitle}[3]{
  \pagenumbering{roman}
  \renewcommand{\LNTitle}{Lecture Note #1, #3}
  \renewcommand{\currentLN}{#1}
  \codeListingstrue
  % number sections, theorems etc., exercises, figures, tables and equations
  % within the lecture note (e.g. Section 8.2, Lemma 8.6), as in the book
  \renewcommand{\thesection}{#1.\arabic{section}}
  \renewcommand{\thelemma}{#1.\arabic{lemma}}
  \renewcommand{\thedefinition}{#1.\arabic{definition}}
  \renewcommand{\thetheorem}{#1.\arabic{theorem}}
  \renewcommand{\thecorollary}{#1.\arabic{corollary}}
  \renewcommand{\theclaim}{#1.\arabic{claim}}
  \renewcommand{\theobservation}{#1.\arabic{observation}}
  \renewcommand{\theconjecture}{#1.\arabic{conjecture}}
  \renewcommand{\theexercise}{#1.\arabic{exercise}}
  \renewcommand{\thefigure}{#1.\arabic{figure}}
  \renewcommand{\thetable}{#1.\arabic{table}}
  \renewcommand{\theequation}{#1.\arabic{equation}}
  \LNtocWidths
  \title{Lecture Note #1\\ #2}
  % a note's \LNabstract, if it defines one before \doTitle, goes right below the title
  \date{\ifdefined\LNabstract\parbox{0.85\textwidth}{\normalsize\normalfont\LNabstract}\\[3ex]\fi
    draft of \today, by \href{\authorURL}{N.~Young}\\[2ex] {\small\copyrightNotice}}
  \maketitle
  {\small
    \hypertarget{TOC}{}
    \tableofcontents
  }
  \thispagestyle{empty}
  % the text starts on page 1, after the title and contents (numbered i, ii, ...)
  \clearpage
  \pagenumbering{arabic}
}

\def\clap#1{\hbox to 0pt{\hss#1\hss}}

\newcommand{\refreshHeader}{}
\newcommand{\header}{}

\newcommand{\doHeader}[1]{      % for lecture notes
  \renewcommand{\header}{%
    \clap{\raisebox{0pt}{\protect\hyperlink{TOC}{$\uparrow$\scriptsize\,\textsc{contents}\hspace{1.1in}}}}%
    {\footnotesize \LNTitle{} / Section \thesection, #1}} %, #1, draft of \today}}
  \markboth{\header}{\header}
}

\newcommand{\SID}{}
\newcommand{\setAuthor}[1]{\renewcommand{\author}{#1}}
\newcommand{\setSID}[1]{\renewcommand{\SID}{#1}}

\newcommand{\doHwHeader}[1]{                                         % for homeworks
  \renewcommand{\refreshHeader}{
    \renewcommand{\header}{%
      {\footnotesize \classSection #1 / Problem~\arabic{problem} / draft of \today \hfill \author\ifGradescope~(\SID)\fi~~}}
    \markboth{\header}{\header}
  }
  \renewcommand{\header}{%
    {\footnotesize \classSection #1 / Instructions / draft of \today \hfill \author\ifGradescope~(\SID)\fi~~}}
  \markboth{\header}{\header}
}

% Homework settings, for instructors using the homeworks in a course:
% \classSection is the course name shown in each homework's running head
% (e.g. \renewcommand{\classSection}{CS 3000\xspace}), and \Gradescopetrue turns on
% the instructions for submitting to Gradescope (in its fixed-length mode),
% the hints saying on which page of the PDF each answer should be, and the
% student ID in the running head.
\newcommand{\classSection}{}
\newif\ifGradescope
\Gradescopefalse
\newcommand{\GradescopeComment}[1]{\ifGradescope\comment{#1}\fi}

\pagestyle{myheadings}
\addtolength{\headsep}{0.2in}
\addtolength{\topmargin}{-0.3in}
\addtolength{\textheight}{0.4in}

%%%%%%%%%%%%%%%%%%%%%%%%%%%%%%%%%%%%%%%%%%%%%%%%
%%%%%%%%%%%%%%%%%%%%%%%%%%%%%%%%%%%%%%%%%%%%%%%%

\newcommand{\MF}{\href
  {https://mfleck.cs.illinois.edu/building-blocks/}
  {MF}\xspace}

\newcommand{\KWSlides}{\href
  {https://www.cs.princeton.edu/~wayne/kleinberg-tardos/}
  {KW slides}\xspace}

\newcommand{\LLM}{\href
  {https://courses.csail.mit.edu/6.042/spring18/mcs.pdf}
  {LLM}\xspace}

% \DPV, \KT, \CLRS: textbooks listed (without links) in the table of sources in
% the Preliminaries note; each prints as a link to its entry in that table.
\DeclareRobustCommand{\DPV}{\noteLink{prelim}{source.DPV}{DPV}\xspace}
\DeclareRobustCommand{\KT}{\noteLink{prelim}{source.KT}{KT}\xspace}
\DeclareRobustCommand{\CLRS}{\noteLink{prelim}{source.CLRS}{CLRS}\xspace}
\AtBeginDocument{\pdfstringdefDisableCommands{\def\DPV{DPV }\def\KT{KT }\def\CLRS{CLRS }}}

\newcommand{\JE}{\href
  {https://jeffe.cs.illinois.edu/teaching/algorithms/}
  {JE}\xspace}

\newcommand{\JEGreedy}{\href
  {https://jeffe.cs.illinois.edu/teaching/algorithms/book/04-greedy.pdf}
  {JE Chapter 4}\xspace}

\newcommand{\BBCode}{\href
{https://northeastern.blackboard.com/webapps/cmsmain/webui/courses/CS3000.MERGED.202010/code?action=frameset&subaction=view&uniq=-jj2max&course_id=_2597392_1&mask=\%2Fcourses\%2FCS3000.MERGED.202010}
  {code}\xspace}

% The Python code (src/code).  Each lecture note that has code ends with a
% section listing its files: \codeSection{path, path, ...}.  In the text,
% \pycode{greedy_algorithms/Dijkstra.py} prints the path as a link to its listing
% (in the homeworks, which have no listings, to the file on GitHub), and
% \pycode[text]{...} uses the given text instead.  \pythonCode{path, ...} is the
% standard sentence ending a section whose algorithms have code:
% "Python code: path, ... (Section N.k)."
\newcommand{\codeFolder}{https://github.com/nealeyoung/algs101/tree/main/src/code}
\newcommand{\codeRepo}{https://github.com/nealeyoung/algs101/blob/main/src/code}
\newif\ifcodeListings   % set by \doTitle: this document has the listings
\newcommand{\githubCode}[2][]{\href{\codeRepo/#2}{\ifx\relax#1\relax\nolinkurl{#2}\else#1\fi}}
\newcommand{\pycode}[2][]{%
  \ifcodeListings\hyperref[code: #2]{\ifx\relax#1\relax\nolinkurl{#2}\else#1\fi}%
  \else\githubCode[#1]{#2}\fi}
\makeatletter
\newcommand{\pythonCode}[1]{%
  \def\pythonCode@first{}\def\pythonCode@list{}%
  \forcsvlist{\pythonCode@add}{#1}%
  Python code: \pythonCode@list\ (Section~\ref{code: \pythonCode@first}).}
\newcommand{\pythonCode@add}[1]{%
  \ifx\pythonCode@first\@empty\def\pythonCode@first{#1}\appto\pythonCode@list{\pycode{#1}}%
  \else\appto\pythonCode@list{, \pycode{#1}}\fi}
\makeatother

% Code listings that copy and paste exactly (in PDF viewers that honor the
% ToUnicode maps, e.g. MuPDF-based ones; PDFKit viewers such as Preview and
% Skim drop leading spaces, so each listing's heading links to the file on
% GitHub).  Spaces are typeset as white visible-space glyphs that copy as
% spaces, and hyphens and asterisks as the T1 glyphs, so they copy as themselves.
\usepackage[T1,OT1]{fontenc}   % T1 for the code listings only; OT1 stays the default
\usepackage{listings}
\pdfgentounicode=1
\input glyphtounicode
\pdfglyphtounicode{uni2423}{0020}
\makeatletter
\def\lst@visiblespace{\textcolor{white}{\char32}}
% put an invisible space on each blank line, so that blank lines copy too
% (listings' OnEmptyLine hook runs before the line break, so patch the loop)
\def\lst@DoNewLines{%
    \@whilenum\lst@newlines>\@ne \do {\lst@NewLine \lst@visiblespace}%
    \ifnum\lst@newlines>\z@ \lst@NewLine \fi}
\makeatother
\lstset{language=Python,
  basicstyle=\fontencoding{T1}\fontfamily{lmtt}\selectfont\footnotesize,
  upquote=true, showstringspaces=false, columns=fullflexible, keepspaces=true,
  showspaces=true, literate={-}{{\char45}}1 {*}{{\char42}}1,
  frame=single, framesep=1ex, rulecolor=\color{gray}}
% the listings' paths are relative to the lecture note's folder (in the book, to src/)
\ifdefined\BookMode\newcommand{\codeDir}{code/}\else\newcommand{\codeDir}{../code/}\fi
% \codefile{path}: a listing of src/code/path, headed by a link to it on GitHub
\newcommand{\codefile}[1]{%
  \subsection*{\githubCode{#1}}\label{code: #1}%
  \lstinputlisting{\codeDir#1}}
% \codeSection{path, path, ...}: the lecture note's last section, listing its code
\newcommand{\codeSection}[1]{%
  \newpage
  \section{Python code}\label{\LNtag{\currentLN}: sec: code}\doHeader{Python code}
  This section lists the Python code for this lecture note.
  Click on a file's name to open the file on GitHub.
  Copying code from GitHub is more reliable than copying it from this PDF,
  since some PDF viewers drop the leading spaces (the indentation) of copied lines.
  \forcsvlist{\codefile}{#1}}
% the copyright notice, for title pages and the homework first pages
\newcommand{\licenseURL}{https://github.com/nealeyoung/algs101/blob/main/LICENSE}
\newcommand{\authorURL}{https://www.cs.ucr.edu/\string~neal}
\newcommand{\copyrightNotice}{%
  \copyright\ 2018--2026 \href{\authorURL}{Neal E. Young}.
  Text: CC BY-NC-SA 4.0; code: MIT; except where noted.
  See \url{\licenseURL}.}
% the published site: PDFs of the notes, and the homework assignments and templates
\newcommand{\siteURL}{https://nealeyoung.github.io/algs101/}

\newcommand{\blackboard}{\href
  {https://northeastern.blackboard.com/webapps/blackboard/execute/launcher?type=Course&id=_2597392_1&url=}
  {Blackboard}\xspace}

\newcommand{\JUSlides}{\href
  {http://www.khoury.neu.edu/home/jullman/cs3000f18/schedule.html}
  {Prof.~Ullman's slides}\xspace}

% a URL on its own line, in three layouts
\newcommand{\URL}[1]{\par-- \parbox[t]{0.8\textwidth}{\url{#1}} \medskip\par}
\newcommand{\URLsmall}[1]{\par-- \parbox[t]{0.95\textwidth}{\small\url{#1}} \medskip\par}
\newcommand{\URLplain}[1]{\par\parbox[t]{\textwidth}{\url{#1}}\par}

%%%%%%%%%%%%%%%%%%%%%%%%%%%%%%%%%%%%%%%%%%%%%%%%
%%%%%%%%%%%%%%%%%%%%%%%%%%%%%%%%%%%%%%%%%%%%%%%%

\newcommand{\continued}{\vfill\centerline{\emph{(continued)}}\vfill\newpage}

% load hyperref last to avoid conflicts; long URLs may also break after hyphens
\PassOptionsToPackage{hyphens}{url}
\ifdefined\BookMode\else\usepackage{xr-hyper}\fi   % for \lectureNoteRefs (homeworks), below
\usepackage{hyperref}
\hypersetup{
    colorlinks,
    linkcolor={blue!50!black},
    citecolor={blue!50!black},
    urlcolor={blue!80!black}
}

\providecommand{\answerDirectory}{.}
\newcommand{\inputAnswer}[1]{\refreshHeader\input{\answerDirectory/#1}}
\newcommand{\inputSection}[1]{\input{#1}}

% Refer to a lecture note by its tag (the prefix of its labels), e.g.
% \LN{proofs} prints "Lecture Note 2" (in the book, a link to that chapter).  The table below is the only place the numbers appear.
\makeatletter
\newcommand{\defineLN}[3]{\@namedef{LNnumber@#2}{#1}\@namedef{LNtag@#1}{#2}\@namedef{LNfile@#2}{#3}}
\newcommand{\LNnumber}[1]{%
  \ifcsname LNnumber@#1\endcsname\csname LNnumber@#1\endcsname
  \else\PackageError{macros}{Unknown lecture-note tag `#1'}{}\fi}
\newcommand{\LNtag}[1]{\csname LNtag@#1\endcsname}
\newcommand{\LNfile}[1]{\csname LNfile@#1\endcsname}
\makeatother
\defineLN{1}{prelim}{1_preliminaries}
\defineLN{2}{proofs}{2_proofs}
\defineLN{3}{matching}{3_stable_matching}
\defineLN{4}{dc}{4_divide_and_conquer}
\defineLN{5}{greedy}{5_greedy}
\defineLN{6}{graphs}{6_graph_traversal}
\defineLN{7}{dp}{7_dynamic_programming}
\defineLN{8}{flow}{8_flow_reductions_LP}
\defineLN{9}{np}{9_NP}

% The text is the same in the book and in the standalone lecture notes; only
% links differ.  In the book, links to other lecture notes stay in the book.
% In a standalone note, they open the other note's published PDF.
% \noteLink{tag}{target}{text} links text to \hypertarget{target} in lecture note tag.
% \LN{proofs} prints "Lecture Note 2", linked to that lecture note.
\newcommand{\currentLN}{0}   % in a standalone note, its number (set by \doTitle)
\newcommand{\notePDF}[1]{\siteURL notes/\LNfile{#1}.pdf}
\newcommand{\noteLink}[3]{%
  \ifnum\LNnumber{#1}=\currentLN\relax\hyperlink{#2}{#3}%
  \else\href{\notePDF{#1}\#nameddest=#2}{#3}\fi}
% \bookLink{target}{text} links text to \hypertarget{target} in the book's
% appendices (from a standalone note, in the published book).
\DeclareRobustCommand{\bookLink}[2]{%
  \ifdefined\BookMode\hyperlink{#1}{#2}%
  \else\href{\siteURL lecture_notes_on_algorithms.pdf\#nameddest=#1}{#2}\fi}
\DeclareRobustCommand{\LN}[1]{%
  \ifnum\LNnumber{#1}=\currentLN\relax Lecture Note~\LNnumber{#1}%
  \else\href{\notePDF{#1}}{Lecture Note~\LNnumber{#1}}\fi}
\pdfstringdefDisableCommands{\def\LN#1{Lecture Note \LNnumber{#1}}}

% Homeworks refer to the notes' sections, exercises and lemmas by \ref: each
% homework's main file says \lectureNoteRefs after \input{macros}.  A homework
% folder has a link refs/ to src/refs/, which holds each note's label records
% (made by tools/refs.sh from the notes' builds).  Read with xr-hyper, they make
% \ref{greedy: exercise: ...} give the exercise's number in the published note,
% linked to it there; labels only the book has (its appendices) link to the book.
\ifdefined\BookMode\else
  \makeatletter
  \newcommand{\lectureNoteRefs}{%
    \@for\LN@tag:=prelim,proofs,matching,dc,greedy,graphs,dp,flow,np\do{%
      \externaldocument{refs/\LNfile{\LN@tag}}[\notePDF{\LN@tag}]}%
    \externaldocument{refs/book}[\siteURL lecture_notes_on_algorithms.pdf]}
  % hyperref links a \ref to another document's label as URL#anchor; browsers'
  % PDF viewers need URL#nameddest=anchor (the form \noteLink and \bookLink use)
  \def\Hy@setref@link#1#2#3#4#5#6\@nil#7{%
    \begingroup
      \ifx\\#5\\\toks0={\hyper@@link{#5}{#4}}\else\toks0={\hyper@@link{#5}{nameddest=#4}}\fi
      \toks1=\expandafter{#7{#1}{#2}{#3}{#4}{#5}}%
      \edef\x{\endgroup\the\toks0 {\the\toks1 }}%
    \x}
  \makeatother
\fi

% Book mode: main.tex defines \BookMode before \input{macros} to combine the
% lecture notes into one report-class document, one chapter per note.
% Without \BookMode (each lecture note on its own, and the homeworks) none
% of this applies.
\ifdefined\BookMode
  \usepackage{etoc}

  % links to other lecture notes stay in the book
  \renewcommand{\noteLink}[3]{\hyperlink{#2}{#3}}

  % each chapter is a lecture note, and is called one, as in the standalone notes
  \renewcommand{\chaptername}{Lecture Note}
  \setcounter{secnumdepth}{3}   % number subsubsections, as in the standalone notes

  % global table of contents, for main.tex
  \newcommand{\bookContents}{%
    \hypertarget{TOC}{}%
    \etocsetnexttocdepth{section}%
    \tableofcontents}

  % \doTitle{N}{title}{short title} starts chapter N, with a link to the global
  % table of contents and the chapter's own table of contents
  \renewcommand{\doTitle}[3]{%
    \setcounter{chapter}{\numexpr#1-1\relax}%
    \chapter{\texorpdfstring{#2}{#3}}%
    \label{\LNtag{#1}: chapter}%
    \csgdef{LNshorttitle@#1}{#3}%
    \csgdef{LNtitle@#1}{#2}%
    \renewcommand{\LNTitle}{Lecture Note #1, #3}%
    \codeListingstrue
    \hyperlink{TOC}{$\uparrow$\scriptsize\,\textsc{contents}}%
    {\small
      \hypertarget{TOC-#1}{}%
      \etocsettocstyle{\section*{\contentsname}}{}%
      \etocsetnexttocdepth{subsubsection}%
      \localtableofcontents
    }}

  % the code: each lecture note's \codeSection just records its files, and
  % \codeAppendix (in the book's appendix) lists them all, a section per note
  \makeatletter   % (file names have no spaces, so drop the spaces around them)
  \renewcommand{\codeSection}[1]{\csxdef{LNcode@\arabic{chapter}}{\zap@space#1 \@empty}}
  \makeatother
  % \codeAppendix: first a contents list (each note's topic, then its files,
  % linked to the listings), then a section per note listing its files
  \makeatletter
  \def\LNbasename#1{\LN@base#1/\LN@end}   % a/b/c.py -> c.py
  \def\LN@base#1/#2\LN@end{\ifx\relax#2\relax#1\else\LN@base#2\LN@end\fi}
  \makeatother
  \newcommand{\codeContentsFile}[1]{%
    \edef\LNfilename{\LNbasename{#1}}%
    \item \pycode[\texttt{\expandafter\detokenize\expandafter{\LNfilename}}]{#1}}
  \newcommand{\codeAppendix}{%
    \begin{itemize}
    \foreach \n in {1,...,9} {%
      \ifcsdef{LNcode@\n}{%
        \item \hyperref[code: note \n]{\csuse{LNtitle@\n}} (\LN{\LNtag{\n}})
        \begin{itemize}
        \edef\LNcodelist{\csuse{LNcode@\n}}%
        \expandafter\forcsvlist\expandafter\codeContentsFile\expandafter{\LNcodelist}%
        \end{itemize}}{}}
    \end{itemize}
    \newpage
    \foreach \n in {1,...,9} {%
      \ifcsdef{LNcode@\n}{%
        \section[Lecture Note \n: \csuse{LNtitle@\n}]{%
          \texorpdfstring{\LN{\LNtag{\n}}: \csuse{LNtitle@\n}}{Lecture Note \n: \csuse{LNtitle@\n}}}%
        \label{code: note \n}%
        \edef\LNcodelist{\csuse{LNcode@\n}}%
        \expandafter\forcsvlist\expandafter\codefile\expandafter{\LNcodelist}%
        \newpage}{}}}

  % a lecture note is a chapter: \LN{proofs} prints a link "Lecture Note 2"
  \DeclareRobustCommand{\LN}[1]{\hyperref[#1: chapter]{Lecture Note~\ref*{#1: chapter}}}
  \pdfstringdefDisableCommands{\def\LN#1{Lecture Note \LNnumber{#1}}}

  % running head; "contents" links to the current chapter's table of contents
  \renewcommand{\doHeader}[1]{%
    \renewcommand{\header}{%
      \clap{\raisebox{0pt}{\protect\hyperlink{TOC-\arabic{chapter}}{$\uparrow$\scriptsize\,\textsc{contents}\hspace{1.1in}}}}%
      {\footnotesize \LNTitle{} / Section \thesection, #1}}%
    \markboth{\header}{\header}}

  % number theorems, lemmas, etc. and exercises within chapters (e.g. Lemma 8.1)
  \counterwithin{lemma}{chapter}
  \counterwithin{definition}{chapter}
  \counterwithin{theorem}{chapter}
  \counterwithin{corollary}{chapter}
  \counterwithin{claim}{chapter}
  \counterwithin{observation}{chapter}
  \counterwithin{conjecture}{chapter}
  \counterwithin{exercise}{chapter}
\fi

% %%%%%%%%%%%%%%%%%%%%%%%%%%%%%%%%%%%%%%%%%%%%%%%%%%%%%%%%%%%%%%%%%%%%%%%%
% %%%%%%%%%%%%%%%%%%%%%%%%%%%%%%%%%%%%%%%%%%%%%%%%%%%%%%%%%%%%%%%%%%%%%%%%

\ExplSyntaxOn
\seq_new:N \l_itemize_seq
% \NewDocumentEnvironment{REITEMIZE}{ +b }{
\NewDocumentCommand{\REITEMIZECMD}{ +m }{
  \regex_split:nnN {\c{item}} {#1} \l_itemize_seq
  \seq_pop:NN \l_itemize_seq \l_tmpa_tl% remove the (hopefully) empty first item
  \seq_if_empty:NF \l_itemize_seq {
    \l_tmpa_tl
    \seq_map_inline:Nn \l_itemize_seq {
      \regex_split:nnN {\c{item}} {#1} \l_itemize_seq {
        \regex_extract_once:nnNTF { ^\[([^]]*)\](.*) } { ##1} \l_foo_seq
      { \item[\seq_item:Nn \l_foo_seq {2}]\ITEMIZER{\seq_item:Nn \l_foo_seq {3}} } { \item\ITEMIZER{##1}  }
    }
  }
}
}
\ExplSyntaxOff

\newlength{\ITEMHT}
\setlength{\ITEMHT}{1in}

% \setlength{\fboxsep}{1pt}
\newcommand{\ITEMIZER}[1]{%
  % \framebox{
\adjustbox{valign=t}{\parbox{\textwidth}{\linespread{1.15}\selectfont #1}}~~\adjustbox{valign=t}{\parbox{0pt}{\color{gray}\rule[-\ITEMHT]{0.01pt}{\ITEMHT}}}
% }
  % \begin{tabular}{@{}p{\textwidth}@{\rule[-\ITEMHT]{0pt}{\ITEMHT}}}
  %   #1
  % \end{tabular}
}

\newsavebox\thesmashminipage
\newlength{\thesmashminipagelength}
\newlength{\answerboxlength}
\newlength{\answerboxoverrun}
\newenvironment{smashminipage}[2]
{\setlength\answerboxlength{#2}
\begin{lrbox}{\thesmashminipage}\begin{minipage}[t]{#1}}
{\end{minipage}\end{lrbox}%
\setlength{\thesmashminipagelength}{\dimexpr\ht\thesmashminipage+\dp\thesmashminipage}%
\setlength{\answerboxoverrun}{\dimexpr\thesmashminipagelength-\answerboxlength}%
\ifdim\the\thesmashminipagelength>\the\answerboxlength%
\usebox{\thesmashminipage}%
\smash{\color{red}\hspace*{-0.98\textwidth}\rule[-\answerboxlength]{0.98\textwidth}{1pt}}
\errmessage{Answer exceeds allotted length, please shorten it.}
\else
\usebox{\thesmashminipage}%
\fi
}



\newenvironment{answerBox}[1]{%
  \def\ANSWERBOXHEIGHT{#1}
  \begin{adjustbox}{valign=t}
  \begin{smashminipage}{\textwidth}{\ANSWERBOXHEIGHT}
}{\end{smashminipage}\end{adjustbox}~~~\adjustbox{valign=t}{\parbox{0pt}{\color{gray}\rule[-\ANSWERBOXHEIGHT]{0.01pt}{\ANSWERBOXHEIGHT}\rule[-\ANSWERBOXHEIGHT]{10pt}{0.01pt}}}}

% % \setlist[enumerate]{itemsep=0pt,parsep=0pt,topsep=0pt,leftmargin=100pt}

% \makeatletter 
\LetLtxMacro\SVenumerate\enumerate
\let\endSVenumerate\endenumerate

\newcommand{\injectEnumerate}{
  \let\SValph\alph
  \def\alph{}
  \renewenvironment{enumerate}[1][]
  {\let\alph\SValph \SVenumerate[label=(\alph*),itemsep=2pt,parsep=0pt]\REITEMIZE}{\endREITEMIZE\endSVenumerate}
  % \setlist[enumerate]{before=\Collect@Body\REITEMIZECMD}
}
% \makeatother
\newcommand{\restoreEnumerate}{
  \renewenvironment{enumerate}[1][]
  \let\SValph\alph
  \def\alph{}
  {\let\alph\SValph \SVenumerate[label=(\alph*)]}{\endSVenumerate}
  % \setlist[enumerate]{before=,after=}
}
\newcommand{\injectBlue}{
  \let\SVblue\blue
  \let\endSVblue\endblue
  \def\blue{\SVblue\REITEMIZE}
  \def\endblue{\endREITEMIZE\endSVblue}
}
\newcommand{\restoreBlue}{
  \let\blue\SVblue
  \let\endblue\endSVblue
}
\newcommand{\injectEnumerateBlue}{
  \setlist[enumerate]{before=\injectBlue,after=\restoreBlue,itemsep=2pt,parsep=0pt}
}
\newcommand{\restoreEnumerateBlue}{
  \setlist[enumerate]{before=,after=}
}


% %%%%%%%%%%%%%%%%%%%%%%%%%%%%%%%%%%%%%%%%%%%%%%%%%%%%%%%%%%%%%%%%%%%%%%%%
% %%%%%%%%%%%%%%%%%%%%%%%%%%%%%%%%%%%%%%%%%%%%%%%%%%%%%%%%%%%%%%%%%%%%%%%%

% \NewDocumentEnvironment{env}{ }
% {\Environenv}
% {\endEnvironenv}

% \NewEnviron{REITEMIZE}{\REITEMIZECMD\BODY}

\makeatletter 
\newenvironment{REITEMIZE}{\Collect@Body\REITEMIZECMD}{}
\makeatother 

\newcommand{\NODEone}[1]{\ensuremath{v_{#1}}\xspace}
\newcommand{\NODE}[2]{{\NODEone {#1,#2}}\xspace}
\newcommand{\EDGE}[2]{\ensuremath{(\textsf{#1}, \textsf{#2})}}
\newcommand{\ITEM}{\item}

%%% Local Variables:
%%% mode: latex
%%% End:
 to combine the
% lecture notes into one report-class document, one chapter per note.
% Without \BookMode (each lecture note on its own, and the homeworks) none
% of this applies.
\ifdefined\BookMode
  \usepackage{etoc}

  % links to other lecture notes stay in the book
  \renewcommand{\noteLink}[3]{\hyperlink{#2}{#3}}

  % each chapter is a lecture note, and is called one, as in the standalone notes
  \renewcommand{\chaptername}{Lecture Note}
  \setcounter{secnumdepth}{3}   % number subsubsections, as in the standalone notes

  % global table of contents, for main.tex
  \newcommand{\bookContents}{%
    \hypertarget{TOC}{}%
    \etocsetnexttocdepth{section}%
    \tableofcontents}

  % \doTitle{N}{title}{short title} starts chapter N, with a link to the global
  % table of contents and the chapter's own table of contents
  \renewcommand{\doTitle}[3]{%
    \setcounter{chapter}{\numexpr#1-1\relax}%
    \chapter{\texorpdfstring{#2}{#3}}%
    \label{\LNtag{#1}: chapter}%
    \csgdef{LNshorttitle@#1}{#3}%
    \csgdef{LNtitle@#1}{#2}%
    \renewcommand{\LNTitle}{Lecture Note #1, #3}%
    \codeListingstrue
    \hyperlink{TOC}{$\uparrow$\scriptsize\,\textsc{contents}}%
    {\small
      \hypertarget{TOC-#1}{}%
      \etocsettocstyle{\section*{\contentsname}}{}%
      \etocsetnexttocdepth{subsubsection}%
      \localtableofcontents
    }}

  % the code: each lecture note's \codeSection just records its files, and
  % \codeAppendix (in the book's appendix) lists them all, a section per note
  \makeatletter   % (file names have no spaces, so drop the spaces around them)
  \renewcommand{\codeSection}[1]{\csxdef{LNcode@\arabic{chapter}}{\zap@space#1 \@empty}}
  \makeatother
  % \codeAppendix: first a contents list (each note's topic, then its files,
  % linked to the listings), then a section per note listing its files
  \makeatletter
  \def\LNbasename#1{\LN@base#1/\LN@end}   % a/b/c.py -> c.py
  \def\LN@base#1/#2\LN@end{\ifx\relax#2\relax#1\else\LN@base#2\LN@end\fi}
  \makeatother
  \newcommand{\codeContentsFile}[1]{%
    \edef\LNfilename{\LNbasename{#1}}%
    \item \pycode[\texttt{\expandafter\detokenize\expandafter{\LNfilename}}]{#1}}
  \newcommand{\codeAppendix}{%
    \begin{itemize}
    \foreach \n in {1,...,9} {%
      \ifcsdef{LNcode@\n}{%
        \item \hyperref[code: note \n]{\csuse{LNtitle@\n}} (\LN{\LNtag{\n}})
        \begin{itemize}
        \edef\LNcodelist{\csuse{LNcode@\n}}%
        \expandafter\forcsvlist\expandafter\codeContentsFile\expandafter{\LNcodelist}%
        \end{itemize}}{}}
    \end{itemize}
    \newpage
    \foreach \n in {1,...,9} {%
      \ifcsdef{LNcode@\n}{%
        \section[Lecture Note \n: \csuse{LNtitle@\n}]{%
          \texorpdfstring{\LN{\LNtag{\n}}: \csuse{LNtitle@\n}}{Lecture Note \n: \csuse{LNtitle@\n}}}%
        \label{code: note \n}%
        \edef\LNcodelist{\csuse{LNcode@\n}}%
        \expandafter\forcsvlist\expandafter\codefile\expandafter{\LNcodelist}%
        \newpage}{}}}

  % a lecture note is a chapter: \LN{proofs} prints a link "Lecture Note 2"
  \DeclareRobustCommand{\LN}[1]{\hyperref[#1: chapter]{Lecture Note~\ref*{#1: chapter}}}
  \pdfstringdefDisableCommands{\def\LN#1{Lecture Note \LNnumber{#1}}}

  % running head; "contents" links to the current chapter's table of contents
  \renewcommand{\doHeader}[1]{%
    \renewcommand{\header}{%
      \clap{\raisebox{0pt}{\protect\hyperlink{TOC-\arabic{chapter}}{$\uparrow$\scriptsize\,\textsc{contents}\hspace{1.1in}}}}%
      {\footnotesize \LNTitle{} / Section \thesection, #1}}%
    \markboth{\header}{\header}}

  % number theorems, lemmas, etc. and exercises within chapters (e.g. Lemma 8.1)
  \counterwithin{lemma}{chapter}
  \counterwithin{definition}{chapter}
  \counterwithin{theorem}{chapter}
  \counterwithin{corollary}{chapter}
  \counterwithin{claim}{chapter}
  \counterwithin{observation}{chapter}
  \counterwithin{conjecture}{chapter}
  \counterwithin{exercise}{chapter}
\fi

% %%%%%%%%%%%%%%%%%%%%%%%%%%%%%%%%%%%%%%%%%%%%%%%%%%%%%%%%%%%%%%%%%%%%%%%%
% %%%%%%%%%%%%%%%%%%%%%%%%%%%%%%%%%%%%%%%%%%%%%%%%%%%%%%%%%%%%%%%%%%%%%%%%

\ExplSyntaxOn
\seq_new:N \l_itemize_seq
% \NewDocumentEnvironment{REITEMIZE}{ +b }{
\NewDocumentCommand{\REITEMIZECMD}{ +m }{
  \regex_split:nnN {\c{item}} {#1} \l_itemize_seq
  \seq_pop:NN \l_itemize_seq \l_tmpa_tl% remove the (hopefully) empty first item
  \seq_if_empty:NF \l_itemize_seq {
    \l_tmpa_tl
    \seq_map_inline:Nn \l_itemize_seq {
      \regex_split:nnN {\c{item}} {#1} \l_itemize_seq {
        \regex_extract_once:nnNTF { ^\[([^]]*)\](.*) } { ##1} \l_foo_seq
      { \item[\seq_item:Nn \l_foo_seq {2}]\ITEMIZER{\seq_item:Nn \l_foo_seq {3}} } { \item\ITEMIZER{##1}  }
    }
  }
}
}
\ExplSyntaxOff

\newlength{\ITEMHT}
\setlength{\ITEMHT}{1in}

% \setlength{\fboxsep}{1pt}
\newcommand{\ITEMIZER}[1]{%
  % \framebox{
\adjustbox{valign=t}{\parbox{\textwidth}{\linespread{1.15}\selectfont #1}}~~\adjustbox{valign=t}{\parbox{0pt}{\color{gray}\rule[-\ITEMHT]{0.01pt}{\ITEMHT}}}
% }
  % \begin{tabular}{@{}p{\textwidth}@{\rule[-\ITEMHT]{0pt}{\ITEMHT}}}
  %   #1
  % \end{tabular}
}

\newsavebox\thesmashminipage
\newlength{\thesmashminipagelength}
\newlength{\answerboxlength}
\newlength{\answerboxoverrun}
\newenvironment{smashminipage}[2]
{\setlength\answerboxlength{#2}
\begin{lrbox}{\thesmashminipage}\begin{minipage}[t]{#1}}
{\end{minipage}\end{lrbox}%
\setlength{\thesmashminipagelength}{\dimexpr\ht\thesmashminipage+\dp\thesmashminipage}%
\setlength{\answerboxoverrun}{\dimexpr\thesmashminipagelength-\answerboxlength}%
\ifdim\the\thesmashminipagelength>\the\answerboxlength%
\usebox{\thesmashminipage}%
\smash{\color{red}\hspace*{-0.98\textwidth}\rule[-\answerboxlength]{0.98\textwidth}{1pt}}
\errmessage{Answer exceeds allotted length, please shorten it.}
\else
\usebox{\thesmashminipage}%
\fi
}



\newenvironment{answerBox}[1]{%
  \def\ANSWERBOXHEIGHT{#1}
  \begin{adjustbox}{valign=t}
  \begin{smashminipage}{\textwidth}{\ANSWERBOXHEIGHT}
}{\end{smashminipage}\end{adjustbox}~~~\adjustbox{valign=t}{\parbox{0pt}{\color{gray}\rule[-\ANSWERBOXHEIGHT]{0.01pt}{\ANSWERBOXHEIGHT}\rule[-\ANSWERBOXHEIGHT]{10pt}{0.01pt}}}}

% % \setlist[enumerate]{itemsep=0pt,parsep=0pt,topsep=0pt,leftmargin=100pt}

% \makeatletter 
\LetLtxMacro\SVenumerate\enumerate
\let\endSVenumerate\endenumerate

\newcommand{\injectEnumerate}{
  \let\SValph\alph
  \def\alph{}
  \renewenvironment{enumerate}[1][]
  {\let\alph\SValph \SVenumerate[label=(\alph*),itemsep=2pt,parsep=0pt]\REITEMIZE}{\endREITEMIZE\endSVenumerate}
  % \setlist[enumerate]{before=\Collect@Body\REITEMIZECMD}
}
% \makeatother
\newcommand{\restoreEnumerate}{
  \renewenvironment{enumerate}[1][]
  \let\SValph\alph
  \def\alph{}
  {\let\alph\SValph \SVenumerate[label=(\alph*)]}{\endSVenumerate}
  % \setlist[enumerate]{before=,after=}
}
\newcommand{\injectBlue}{
  \let\SVblue\blue
  \let\endSVblue\endblue
  \def\blue{\SVblue\REITEMIZE}
  \def\endblue{\endREITEMIZE\endSVblue}
}
\newcommand{\restoreBlue}{
  \let\blue\SVblue
  \let\endblue\endSVblue
}
\newcommand{\injectEnumerateBlue}{
  \setlist[enumerate]{before=\injectBlue,after=\restoreBlue,itemsep=2pt,parsep=0pt}
}
\newcommand{\restoreEnumerateBlue}{
  \setlist[enumerate]{before=,after=}
}


% %%%%%%%%%%%%%%%%%%%%%%%%%%%%%%%%%%%%%%%%%%%%%%%%%%%%%%%%%%%%%%%%%%%%%%%%
% %%%%%%%%%%%%%%%%%%%%%%%%%%%%%%%%%%%%%%%%%%%%%%%%%%%%%%%%%%%%%%%%%%%%%%%%

% \NewDocumentEnvironment{env}{ }
% {\Environenv}
% {\endEnvironenv}

% \NewEnviron{REITEMIZE}{\REITEMIZECMD\BODY}

\makeatletter 
\newenvironment{REITEMIZE}{\Collect@Body\REITEMIZECMD}{}
\makeatother 

\newcommand{\NODEone}[1]{\ensuremath{v_{#1}}\xspace}
\newcommand{\NODE}[2]{{\NODEone {#1,#2}}\xspace}
\newcommand{\EDGE}[2]{\ensuremath{(\textsf{#1}, \textsf{#2})}}
\newcommand{\ITEM}{\item}

%%% Local Variables:
%%% mode: latex
%%% End:
 to combine the
% lecture notes into one report-class document, one chapter per note.
% Without \BookMode (each lecture note on its own, and the homeworks) none
% of this applies.
\ifdefined\BookMode
  \usepackage{etoc}

  % links to other lecture notes stay in the book
  \renewcommand{\noteLink}[3]{\hyperlink{#2}{#3}}

  % each chapter is a lecture note, and is called one, as in the standalone notes
  \renewcommand{\chaptername}{Lecture Note}
  \setcounter{secnumdepth}{3}   % number subsubsections, as in the standalone notes

  % global table of contents, for main.tex
  \newcommand{\bookContents}{%
    \hypertarget{TOC}{}%
    \etocsetnexttocdepth{section}%
    \tableofcontents}

  % \doTitle{N}{title}{short title} starts chapter N, with a link to the global
  % table of contents and the chapter's own table of contents
  \renewcommand{\doTitle}[3]{%
    \setcounter{chapter}{\numexpr#1-1\relax}%
    \chapter{\texorpdfstring{#2}{#3}}%
    \label{\LNtag{#1}: chapter}%
    \csgdef{LNshorttitle@#1}{#3}%
    \csgdef{LNtitle@#1}{#2}%
    \renewcommand{\LNTitle}{Lecture Note #1, #3}%
    \codeListingstrue
    \hyperlink{TOC}{$\uparrow$\scriptsize\,\textsc{contents}}%
    {\small
      \hypertarget{TOC-#1}{}%
      \etocsettocstyle{\section*{\contentsname}}{}%
      \etocsetnexttocdepth{subsubsection}%
      \localtableofcontents
    }}

  % the code: each lecture note's \codeSection just records its files, and
  % \codeAppendix (in the book's appendix) lists them all, a section per note
  \makeatletter   % (file names have no spaces, so drop the spaces around them)
  \renewcommand{\codeSection}[1]{\csxdef{LNcode@\arabic{chapter}}{\zap@space#1 \@empty}}
  \makeatother
  % \codeAppendix: first a contents list (each note's topic, then its files,
  % linked to the listings), then a section per note listing its files
  \makeatletter
  \def\LNbasename#1{\LN@base#1/\LN@end}   % a/b/c.py -> c.py
  \def\LN@base#1/#2\LN@end{\ifx\relax#2\relax#1\else\LN@base#2\LN@end\fi}
  \makeatother
  \newcommand{\codeContentsFile}[1]{%
    \edef\LNfilename{\LNbasename{#1}}%
    \item \pycode[\texttt{\expandafter\detokenize\expandafter{\LNfilename}}]{#1}}
  \newcommand{\codeAppendix}{%
    \begin{itemize}
    \foreach \n in {1,...,9} {%
      \ifcsdef{LNcode@\n}{%
        \item \hyperref[code: note \n]{\csuse{LNtitle@\n}} (\LN{\LNtag{\n}})
        \begin{itemize}
        \edef\LNcodelist{\csuse{LNcode@\n}}%
        \expandafter\forcsvlist\expandafter\codeContentsFile\expandafter{\LNcodelist}%
        \end{itemize}}{}}
    \end{itemize}
    \newpage
    \foreach \n in {1,...,9} {%
      \ifcsdef{LNcode@\n}{%
        \section[Lecture Note \n: \csuse{LNtitle@\n}]{%
          \texorpdfstring{\LN{\LNtag{\n}}: \csuse{LNtitle@\n}}{Lecture Note \n: \csuse{LNtitle@\n}}}%
        \label{code: note \n}%
        \edef\LNcodelist{\csuse{LNcode@\n}}%
        \expandafter\forcsvlist\expandafter\codefile\expandafter{\LNcodelist}%
        \newpage}{}}}

  % a lecture note is a chapter: \LN{proofs} prints a link "Lecture Note 2"
  \DeclareRobustCommand{\LN}[1]{\hyperref[#1: chapter]{Lecture Note~\ref*{#1: chapter}}}
  \pdfstringdefDisableCommands{\def\LN#1{Lecture Note \LNnumber{#1}}}

  % running head; "contents" links to the current chapter's table of contents
  \renewcommand{\doHeader}[1]{%
    \renewcommand{\header}{%
      \clap{\raisebox{0pt}{\protect\hyperlink{TOC-\arabic{chapter}}{$\uparrow$\scriptsize\,\textsc{contents}\hspace{1.1in}}}}%
      {\footnotesize \LNTitle{} / Section \thesection, #1}}%
    \markboth{\header}{\header}}

  % number theorems, lemmas, etc. and exercises within chapters (e.g. Lemma 8.1)
  \counterwithin{lemma}{chapter}
  \counterwithin{definition}{chapter}
  \counterwithin{theorem}{chapter}
  \counterwithin{corollary}{chapter}
  \counterwithin{claim}{chapter}
  \counterwithin{observation}{chapter}
  \counterwithin{conjecture}{chapter}
  \counterwithin{exercise}{chapter}
\fi

% %%%%%%%%%%%%%%%%%%%%%%%%%%%%%%%%%%%%%%%%%%%%%%%%%%%%%%%%%%%%%%%%%%%%%%%%
% %%%%%%%%%%%%%%%%%%%%%%%%%%%%%%%%%%%%%%%%%%%%%%%%%%%%%%%%%%%%%%%%%%%%%%%%

\ExplSyntaxOn
\seq_new:N \l_itemize_seq
% \NewDocumentEnvironment{REITEMIZE}{ +b }{
\NewDocumentCommand{\REITEMIZECMD}{ +m }{
  \regex_split:nnN {\c{item}} {#1} \l_itemize_seq
  \seq_pop:NN \l_itemize_seq \l_tmpa_tl% remove the (hopefully) empty first item
  \seq_if_empty:NF \l_itemize_seq {
    \l_tmpa_tl
    \seq_map_inline:Nn \l_itemize_seq {
      \regex_split:nnN {\c{item}} {#1} \l_itemize_seq {
        \regex_extract_once:nnNTF { ^\[([^]]*)\](.*) } { ##1} \l_foo_seq
      { \item[\seq_item:Nn \l_foo_seq {2}]\ITEMIZER{\seq_item:Nn \l_foo_seq {3}} } { \item\ITEMIZER{##1}  }
    }
  }
}
}
\ExplSyntaxOff

\newlength{\ITEMHT}
\setlength{\ITEMHT}{1in}

% \setlength{\fboxsep}{1pt}
\newcommand{\ITEMIZER}[1]{%
  % \framebox{
\adjustbox{valign=t}{\parbox{\textwidth}{\linespread{1.15}\selectfont #1}}~~\adjustbox{valign=t}{\parbox{0pt}{\color{gray}\rule[-\ITEMHT]{0.01pt}{\ITEMHT}}}
% }
  % \begin{tabular}{@{}p{\textwidth}@{\rule[-\ITEMHT]{0pt}{\ITEMHT}}}
  %   #1
  % \end{tabular}
}

\newsavebox\thesmashminipage
\newlength{\thesmashminipagelength}
\newlength{\answerboxlength}
\newlength{\answerboxoverrun}
\newenvironment{smashminipage}[2]
{\setlength\answerboxlength{#2}
\begin{lrbox}{\thesmashminipage}\begin{minipage}[t]{#1}}
{\end{minipage}\end{lrbox}%
\setlength{\thesmashminipagelength}{\dimexpr\ht\thesmashminipage+\dp\thesmashminipage}%
\setlength{\answerboxoverrun}{\dimexpr\thesmashminipagelength-\answerboxlength}%
\ifdim\the\thesmashminipagelength>\the\answerboxlength%
\usebox{\thesmashminipage}%
\smash{\color{red}\hspace*{-0.98\textwidth}\rule[-\answerboxlength]{0.98\textwidth}{1pt}}
\errmessage{Answer exceeds allotted length, please shorten it.}
\else
\usebox{\thesmashminipage}%
\fi
}



\newenvironment{answerBox}[1]{%
  \def\ANSWERBOXHEIGHT{#1}
  \begin{adjustbox}{valign=t}
  \begin{smashminipage}{\textwidth}{\ANSWERBOXHEIGHT}
}{\end{smashminipage}\end{adjustbox}~~~\adjustbox{valign=t}{\parbox{0pt}{\color{gray}\rule[-\ANSWERBOXHEIGHT]{0.01pt}{\ANSWERBOXHEIGHT}\rule[-\ANSWERBOXHEIGHT]{10pt}{0.01pt}}}}

% % \setlist[enumerate]{itemsep=0pt,parsep=0pt,topsep=0pt,leftmargin=100pt}

% \makeatletter 
\LetLtxMacro\SVenumerate\enumerate
\let\endSVenumerate\endenumerate

\newcommand{\injectEnumerate}{
  \let\SValph\alph
  \def\alph{}
  \renewenvironment{enumerate}[1][]
  {\let\alph\SValph \SVenumerate[label=(\alph*),itemsep=2pt,parsep=0pt]\REITEMIZE}{\endREITEMIZE\endSVenumerate}
  % \setlist[enumerate]{before=\Collect@Body\REITEMIZECMD}
}
% \makeatother
\newcommand{\restoreEnumerate}{
  \renewenvironment{enumerate}[1][]
  \let\SValph\alph
  \def\alph{}
  {\let\alph\SValph \SVenumerate[label=(\alph*)]}{\endSVenumerate}
  % \setlist[enumerate]{before=,after=}
}
\newcommand{\injectBlue}{
  \let\SVblue\blue
  \let\endSVblue\endblue
  \def\blue{\SVblue\REITEMIZE}
  \def\endblue{\endREITEMIZE\endSVblue}
}
\newcommand{\restoreBlue}{
  \let\blue\SVblue
  \let\endblue\endSVblue
}
\newcommand{\injectEnumerateBlue}{
  \setlist[enumerate]{before=\injectBlue,after=\restoreBlue,itemsep=2pt,parsep=0pt}
}
\newcommand{\restoreEnumerateBlue}{
  \setlist[enumerate]{before=,after=}
}


% %%%%%%%%%%%%%%%%%%%%%%%%%%%%%%%%%%%%%%%%%%%%%%%%%%%%%%%%%%%%%%%%%%%%%%%%
% %%%%%%%%%%%%%%%%%%%%%%%%%%%%%%%%%%%%%%%%%%%%%%%%%%%%%%%%%%%%%%%%%%%%%%%%

% \NewDocumentEnvironment{env}{ }
% {\Environenv}
% {\endEnvironenv}

% \NewEnviron{REITEMIZE}{\REITEMIZECMD\BODY}

\makeatletter 
\newenvironment{REITEMIZE}{\Collect@Body\REITEMIZECMD}{}
\makeatother 

\newcommand{\NODEone}[1]{\ensuremath{v_{#1}}\xspace}
\newcommand{\NODE}[2]{{\NODEone {#1,#2}}\xspace}
\newcommand{\EDGE}[2]{\ensuremath{(\textsf{#1}, \textsf{#2})}}
\newcommand{\ITEM}{\item}

%%% Local Variables:
%%% mode: latex
%%% End:
 to combine the
% lecture notes into one report-class document, one chapter per note.
% Without \BookMode (each lecture note on its own, and the homeworks) none
% of this applies.
\ifdefined\BookMode
  \usepackage{etoc}

  % links to other lecture notes stay in the book
  \renewcommand{\noteLink}[3]{\hyperlink{#2}{#3}}

  % each chapter is a lecture note, and is called one, as in the standalone notes
  \renewcommand{\chaptername}{Lecture Note}
  \setcounter{secnumdepth}{3}   % number subsubsections, as in the standalone notes

  % global table of contents, for main.tex
  \newcommand{\bookContents}{%
    \hypertarget{TOC}{}%
    \etocsetnexttocdepth{section}%
    \tableofcontents}

  % \doTitle{N}{title}{short title} starts chapter N, with a link to the global
  % table of contents and the chapter's own table of contents
  \renewcommand{\doTitle}[3]{%
    \setcounter{chapter}{\numexpr#1-1\relax}%
    \chapter{\texorpdfstring{#2}{#3}}%
    \label{\LNtag{#1}: chapter}%
    \csgdef{LNshorttitle@#1}{#3}%
    \csgdef{LNtitle@#1}{#2}%
    \renewcommand{\LNTitle}{Lecture Note #1, #3}%
    \codeListingstrue
    \hyperlink{TOC}{$\uparrow$\scriptsize\,\textsc{contents}}%
    {\small
      \hypertarget{TOC-#1}{}%
      \etocsettocstyle{\section*{\contentsname}}{}%
      \etocsetnexttocdepth{subsubsection}%
      \localtableofcontents
    }}

  % the code: each lecture note's \codeSection just records its files, and
  % \codeAppendix (in the book's appendix) lists them all, a section per note
  \makeatletter   % (file names have no spaces, so drop the spaces around them)
  \renewcommand{\codeSection}[1]{\csxdef{LNcode@\arabic{chapter}}{\zap@space#1 \@empty}}
  \makeatother
  % \codeAppendix: first a contents list (each note's topic, then its files,
  % linked to the listings), then a section per note listing its files
  \makeatletter
  \def\LNbasename#1{\LN@base#1/\LN@end}   % a/b/c.py -> c.py
  \def\LN@base#1/#2\LN@end{\ifx\relax#2\relax#1\else\LN@base#2\LN@end\fi}
  \makeatother
  \newcommand{\codeContentsFile}[1]{%
    \edef\LNfilename{\LNbasename{#1}}%
    \item \pycode[\texttt{\expandafter\detokenize\expandafter{\LNfilename}}]{#1}}
  \newcommand{\codeAppendix}{%
    \begin{itemize}
    \foreach \n in {1,...,9} {%
      \ifcsdef{LNcode@\n}{%
        \item \hyperref[code: note \n]{\csuse{LNtitle@\n}} (\LN{\LNtag{\n}})
        \begin{itemize}
        \edef\LNcodelist{\csuse{LNcode@\n}}%
        \expandafter\forcsvlist\expandafter\codeContentsFile\expandafter{\LNcodelist}%
        \end{itemize}}{}}
    \end{itemize}
    \newpage
    \foreach \n in {1,...,9} {%
      \ifcsdef{LNcode@\n}{%
        \section[Lecture Note \n: \csuse{LNtitle@\n}]{%
          \texorpdfstring{\LN{\LNtag{\n}}: \csuse{LNtitle@\n}}{Lecture Note \n: \csuse{LNtitle@\n}}}%
        \label{code: note \n}%
        \edef\LNcodelist{\csuse{LNcode@\n}}%
        \expandafter\forcsvlist\expandafter\codefile\expandafter{\LNcodelist}%
        \newpage}{}}}

  % a lecture note is a chapter: \LN{proofs} prints a link "Lecture Note 2"
  \DeclareRobustCommand{\LN}[1]{\hyperref[#1: chapter]{Lecture Note~\ref*{#1: chapter}}}
  \pdfstringdefDisableCommands{\def\LN#1{Lecture Note \LNnumber{#1}}}

  % running head; "contents" links to the current chapter's table of contents
  \renewcommand{\doHeader}[1]{%
    \renewcommand{\header}{%
      \clap{\raisebox{0pt}{\protect\hyperlink{TOC-\arabic{chapter}}{$\uparrow$\scriptsize\,\textsc{contents}\hspace{1.1in}}}}%
      {\footnotesize \LNTitle{} / Section \thesection, #1}}%
    \markboth{\header}{\header}}

  % number theorems, lemmas, etc. and exercises within chapters (e.g. Lemma 8.1)
  \counterwithin{lemma}{chapter}
  \counterwithin{definition}{chapter}
  \counterwithin{theorem}{chapter}
  \counterwithin{corollary}{chapter}
  \counterwithin{claim}{chapter}
  \counterwithin{observation}{chapter}
  \counterwithin{conjecture}{chapter}
  \counterwithin{exercise}{chapter}
\fi

% %%%%%%%%%%%%%%%%%%%%%%%%%%%%%%%%%%%%%%%%%%%%%%%%%%%%%%%%%%%%%%%%%%%%%%%%
% %%%%%%%%%%%%%%%%%%%%%%%%%%%%%%%%%%%%%%%%%%%%%%%%%%%%%%%%%%%%%%%%%%%%%%%%

\ExplSyntaxOn
\seq_new:N \l_itemize_seq
% \NewDocumentEnvironment{REITEMIZE}{ +b }{
\NewDocumentCommand{\REITEMIZECMD}{ +m }{
  \regex_split:nnN {\c{item}} {#1} \l_itemize_seq
  \seq_pop:NN \l_itemize_seq \l_tmpa_tl% remove the (hopefully) empty first item
  \seq_if_empty:NF \l_itemize_seq {
    \l_tmpa_tl
    \seq_map_inline:Nn \l_itemize_seq {
      \regex_split:nnN {\c{item}} {#1} \l_itemize_seq {
        \regex_extract_once:nnNTF { ^\[([^]]*)\](.*) } { ##1} \l_foo_seq
      { \item[\seq_item:Nn \l_foo_seq {2}]\ITEMIZER{\seq_item:Nn \l_foo_seq {3}} } { \item\ITEMIZER{##1}  }
    }
  }
}
}
\ExplSyntaxOff

\newlength{\ITEMHT}
\setlength{\ITEMHT}{1in}

% \setlength{\fboxsep}{1pt}
\newcommand{\ITEMIZER}[1]{%
  % \framebox{
\adjustbox{valign=t}{\parbox{\textwidth}{\linespread{1.15}\selectfont #1}}~~\adjustbox{valign=t}{\parbox{0pt}{\color{gray}\rule[-\ITEMHT]{0.01pt}{\ITEMHT}}}
% }
  % \begin{tabular}{@{}p{\textwidth}@{\rule[-\ITEMHT]{0pt}{\ITEMHT}}}
  %   #1
  % \end{tabular}
}

\newsavebox\thesmashminipage
\newlength{\thesmashminipagelength}
\newlength{\answerboxlength}
\newlength{\answerboxoverrun}
\newenvironment{smashminipage}[2]
{\setlength\answerboxlength{#2}
\begin{lrbox}{\thesmashminipage}\begin{minipage}[t]{#1}}
{\end{minipage}\end{lrbox}%
\setlength{\thesmashminipagelength}{\dimexpr\ht\thesmashminipage+\dp\thesmashminipage}%
\setlength{\answerboxoverrun}{\dimexpr\thesmashminipagelength-\answerboxlength}%
\ifdim\the\thesmashminipagelength>\the\answerboxlength%
\usebox{\thesmashminipage}%
\smash{\color{red}\hspace*{-0.98\textwidth}\rule[-\answerboxlength]{0.98\textwidth}{1pt}}
\errmessage{Answer exceeds allotted length, please shorten it.}
\else
\usebox{\thesmashminipage}%
\fi
}



\newenvironment{answerBox}[1]{%
  \def\ANSWERBOXHEIGHT{#1}
  \begin{adjustbox}{valign=t}
  \begin{smashminipage}{\textwidth}{\ANSWERBOXHEIGHT}
}{\end{smashminipage}\end{adjustbox}~~~\adjustbox{valign=t}{\parbox{0pt}{\color{gray}\rule[-\ANSWERBOXHEIGHT]{0.01pt}{\ANSWERBOXHEIGHT}\rule[-\ANSWERBOXHEIGHT]{10pt}{0.01pt}}}}

% % \setlist[enumerate]{itemsep=0pt,parsep=0pt,topsep=0pt,leftmargin=100pt}

% \makeatletter 
\LetLtxMacro\SVenumerate\enumerate
\let\endSVenumerate\endenumerate

\newcommand{\injectEnumerate}{
  \let\SValph\alph
  \def\alph{}
  \renewenvironment{enumerate}[1][]
  {\let\alph\SValph \SVenumerate[label=(\alph*),itemsep=2pt,parsep=0pt]\REITEMIZE}{\endREITEMIZE\endSVenumerate}
  % \setlist[enumerate]{before=\Collect@Body\REITEMIZECMD}
}
% \makeatother
\newcommand{\restoreEnumerate}{
  \renewenvironment{enumerate}[1][]
  \let\SValph\alph
  \def\alph{}
  {\let\alph\SValph \SVenumerate[label=(\alph*)]}{\endSVenumerate}
  % \setlist[enumerate]{before=,after=}
}
\newcommand{\injectBlue}{
  \let\SVblue\blue
  \let\endSVblue\endblue
  \def\blue{\SVblue\REITEMIZE}
  \def\endblue{\endREITEMIZE\endSVblue}
}
\newcommand{\restoreBlue}{
  \let\blue\SVblue
  \let\endblue\endSVblue
}
\newcommand{\injectEnumerateBlue}{
  \setlist[enumerate]{before=\injectBlue,after=\restoreBlue,itemsep=2pt,parsep=0pt}
}
\newcommand{\restoreEnumerateBlue}{
  \setlist[enumerate]{before=,after=}
}


% %%%%%%%%%%%%%%%%%%%%%%%%%%%%%%%%%%%%%%%%%%%%%%%%%%%%%%%%%%%%%%%%%%%%%%%%
% %%%%%%%%%%%%%%%%%%%%%%%%%%%%%%%%%%%%%%%%%%%%%%%%%%%%%%%%%%%%%%%%%%%%%%%%

% \NewDocumentEnvironment{env}{ }
% {\Environenv}
% {\endEnvironenv}

% \NewEnviron{REITEMIZE}{\REITEMIZECMD\BODY}

\makeatletter 
\newenvironment{REITEMIZE}{\Collect@Body\REITEMIZECMD}{}
\makeatother 

\newcommand{\NODEone}[1]{\ensuremath{v_{#1}}\xspace}
\newcommand{\NODE}[2]{{\NODEone {#1,#2}}\xspace}
\newcommand{\EDGE}[2]{\ensuremath{(\textsf{#1}, \textsf{#2})}}
\newcommand{\ITEM}{\item}

%%% Local Variables:
%%% mode: latex
%%% End:
